\documentclass[11pt,openany]{book}

\usepackage[a4paper,inner=31mm,outer=27mm,top=28mm,bottom=30mm,headheight=15pt]{geometry}
\usepackage[T1]{fontenc}
\usepackage[utf8]{inputenc}
\usepackage{newtxtext}
\usepackage{mathtools,amssymb,amsthm,mathrsfs}
\usepackage{newtxmath}
\usepackage{microtype}
\usepackage{bm}
\usepackage{enumitem}
\usepackage{booktabs,longtable,array,tabularx}
\usepackage{xcolor}
\usepackage{tikz-cd}
\usepackage[most]{tcolorbox}
\usepackage{fancyhdr}
\usepackage{titlesec}
\usepackage{hyperref}
\usepackage{bookmark}
\usepackage[nameinlink,capitalise,noabbrev]{cleveref}

\definecolor{deepblue}{HTML}{173B57}
\definecolor{midblue}{HTML}{2F6685}
\definecolor{paleblue}{HTML}{EFF6FA}
\definecolor{deepgreen}{HTML}{285943}
\definecolor{palegreen}{HTML}{F0F7F3}
\definecolor{deepred}{HTML}{8B2E2E}
\definecolor{palered}{HTML}{FFF3F1}
\definecolor{softgray}{HTML}{F4F5F6}
\definecolor{perverseviolet}{HTML}{5B3F72}
\definecolor{paleviolet}{HTML}{F6F1FA}
\definecolor{proetaleteal}{HTML}{176B68}
\definecolor{paleproetale}{HTML}{EEF9F8}
\definecolor{goldbrown}{HTML}{80611A}
\definecolor{palegold}{HTML}{FFF9E9}

\hypersetup{
  colorlinks=true,
  linkcolor=deepblue,
  citecolor=deepgreen,
  urlcolor=midblue,
  pdftitle={The Weil Conjectures on the Pro-etale Site through Perverse Sheaves},
  pdfauthor={Prepared for a Lean 4 formalization project},
  pdfsubject={A noncircular proof using integral pro-etale coefficients and the perverse weight engine}
}

\setlist[itemize]{leftmargin=2em,itemsep=0.25em,topsep=0.4em}
\setlist[enumerate]{leftmargin=2.2em,itemsep=0.25em,topsep=0.4em}
\setcounter{tocdepth}{2}
\setcounter{secnumdepth}{3}
\allowdisplaybreaks

\titleformat{\chapter}[display]
  {\normalfont\huge\bfseries\color{deepblue}}
  {\filleft\Large\color{midblue}\chaptername\ \thechapter}
  {1ex}
  {\titlerule\vspace{1.2ex}\filright}
  [\vspace{0.8ex}\titlerule]
\titleformat{\section}
  {\normalfont\Large\bfseries\color{deepblue}}
  {\thesection}{0.7em}{}
\titleformat{\subsection}
  {\normalfont\large\bfseries\color{midblue}}
  {\thesubsection}{0.7em}{}

\pagestyle{fancy}
\fancyhf{}
\fancyhead[LE]{\small\color{midblue}\nouppercase{\leftmark}}
\fancyhead[RO]{\small\color{midblue}\nouppercase{\rightmark}}
\fancyfoot[C]{\small\thepage}
\renewcommand{\headrulewidth}{0.4pt}

\newtheorem{theorem}{Theorem}[chapter]
\newtheorem{proposition}[theorem]{Proposition}
\newtheorem{lemma}[theorem]{Lemma}
\newtheorem{corollary}[theorem]{Corollary}
\theoremstyle{definition}
\newtheorem{definition}[theorem]{Definition}
\newtheorem{example}[theorem]{Example}
\theoremstyle{remark}
\newtheorem{remark}[theorem]{Remark}
\newtheorem{warning}[theorem]{Warning}
\crefname{warning}{warning}{warnings}
\Crefname{warning}{Warning}{Warnings}

\newtcolorbox{foundation}[1][]{
  enhanced,breakable,colback=paleblue,colframe=midblue,
  boxrule=0.6pt,arc=1.2mm,title={Foundational proof capsule},
  fonttitle=\bfseries,#1}
\newtcolorbox{convention}[1][]{
  enhanced,breakable,colback=palegreen,colframe=deepgreen,
  boxrule=0.6pt,arc=1.2mm,title={Convention check},
  fonttitle=\bfseries,#1}
\newtcolorbox{formalization}[1][]{
  enhanced,breakable,colback=softgray,colframe=black!55,
  boxrule=0.5pt,arc=1.2mm,title={Formalization interface},
  fonttitle=\bfseries,#1}
\newtcolorbox{perversebox}[1][]{
  enhanced,breakable,colback=paleviolet,colframe=perverseviolet,
  boxrule=0.6pt,arc=1.2mm,title={Perverse-sheaf checkpoint},
  fonttitle=\bfseries,#1}
\newtcolorbox{proetalebox}[1][]{
  enhanced,breakable,colback=paleproetale,colframe=proetaleteal,
  boxrule=0.6pt,arc=1.2mm,title={Pro-etale simplification},
  fonttitle=\bfseries,#1}
\newtcolorbox{powerbox}[1][]{
  enhanced,breakable,colback=palegold,colframe=goldbrown,
  boxrule=0.6pt,arc=1.2mm,title={Where perversity does the work},
  fonttitle=\bfseries,#1}
\newtcolorbox{auditbox}[1][]{
  enhanced,breakable,colback=palered,colframe=deepred,
  boxrule=0.6pt,arc=1.2mm,title={Logical audit},
  fonttitle=\bfseries,#1}

\newcommand{\Fq}{\mathbf F_q}
\newcommand{\Fbar}{\overline{\mathbf F}_q}
\newcommand{\Ql}{\mathbf Q_\ell}
\newcommand{\Zl}{\mathbf Z_\ell}
\newcommand{\Fl}{\mathbf F_\ell}
\newcommand{\Qbarl}{\overline{\mathbf Q}_\ell}
\newcommand{\Qbar}{\overline{\mathbf Q}}
\newcommand{\A}{\mathbf A}
\newcommand{\PP}{\mathbf P}
\newcommand{\Gm}{\mathbf G_m}
\newcommand{\Ga}{\mathbf G_a}
\newcommand{\Spec}{\operatorname{Spec}}
\newcommand{\Gal}{\operatorname{Gal}}
\newcommand{\Frob}{\operatorname{Frob}}
\newcommand{\Fr}{\operatorname{Fr}}
\newcommand{\Tr}{\operatorname{Tr}}
\newcommand{\rk}{\operatorname{rk}}
\newcommand{\im}{\operatorname{im}}
\newcommand{\id}{\operatorname{id}}
\newcommand{\coker}{\operatorname{coker}}
\newcommand{\End}{\operatorname{End}}
\newcommand{\Hom}{\operatorname{Hom}}
\newcommand{\Ext}{\operatorname{Ext}}
\newcommand{\Supp}{\operatorname{Supp}}
\newcommand{\et}{\mathrm{\acute et}}
\newcommand{\proet}{\mathrm{pro\acute et}}
\newcommand{\perf}{\mathrm{perf}}
\newcommand{\lis}{\mathrm{lisse}}
\newcommand{\cons}{\mathrm{cons}}
\newcommand{\comp}{\mathrm{comp}}
\newcommand{\op}{\mathrm{op}}
\newcommand{\D}{\mathrm D}
\newcommand{\R}{\mathrm R}
\newcommand{\cH}{\mathcal H}
\newcommand{\cF}{\mathcal F}
\newcommand{\cG}{\mathcal G}
\newcommand{\cL}{\mathcal L}
\newcommand{\cK}{\mathcal K}
\newcommand{\cE}{\mathcal E}
\newcommand{\cM}{\mathcal M}
\newcommand{\cP}{\mathcal P}
\newcommand{\cQ}{\mathcal Q}
\newcommand{\one}{\mathbf 1}
\newcommand{\wt}{\operatorname{wt}}
\newcommand{\FT}{\operatorname{FT}}
\newcommand{\IC}{\operatorname{IC}}
\newcommand{\eps}{\varepsilon}
\newcommand{\Perv}{\operatorname{Perv}}
\newcommand{\pH}{{}^{p}\!\mathcal H}
\newcommand{\ptau}{{}^{p}\!\tau}
\newcommand{\Dbc}{\mathrm D^b_c}
\newcommand{\Dcons}{\mathrm D_{\mathrm{cons}}}
\newcommand{\Dcomp}{\mathrm D_{\mathrm{comp}}}
\newcommand{\Dnaive}{\mathrm D_{\mathrm{naive}}}
\newcommand{\Dsolid}{\mathrm D^{\blacksquare}}
\newcommand{\Dindcons}{\mathrm D_{\mathrm{indcons}}}
\newcommand{\Lhat}{\widehat{\mathbf L}}
\newcommand{\Rlim}{\mathrm R\!\varprojlim}
\DeclareMathOperator{\Fix}{Fix}
\DeclareMathOperator{\rad}{rad}
\DeclareMathOperator{\charpoly}{charpoly}

% Compatibility aliases for the independently reviewed chapters retained in
% this edition.  The new pro-etale chapters replace the old adic-system layer.
\makeatletter
\AtBeginDocument{
  \csletcs{r@chap:etale-foundations}{r@chap:proetale-site}
  \csletcs{r@chap:etale-foundations@cref}{r@chap:proetale-site@cref}
  \csletcs{r@ch:etale-foundations}{r@chap:proetale-site}
  \csletcs{r@ch:etale-foundations@cref}{r@chap:proetale-site@cref}
  \csletcs{r@chap:six-operations}{r@chap:proetale-six}
  \csletcs{r@chap:six-operations@cref}{r@chap:proetale-six@cref}
  \csletcs{r@ch:six-operations}{r@chap:proetale-six}
  \csletcs{r@ch:six-operations@cref}{r@chap:proetale-six@cref}
  \csletcs{r@chap:adic-weil}{r@chap:proetale-weil}
  \csletcs{r@chap:adic-weil@cref}{r@chap:proetale-weil@cref}
  \csletcs{r@ch:adic-weil}{r@chap:proetale-weil}
  \csletcs{r@ch:adic-weil@cref}{r@chap:proetale-weil@cref}
  \csletcs{r@chap:trace-formula}{r@ch:trace-formula}
  \csletcs{r@chap:trace-formula@cref}{r@ch:trace-formula@cref}
  \csletcs{r@chap:duality}{r@ch:duality}
  \csletcs{r@chap:duality@cref}{r@ch:duality@cref}
  \csletcs{r@chap:algebra}{r@chap:algebra-core}
  \csletcs{r@chap:algebra@cref}{r@chap:algebra-core@cref}
}
\makeatother

\begin{document}

\frontmatter
\begin{titlepage}
  \centering
  \vspace*{1.7cm}
  {\Huge\bfseries\color{deepblue}
    The Weil Conjectures\\[0.25em]
    on the Pro-etale Site\par}
  \vspace{0.55cm}
  {\LARGE\bfseries\color{perverseviolet}
    through Perverse Sheaves\par}
  \vspace{0.9cm}
  {\Large Integral \(\ell\)-adic coefficients, six operations, weights,
    and Fourier transform\par}
  \vspace{1.2cm}
  {\large A detailed and noncircular natural-language proof\par}
  \vfill
  \begin{tcolorbox}[
    width=0.88\textwidth,colback=paleproetale,colframe=proetaleteal,
    boxrule=0.7pt,arc=1.5mm]
  \centering
  The pro-etale site makes \(\Zl\) an honest coefficient sheaf.\par
  Route B builds the six operations on the solid ambient category,\par
  then proves their restriction to constructible complexes.\par
  Perverse sheaves turn the hard weight argument into exact
  local-to-global operations.
  \end{tcolorbox}
  \vfill
  {\large Prepared as a mathematical specification for reusable\\
  Lean 4 and mathlib formalization\par}
  \vspace{0.7cm}
  {\large August 2026\par}
\end{titlepage}

\chapter*{Preface: two changes of language}
\addcontentsline{toc}{chapter}{Preface: two changes of language}

These notes prove the Weil conjectures by making two changes of language
before beginning the arithmetic argument.  The first change is
coefficient-theoretic:
\[
  X_{\et}\text{ with towers }
  (\mathbf Z/\ell^n)_{n\geq1}
  \quad\rightsquigarrow\quad
  X_{\proet}\text{ with an honest }\Zl\text{-sheaf}.
\tag{P.1}\label{eq:preface-proetale-change}
\]
The second is cohomological:
\[
  \text{ordinary cohomological degree}
  \quad\rightsquigarrow\quad
  \text{middle perverse degree}.
\tag{P.2}\label{eq:preface-perverse-change}
\]
The first replaces the Artin--Rees quotient presentation in the proof spine
by a comparison-equivalent honest-sheaf working model.  The second makes
dimension, extension across a boundary, duality, and the Fourier transform
interact in the correct abelian category.
Together they give a cleaner proof and, more importantly for formalization, a
cleaner collection of reusable interfaces.

\begin{proetalebox}
On \(X_{\proet}\), the topological ring \(\Zl\) determines a genuine sheaf
\[
 \widehat{\mathbf Z}_{\ell,X}
   =\varprojlim_n\mathbf Z/\ell^n\mathbf Z.
\]
Thus coefficients and cohomology are operations on actual pro-etale
sheaves.  The naive category of all such sheaves is nevertheless too large:
proper base change fails there even for finite coefficients.  Route B first
passes to the reflective localization of solid pro-etale complexes,
\[
 \Dsolid(X,\Zl)\subset\Dnaive(X,\Zl),
\]
constructs the coherent six operations on that presentable ambient category,
and only then restricts them to \(\Dcons(X,\Zl)\).  Solidity is a continuity
condition, not a finite-stratification condition.  The historical Verdier
quotient of coefficient towers remains a comparison-equivalent model on
constructible objects.
\end{proetalebox}

\begin{powerbox}
Perverse sheaves are not merely used in the final three-line shift
calculation.  They perform three essential tasks in the proof:
\begin{enumerate}
\item affine Artin vanishing becomes one-sided perverse \(t\)-exactness;
\item Fourier--Deligne transform becomes an exact equivalence of perverse
      hearts;
\item the global weight theorem is expressed by the canonical objects
      \(\pH^a(Rf_!\cP)\), so boundary and support terms are controlled by
      subobjects, quotients, and intermediate extension.
\end{enumerate}
\end{powerbox}

\section*{What the pro-etale site simplifies}

Fix a prime \(\ell\) different from the characteristic of the ground field.
More generally, let \(E/\Ql\) be finite, let \(\mathcal O_E\) be its ring of
integers, and let \(\varpi\) be a uniformizer.  The coefficient ladder used in
the notes is
\[
\begin{tikzcd}[column sep=large]
 \widehat{\mathbf Z}_{\ell,X}
   \arrow[r] &
 \widehat{\mathcal O}_{E,X}
   \arrow[r,"\varpi^{-1}"] &
 E_X
   \arrow[r] &
 (\Qbarl)_X .
\end{tikzcd}
\tag{P.3}\label{eq:preface-coefficient-ladder}
\]
The last term is shorthand for descent to some finite \(E\), followed by
extension of scalars.  It is not treated as a complete noetherian coefficient
ring.

For integral coefficients, the Route B ambient category is
\(\Dsolid(X,\widehat{\mathcal O}_E)\).  Inside it, the category used by the
arithmetic proof consists of constructible, derived
\(\varpi\)-complete complexes.  For such a complex,
\[
  K\simeq
  \Rlim_n\bigl(K\otimes^L_{\widehat{\mathcal O}_E}
                    \mathcal O_E/\varpi^n\bigr).
\tag{P.4}\label{eq:preface-derived-completion}
\]
Weakly contractible pro-etale affines make this formula behave as the
geometric intuition suggests.  On a topologically noetherian scheme,
constructibility is equivalently local perfectness along a finite
stratification.

Formula~\eqref{eq:preface-derived-completion} is a compatibility and
recognition theorem; it is not the definition of an ambient operation.
Solid rigidity constructs the six operations before constructibility is
imposed.  Finiteness and constructible duality are then used to prove that
those operations preserve the constructible subcategory.

This is a simplification of the categorical architecture, not a magical
proof of the geometric theorems.  Proper and smooth base change, generic
local acyclicity, constructibility of direct images, absolute purity,
Verdier duality, and the trace formula remain deep.  The known proof of
ambient solid base change still uses torsion constructible generators and
classical torsion base change; Route B removes the levelwise definition, not
the finite theorem's mathematical content.

\begin{auditbox}
For a separated finite-presentation map, the exceptional direct image
\(Rf_!\) is first constructed in the ambient solid category by the coherent
compactification formalism and then restricted to constructibles.  It is
generally \emph{not} the derived functor of an abelian sheaf-level \(f_!\).
Pullback on the constructible complete subcategory agrees with completed
pullback.  Neither completion alone nor the unrestricted naive pro-etale
category satisfies the required ambient axioms; the counterexample and the
solid correction are proved in \cref{chap:proetale-six}.
\end{auditbox}

\section*{Why the coefficients change during the proof}

The integral category supplies lattices, reduction maps, continuous
cohomology, and perfect complexes.  The arithmetic weight argument is carried
out after passing to \(E\) or \(\Qbarl\).  There are two reasons.

First, an Artin--Schreier character
\(\psi:\mathbf F_p\to E^\times\) may require enlarging \(E\) so that it
contains the relevant \(p\)-th roots of unity.  Its rank-one kernel naturally
has an \(\mathcal O_E\)-lattice, but need not be defined over \(\Zl\).

Second, the familiar middle perverse \(t\)-structure is Verdier self-dual over
a field.  Over the discrete valuation ring \(\mathcal O_E\), duality exchanges
the two integral perversities \(p\) and \(p^+\).  Consequently integral
intermediate extensions come in two variants, and an integral decomposition
theorem is false in general.  These notes develop enough of the integral
pair to make reduction and lattices honest, then rationalize before invoking
self-duality, pure semisimplicity, or decomposition.

\section*{The perverse proof engine}

Let \(f_0:X_0\to Y_0\) be separated of finite type over \(\Fq\).  The central
theorem is first stated for an \(E\)-perverse Weil sheaf:
\[
 \cP_0\text{ has weights }\leq w
 \quad\Longrightarrow\quad
 \pH^a(Rf_{0!}\cP_0)
 \text{ has weights }\leq w+a.
\tag{P.5}\label{eq:preface-perverse-direct-image}
\]
The derived version follows from the perverse truncation tower.  This order
matters: perversity is the proof language of the theorem, not a restatement
attached after an ordinary-cohomology proof.

The difficult step is the relative curve estimate.  For a pure lisse sheaf
\(\cL\) on a smooth open curve \(j:U\hookrightarrow C\), one forms
\[
  \cP=j_{!*}(\cL[1]).
\tag{P.6}\label{eq:preface-curve-perverse-object}
\]
The Fourier--Deligne transform, normalized by a shift \([1]\), is exact on
the perverse heart.  Its stalk traces are exponential sums, and Plancherel
turns their second moment into a geometric norm.  Local Fourier transform
records the boundary contribution.  The resulting stalk estimate gives the
weight bound for \(H_c^1(U,\cL)\).  Relative Noether normalization,
recollement, and affine one-sided exactness then promote the curve result to
\eqref{eq:preface-perverse-direct-image}.

The complete noncircular chain is
\[
\begin{tikzcd}[column sep=small,row sep=large]
\substack{\text{pro-etale coefficients}\\
          \text{and six operations}}
 \arrow[r] &
\substack{\text{perverse recollement}\\
          \text{and Artin vanishing}}
 \arrow[r] &
\substack{\text{Fourier exactness}\\
          \text{and Plancherel}}
 \arrow[d] \\
\substack{\text{proper purity}\\
          \text{and IC purity}}
 \arrow[d] &
\substack{Rf_!\text{ upper}\\
          \text{weights}}
 \arrow[l] &
\substack{\text{perverse curve}\\
          \text{estimate}}
 \arrow[l] \\
\substack{\text{pure perverse}\\
          \text{semisimplicity}}
 \arrow[r] &
\substack{\text{decomposition}\\
          \text{theorem}}
 &
\\[-0.4em]
\substack{\text{smooth proper}\\
          H^i\text{ pure of weight }i}
 \arrow[r] &
\substack{\text{trace, duality,}\\
          \text{weight separation}}
 \arrow[r] &
\text{Weil conjectures.}
\end{tikzcd}
\tag{P.7}\label{eq:preface-dependency-chain}
\]
The smooth-proper purity node is obtained directly from proper purity; it
does not depend on the decomposition node.  Both are consequences of the
same stronger perverse theorem.

\section*{The stronger theorem proved first}

The natural endpoint of the perverse argument is stronger than the classical
smooth statement.  If \(X_0\) is proper and pure dimensional, and if
\(\cL_0\) is a pure lisse sheaf on a smooth dense stratum, then
\[
  \IC_{X_0}(\cL_0)=j_{!*}(\cL_0[d])
\]
is pure, and its intersection cohomology is pure with the predicted shift.
For a proper map \(f_0\), every
\(\pH^a(Rf_{0*}\IC_{X_0}(\cL_0))\) is pure and becomes semisimple after
geometric base change.  After base change to \(\Fbar\) and after forgetting
the Weil structure, the derived direct image splits noncanonically into these
canonical perverse cohomology objects.  No Frobenius-equivariant splitting
over \(\Fq\) is asserted; the perverse cohomology objects themselves, unlike
a chosen splitting, remain canonical and carry their Weil structures.

For a smooth \(d\)-fold,
\(\IC_{X_0}(E)=E[d]\).  Taking the structural morphism
\(a_0:X_0\to\Spec\Fq\) gives
\[
\begin{aligned}
 E[d]\text{ is pure of weight }d
 &\Longrightarrow Ra_{0*}E[d]\text{ is pure of weight }d,\\
 &\Longrightarrow
 H^{i-d}(Ra_{0*}E[d])
   =H^i(X,E)\text{ is pure of weight }i.
\end{aligned}
\tag{P.8}\label{eq:preface-terminal-calculation}
\]
This is the Riemann-hypothesis part.  The trace formula gives rationality,
Poincare duality gives reciprocal roots and the functional equation, and
universal-realization purity separates the rational zeta function into
integral, \(\ell\)-independent factors.

\section*{Completeness contract}

“Self-contained” in these notes has the following exact meaning.

\begin{enumerate}
\item Every coefficient category, functor, Weil structure, shift, twist,
      weight convention, perverse truncation, and comparison used in the
      proof is defined.
\item The implications from the foundational theorem dossier to
      \eqref{eq:preface-terminal-calculation} and then to every clause of the
      Weil conjectures are proved.
\item SGA-scale construction theorems are never hidden under the phrase
      “the usual formalism.”  Their exact hypotheses, compatibilities, and
      primary locations are recorded in the appendices.
\item The pro-etale/classical-adic equivalence is an explicit comparison
      theorem.  It is not treated as definitional equality.
\item No smooth-proper purity, decomposition theorem, integral Frobenius
      factor, or \(\ell\)-independence statement is used upstream of its
      proof.
\end{enumerate}

The foundational dossier distinguishes four statuses:

\begin{center}
\begin{tabular}{cl}
\toprule
\(\mathbf F\) & deep foundational theorem used with an exact source and
                interface,\\
\(\mathbf C\) & construction or comparison proved in the pro-etale chapters,\\
\(\mathbf P\) & theorem proved in full in the mathematical proof spine,\\
\(\mathbf A\) & final assembly from already proved results.\\
\bottomrule
\end{tabular}
\end{center}

These letters are expository metadata only.  They never occur as hypotheses
of the final theorem.

\section*{How to read the notes}

Readers interested in the new coefficient architecture should begin with
\cref{chap:proetale-site,chap:proetale-coefficients,chap:proetale-six,%
chap:proetale-weil}.  Readers who know the six operations may begin with the
middle perverse \(t\)-structure and the integral-lattice warning, then read
the perverse-engine and Fourier chapters.  The shortest logically complete
route to the final theorem is:
\[
 \text{perverse direct-image theorem}
 \longrightarrow
 \text{proper purity}
 \longrightarrow
 \text{terminal calculation}
 \longrightarrow
 \text{algebraic endgame}.
\]

The principal primary sources are Bhatt--Scholze and
Hemo--Richarz--Scholbach for pro-etale constructibles,
Ruimy--Tubach--Wolf \cite{RuimyTubachWolfSolid} for the solid Route B
ambient formalism, SGA and the finiteness/duality literature for the
geometric inputs, BBD \cite{BBD} for perverse sheaves,
Deligne \cite{DeligneWeilII} for weights, and Laumon
\cite{LaumonFourier} for the Fourier proof.  The exact theorem crosswalk is
given in the final appendix.

\tableofcontents

\mainmatter
\part{The target and the algebraic endgame}
\chapter{The Weil conjectures: exact statement and conventions}
\label{chap:statement}

\section{Arithmetic data attached to a variety}

Let \(k_0=\Fq\), let \(k=\Fbar\), and let \(X_0\) be a scheme of finite
type over \(k_0\).  For every \(n\geq1\), the set
\(X_0(\mathbf F_{q^n})\) is finite.  We write
\[
  N_n(X_0)=\#X_0(\mathbf F_{q^n}).
\]
The Hasse--Weil zeta function packages all these point counts into one
formal series.

\begin{definition}
\label{def:hasse-weil-zeta}
The zeta function of \(X_0\) is
\[
 Z(X_0,T)
 =\exp\left(\sum_{n\geq1}N_n(X_0)\frac{T^n}{n}\right)
 \in1+T\mathbf Q[[T]].
\]
\end{definition}

The denominators \(n\) make it nonobvious from this definition that the
coefficients are integers.  Integrality becomes transparent after
rewriting the series as an Euler product.  If \(x\) is a closed point,
its residue field is finite over \(k_0\), and we set
\[
  \deg(x)=[\kappa(x):k_0].
\]
The closed-point formula, proved without cohomology in
\cref{prop:euler-product}, is
\[
  Z(X_0,T)=\prod_{x\in|X_0|}
       (1-T^{\deg(x)})^{-1}\in1+T\mathbf Z[[T]].
  \tag{1.1}\label{eq:zeta-euler-intro}
\]
This is a formal product: only finitely many factors can affect a fixed
coefficient.  Analytic convergence is neither assumed nor needed.

\section{Frobenius and Tate twists}

Put \(X=X_0\times_{k_0}k\).  The symbol
\[
 F\colon X\longrightarrow X
\]
always denotes geometric Frobenius, inverse to the \(q\)-power arithmetic
Frobenius.  On geometric points it generates the same finite orbits as
arithmetic Frobenius, but in the opposite direction.  On etale
cohomology, taking the inverse changes every eigenvalue, so the
distinction is indispensable.

\begin{convention}
Our Tate-twist convention is
\[
 F\mid\Ql(1)=q^{-1},\qquad
 F\mid\Ql(-1)=q.
\]
Thus \(F\) acts by \(q^d\) on \(\Ql(-d)\).  If \(X_0\) is smooth and
proper of dimension \(d\), the trace map identifies
\[
 H^{2d}_{\et}(X,\Ql)\simeq\Ql(-d)
\]
when \(X\) is connected.  The top reciprocal Frobenius root is therefore
\(q^d\), not \(q^{-d}\).
\end{convention}

Fix a prime \(\ell\neq\operatorname{char}(k_0)\).  For a smooth proper
\(X_0\), define
\[
 V_{i,\ell}(X_0)=H^i_{\et}(X,\Ql),\qquad
 P_{i,\ell}(T)=\det(1-TF\mid V_{i,\ell}(X_0)).
 \tag{1.2}\label{eq:def-Pi}
\]
These spaces are finite-dimensional and vanish outside
\(0\leq i\leq2d\).  If, over \(\Qbarl\),
\[
 P_{i,\ell}(T)=\prod_{j=1}^{b_{i,\ell}}
       (1-\alpha_{ij}T),
\]
then the \(\alpha_{ij}\) are called the reciprocal roots, or Frobenius
eigenvalues.  The ordinary roots of the polynomial are
\(\alpha_{ij}^{-1}\).

\section{Weights: the initial and final formulations}

There are two logically different formulations of purity.  The
sheaf-theoretic theorem first says that an eigenvalue
\(\alpha\in\Qbarl^\times\) has
\[
 |\iota(\alpha)|=q^{w/2}
 \tag{1.3}\label{eq:universal-iota-weight}
\]
for every abstract field isomorphism
\(\iota\colon\Qbarl\xrightarrow{\sim}\mathbf C\).  It does not initially
assert that \(\alpha\) is algebraic over \(\mathbf Q\).  After
no-cancellation and rational factor recovery are proved, \(\alpha\) is
known to be an algebraic number, and \eqref{eq:universal-iota-weight}
becomes the classical assertion
\[
 |\tau(\alpha)|=q^{w/2}
\]
for every embedding
\(\tau\colon\mathbf Q(\alpha)\hookrightarrow\mathbf C\).

\begin{warning}
It is not enough to choose one isomorphism
\(\Qbarl\simeq\mathbf C\).  The final rationality of each weight factor
uses stability under \(\operatorname{Aut}(\Qbarl/\mathbf Q)\).  That
stability follows because \eqref{eq:universal-iota-weight} is quantified
over every \(\iota\): composing \(\iota\) with an automorphism preserves
the same absolute value.  This quantifier is a mathematical input, not a
matter of exposition.
\end{warning}

\section{The classical theorem}

\begin{theorem}[The Weil conjectures]
\label{thm:weil-conjectures-statement}
Let \(X_0\) be a smooth, projective, geometrically connected variety of
pure dimension \(d\) over \(\Fq\).  Then:

\begin{enumerate}
\item \emph{Cohomological rationality.}  For every auxiliary prime
      \(\ell\neq\operatorname{char}(\Fq)\),
      \[
        Z(X_0,T)
        =\prod_{i=0}^{2d}P_{i,\ell}(T)^{(-1)^{i+1}}.
        \tag{1.4}\label{eq:weil-cohomological-factorization}
      \]
      In particular \(Z(X_0,T)\) is a rational function.

\item \emph{Integral, realization-independent factors.}  There are
      polynomials \(P_i(T)\in\mathbf Z[T]\), independent of \(\ell\),
      such that \(P_i=P_{i,\ell}\) for every auxiliary prime.  They have
      constant term \(1\), and
      \[
        P_0(T)=1-T,\qquad P_{2d}(T)=1-q^dT.
        \tag{1.5}\label{eq:weil-endpoints}
      \]
      Consequently
      \[
        Z(X_0,T)=
        \frac{P_1(T)P_3(T)\cdots}
        {(1-T)P_2(T)P_4(T)\cdots(1-q^dT)},
      \]
      with only degrees in the interval \([0,2d]\) included.

\item \emph{Riemann hypothesis.}  Every reciprocal root \(\alpha\) of
      \(P_i\) is an algebraic integer, and for every embedding
      \(\tau\colon\mathbf Q(\alpha)\hookrightarrow\mathbf C\),
      \[
        |\tau(\alpha)|=q^{i/2}.
        \tag{1.6}\label{eq:weil-rh}
      \]
      Equivalently, every ordinary root of \(P_i(T)\) has absolute value
      \(q^{-i/2}\) under every complex embedding.

\item \emph{Betti numbers.}  The integer
      \[
        b_i=\deg P_i=\dim_{\Ql}H^i_{\et}(X,\Ql)
        \tag{1.7}\label{eq:weil-betti}
      \]
      is independent of \(\ell\).  Poincare duality gives
      \(b_i=b_{2d-i}\), and geometric connectedness gives
      \(b_0=b_{2d}=1\).  If \(X_0\) is obtained by good reduction from a
      smooth proper variety in characteristic zero, then \(b_i\) equals
      the \(i\)-th singular Betti number of the associated compact
      complex manifold.

\item \emph{Functional equation.}  If
      \[
        \chi=\sum_{i=0}^{2d}(-1)^ib_i,
      \]
      then a sign \(\varepsilon\in\{1,-1\}\) exists such that
      \[
        Z(X_0,q^{-d}T^{-1})
        =\varepsilon q^{d\chi/2}T^\chi Z(X_0,T).
        \tag{1.8}\label{eq:weil-functional-equation}
      \]
      If \(D_i=\det(F\mid H^i_{\et}(X,\Ql))\), then
      \[
        \varepsilon=(-1)^\chi q^{-d\chi/2}
          \prod_{i=0}^{2d}D_i^{(-1)^i}.
        \tag{1.9}\label{eq:epsilon-determinant}
      \]
      Poincare duality proves that the right side is a sign.
\end{enumerate}
\end{theorem}

\section{What the Betti-number assertion means}

In characteristic zero, a smooth projective complex variety has
topological cohomology groups whose dimensions are its classical Betti
numbers.  An arbitrary variety over a finite field need not possess a
lift to characteristic zero, so one cannot define its \(b_i\) by first
choosing such a lift.  The unconditional definition in
\eqref{eq:weil-betti} is etale.  The substantive assertion is that the
dimension is the same for every \(\ell\).  When a smooth proper lift
exists, smooth proper base change and the etale--singular comparison
theorem prove agreement with topology.

The theorem does not assert hard Lefschetz as part of its statement.
Hard Lefschetz gives further structure on the \(b_i\), but it is not
needed in the Laumon route used in these notes.  Poincare duality alone
is enough for the symmetry \(b_i=b_{2d-i}\) and for the functional
equation.

\section{The geometric-connectedness caveat}

For any smooth proper \(X_0\),
\[
 H^0_{\et}(X,\Ql)
\]
is the vector space of functions on the finite set \(\pi_0(X)\), and
\(F\) permutes this set.  Hence in general
\[
 P_{0,\ell}(T)=\det(1-TF\mid\Ql^{\pi_0(X)}).
 \tag{1.10}\label{eq:general-P0}
\]
It is \(1-T\) precisely when \(X\) is connected.  Duality makes the top
factor the same permutation determinant with each eigenvalue multiplied
by \(q^d\).

For instance, \(\Spec\mathbf F_{q^m}\) is connected over \(\Fq\), but
its geometric fiber consists of \(m\) points cyclically permuted by
Frobenius.  Its \(P_0\) is \(1-T^m\).  Thus ``connected'' cannot replace
``geometrically connected'' in the endpoint assertion.  The
cohomological trace formula, purity, and factor arguments remain true
without this hypothesis; only the simplified endpoints
\eqref{eq:weil-endpoints} must be replaced by
\eqref{eq:general-P0} and its dual.

For a smooth scheme, geometric connectedness implies geometric
irreducibility.  Indeed, a regular scheme has disjoint irreducible
components: if two components met, the local ring at an intersection
point would have more than one minimal prime and would not be a domain,
whereas a regular local ring is a domain.  The components are therefore
open and closed.  This explains why older statements often say
``nonsingular projective absolutely irreducible variety.''

\section{Logical order of the proof}

The proof does not establish the five clauses of
\cref{thm:weil-conjectures-statement} independently.
It follows this directed chain:
\[
\begin{array}{c}
\text{trace formula}+\text{finite-dimensional cohomology}\\
\Downarrow\\
\text{\(\Ql\)-rational determinant expression for }Z\\
\Downarrow\\
\text{Hankel descent gives }Z\in\mathbf Q(T),\\[1mm]
\text{while Poincare duality gives the functional equation};
\end{array}
\]
\[
\begin{array}{c}
\text{Fourier--Deligne transform}+\text{curve estimate}\\
\Downarrow\\
\text{Deligne--Laumon direct-image theorem}\\
\Downarrow\\
\text{smooth-proper universal realization-purity}\\
\Downarrow\\
\text{no cancellation between weights}\\
\Downarrow\\
\text{algebraicity, rational factors, integral factors,
and \(\ell\)-independence}.
\end{array}
\]
In particular, the individual \(P_{i,\ell}\) are not assumed integral,
rational over \(\mathbf Q\), or \(\ell\)-independent during the proof of
purity.  Their alternating product is known earlier, but it is purity
which makes the factors recoverable from that product.

The route is exclusively the Laumon--Deligne Weil II route.  We do not
use the Weil I chain through Lefschetz pencils, Picard--Lefschetz
monodromy, Chebotarev, or tensor-power amplification.  Keeping one route
prevents deep inputs from being imported twice under different names and
makes the dependency audit suitable for formalization.

\begin{formalization}
A reusable theorem statement should separate the following structures:
\begin{itemize}
\item a counting sequence \(n\mapsto N_n\) and a locally finite
      collection of closed-point degrees;
\item a bounded graded family of finite-dimensional Frobenius modules;
\item a trace-formula equality connecting the first two;
\item universal realization-purity, arithmetic algebraicity, and
      algebraic integrality as three different predicates;
\item a perfect graded similitude pairing with multiplier \(q^d\);
\item geometric connectedness as the hypothesis which simplifies the
      endpoint permutation representation.
\end{itemize}
The polynomial \(P_i=\det(1-TF_i)\) should be constructed from the
Frobenius module rather than stored as redundant data.  The final
\(\ell\)-independent polynomial should be returned by a uniqueness
theorem for weight factors of the integral zeta function.
\end{formalization}

\chapter{The algebraic core}
\label{chap:algebra-core}

This chapter contains the part of the proof which is independent of
schemes and cohomology.  Its inputs are finite Frobenius sets, formal
power series, sequences, finite-dimensional vector spaces, and integral
lattices.  Separating this material has two benefits.  It makes clear
which conclusions really use the geometric theorems, and it produces a
library reusable for other zeta functions and compatible systems.

\section{Formal logarithms, exponentials, and products}

Let \(R\) be a commutative \(\mathbf Q\)-algebra.  The \(T\)-adic
filtration makes \(R[[T]]\) complete and separated.  Therefore
\[
 \exp(U)=\sum_{n\geq0}\frac{U^n}{n!}
 \quad(U\in TR[[T]])
\]
and
\[
 \log(1+U)=\sum_{n\geq1}(-1)^{n+1}\frac{U^n}{n}
 \quad(U\in TR[[T]])
\]
are defined coefficientwise.  They are inverse bijections between
\(TR[[T]]\) under addition and \(1+TR[[T]]\) under multiplication.
Every coefficient depends on finitely many terms, so no analytic
convergence is involved.

An infinite product \(\prod_{\lambda}(1+u_\lambda)\), with
\(u_\lambda\in TR[[T]]\), is called \(T\)-adically locally finite if, for
each \(N\), only finitely many \(u_\lambda\) are nonzero modulo
\(T^{N+1}\).  Such a product is well defined: its image in each quotient
\(R[[T]]/(T^{N+1})\) is a finite product, and these finite products are
compatible as \(N\) varies.  Formal logarithm then converts it into the
locally finite sum of the logarithms.  This elementary completeness
argument is what justifies the closed-point Euler product.

\section{Closed points as Frobenius orbits}

Let \(Y_0\) be of finite type over \(\Fq\), and put
\(Y=Y_0\times_{\Fq}\Fbar\).  A geometric point of \(Y\) is defined over
some finite extension of \(\Fq\), so its orbit under geometric
Frobenius \(F\) is finite.  If the least positive integer fixing the
point is \(m\), the orbit has \(m\) elements and descends to a closed
point with residue field \(\mathbf F_{q^m}\).  Conversely, the geometric
fiber of a closed point \(y\) with
\(\kappa(y)=\mathbf F_{q^m}\) consists of the \(m\) embeddings of its
residue field into \(\Fbar\), and Frobenius acts transitively on them.
Thus:
\[
 \{\text{closed points of }Y_0\}
 \longleftrightarrow
 \{\text{finite Frobenius orbits in }Y(\Fbar)\},
 \tag{2.1}\label{eq:closed-orbit-bijection}
\]
with degree equal to orbit length.

Let \(a_m\) be the number of closed points of degree \(m\), and let
\(N_n=\#Y_0(\mathbf F_{q^n})\).

\begin{lemma}[Counting fixed points orbit by orbit]
\label{lem:point-count-orbits}
For every \(n\geq1\),
\[
  N_n=\sum_{m\mid n}ma_m.
  \tag{2.2}\label{eq:point-orbit-count}
\]
\end{lemma}

\begin{proof}
An \(\mathbf F_{q^n}\)-point is a geometric point fixed by \(F^n\).  On
an orbit of length \(m\), the permutation \(F^n\) is the identity if
\(m\mid n\) and has no fixed point otherwise.  In the first case all
\(m\) points contribute.  Summation over the disjoint orbits gives
\eqref{eq:point-orbit-count}.
\end{proof}

There are finitely many closed points of every bounded degree.  Indeed,
if \(m\leq r\), all geometric points above a degree-\(m\) point are
rational over \(\mathbf F_{q^{r!}}\).  They belong to the finite set
\(Y_0(\mathbf F_{q^{r!}})\).  This verifies the local finiteness needed
for the next product.

\begin{proposition}[Closed-point Euler product]
\label{prop:euler-product}
For every finite-type \(\Fq\)-scheme \(Y_0\),
\[
 \exp\left(\sum_{n\geq1}N_n\frac{T^n}{n}\right)
 =\prod_{y\in|Y_0|}(1-T^{\deg(y)})^{-1}
 \quad\text{in }1+T\mathbf Z[[T]].
 \tag{2.3}\label{eq:closed-point-product}
\]
\end{proposition}

\begin{proof}
The preceding bounded-degree argument makes the right side a
\(T\)-adically locally finite product.  Each factor has integral
coefficients, so the product belongs to \(1+T\mathbf Z[[T]]\).

Compute its logarithm in \(\mathbf Q[[T]]\).  Grouping closed points by
degree and using
\(-\log(1-U)=\sum_{s\geq1}U^s/s\), we obtain
\[
\begin{aligned}
 \log\prod_{y\in|Y_0|}(1-T^{\deg(y)})^{-1}
 &=\sum_{m\geq1}a_m\sum_{s\geq1}\frac{T^{ms}}s\\
 &=\sum_{n\geq1}\left(\sum_{m\mid n}\frac{a_m}{n/m}\right)T^n\\
 &=\sum_{n\geq1}\frac1n
        \left(\sum_{m\mid n}ma_m\right)T^n\\
 &=\sum_{n\geq1}N_n\frac{T^n}{n},
\end{aligned}
\]
where the last equality is \cref{lem:point-count-orbits}.  Formal
exponential is inverse to formal logarithm, so
\eqref{eq:closed-point-product} follows.
\end{proof}

\begin{corollary}[M\"obius recovery of closed points]
\label{cor:mobius-closed-points}
The point counts determine the number of degree-\(m\) closed points:
\[
 a_m=\frac1m\sum_{r\mid m}\mu(r)N_{m/r}.
\]
\end{corollary}

\begin{proof}
Apply ordinary M\"obius inversion to
\(N_n=\sum_{m\mid n}ma_m\).
\end{proof}

The quotient in this formula is automatically a nonnegative integer
because it counts orbits.  The integrality is much less evident if one
starts with an arbitrary numerical sequence \(N_n\); it is a geometric
constraint encoded by the Euler product.

\section{The determinant--trace identity}

Let \(K\) be a characteristic-zero field, let \(V\) be a
finite-dimensional \(K\)-vector space, and let \(u\in\End_K(V)\).

\begin{theorem}[Determinant as the exponential of traces]
\label{prop:det-trace}
In \(K[[T]]\),
\[
 \det(1-Tu)^{-1}
 =\exp\left(\sum_{n\geq1}\Tr(u^n)\frac{T^n}{n}\right).
 \tag{2.4}\label{eq:det-trace}
\]
Equivalently,
\[
 -\frac{d}{dT}\log\det(1-Tu)
 =\sum_{n\geq1}\Tr(u^n)T^{n-1}.
 \tag{2.5}\label{eq:log-derivative-trace}
\]
\end{theorem}

\begin{proof}
Scalar extension preserves determinant, trace, and equality of formal
series, so extend \(K\) to an algebraic closure.  Choose a basis in which
\(u\) is upper triangular.  Let the diagonal entries, with algebraic
multiplicity, be \(\lambda_1,\ldots,\lambda_r\).  Nilpotent entries above
the diagonal affect neither the determinant nor the traces of powers:
\[
 \det(1-Tu)=\prod_{j=1}^r(1-\lambda_jT),\qquad
 \Tr(u^n)=\sum_{j=1}^r\lambda_j^n.
\]
Thus
\[
\begin{aligned}
 -\log\det(1-Tu)
 &=-\sum_{j=1}^r\log(1-\lambda_jT)\\
 &=\sum_{j=1}^r\sum_{n\geq1}
       \frac{\lambda_j^nT^n}{n}\\
 &=\sum_{n\geq1}\Tr(u^n)\frac{T^n}{n}.
\end{aligned}
\]
Exponentiating proves \eqref{eq:det-trace}; differentiating proves
\eqref{eq:log-derivative-trace}.  Since both sides before scalar
extension have coefficients in \(K\), the identity descends to \(K\).
\end{proof}

\begin{remark}
There is also a basis-free proof.  Jacobi's determinant formula gives
\[
 \frac{d}{dT}\det(1-Tu)
 =-\det(1-Tu)\Tr((1-Tu)^{-1}u).
\]
The geometric series
\((1-Tu)^{-1}=\sum_{n\geq0}T^nu^n\) then yields
\eqref{eq:log-derivative-trace}.  The triangular proof was chosen because
it simultaneously records the reciprocal roots and makes clear that
semisimplicity is unnecessary.
\end{remark}

\begin{corollary}[Alternating determinant formula]
\label{cor:alternating-det-trace}
For a bounded graded family \((V_i,u_i)\),
\[
 \prod_i\det(1-Tu_i)^{(-1)^{i+1}}
 =\exp\left(\sum_{n\geq1}
   \left(\sum_i(-1)^i\Tr(u_i^n)\right)\frac{T^n}{n}\right).
 \tag{2.6}\label{eq:alternating-det-trace}
\]
\end{corollary}

\begin{proof}
Apply \cref{prop:det-trace} in each degree.  The factor
\(\det(1-Tu_i)^{-1}\) has logarithm
\(\sum_n\Tr(u_i^n)T^n/n\).  Raising it to \((-1)^i\) and multiplying
over \(i\) gives exactly \eqref{eq:alternating-det-trace}.  Boundedness
makes the degree product finite.
\end{proof}

\begin{lemma}[Additivity through a stable filtration]
\label{lem:det-trace-filtration}
If \(u\) preserves a finite filtration of \(V\), then
\[
 \Tr(u^n\mid V)=\sum_a\Tr(u^n\mid\operatorname{gr}_aV)
\]
and
\[
 \det(1-Tu\mid V)
 =\prod_a\det(1-Tu\mid\operatorname{gr}_aV).
\]
\end{lemma}

\begin{proof}
Choose a basis adapted to the filtration.  The matrix of \(u\) is block
upper triangular, with diagonal blocks the induced maps on the graded
pieces.  The same is true of every power.  Trace adds and determinant
multiplies across those diagonal blocks.
\end{proof}

This lemma is the linear-algebra reason spectral-sequence arguments need
only a finite Frobenius-stable filtration and subquotient control; they
do not require a noncanonical direct-sum splitting.

\section{Rational series and eventual recurrences}

Let \(a=(a_0,a_1,\ldots)\) be a sequence in a field \(K\).  Write
\[
 Sa=(a_1,a_2,\ldots)
\]
for the left shift, and let
\[
 H(a)=(a_{r+s})_{r,s\geq0}
\]
be its infinite Hankel matrix.  The rank of an infinite matrix is the
supremum of the ranks of all its finite submatrices.  Equivalently, it is
the dimension of the span of its columns when that dimension is finite.

\begin{proposition}[Rationality, recurrence, shift span, and Hankel rank]
\label{prop:rational-recurrence-hankel}
For \(f(T)=\sum_{n\geq0}a_nT^n\in K[[T]]\), the following conditions are
equivalent.
\begin{enumerate}
\item \(f\) is the Taylor expansion at zero of an element of \(K(T)\)
      regular at zero.
\item There are \(r,N\geq0\) and \(b_1,\ldots,b_r\in K\) such that
      \[
       a_n+b_1a_{n-1}+\cdots+b_ra_{n-r}=0
       \qquad(n\geq N).
       \tag{2.7}\label{eq:eventual-recurrence}
      \]
\item The \(K\)-span of \(S^na\), \(n\geq0\), is finite-dimensional.
\item The Hankel matrix \(H(a)\) has finite rank over \(K\).
\end{enumerate}
\end{proposition}

\begin{proof}
Suppose \(f=A/B\), with \(A,B\in K[T]\) and \(B(0)=1\).  Write
\[
 B(T)=1+b_1T+\cdots+b_rT^r.
\]
The coefficient of \(T^n\) in \(Bf\) is
\[
 a_n+b_1a_{n-1}+\cdots+b_ra_{n-r},
\]
where negative subscripts are interpreted as zero.  Since \(Bf=A\), this
coefficient vanishes when \(n>\deg A\).  Hence (1) implies (2).

Assume (2).  Discarding the first \(N\) coordinates, the recurrence
expresses every tail after the first \(r\) tails as a linear combination
of those \(r\).  Restoring the discarded coordinates adds only the
finite-dimensional space of sequences supported in
\(\{0,\ldots,N+r-1\}\).  Therefore all shifts lie in a
finite-dimensional span, and (2) implies (3).

The \(s\)-th column of \(H(a)\) is the sequence
\[
 (a_s,a_{s+1},a_{s+2},\ldots)=S^sa.
\]
Thus the column span of the Hankel matrix is precisely the shift span.
This proves the equivalence of (3) and (4).

Finally suppose
\[
 W=\operatorname{span}_K\{S^na:n\geq0\}
\]
is finite-dimensional.  It is stable under \(S\).  If
\[
 c(t)=t^r+c_{r-1}t^{r-1}+\cdots+c_0
\]
is the characteristic polynomial of \(S\mid W\), Cayley--Hamilton gives
\(c(S)a=0\).  In the \(n\)-th coordinate,
\[
 a_{n+r}+c_{r-1}a_{n+r-1}+\cdots+c_0a_n=0
 \qquad(n\geq0).
 \tag{2.8}\label{eq:CH-recurrence}
\]
Set
\[
 C(T)=1+c_{r-1}T+\cdots+c_0T^r.
\]
For every \(m\geq r\), the coefficient of \(T^m\) in \(Cf\) is the
left side of \eqref{eq:CH-recurrence} with \(n=m-r\), and is zero.
Therefore \(Cf\) is a polynomial of degree less than \(r\), proving that
\(f\in K(T)\).  Thus (3) implies (1).
\end{proof}

\begin{remark}
The recurrence obtained from a chosen rational expression need only be
eventual because its numerator creates finitely many exceptional
coefficients.  Passing to the full shift span absorbs those exceptions,
and Cayley--Hamilton returns a recurrence valid from \(n=0\).  This
distinction is helpful in formal code: an eventual recurrence is easy to
extract from a denominator, while a global recurrence is easy to use as
an endomorphism identity.
\end{remark}

\section{Hankel descent}

\begin{lemma}[Rank is unchanged by scalar extension]
\label{lem:matrix-rank-extension}
Let \(K\subset L\) be fields and let \(M\) be a possibly infinite matrix
with entries in \(K\).  Its rank over \(K\) equals its rank over \(L\).
\end{lemma}

\begin{proof}
For each \(r\), the assertion that the rank is at least \(r\) means that
some finite \(r\times r\) minor has nonzero determinant.  Such a
determinant lies in \(K\), and it is nonzero in \(K\) if and only if its
image is nonzero in \(L\).  The sets of possible \(r\) are therefore the
same over both fields.  Taking their supremum proves the result.
\end{proof}

\begin{theorem}[Hankel rationality descent]
\label{thm:hankel-descent}
For every field extension \(K\subset L\),
\[
  K[[T]]\cap L(T)=K(T)
\]
inside \(L[[T]]\).
\end{theorem}

\begin{proof}
Let \(f=\sum a_nT^n\) have \(a_n\in K\), and suppose
\(f\in L(T)\).  By
\cref{prop:rational-recurrence-hankel}, the Hankel matrix \(H(a)\) has
finite rank over \(L\).  Its entries lie in \(K\), so
\cref{lem:matrix-rank-extension} gives the same finite rank over \(K\).
Applying \cref{prop:rational-recurrence-hankel} again, now over \(K\),
gives \(f\in K(T)\).
\end{proof}

This argument is preferable to averaging a denominator over field
automorphisms.  A sum of denominators need not remain a denominator, and
an infinite extension may have no useful finite averaging group.  Hankel
rank is canonical and is defined by the vanishing of minors with
coefficients in the smaller field.

\section{The integral Fatou lattice}

Hankel descent gives rational coefficients in a rational expression, but
the Weil conjectures require integral polynomial factors.  The next
theorem upgrades an integral rational series to an integral numerator and
denominator.

\begin{lemma}[Finite coordinates detect a finite-dimensional sequence space]
\label{lem:finite-coordinate-detection}
Let \(V\subset\mathbf Q^{\mathbf N}\) be finite-dimensional.  There is an
\(N\) such that projection to the first \(N\) coordinates
\[
 \pi_N\colon V\longrightarrow\mathbf Q^N
\]
is injective.
\end{lemma}

\begin{proof}
The kernels \(\ker\pi_N\) form a descending sequence of subspaces of the
finite-dimensional space \(V\).  Their intersection is zero, because a
sequence all of whose coordinates vanish is the zero sequence.  The
descending sequence of dimensions eventually stabilizes.  Once it
stabilizes, the subspaces themselves stabilize, and their stable value is
their intersection.  Hence that stable value is zero.
\end{proof}

\begin{theorem}[Integral Fatou lemma]
\label{thm:integral-fatou}
If
\[
 f(T)=\sum_{n\geq0}a_nT^n\in\mathbf Z[[T]]\cap\mathbf Q(T),
\]
then \(f=A/B\) for polynomials \(A,B\in\mathbf Z[T]\) with \(B(0)=1\).
If \(a_0=1\), one also has \(A(0)=1\).
\end{theorem}

\begin{proof}
Let \(a=(a_0,a_1,\ldots)\).  Rationality and
\cref{prop:rational-recurrence-hankel} show that
\[
 V=\operatorname{span}_{\mathbf Q}\{S^na:n\geq0\}
\]
is finite-dimensional.  Choose \(N\) so that \(\pi_N\) is injective on
\(V\), using \cref{lem:finite-coordinate-detection}.

Inside \(V\), form the shift-stable abelian group
\[
 M=\sum_{n\geq0}\mathbf Z\,S^na.
 \tag{2.9}\label{eq:fatou-lattice}
\]
Every generator is an integral sequence.  Hence
\(\pi_N(M)\subset\mathbf Z^N\).  A submodule of the finitely generated
\(\mathbf Z\)-module \(\mathbf Z^N\) is finitely generated because
\(\mathbf Z\) is noetherian.  Since \(\pi_N\) is injective, \(M\) is
finitely generated as well.  It is torsion-free as a subgroup of a
\(\mathbf Q\)-vector space, so it is a finite free
\(\mathbf Z\)-module.

The shift maps \(M\) to itself.  Let
\[
 c(t)=t^r+c_{r-1}t^{r-1}+\cdots+c_0\in\mathbf Z[t]
\]
be the characteristic polynomial of the integral endomorphism
\(S\mid M\).  Cayley--Hamilton holds for endomorphisms of finite free
modules, and yields \(c(S)a=0\).  Coordinatewise,
\[
 a_{n+r}+c_{r-1}a_{n+r-1}+\cdots+c_0a_n=0
 \quad(n\geq0).
\]
Define
\[
 B(T)=1+c_{r-1}T+\cdots+c_0T^r\in\mathbf Z[T].
\]
The recurrence says that \(A(T)=B(T)f(T)\) has degree less than \(r\).
Its coefficients are integral because both \(B\) and the series \(f\)
are integral.  The constant coefficients satisfy
\[
 B(0)=1,\qquad A(0)=B(0)a_0=a_0.
\]
This proves every assertion.
\end{proof}

\begin{remark}
The essential lattice is not an arbitrarily chosen lattice in the
rational shift space.  It is the group generated by all integral shifts.
Finite-coordinate detection embeds this possibly infinite-looking group
into \(\mathbf Z^N\), where noetherianity proves finite generation.
This is the step which turns a rational recurrence into an integral monic
characteristic polynomial.
\end{remark}

\section{Normalized Gauss factors}

\begin{lemma}[Normalized rational factors are integral]
\label{lem:normalized-gauss-factor}
Suppose \(C\in\mathbf Z[T]\), \(C(0)=1\), and
\[
 C(T)=P(T)Q(T)
\]
with \(P,Q\in\mathbf Q[T]\) and \(P(0)=Q(0)=1\).  Then
\(P,Q\in\mathbf Z[T]\).
\end{lemma}

\begin{proof}
The polynomial \(C\) is primitive because its constant coefficient is
\(1\).  Write
\[
 P=aP_0,\qquad Q=bQ_0,
\]
where \(a,b\in\mathbf Q^\times\) and \(P_0,Q_0\in\mathbf Z[T]\) are
primitive.  Gauss's lemma says \(P_0Q_0\) is primitive.  Since
\[
 C=abP_0Q_0
\]
and both \(C\) and \(P_0Q_0\) are primitive, comparison of contents gives
\(ab=\pm1\), after harmlessly changing the sign of a primitive factor.
The normalizations at zero say
\[
 aP_0(0)=1,\qquad bQ_0(0)=1.
\]
Thus \(P_0(0)Q_0(0)=\pm1\).  These are integers whose product is a unit,
so each is \(\pm1\).  Hence \(a,b\in\{1,-1\}\), and \(P,Q\) are integral.
\end{proof}

\begin{corollary}[Reduced integral presentation]
\label{cor:reduced-integral-presentation}
If \(f\in1+T\mathbf Z[[T]]\cap\mathbf Q(T)\), then
\[
 f=A/B
\]
for coprime \(A,B\in\mathbf Z[T]\) satisfying
\(A(0)=B(0)=1\).
\end{corollary}

\begin{proof}
\Cref{thm:integral-fatou} gives one integral presentation with both
constant terms \(1\).  Let \(G\) be the greatest common divisor in
\(\mathbf Q[T]\), normalized by \(G(0)=1\).  Then
\[
 A=GA_1,\qquad B=GB_1,
\]
where every displayed factor has constant term \(1\).  Apply
\cref{lem:normalized-gauss-factor} to each equality.  It shows that
\(G,A_1,B_1\) are integral.  Cancel \(G\); the remaining fraction is
coprime and normalized.
\end{proof}

\section{Reciprocal determinant identities}

\begin{lemma}[Inverting an automorphism]
\label{lem:inverse-determinant}
If \(u\) is an automorphism of an \(r\)-dimensional vector space, then
\[
 \det(1-Tu)=(-T)^r\det(u)\det(1-T^{-1}u^{-1})
 \tag{2.10}\label{eq:inverse-det}
\]
in \(K[T,T^{-1}]\).
\end{lemma}

\begin{proof}
The endomorphism identity
\[
 1-Tu=(-Tu)(1-T^{-1}u^{-1})
\]
is checked by multiplication.  Taking determinants gives
\eqref{eq:inverse-det}.
\end{proof}

Let \(V,W\) have dimension \(r\), let
\(\langle\, ,\,\rangle\colon V\times W\to K\) be perfect, and suppose
automorphisms \(u,v\) satisfy
\[
 \langle ux,vw\rangle=\mu\langle x,w\rangle
 \qquad(\mu\in K^\times).
 \tag{2.11}\label{eq:similitude-pair}
\]
Identifying \(W\) with \(V^\vee\), this is equivalent to
\[
 v=\mu(u^{-1})^\vee.
 \tag{2.12}\label{eq:paired-adjoint}
\]

\begin{proposition}[Reciprocal characteristic polynomials]
\label{prop:reciprocal-characteristic}
Put
\[
 P_V(T)=\det(1-Tu\mid V),\qquad
 P_W(T)=\det(1-Tv\mid W).
\]
Then
\[
 \det(v)=\frac{\mu^r}{\det(u)}
 \tag{2.13}\label{eq:paired-determinants}
\]
and
\[
 P_W(\mu^{-1}T^{-1})
 =(-1)^rT^{-r}\det(u)^{-1}P_V(T).
 \tag{2.14}\label{eq:paired-polynomials}
\]
If the eigenvalues of \(u\) are \(\alpha_1,\ldots,\alpha_r\), those of
\(v\) are \(\mu/\alpha_1,\ldots,\mu/\alpha_r\), with multiplicity.
\end{proposition}

\begin{proof}
Equation \eqref{eq:paired-adjoint} immediately gives
\eqref{eq:paired-determinants} and the eigenvalue statement.  Over an
algebraic closure,
\[
\begin{aligned}
 P_W(\mu^{-1}T^{-1})
 &=\prod_{j=1}^r
   \left(1-\frac{\mu}{\alpha_j}\mu^{-1}T^{-1}\right)\\
 &=\prod_{j=1}^r
   \left[-(\alpha_jT)^{-1}(1-\alpha_jT)\right]\\
 &=(-1)^rT^{-r}\det(u)^{-1}P_V(T).
\end{aligned}
\]
Both sides have coefficients in \(K\), so the identity descends.
\end{proof}

\section{The graded functional-equation computation}

Let \(V_0,\ldots,V_{2d}\) be finite-dimensional spaces with
automorphisms \(F_i\).  Suppose there are perfect pairings
\[
 V_i\times V_{2d-i}\longrightarrow K
\]
of common multiplier \(\mu\), and suppose the middle pairing is
\((-1)^d\)-symmetric.  Set
\[
 b_i=\dim V_i,\quad D_i=\det(F_i),\quad
 \chi=\sum_i(-1)^ib_i,
\]
and
\[
 Z_F(T)=\prod_i\det(1-TF_i)^{(-1)^{i+1}}.
\]

\begin{theorem}[Abstract functional equation]
\label{thm:abstract-functional-equation}
One has
\[
 Z_F(\mu^{-1}T^{-1})
 =(-1)^\chi T^\chi
   \left(\prod_{i=0}^{2d}D_i^{(-1)^i}\right)Z_F(T).
 \tag{2.15}\label{eq:abstract-FE-first}
\]
If \(\mu=q^d\), then
\[
 \prod_iD_i^{(-1)^i}
 =\delta^{(-1)^d}q^{d\chi/2}
 \tag{2.16}\label{eq:alt-det-sign}
\]
for a sign \(\delta\in\{1,-1\}\).  Hence
\[
 Z_F(q^{-d}T^{-1})
 =\varepsilon q^{d\chi/2}T^\chi Z_F(T),
 \quad
 \varepsilon=(-1)^\chi\delta^{(-1)^d}.
 \tag{2.17}\label{eq:abstract-FE-final}
\]
\end{theorem}

\begin{proof}
Apply \cref{prop:reciprocal-characteristic} to the pair
\((V_i,V_{2d-i})\).  Since \(i\) and \(2d-i\) have the same parity,
\[
\begin{aligned}
 Z_F(\mu^{-1}T^{-1})
 &=\prod_i
 \det(1-\mu^{-1}T^{-1}F_{2d-i})^{(-1)^{i+1}}\\
 &=\prod_i
 \left((-1)^{b_i}T^{-b_i}D_i^{-1}
       \det(1-TF_i)\right)^{(-1)^{i+1}}.
\end{aligned}
\]
The last determinant factors give \(Z_F(T)\).  The exponent of \(T\) is
\[
 -\sum_i(-1)^{i+1}b_i=\chi.
\]
The product of the signs is \((-1)^\chi\): a sign raised to exponent
\(-1\) is the same sign.  The remaining determinants contribute
\(\prod_iD_i^{(-1)^i}\).  This proves
\eqref{eq:abstract-FE-first}.

For \(i<d\), \eqref{eq:paired-determinants} says
\[
 D_iD_{2d-i}=\mu^{b_i}.
 \tag{2.18}\label{eq:det-pairs}
\]
In degree \(d\), taking determinants in the middle similitude relation
gives
\[
 D_d^2=\mu^{b_d}.
 \tag{2.19}\label{eq:middle-det-square}
\]
If \(d\) is odd, the middle pairing is alternating, so \(b_d\) is even.
If \(d\) is even, \(db_d\) is even without a restriction on \(b_d\).
Thus \(db_d/2\) is always an integer.  With \(\mu=q^d\), put
\[
 \delta=D_dq^{-db_d/2}.
\]
Equation \eqref{eq:middle-det-square} gives \(\delta^2=1\).

Pairing the determinant terms symmetrically about degree \(d\) now gives
\[
\begin{aligned}
 \prod_iD_i^{(-1)^i}
 &=\prod_{i<d}(D_iD_{2d-i})^{(-1)^i}
     D_d^{(-1)^d}\\
 &=\delta^{(-1)^d}
   q^{d\left(\sum_{i<d}(-1)^ib_i+
           \frac12(-1)^db_d\right)}.
\end{aligned}
\]
Dimension symmetry \(b_i=b_{2d-i}\) gives
\[
 \frac\chi2
 =\sum_{i<d}(-1)^ib_i+\frac12(-1)^db_d.
\]
Substitution proves \eqref{eq:alt-det-sign}, and then
\eqref{eq:abstract-FE-final}.
\end{proof}

\begin{remark}
The notation \(q^{d\chi/2}\) hides no fractional power.  If \(d\) is
even, \(d\chi\) is even.  If \(d\) is odd, the middle alternating form
forces \(b_d\), and hence \(\chi\), to be even.  The proof above
constructs the exponent as an integer before using exponential notation.
\end{remark}

\begin{remark}
For an alternating middle pairing, the Pfaffian identity improves
\(D_d^2=\mu^{b_d}\) to \(D_d=\mu^{b_d/2}\), so \(\delta=1\).  For a
symmetric middle pairing both signs can occur.  The Weil conjectures need
only the uniform conclusion \(\varepsilon\in\{1,-1\}\).
\end{remark}

\section{Formalization interfaces}

\begin{formalization}
The algebraic development should be divided into small, geometry-free
modules.
\begin{enumerate}
\item A \(T\)-adically locally finite product API should prove that
      formal logarithm exchanges Euler products and divisor sums.
\item A finite-permutation API should identify fixed points of \(F^n\)
      with orbit lengths dividing \(n\), yielding
      \cref{prop:euler-product}.
\item Determinants and traces should provide
      \cref{prop:det-trace}, its graded form, and filtration
      multiplicativity.
\item A sequence module should identify rational series, eventual
      recurrences, finite shift spans, and finite Hankel rank.  Matrix
      rank under scalar extension then gives
      \cref{thm:hankel-descent}.
\item The Fatou module should construct the integral shift lattice
      \eqref{eq:fatou-lattice}, prove finite generation by finite-coordinate
      detection, and apply Cayley--Hamilton over \(\mathbf Z\).
\item A polynomial-content module should expose the normalized factor
      lemma and reduced integral presentation separately.
\item A `PerfectSimilitudePairing` should derive adjoint, determinant,
      eigenvalue, and reciprocal-polynomial identities from one
      compatibility equation.
\item A bounded graded pairing should return
      \cref{thm:abstract-functional-equation}, including the middle-form
      parity lemma and an integer exponent for \(q\).
\end{enumerate}
These results should import no scheme theory.  The cohomological chapters
instantiate them with closed points, Frobenius modules, trace formula,
and Poincare duality.
\end{formalization}


\part{The pro-etale coefficient machine}
\label{part:proetale-machine}
\chapter{The Pro-etale Site as a Cohomological Workspace}
\label{chap:proetale-site}

The pro-etale site is not a replacement for the etale topology in the
finiteness theorems.  Its advantage is different and very concrete: inverse
limits of finite etale covers become actual objects, complete topological
coefficient rings become honest sheaves, and derived inverse limits can be
controlled internally.  This chapter isolates precisely those gains.

Throughout, schemes are locally noetherian unless explicitly stated.  The
definitions themselves make sense for arbitrary schemes.

\section{Weakly etale morphisms}

\begin{definition}[Weakly etale morphism]
\label{def:weakly-etale}
A morphism \(f:Y\to X\) is \emph{weakly etale} if \(f\) is flat and its
diagonal
\[
  \Delta_f:Y\longrightarrow Y\times_XY
\]
is flat.  It is \emph{pro-etale} in the affine sense if, locally on source
and target, it is a cofiltered inverse limit of etale affine morphisms with
affine transition maps.
\end{definition}

Every etale morphism is weakly etale.  A filtered colimit of etale algebras
is weakly etale, so inverse limits of affine etale schemes provide the basic
new examples.  Conversely, weakly etale morphisms are locally ind-etale in
the topology relevant below; this is
\cite[Theorem~2.3.4]{BhattScholzeProEtale}.

\begin{example}[A profinite set becomes geometric]
Let \(S=\varprojlim_i S_i\) be a profinite set, written as a limit of finite
sets, and let \(X\) be a scheme.  Then
\[
  X\times S:=\varprojlim_i \coprod_{s\in S_i}X
\]
is pro-etale over \(X\).  Thus a continuous action of a profinite group may
be studied on an actual cover rather than only through all of its finite
quotients.
\end{example}

\begin{definition}[Small pro-etale site]
\label{def:small-proetale-site}
The site \(X_{\proet}\) has weakly etale \(X\)-schemes as objects.  Coverings
are generated by jointly surjective fpqc families.  Equivalently, one may
use the basis of affine pro-etale objects and finite jointly surjective
families after refining.  We write
\[
  \nu_X:X_{\proet}\longrightarrow X_{\et}
\]
for the morphism of topoi induced by viewing every etale object as a
pro-etale object.
\end{definition}

The equivalence of the weakly-etale and ind-etale descriptions of the
topology is a theorem, not a definitional equality.  This distinction matters
in a formal development: morphisms, covering families, and the chosen basis
should be separate structures.

\begin{proposition}[Full faithfulness for bounded-below etale complexes]
\label{prop:proetale-etale-comparison}
For \(K\in\D^+(X_{\et})\), the adjunction unit
\[
 K\longrightarrow \R\nu_{X,*}\nu_X^*K
\]
is an equivalence.  Hence
\(\nu_X^*:\D^+(X_{\et})\to\D^+(X_{\proet})\) is fully faithful, and
\[
 \R\Gamma(X_{\et},K)
   \simeq \R\Gamma(X_{\proet},\nu_X^*K).
\]
\end{proposition}

\begin{proof}
If an affine pro-etale object has a presentation \(U=\varprojlim_iU_i\) by
affine etale objects, then
\[
 \R\Gamma(U,\nu_X^*K)
   \simeq\varinjlim_i\R\Gamma(U_i,K).
\]
Refining covers by limits of affine etale covers reduces the assertion to
Cech descent on \(X_{\et}\).  This is precisely
\cite[Corollary~5.1.6]{BhattScholzeProEtale}.  Bounded-below is essential;
\cite[Remark~5.1.8]{BhattScholzeProEtale} gives the unbounded correction by
a Postnikov limit and explains why the displayed unit can then fail.
\end{proof}

\begin{proposition}[Pushforward comparison]
\label{prop:proetale-pushforward-comparison}
Let \(f:Y\to X\) be qcqs and let \(K\in\D^+(Y_{\et})\).  The canonical
exchange morphism
\[
 \nu_X^*\R f_{\et,*}K
   \longrightarrow
 \R f_{\proet,*}\nu_Y^*K
\]
is an equivalence.
\end{proposition}

\begin{proof}
For affine \(X\) and \(Y\), evaluate on a presentation
\(U=\varprojlim_iU_i\) and use the preceding colimit formula on
\(f^{-1}U=\varprojlim_i f^{-1}U_i\).  A finite Mayer--Vietoris argument first
handles quasicompact separated \(Y\), and a second application handles the
qcqs case.  For complexes, finite limits are replaced by finite homotopy
limits.  This is \cite[Lemma~5.4.3]{BhattScholzeProEtale}.
\end{proof}

\begin{warning}
This comparison does not say that every pro-etale sheaf comes from the etale
site.  In particular, the sheaf of continuous \(\Zl\)-valued functions is not
the pullback of the constant discrete sheaf with underlying group \(\Zl\).
That difference is exactly what makes integral \(\ell\)-adic coefficients
work cleanly.
\end{warning}

\section{Weakly contractible affine objects}

\begin{definition}
A qcqs scheme \(U\) is \emph{w-local} if its underlying spectral space
satisfies both:
\begin{enumerate}
\item every open cover \((U_i\hookrightarrow U)_i\) splits, meaning that
      \(\coprod_iU_i\to U\) admits a section;
\item the subspace \(U^c\subset U\) of closed points is closed.
\end{enumerate}
Equivalently, \(U^c\) is closed and the natural map
\[
  U^c\longrightarrow\pi_0(U)
\]
is a bijection, indeed a homeomorphism onto the profinite component space.
An affine w-local scheme is \emph{w-strictly local} if every local ring at a
closed point is strictly henselian.  It is \emph{w-contractible} if every
faithfully flat weakly etale map \(V\to U\) admits a section.
\end{definition}

The equivalence in the definition is
\cite[Definition~2.1.1 and Lemma~2.1.4]{BhattScholzeProEtale}.

The last condition is the site-theoretic one: every covering of the
representable sheaf \(h_U\) splits.  It implies exactness of sections on
epimorphisms of sheaves over \(U\).

\begin{theorem}[Existence of a contractible basis]
\label{thm:proetale-wcontractible-cover}
Every affine scheme admits a faithfully flat ind-etale cover by a
w-contractible affine scheme.  Hence every scheme is locally weakly
contractible on its pro-etale site.
\end{theorem}

\begin{proof}
First pass to a w-local ind-Zariski cover.  Strictly henselize at all closed
points simultaneously, then take the inverse limit over all finite etale
covers equipped with compatible lifts at those points.  The resulting affine
scheme is w-strictly local and has no nontrivial weakly etale covering without
a section.  The construction and the final lifting argument are
\cite[Lemma~2.4.9 and Proposition~4.2.8]{BhattScholzeProEtale}; the global
form is \cite[Theorem~1.5]{BhattScholzeProEtale}.
\end{proof}

\begin{proetalebox}
On the etale site, a geometric point probes a sheaf but is not itself a large
enough neighborhood on which an infinite inverse system becomes constant.
On the pro-etale site, w-contractible affines play the role of projective
test objects.  This turns many derived questions about complete coefficients
into ordinary algebra over the ring of sections on such an affine.
\end{proetalebox}

\section{Repleteness and derived inverse limits}

\begin{definition}
A topos is \emph{replete} if the inverse limit of every countable tower of
epimorphisms is again an epimorphism.
\end{definition}

\begin{theorem}[Homological consequences of the pro-etale basis]
\label{thm:proetale-repleteness}
The pro-etale topos of a scheme is replete.  More precisely:
\begin{enumerate}
\item countable products of abelian sheaves are exact;
\item if \(\cdots\twoheadrightarrow F_2\twoheadrightarrow F_1\) is a tower
      of epimorphisms, then \(\varprojlim F_n\to F_m\) is an epimorphism and
      \(\R\varprojlim F_n\simeq\varprojlim F_n\);
\item for an arbitrary countable tower, \(\R\varprojlim\) has cohomological
      dimension at most one;
\item the unbounded derived category is left complete:
      \[
        K\xrightarrow{\sim}
        \R\varprojlim_n \tau^{\ge -n}K;
      \]
\item complexes satisfy hyperdescent on \(X_{\proet}\).
\end{enumerate}
\end{theorem}

\begin{proof}
Evaluate a tower after pulling it to a w-contractible affine cover.  Since
every epimorphism there splits locally, the elementary set-theoretic lifting
argument for a countable inverse tower applies.  Conservativity of the cover
globalizes the conclusion.  Exact products and the standard two-term
\(1-\mathrm{shift}\) resolution
\[
  \prod_n K_n \xrightarrow{\,1-\mathrm{shift}\,} \prod_n K_n
\]
give the cohomological-dimension bound.  The truncation tower then converges,
which yields left completeness and hyperdescent.  See
\cite[Propositions~3.1.9--3.1.11, 3.3.3, and 3.3.6]{BhattScholzeProEtale}.
\end{proof}

\begin{remark}
Repleteness is not the assertion that inverse limits are exact without
hypotheses.  The only possible obstruction for a countable tower is the
first derived limit, and it vanishes for epimorphic transition maps.  Later
completion arguments will use exactly this statement.
\end{remark}

\section{Honest topological coefficient sheaves}

\begin{definition}[Continuous-function sheaf]
Let \(M\) be a topological abelian group.  On an affine pro-etale object
\(U\), define
\[
  \underline M(U)=\operatorname{Cont}(|U|,M),
\]
and sheafify.  If \(M=\varprojlim_n M_n\) is profinite with finite discrete
quotients, then
\[
  \underline M \simeq \varprojlim_n \underline{M_n}
\]
as sheaves on \(X_{\proet}\).
\end{definition}

\begin{proposition}[The \(\Zl\)-sheaf]
\label{prop:proetale-zl-sheaf}
For \(\ell\) invertible on \(X\), the topological ring \(\Zl\) defines an
honest sheaf of rings
\[
  \underline{\Zl}
     =\varprojlim_n \underline{\mathbf Z/\ell^n\mathbf Z}
\]
on \(X_{\proet}\).  Its restriction to a w-contractible affine \(U\) is the
ring of continuous maps \(|U|\to\Zl\).
\end{proposition}

\begin{proof}
Continuous maps into a profinite limit are precisely compatible systems of
continuous maps into the finite discrete quotients.  Sheaf descent follows
from fpqc descent for locally constant finite functions and repleteness for
the inverse limit.  This is the construction of
\cite[Lemma~4.2.12]{BhattScholzeProEtale}.
\end{proof}

\begin{warning}[Discrete versus continuous constants]
The constant sheaf associated with the discrete group underlying \(\Zl\)
has locally constant sections.  The topological sheaf \(\underline{\Zl}\)
has continuous sections and is the inverse limit of its finite quotients.
They need not agree.  Every later occurrence of \(\Zl\) as a coefficient
ring means the latter sheaf.
\end{warning}

\section{Invariance under universal homeomorphisms}

\begin{theorem}
\label{thm:proetale-universal-homeomorphism}
If \(v:X'\to X\) is a universal homeomorphism, then pullback induces an
equivalence
\[
  v^*:\operatorname{Sh}(X_{\proet})
      \xrightarrow{\sim}
      \operatorname{Sh}(X'_{\proet}).
\]
It is compatible with the comparison to the etale sites and with the
topological coefficient sheaves above.
\end{theorem}

\begin{proof}
A weakly etale morphism descends uniquely across a universal homeomorphism,
and the underlying topological spaces, hence covering families, are
unchanged.  Thus the induced functor of sites is an equivalence.  See
\cite[Lemma~5.4.2]{BhattScholzeProEtale}.
\end{proof}

This theorem is the formal reason geometric Frobenius, although not an
isomorphism of schemes, acts invertibly on the coefficient theory developed
below.

\begin{formalization}
A reusable formalization should expose four layers:
\begin{enumerate}
\item weakly etale morphisms and their stability under composition and base
      change;
\item the pro-etale coverage and the comparison morphism \(\nu_X\);
\item w-contractible affine objects as a basis with a projectivity property;
\item replete topoi and abstract lemmas on \(\R\varprojlim\), independent of
      algebraic geometry.
\end{enumerate}
The last layer belongs in general homological algebra.  The coefficient
chapters should import it, not reprove it for \(\ell\)-adic sheaves.
\end{formalization}

\begin{auditbox}
Nothing in this chapter proves constructibility, finite cohomological
dimension, proper base change, or purity.  The pro-etale site provides the
correct arena for limits and continuous coefficients.  The hard geometric
theorems remain inputs and will be identified explicitly in
\cref{chap:proetale-six}.
\end{auditbox}


\chapter{Complete Coefficients on the Pro-etale Site}
\label{chap:proetale-coefficients}

Fix a prime \(\ell\) invertible on \(X\).  The basic coefficient rings will
be \(\Lambda=\Zl\) and, more generally, the ring of integers
\(\Lambda=\mathcal O_E\) in a finite extension \(E/\Ql\).  Write
\(\mathfrak m=(\varpi)\) and \(\Lambda_n=\Lambda/\mathfrak m^n\).  These are
complete discrete valuation rings, so the ideal of definition is regular.
Regularity permits the unbounded system comparison in
\cite[Lemma~3.5.7(5)]{BhattScholzeProEtale}.  Repleteness gives
\(\R\varprojlim\) cohomological dimension at most one, while surjective
reduction towers satisfy the Mittag--Leffler condition and kill the relevant
\(\varprojlim^1\).  No exactness claim is made for arbitrary inverse systems.

\section{Derived completion}

Let \(\underline\Lambda\) denote the continuous coefficient sheaf of
\cref{prop:proetale-zl-sheaf}, with the identical construction for
\(\mathcal O_E\).

\begin{definition}[Derived complete complex]
\label{def:derived-complete-proetale}
For \(K\in\D(X_{\proet},\underline\Lambda)\), put
\[
  K_n=K\mathbin{\mathop\otimes\limits^{\mathbf L}_{\Lambda}}\Lambda_n.
\]
The complex \(K\) is \emph{derived \(\mathfrak m\)-complete} if the canonical
map
\begin{equation}
\label{eq:derived-completion-map}
 K\longrightarrow
 \Rlim_n K_n
\end{equation}
is an equivalence.  Its derived completion is
\[
  \widehat K:=\Rlim_n
  \bigl(K\mathbin{\mathop\otimes\limits^{\mathbf L}_{\Lambda}}\Lambda_n\bigr).
\]
Write \(\Dcomp(X,\Lambda)\) for the full subcategory of complete objects.
\end{definition}

\begin{proposition}[Completion is a reflector]
\label{prop:completion-reflector}
The inclusion
\[
  \Dcomp(X,\Lambda)\hookrightarrow
  \D(X_{\proet},\underline\Lambda)
\]
admits \(K\mapsto\widehat K\) as a left adjoint.  For complete \(K,L\), define
\[
 K\widehat\otimes^{\mathbf L}_{\Lambda}L
   :=\widehat{K\otimes^{\mathbf L}_{\Lambda}L},
 \qquad
 \R\underline{\operatorname{Hom}}_{\Lambda}(K,L)
\]
with the internal Hom taken in complete objects.  These operations make
\(\Dcomp(X,\Lambda)\) a closed symmetric monoidal stable category.
\end{proposition}

\begin{proof}
For a finitely generated ideal in a replete topos, the abstract completion
reflector is \cite[Lemma~3.4.9]{BhattScholzeProEtale}.  In the noetherian
constant-ring case, \cite[Proposition~3.5.1]{BhattScholzeProEtale} identifies
that reflector with \(K\mapsto\Rlim_n(K\otimes^{\mathbf L}\Lambda_n)\) and
proves the adjunction.  The completed coefficient sheaf and its complete
module category are \cite[Definition~3.5.2]{BhattScholzeProEtale}.
Completed tensor is \cite[Definition~3.4.10]{BhattScholzeProEtale}, while
closedness and tensor--Hom adjunction are
\cite[Lemma~3.4.11]{BhattScholzeProEtale}.  Finally,
\cite[Lemma~3.5.7]{BhattScholzeProEtale} supplies the equivalence with
compatible systems used in the next section; its unbounded part uses that
\((\varpi)\) is a regular ideal.
\end{proof}

\section{Complete objects and compatible torsion systems}

\begin{definition}[Compatible system]
A compatible \(\Lambda_n\)-system is a sequence
\[
 K_n\in\D(X_{\proet},\Lambda_n)
\]
equipped with coherent equivalences
\[
 K_{n+1}\otimes^{\mathbf L}_{\Lambda_{n+1}}\Lambda_n
   \simeq K_n.
\]
Morphisms must respect all transition equivalences and their higher
coherences.
\end{definition}

\begin{theorem}[Completion--system equivalence]
\label{thm:proetale-complete-system-equivalence}
Reduction modulo \(\mathfrak m^n\) induces an equivalence of stable
categories
\begin{equation}
\label{eq:completion-system-equivalence}
 \Dcomp(X,\Lambda)
   \xrightarrow{\sim}
 \varprojlim_n \D(X_{\proet},\Lambda_n).
\end{equation}
Its inverse sends \((K_n)_n\) to \(\Rlim_n K_n\).  The statement remains
valid in the unbounded derived category for \(\Zl\) and \(\mathcal O_E\).
\end{theorem}

\begin{proof}
For complete \(K\), \cref{eq:derived-completion-map} identifies \(K\) with
the inverse limit of its reductions, and the reduction identities supply a
compatible system.  Conversely, let \(K=\Rlim K_n\).  Exactness of products
and the cohomological-dimension-one property of \(\Rlim\) allow derived
tensor with \(\Lambda_m\) to pass through the tower; compatibility makes the
result \(K_m\).  Thus \(K\) is complete and the two constructions are inverse.
The regularity of \((\varpi)\) supplies the unbounded assertion.  This is
\cite[Lemma~3.5.7]{BhattScholzeProEtale}; a coefficient-system formulation is
\cite[Proposition~5.1]{HRSConstructible}.
\end{proof}

\begin{proetalebox}
Traditional adic formalisms begin with a tower
\((K_n)_n\) and then quotient by essentially zero systems.  On the pro-etale
site the tower has a genuine inverse limit, and
\cref{thm:proetale-complete-system-equivalence} says that no information is
lost by replacing the tower with one complete object.  This is the first
major simplification used later.
\end{proetalebox}

\section{Lisse and constructible complexes}

\begin{definition}[Lisse and constructible]
\label{def:proetale-lisse-constructible}
A complete complex \(K\) is \emph{lisse} if it is dualizable for
\(\widehat\otimes^{\mathbf L}_{\Lambda}\).  It is \emph{constructible} if
there is a finite stratification
\[
  X=\coprod_{a=1}^r X_a
\]
by locally closed constructible subschemes such that \(K|_{X_a}\) is lisse
for every \(a\).  We denote the corresponding categories by
\(\D_{\lis}(X,\Lambda)\) and \(\Dcons(X,\Lambda)\).
\end{definition}

For noetherian \(X\), this agrees with the Bhatt--Scholze definition:
\(K\) is complete and \(K\otimes^{\mathbf L}_{\Lambda}\Lambda_1\) has bounded
constructible cohomology.  The equivalence is
\cite[Definition~6.5.1, Lemma~6.5.3, and Proposition~6.6.11]
{BhattScholzeProEtale}.

\begin{theorem}[Local perfect-complex criterion]
\label{thm:proetale-constructible-characterization}
Let \(X\) be locally noetherian.
\begin{enumerate}
\item A complete complex is lisse if and only if pro-etale locally it is the
      constant sheaf attached to a perfect \(\Lambda\)-complex.
\item On a w-contractible affine \(U\), the section functor identifies lisse
      objects with perfect complexes over
      \(\Gamma(U,\underline\Lambda)\).
\item Lisse and constructible complexes satisfy pro-etale hyperdescent.
\item If \(X\) is also qcqs, constructible complexes are bounded.  Completed
      tensor products and internal Homs preserve constructibility.
\end{enumerate}
\end{theorem}

\begin{proof}
On a w-contractible affine the section equivalence is
\cite[Lemma~4.1]{HRSConstructible}, and hyperdescent is
\cite[Corollary~4.7]{HRSConstructible}.  The converse from dualizable to
pro-etale-locally perfect-constant requires that \(X\) have locally finitely
many irreducible components; this is
\cite[Theorem~4.13]{HRSConstructible}.  A locally noetherian scheme satisfies
that hypothesis.  Under qcqs, finitely many affine pieces and their finite
stratifications give a uniform cohomological bound.  Tensor and internal Hom
reduce on each stratum to perfect-complex algebra and then glue through
localization triangles.  Compare
\cite[Lemmas~6.5.3 and 6.5.5]{BhattScholzeProEtale}.
\end{proof}

\begin{example}[The point]
\label{ex:constructible-point}
For a separably closed field \(k\),
\[
  \Dcons(\Spec k,\Lambda)\simeq\operatorname{Perf}(\Lambda).
\]
For a general field, the right-hand side is replaced by perfect
\(\Lambda\)-complexes with continuous action of its absolute Galois group.
The first assertion is \cite[Example~6.5.2]{BhattScholzeProEtale}.
\end{example}

\section{The standard t-structure and a finite-level trap}

\begin{theorem}[Adic t-structure]
\label{thm:proetale-adic-tstructure}
Let \(\Lambda\) be t-admissible: its underlying ring \(\Lambda_*\) is regular
coherent, and
\[
 \Lambda_*\longrightarrow\Gamma(S,\Lambda)
\]
is flat for every extremally disconnected profinite set \(S\).
\begin{enumerate}
\item If \(X\) has Zariski-locally finitely many irreducible components, the
      ambient standard t-structure restricts to
      \(\D_{\lis}(X,\Lambda)\).
\item If every constructible subset of \(X\) has Zariski-locally finitely
      many irreducible components, it restricts to
      \(\Dcons(X,\Lambda)\).
\end{enumerate}
Both geometric conditions hold when \(X\) is locally noetherian.
The topological rings \(\Zl\) and \(\mathcal O_E\) are t-admissible.  Hence,
under the second geometric condition,
\[
  \cH^i(K)\in\Dcons(X,\Lambda)
  \qquad(K\in\Dcons(X,\Lambda)).
\]
\end{theorem}

\begin{proof}
By \cite[Theorem~4.13]{HRSConstructible}, lisse objects are locally
perfect-constant under the first geometric hypothesis.  Regular coherence of
\(\Lambda_*\) keeps truncations of the perfect coefficient complex perfect.
The second part of t-admissibility is equally necessary: flatness of
\(\Lambda_*\to\Gamma(S,\Lambda)\) makes truncation commute with evaluation on
w-contractible affines, whose component spaces are extremally disconnected.
The constructible assertion follows by applying this calculation on a
stratification; the stated condition on every constructible subset ensures
that the calculation remains available on all strata.  This is
\cite[Theorem~6.2]{HRSConstructible}.  Complete discrete valuation rings and
their fraction fields are covered by
\cite[Lemma~6.10 and Corollary~6.11]{HRSConstructible}.
\end{proof}

\begin{warning}[Do not take the inverse limit of finite-level hearts]
For \(n>1\), the ring \(\mathbf Z/\ell^n\mathbf Z\) is not regular.  A truncation of a
perfect complex over it need not be perfect; for example the residue module
\(\Fl\) has infinite projective dimension.  Hence
\(\operatorname{Perf}(\mathbf Z/\ell^n\mathbf Z)\) is not preserved by the standard truncation functors.
The integral t-structure must be constructed on the complete
\(\Zl\)-category and cannot be defined naively as the inverse limit of
finite-level hearts.  This is the warning in
\cite[Section~6]{HRSConstructible}.
\end{warning}

\section{The coefficient tower}

Choose a finite extension \(E/\Ql\), with ring of integers \(\mathcal O_E\).

\begin{theorem}[Integral, rational, and algebraically closed coefficients]
\label{thm:proetale-coefficient-tower}
The constructible categories are connected by the following operations.
\begin{enumerate}
\item Derived completed scalar extension gives
\[
  -\widehat\otimes^{\mathbf L}_{\Zl}\mathcal O_E:
  \Dcons(X,\Zl)\longrightarrow\Dcons(X,\mathcal O_E).
\]
This is a change of coefficients, not an equivalence.
\item Inverting \(\varpi\) induces an equivalence
\begin{equation}
\label{eq:oe-to-e}
  \Dcons(X,\mathcal O_E)[1/\ell]
     \xrightarrow{\sim}\Dcons(X,E).
\end{equation}
\item The algebraic-closure coefficient category is the filtered colimit
\begin{equation}
\label{eq:qbarl-colimit}
  \Dcons(X,\Qbarl)
    :=\varinjlim_{E/\Ql\,, [E:\Ql]<\infty}\Dcons(X,E),
\end{equation}
with transition functors given by scalar extension.
\end{enumerate}
Every constructible \(E\)-complex admits an \(\mathcal O_E\)-model after
possibly enlarging \(E\); such a lattice is not canonical.
\end{theorem}

\begin{proof}
Completed scalar extension may be checked after reduction modulo powers of
the maximal ideals, where it preserves perfectness on strata.  For the second
assertion, an integral object killed after inverting \(\ell\) is precisely
\(\ell\)-power torsion, and every rational constructible object admits a
lattice locally on a finite stratification.  Gluing by localization triangles
gives essential surjectivity; clearing denominators gives full faithfulness.
These are \cite[Proposition~6.8.4, Lemma~6.8.13, and
Proposition~6.8.14]{BhattScholzeProEtale}.  The colimit description is
\cite[Lemma~7.9]{HRSConstructible}.
\end{proof}

\begin{remark}[What \(\Qbarl\) means]
Equation~\eqref{eq:qbarl-colimit} is deliberately categorical.  It avoids
treating the non-discrete, non-locally compact field \(\Qbarl\) as though it
were another complete discrete valuation ring.  Every individual
constructible object and every morphism is defined over some finite
extension \(E/\Ql\).
\end{remark}

\begin{theorem}[Comparison with the classical adic category]
\label{thm:proetale-classical-adic-comparison}
For a noetherian scheme on which \(\ell\) is invertible, the constructible
pro-etale categories with \(\mathcal O_E\)- and \(E\)-coefficients agree with
the classical bounded constructible adic derived categories.  Under this
equivalence, reduction modulo \(\mathfrak m^n\), tensor products, pullback,
and cohomology agree.
\end{theorem}

\begin{proof}
Apply \cref{thm:proetale-complete-system-equivalence} to replace a complete
object by its compatible torsion tower.  The comparison of discrete
coefficients in \cref{prop:proetale-etale-comparison} identifies every level
with the classical etale category.  The adic normalization condition is
exactly the derived reduction compatibility of the tower.  The complete
comparison, including rational localization, is
\cite[Theorem~7.7]{HRSConstructible}.
\end{proof}

\begin{formalization}
The reusable API should distinguish:
\[
\begin{array}{c}
\text{all pro-etale complexes}\\
\downarrow;\text{solidification}\\
\text{solid ambient complexes}\\
\supset \text{ind-constructible complexes}\\
\supset \text{constructible complexes}\\
\supset \text{lisse complexes}.
\end{array}
\]
Derived completion remains a separate reflector on the integral coefficient
side.  The system comparison is an equivalence of stable categories, not an
equality of objects.  Solidification, completion, scalar extension, rational
coefficients, and choice of a lattice must be separate constructions; no
theorem may depend on a definitional choice of lattice.
\end{formalization}

\begin{auditbox}
This chapter constructs the coefficient categories and proves their
comparison with traditional adic systems.  It does not prove that geometric
direct images preserve constructibility.  The next chapter first constructs
the operations on the solid ambient category and then proves their
constructible restriction.  The latter still depends on Grothendieck--Gabber
finiteness and constructible duality.
\end{auditbox}


\chapter{Route B: Ambient Solid Six Operations and Constructible Descent}
\label{chap:proetale-six}

Let \(P\) be a set of primes and put

\[
  \widehat{\mathbf Z}_{P}:=\prod_{r\in P}\mathbf Z_r.
\]

For the Weil conjectures we take \(P=\{\ell\}\), so that
\(\widehat{\mathbf Z}_{P}=\Zl\), and every scheme is a qcqs
\(\mathbf Z[1/\ell]\)-scheme.  This chapter adopts \emph{Route B}:

\[
\boxed{
 \begin{gathered}
 \text{one presentable ambient category of pro-etale sheaves}\\
 \Downarrow\\[-0.4em]
 \text{a coherent six-functor formalism}\\
 \Downarrow\\[-0.4em]
 \text{restriction to the constructible subcategory}.
 \end{gathered}}
\tag{4.B}\label{eq:route-b-overview}
\]

No operation is \emph{defined} by first constructing a compatible family
modulo \(\ell^n\) and then taking its inverse limit.  There is nevertheless
an essential correction to the phrase ``all pro-etale sheaves'': the naive
unrestricted derived category is too large.  The correct ambient category is
the reflective localization of \emph{solid} pro-etale complexes.  Solidity
is an ambient continuity condition, not a constructibility or finite
stratification condition.

\section{The unrestricted category is impossible}

Write
\[
 \Dnaive(X,\Lambda)
   :=\D\bigl(\operatorname{Sh}(X_{\proet},\Lambda)\bigr)
\]
for the unbounded derived category of all pro-etale \(\Lambda\)-module
sheaves (or its hypersheaf enhancement).

\begin{theorem}[Failure of unrestricted proper base change]
\label{thm:naive-proetale-pbc-fails}
Let \(r\) be a prime.  In the cartesian square
\begin{equation}
\label{eq:naive-pbc-counterexample}
\begin{tikzcd}
 \PP^1_{\mathbf C} \arrow[r,"q"] \arrow[d,"g"'] &
 \PP^1_{\mathbf Q} \arrow[d,"f"] \\
 \Spec(\mathbf C) \arrow[r,"p"'] & \Spec(\mathbf Q),
\end{tikzcd}
\end{equation}
the proper base-change transformation
\[
 p^*f_*\longrightarrow g_*q^*
\]
is not an equivalence on
\(\Dnaive(\PP^1_{\mathbf Q},\mathbf F_r)\).
\end{theorem}

\begin{proof}
This is \cite[Remark~7.15]{HaineEtAlCondensedHomotopy}.  We recall the
contradiction because it identifies exactly what goes wrong.  If the displayed
map were an equivalence, then its mate under the left adjoints would be an
equivalence.  Since pullback from \(\Spec(\mathbf Q)\) to
\(\Spec(\mathbf C)\) has the expected left adjoint on these pro-etale
hypersheaf categories, applying the mate to the unit and then taking relative
homology would identify the \(\mathbf F_r\)-homology objects of the condensed
homotopy types of \(\PP^1_{\mathbf C}\) and \(\PP^1_{\mathbf Q}\).

The first homology of the complex-side object has a quotient of cardinality
\(2^{2^{\aleph_0}}\): it receives the quotient of a product of copies of
\(\mathbf F_r\), indexed by \(\mathbf C\), by the corresponding direct sum.
The rational-side first homology is a quotient of a free profinite group on a
countable set and has cardinality at most \(2^{\aleph_0}\).  The two homology
objects therefore cannot be isomorphic.
\end{proof}

\begin{corollary}
There is no Grothendieck six-functor formalism on
\(X\mapsto\Dnaive(X,\Zl)\) having the usual pro-etale pullback and direct
image, arbitrary proper base change, and \(p_!\simeq p_*\) for proper maps.
The failure is not repaired by restricting merely to derived
\(\ell\)-complete objects.
\end{corollary}

\begin{proof}
Such a formalism would imply the proper base-change statement contradicted by
\cref{thm:naive-proetale-pbc-fails}.  Moreover \(\mathbf F_\ell\) is already
derived \(\ell\)-complete, so the same counterexample belongs to the proposed
complete subcategory.
\end{proof}

\begin{auditbox}[title={Why the ambient category must change}]
The obstruction is not that \(\Zl\) fails to be an honest pro-etale sheaf,
nor that completion has been omitted.  It is that arbitrary pro-etale
hypersheaves retain condensed-homotopical variation not controlled by
ordinary etale geometry.  Proper pushforward does not satisfy the base-change
axiom on all of them.  Thus ``all derived-complete pro-etale sheaves'' is not
a valid Route B coefficient system.
\end{auditbox}

\section{Solid pro-etale complexes}

Let \(X\) be qcqs.  If
\(j:U=\varprojlim_\alpha U_\alpha\to X\) is pro-etale with every
\(U_\alpha\to X\) affine etale, set
\[
 j^{\blacksquare}_{\sharp}\widehat{\mathbf Z}_{P}
   :=\varprojlim_\alpha
      (j_\alpha)_\sharp\widehat{\mathbf Z}_{P}.
\]
The solidification maps compare this continuous object with the free
pro-etale sheaf represented by \(U\).

\begin{definition}[Solid ambient category]
\label{def:solid-ambient-category}
A pro-etale sheaf \(F\) of \(\widehat{\mathbf Z}_{P}\)-modules is
\emph{solid} if, for every \(j\) as above, the natural map
\[
 \operatorname{Hom}
  \bigl(j^{\blacksquare}_{\sharp}\widehat{\mathbf Z}_{P},F\bigr)
   \longrightarrow F(U)
\]
is an equivalence.  Write
\[
 \operatorname{Sh}(X,\widehat{\mathbf Z}_{P})^{\blacksquare}
  \subset
 \operatorname{Sh}(X_{\proet},\widehat{\mathbf Z}_{P})
\]
for the full subcategory of solid sheaves.  A complex is solid if all its
cohomology sheaves are solid, and
\[
 \Dsolid(X,\widehat{\mathbf Z}_{P})
  \subset
 \Dnaive(X,\widehat{\mathbf Z}_{P})
\]
denotes the resulting full stable subcategory.  If \(\Lambda\) is a solid
\(\widehat{\mathbf Z}_{P}\)-algebra, define
\[
 \Dsolid(X,\Lambda)
   :=\operatorname{Mod}_{\Lambda}
      \bigl(\Dsolid(X,\widehat{\mathbf Z}_{P})\bigr).
\]
\end{definition}

\begin{theorem}[Solidification]
\label{thm:solidification-properties}
The inclusion
\[
 \Dsolid(X,\Lambda)\hookrightarrow\Dnaive(X,\Lambda)
\]
has a colimit-preserving left adjoint \((-)^{\blacksquare}\).  The solid
category is stable, presentable, closed symmetric monoidal, and compatible
with pullback.  It contains the pro-category generated by constructible
etale \(P\)-torsion sheaves, but it is much larger than the bounded
constructible category.
\end{theorem}

\begin{proof}
At the abelian level, solidity is precisely the reflective localization at
the solidification maps above.  The localized subcategory is closed under
extensions, all limits, and all colimits; compact generators are obtained
from cofiltered limits of constructible etale \(P\)-torsion sheaves.  Passing
to complexes gives the stated presentable stable category, and module objects
give general solid coefficients.  This is
\cite[Theorem~1.2.8 and Corollary~1.3.2]
{RuimyTubachWolfSolid}.
\end{proof}

\begin{warning}
``Solid'' is neither synonymous with ``constructible'' nor with ``derived
\(\ell\)-complete''.  For example, the topological field \(\Ql\) is a solid
\(\Zl\)-algebra.  Conversely, merely imposing derived completion does not
remove the counterexample of \cref{thm:naive-proetale-pbc-fails}.
\end{warning}

\begin{proetalebox}[title={What Route B gains}]
The coefficient object \(\Zl\), completed tensor product, scalar extension
to \(\Ql\), and all ambient adjunctions live in one presentable category.
No Artin--Rees quotient and no compatible-system object occur in the
definition.  Solidification is the additional continuity condition required
by proper geometry.
\end{proetalebox}

\section{The ambient six operations}

\begin{theorem}[Solid proper and smooth base change]
\label{thm:solid-ambient-base-change}
Consider a cartesian square of qcqs schemes
\begin{equation}
\label{eq:solid-base-change-square}
\begin{tikzcd}
X' \arrow[r,"g"] \arrow[d,"q"'] & X \arrow[d,"p"] \\
Y' \arrow[r,"f"'] & Y.
\end{tikzcd}
\end{equation}
For a solid coefficient algebra \(\Lambda\), the canonical map
\[
 f^*p_*K\longrightarrow q_*g^*K
\]
is an equivalence in either of the following cases:
\begin{enumerate}
\item \(p\) is proper;
\item \(f\) is smooth and all primes in \(P\) are invertible on \(Y\).
\end{enumerate}
\end{theorem}

\begin{proof}
We give the published proof because it answers exactly whether finite
coefficients are still needed.  The forgetful functor to
\(\widehat{\mathbf Z}_{P}\)-coefficients is conservative and commutes with
pullback and direct image, so one first reduces to that coefficient ring.
The relevant functors commute with limits and the solid standard
\(t\)-structure is left complete.  Postnikov truncation therefore reduces
to bounded-below objects.  Uniform cohomological-dimension bounds and
commutation with filtered colimits of uniformly bounded-below objects reduce
further to bounded objects, then devissage reduces to the heart.

The solid heart is generated under filtered colimits by torsion
constructible etale sheaves.  On those generators the assertion is the
ordinary torsion proper or smooth base-change theorem.  Colimit preservation
extends the equivalence to every solid object.  This is
\cite[Proposition~1.3.6]{RuimyTubachWolfSolid}.
\end{proof}

\begin{foundation}[title={Finite coefficients: definition versus proof}]
The preceding theorem is an ambient theorem: it is not defined as the limit
of the maps for \(\mathbf Z/\ell^n\).  Nevertheless, its known proof uses
torsion constructible generators and classical torsion base change.  This
cannot disappear logically: substituting \(K=\mathbf Z/\ell^n\) into the
ambient theorem recovers finite-coefficient base change.  Route B removes
the layer-by-layer \emph{construction}; it does not remove the mathematical
content of the finite theorem.
\end{foundation}

The full ambient formalism is obtained by solid rigidity.

\begin{theorem}[Ambient solid six-functor formalism]
\label{thm:ambient-solid-six-functors}
Assume all primes in \(P\) are invertible.  The assignment
\[
 X\longmapsto\Dsolid(X,\Lambda)
\]
on qcqs schemes has the following structure.
\begin{enumerate}
\item Every morphism has an adjunction \(f^*\dashv f_*\).
\item Every finitely presented morphism between qcqs schemes has an
      adjunction \(f_!\dashv f^!\).
\item If \(f\) is proper and finitely presented, then \(f_!\simeq f_*\).
\item Exceptional direct image satisfies arbitrary base change and the
      projection formula
      \[
       f_!(K\otimes f^*L)\simeq f_!K\otimes L.
      \]
\item Each fiber is closed symmetric monoidal, with tensor product and
      internal Hom.
\item Open--closed decompositions give both localization triangles.
\item The operations assemble into a lax symmetric monoidal functor on a
      category of correspondences, so all composition and base-change
      compatibilities are coherent.
\item If \(f\) is smooth, then
      \[
       f^!\simeq\Sigma_{\Omega_f}f^*.
      \]
      For a smooth map of pure relative dimension \(d\), the etale
      orientation identifies this with \(f^*(d)[2d]\).
\end{enumerate}
\end{theorem}

\begin{proof}
Ruimy--Tubach--Wolf first construct the coefficient system of solid
pro-etale motives.  Homotopy invariance, Nisnevich descent, localization,
smooth ambidexterity, and Tate stability make it a motivic coefficient
system.  The abstract correspondence theorem then supplies its six
operations; in particular, exceptional direct image is a colimit-preserving
functor and hence has the right adjoint \(f^!\).

The solid rigidity functor
\[
 \rho^{\blacksquare}_{\sharp}:
 \Dsolid(X,\Lambda)\longrightarrow
 \operatorname{DM}^{\mathrm{eff}}(X,\Lambda)^{\blacksquare}
\]
is an equivalence when all primes of \(P\) are invertible.  Full faithfulness
is proved from the solid generators, torsion rigidity, and compactness of the
Tate comparison object; essential surjectivity follows from the corresponding
motivic generators.  Transporting the motivic formalism across this
equivalence gives the theorem.  See
\cite[Theorems~2.5.4 and~3.3.7]{RuimyTubachWolfSolid}.
\end{proof}

\begin{warning}[Preprint status]
The solid rigidity theorem used here is a 2026 preprint.  It is presently the
most direct published Route B model on the pro-etale site and is especially
useful for a formalization blueprint, but its preprint status must remain
visible in the dependency ledger.  Scholze's profinitely complete
etale-valued formalism \cite[Theorems~7.15--7.16]{ScholzeSixFunctors}
provides a closely related established ambient alternative.
\end{warning}

\subsection{Compactification inside the ambient category}

Let \(f:X\to Y\) be separated and finitely presented.  Choose a Nagata
compactification
\[
 X\xrightarrow{j}\overline X\xrightarrow{p}Y
\]
with \(j\) a qcqs open immersion and \(p\) proper.

\begin{definition}[Exceptional direct image]
\label{def:proetale-fshriek}
In the ambient solid category set
\begin{equation}
\label{eq:fshriek-compactification}
 f_!K:=p_*j_!K.
\end{equation}
\end{definition}

\begin{theorem}[Independence and coherence]
\label{thm:fshriek-independent}
Definition~\eqref{eq:fshriek-compactification} is independent, through a
contractible space of coherent equivalences, of the compactification.  It
agrees with the functor in
\cref{thm:ambient-solid-six-functors}, is functorial in \(f\), satisfies
arbitrary base change, and satisfies the projection formula.
\end{theorem}

\begin{proof}
Two compactifications are dominated by the closure of the diagonal image in
their fiber product.  Proper base change compares the two functors
\(p_*j_!\).  Rather than checking an unbounded list of cocycles by hand, the
correspondence-level construction in
\cite[Theorem~2.5.4]{RuimyTubachWolfSolid} simultaneously supplies all higher
compatibilities.  This is the categorical content behind independence of
compactification.
\end{proof}

\begin{warning}
For a general \(f\), this \(f_!\) is not obtained by deriving an abelian-level
pro-etale functor.  It is a derived geometric operation produced by open
extension, proper pushforward, and coherent compactification.
\end{warning}

\section{The constructible subcategory inside the ambient theory}

We now specialize to \(P=\{\ell\}\), and let
\(\Lambda=\Zl\) or \(\mathcal O_E\).  Recall the intrinsic definition from
\cref{chap:proetale-coefficients}: lisse means tensor-dualizable, and
constructible means lisse after restriction to each member of a finite
constructible stratification.

\begin{theorem}[Constructibles are solid]
\label{thm:constructibles-embed-solid}
For every qcqs \(X\), there are full inclusions
\[
 \Dcons(X,\Zl)\hookrightarrow\Dsolid(X,\Zl),
 \qquad
 \Dcons(X,\Ql)\hookrightarrow\Dsolid(X,\Ql).
\]
The same holds over \(\mathcal O_E\) and \(E\) after scalar extension.
\end{theorem}

\begin{proof}
If \(K\) is lisse, dualizability shows that tensoring with \(K\) commutes
with limits.  Writing \(\Zl=\varprojlim_n\mathbf Z/\ell^n\), one obtains
\(K\simeq\varprojlim_nK/\ell^n\); the reductions are bounded torsion etale
complexes and hence solid.  Solidity is local and stable under the
open--closed gluing used by a finite stratification, proving the integral
claim.  A constructible \(\Ql\)-complex has a constructible \(\Zl\)-lattice,
so the rational claim follows by scalar extension.  This is
\cite[Proposition~1.3.11 and Corollary~1.3.12]
{RuimyTubachWolfSolid}.
\end{proof}

The full solid category has its own compact generators, and they must not be
confused with constructible complexes.  The correct compact recognition
takes place in the ind-constructible intermediate category.

\begin{definition}[Ind-constructible objects]
\label{def:indconstructible}
Let \(\Dindcons(X,\Lambda)\) be the full subcategory of pro-etale complexes
which are filtered colimits of objects of \(\Dcons(X,\Lambda)\).  Since the
solid category is closed under colimits and contains constructibles,
\[
 \Dcons(X,\Lambda)
  \subset\Dindcons(X,\Lambda)
  \subset\Dsolid(X,\Lambda).
\]
\end{definition}

\begin{theorem}[Compact recognition]
\label{thm:indconstructible-compact-recognition}
Suppose \(X\) has finite \(\Lambda\)-cohomological dimension in the sense of
Hemo--Richarz--Scholbach.  Then
\[
 \Dindcons(X,\Lambda)^\omega=\Dcons(X,\Lambda),
 \qquad
 \operatorname{Ind}\bigl(\Dcons(X,\Lambda)\bigr)
   \simeq\Dindcons(X,\Lambda).
\]
Finite-type schemes over finite or separably closed fields satisfy this
hypothesis for \(\Zl,\mathcal O_E,E\), and \(\Ql\).
\end{theorem}

\begin{proof}
For a constructible \(M\), the functor
\(\R\underline{\operatorname{Hom}}(M,-)\), followed by global sections,
commutes with filtered colimits of constructible terms.  The point is a
uniform cohomological-dimension bound; it removes the uniform lower-bound
hypothesis otherwise needed for direct image.  Thus every constructible
object is compact in \(\Dindcons\).

Conversely, write a compact \(M\in\Dindcons\) as a filtered colimit
\(M\simeq\varinjlim_iM_i\) of constructible objects.  Compactness makes
\(\id_M\) factor through some \(M_i\), so \(M\) is a retract of \(M_i\).
The constructible category is idempotent complete; hence \(M\) is
constructible.  The universal property of Ind-completion now gives the
second equivalence.  This is
\cite[Proposition~8.2 and Corollary~8.3]{HRSConstructible}.
\end{proof}

\begin{warning}
The theorem says
\(\Dcons=\Dindcons^\omega\), not
\(\Dcons=(\Dsolid)^\omega\).  The latter assertion is neither used nor made.
Solid compact generators include pro-torsion objects of a different nature.
\end{warning}

\section{The formal half of constructible descent}

Three operations restrict without a geometric finiteness theorem.

\begin{proposition}[Formal preservation]
\label{prop:formal-constructible-preservation}
Let \(f:X\to Y\) be a morphism of noetherian schemes and let
\(j:U\hookrightarrow X\) be a quasi-compact open immersion.
\begin{enumerate}
\item \(f^*\) preserves constructible complexes.
\item Completed tensor and internal Hom preserve pairs of constructible
      complexes.
\item \(j_!\) preserves constructible complexes.
\end{enumerate}
\end{proposition}

\begin{proof}
Pullback is strong symmetric monoidal and therefore preserves dualizable
objects.  Pulling back a finite stratification proves the first assertion.
For the second, choose a common refinement of the two stratifications.  On
each stratum \(K\) and \(L\) are dualizable, so
\[
 \R\underline{\operatorname{Hom}}(K,L)
   \simeq K^\vee\widehat\otimes L,
\]
which is again dualizable.  For the last assertion, adjoin the closed
complement to a finite stratification of \(U\).  The restriction of \(j_!K\)
to a stratum in \(U\) is the corresponding restriction of \(K\), while its
restriction to the complement is zero.  Equivalently, use
\cite[Corollary~4.6]{HRSConstructible}.
\end{proof}

\subsection{Completed pullback is now a restriction theorem}

\begin{definition}[Completed pullback]
\label{def:completed-pullback}
For a derived-complete constructible object \(K\), write
\begin{equation}
\label{eq:completed-pullback}
 \Lhat f^*K:=\bigl(\mathbf Lf^*K\bigr)^\wedge.
\end{equation}
\end{definition}

\begin{proposition}[Comparison with ambient pullback]
\label{prop:completed-pullback-adjunction}
Under the inclusion of
\cref{thm:constructibles-embed-solid}, the restriction of the ambient solid
pullback is canonically equivalent to \(\Lhat f^*\).  It preserves
constructibility and is left adjoint to the restricted \(Rf_*\) once the
latter is shown to preserve constructibility.
\end{proposition}

\begin{proof}
The ambient pullback is obtained by pullback followed by the relevant
coefficient localization; on a complete constructible object this is exactly
derived completion.  The first part of
\cref{prop:formal-constructible-preservation} gives preservation.  The
ambient adjunction restricts after the direct-image preservation theorem
below.  Reduction modulo \(\ell^n\) verifies the familiar comparison
\[
 (\Lhat f^*K)/\ell^n\simeq
 \mathbf Lf^*(K/\ell^n),
\]
but this identity is a compatibility theorem, not the definition of the
ambient operation.
\end{proof}

\section{The geometric half of constructible descent}

Ambient existence alone cannot force a functor to preserve a small
geometric subcategory.  The remaining restriction uses exactly the
Grothendieck--Gabber finiteness and constructible-duality content.

\begin{theorem}[Geometric restriction core]
\label{thm:torsion-finiteness-input}
Let \(k\) be finite or separably closed, let \(X,Y\) be finite-type
\(k\)-schemes, and let coefficients be \(\Zl\) or \(\mathcal O_E\), with
\(\ell\ne\operatorname{char}k\).
\begin{enumerate}
\item If \(p:X\to Y\) is proper and \(K\in\Dcons(X,\Lambda)\), then on each
      irreducible component of \(Y\) there is a dense open \(V\) on which
      \((p_*K)|_V\) is lisse and perfect.  Its formation is compatible with
      geometric base change on \(V\), and its cohomological amplitude is
      bounded in terms of \(p\) and the amplitude of \(K\).
\item If \(j:U\hookrightarrow X\) is a quasi-compact open immersion, then
      \(j_*\) preserves constructibility.
\item For \(a_X:X\to\Spec k\), the object
      \(\omega_X:=a_X^!\Lambda\) is constructible, the functor
      \[
       \mathbb D_XK:=
       \R\underline{\operatorname{Hom}}(K,\omega_X)
      \]
      preserves constructibles, and
      \(K\to\mathbb D_X\mathbb D_XK\) is an equivalence for constructible
      \(K\).
\end{enumerate}
\end{theorem}

\begin{proof}[Proof capsule and exact foundational status]
The first two assertions are the generic-local-acyclicity and boundary
pieces of the Grothendieck--Gabber constructibility theorem.  The proof is a
simultaneous noetherian induction on the support.  After shrinking an
irreducible base, generic local acyclicity makes the proper pushforward
locally constant and base-change compatible.  Cohomological-dimension bounds
give boundedness; finite generation on geometric fibers gives perfectness.
On the closed complement, the localization triangle and the induction
hypothesis control the boundary.  The same induction, applied to a
compactification of an open immersion, gives boundary finiteness.

For duality, compactify the structural map, use proper duality and the two
localization triangles, and induct on a regular stratification.  On a smooth
stratum of dimension \(d\), purity gives
\(\omega\simeq\Lambda(d)[2d]\); biduality is then perfect-module biduality.
The localization triangles extend it across the complement.

These are deep geometric theorems, assembled for complete coefficients in
\cite[Section~6.7]{BhattScholzeProEtale} from the classical finiteness,
local acyclicity, duality, and purity theorems.  Their currently available
proof uses finite torsion coefficients and uniform bounds.  The present notes
do not relabel that proof as ``formal solidification''.  A proof avoiding
every torsion theorem would be a new proof of the torsion special cases as
well.
\end{proof}

The global preservation theorem is now a short Route B argument in the
ambient category.

\begin{proposition}[Proper direct image preserves constructibility]
\label{prop:route-b-proper-constructible}
If \(p:X\to Y\) is proper between finite-type \(k\)-schemes, then
\[
 p_*:\Dcons(X,\Lambda)\longrightarrow\Dcons(Y,\Lambda).
\]
\end{proposition}

\begin{proof}
Proceed by noetherian induction on closed subsets of \(Y\).  For
\(K\in\Dcons(X,\Lambda)\), choose, by
\cref{thm:torsion-finiteness-input}, a dense open immersion
\(v:V\hookrightarrow Y\) such that \(v^*p_*K\) is lisse.  Let
\(i:Z\hookrightarrow Y\) be the reduced closed complement and form
\[
\begin{tikzcd}
X_Z \arrow[r,"i'"] \arrow[d,"p_Z"'] & X \arrow[d,"p"] \\
Z \arrow[r,"i"'] & Y.
\end{tikzcd}
\]
Ambient solid proper base change gives
\[
 i^*p_*K\simeq p_{Z*}i'^*K.
\]
The right side is constructible by the induction hypothesis.  Thus \(p_*K\)
is lisse on \(V\) and constructible on \(Z\); the finite-stratification
gluing criterion makes it constructible on \(Y\).
\end{proof}

\begin{corollary}[Integral direct and exceptional images]
\label{cor:integral-direct-image}
For every separated finite-type morphism \(f:X\to Y\),
\[
 f_!:\Dcons(X,\Lambda)\longrightarrow\Dcons(Y,\Lambda).
\]
Moreover, after constructible Verdier duality is applied below, the ambient
\(f_*\) also restricts to constructibles.
\end{corollary}

\begin{proof}
Choose \(f=pj\) with \(j\) open and \(p\) proper.  By
\cref{prop:formal-constructible-preservation}, \(j_!K\) is constructible;
by \cref{prop:route-b-proper-constructible}, \(p_*j_!K\) is constructible.
The ambient identity \(f_!\simeq p_*j_!\) proves the claim.
\end{proof}

\section{Verdier duality descends the remaining functors}

The ambient formalism supplies natural Verdier exchange transformations.
On reflexive constructible objects they are equivalences.

\begin{theorem}[Constructible exceptional pullback and direct image]
\label{thm:proetale-fupper-shriek}
Let \(f:X\to Y\) be separated and of finite type over \(k\).  For
constructible \(K\) there are natural ambient equivalences
\begin{equation}
\label{eq:six-duality-identities}
 f^!K\simeq\mathbb D_Xf^*\mathbb D_YK,
 \qquad
 f_*K\simeq\mathbb D_Yf_!\mathbb D_XK.
\end{equation}
Consequently both \(f^!\) and \(f_*\) preserve constructibility.
\end{theorem}

\begin{proof}
Since \(a_X=a_Yf\), transitivity gives
\(\omega_X\simeq f^!\omega_Y\).  For arbitrary ambient \(M\) and constructible \(K\), the first
Verdier exchange is verified by the chain
\begin{align*}
 \operatorname{Map}_X
  \bigl(M,\mathbb D_Xf^*\mathbb D_YK\bigr)
 &\simeq
 \operatorname{Map}_X
  \bigl(M\widehat\otimes f^*\mathbb D_YK,\omega_X\bigr)\\
 &\simeq
 \operatorname{Map}_Y
  \bigl(f_!M\widehat\otimes\mathbb D_YK,\omega_Y\bigr)\\
 &\simeq
 \operatorname{Map}_Y
  \bigl(f_!M,\mathbb D_Y\mathbb D_YK\bigr)\\
 &\simeq \operatorname{Map}_Y(f_!M,K)\\
 &\simeq \operatorname{Map}_X(M,f^!K).
\end{align*}
The second line uses the projection formula and \(f_!\dashv f^!\), and the
fourth uses constructible biduality.  This calculation is the mate
verification of the ambient Verdier exchange map, hence proves the first
equivalence.  Applying the same calculation to the other adjunction, or
taking Verdier duals, gives the second.

Every functor on the right sides of
\eqref{eq:six-duality-identities} preserves constructibility by
\cref{prop:formal-constructible-preservation,%
thm:torsion-finiteness-input,cor:integral-direct-image}.  Therefore the two
left sides are constructible.
\end{proof}

\begin{proposition}[Tensor, internal Hom, and projection]
\label{prop:proetale-tensor-hom}
For constructible \(K,L\), completed tensor product and internal Hom are
constructible, and for constructible \(M\) on \(Y\) there is an equivalence
\begin{equation}
\label{eq:fshriek-projection}
 f_!(K\widehat\otimes\Lhat f^*M)
   \simeq f_!K\widehat\otimes M.
\end{equation}
\end{proposition}

\begin{proof}
Preservation was proved in
\cref{prop:formal-constructible-preservation}.  The projection formula is an
ambient equivalence from \cref{thm:ambient-solid-six-functors}; since both
sides are constructible, it restricts without a new proof.
\end{proof}

\begin{theorem}[Purity inputs]
\label{thm:proetale-purity-inputs}
Let coefficients be as above.
\begin{enumerate}
\item If \(f:X\to Y\) is smooth of pure relative dimension \(d\), the
      etale orientation gives
      \[
       f^!K\simeq\Lhat f^*K(d)[2d].
      \]
\item If \(i:Z\hookrightarrow X\) is a regular immersion of pure codimension
      \(c\) between regular noetherian schemes, then
      \[
       i^!\Lambda\simeq\Lambda(-c)[-2c].
      \]
\end{enumerate}
\end{theorem}

\begin{proof}
The first assertion is the oriented form of smooth ambidexterity in the
ambient formalism.  The second is Gabber absolute purity.  Solidification
does not prove absolute purity: its proof uses deformation to the normal
cone, the zero-section calculation, and the corresponding torsion purity
theorem.  See \cite{ILOTravauxGabber}.
\end{proof}

\section{The restricted coherent package}

Return to a cartesian square
\begin{equation}
\label{eq:six-base-change-square}
\begin{tikzcd}
X' \arrow[r,"g'"] \arrow[d,"f'"'] & X \arrow[d,"f"] \\
Y' \arrow[r,"g"'] & Y.
\end{tikzcd}
\end{equation}

\begin{theorem}[Constructible base change]
\label{thm:proetale-base-change}
For constructible complexes:
\begin{enumerate}
\item if \(f\) is proper, then
\[
 \Lhat g^*f_*K\xrightarrow{\sim}f'_*\Lhat{g'}^*K;
\]
\item for separated finite-type \(f\), exceptional direct image satisfies
\[
 \Lhat g^*f_!K\xrightarrow{\sim}f'_!\Lhat{g'}^*K;
\]
\item the smooth base-change and projection-formula equivalences also
restrict to constructible objects.
\end{enumerate}
\end{theorem}

\begin{proof}
These are the ambient equivalences of
\cref{thm:solid-ambient-base-change,%
thm:ambient-solid-six-functors}.  The preceding sections prove that every
object appearing on either side is constructible.  Fullness of the
constructible inclusion therefore restricts the equivalences.  No separate
coherence check modulo \(\ell^n\) is necessary.
\end{proof}

\begin{theorem}[Route B constructible six operations]
\label{thm:proetale-six-functor-package}
For finite-type schemes over a finite field and separated finite-type
morphisms, the ambient solid operations restrict to
\[
 X\longmapsto\Dcons(X,\Lambda)
\]
and give
\[
 \Lhat f^*,\quad f_*,\quad f_!,\quad f^!,\quad
 -\widehat\otimes-,\quad
 \R\underline{\operatorname{Hom}}(-,-).
\]
They carry both adjunctions, composition equivalences, proper and
exceptional base change, projection formulas, localization triangles,
smooth and absolute purity where applicable, and constructible Verdier
duality.  Every higher coherence diagram is inherited from the ambient
correspondence formalism.
\end{theorem}

\begin{proof}
Pullback, tensor, internal Hom, and open extension by zero preserve
constructibility by
\cref{prop:formal-constructible-preservation}.  Proper direct image and
exceptional direct image preserve it by
\cref{prop:route-b-proper-constructible,cor:integral-direct-image}.
Exceptional pullback and ordinary direct image preserve it by
\cref{thm:proetale-fupper-shriek}.  Hence all ambient functors restrict.

The constructible categories are full subcategories.  Units, counits,
exchange maps, and projection maps are therefore the restrictions of the
ambient natural transformations.  Every equation between their composites
already holds in the ambient correspondence category, so all coherence
diagrams continue to commute after restriction.
\end{proof}

\subsection{Rational coefficients}

For a finite extension \(E/\Ql\), the safe constructible relation is
\[
 \Dcons(X,\mathcal O_E)[1/\ell]
   \simeq\Dcons(X,E).
\]
Define
\[
 \Dcons(X,\Qbarl)
   :=\varinjlim_{E/\Ql\ \mathrm{finite}}\Dcons(X,E).
\]
Linearity and finite coefficient-field descent extend the restricted six
operations to \(E\) and \(\Qbarl\).  We do not identify the rationalization
of an arbitrary derived-complete naive category with all
\(\Ql\)-sheaves; the solid topological coefficient algebra is the ambient
rational object.

\begin{powerbox}[title={Why this architecture helps the perverse proof}]
The perverse \(t\)-structure is imposed only after the ambient operations
have been constructed and after their constructible preservation has been
proved.  Consequently recollement, proper pushforward, exceptional
pullback, duality, Fourier kernels, and all base-change maps used later are
restrictions of one coherent formalism.  Perverse truncation need not carry
compactification choices or compatible \(\ell^n\)-systems in its data.
\end{powerbox}

\begin{formalization}
The reusable implementation should expose four nested layers:
\[
\begin{tikzcd}[row sep=small]
\Dnaive(X,\Lambda)
 \arrow[d,"(-)^{\blacksquare}"]\\
\Dsolid(X,\Lambda)
 \arrow[d,phantom,"\supset"]\\
\Dindcons(X,\Lambda)
 \arrow[d,phantom,"\supset"]\\
\Dcons(X,\Lambda)=\Dindcons(X,\Lambda)^\omega.
\end{tikzcd}
\]
The ambient correspondence functor and its coherence belong to the solid
layer.  Constructibility is a predicate with preservation theorems, not a
field in the ambient functor record.  The code must not assert
\(\Dcons=(\Dsolid)^\omega\), must not identify solidity with completion, and
must retain the preprint/foundational status of solid rigidity and geometric
finiteness separately.
\end{formalization}

\begin{auditbox}[title={Final Route B dependency audit}]
The exact dependency chain is
\[
\begin{tikzcd}[column sep=small,row sep=large]
\substack{\text{solidification of the}\\
          \text{naive pro-etale category}}
 \arrow[r] &
\substack{\text{solid pro-etale motives}\\
          \text{and rigidity}}
 \arrow[r] &
\substack{\text{ambient coherent}\\
          \text{six operations}}\\
\substack{\text{torsion proper base change}\\
          \text{and torsion rigidity}}
 \arrow[ur] \arrow[r] &
\substack{\text{Grothendieck--Gabber}\\
          \text{finiteness and duality}}
 \arrow[r] &
\substack{\text{restriction to}\\
          \Dcons}
 \arrow[u].
\end{tikzcd}
\]
Thus Route B genuinely constructs the operations before constructibility and
does not define them level by level.  It does \emph{not} prove that classical
finite-coefficient geometry was unnecessary.  In the current proof that
geometry enters twice: to establish ambient solid base change on generators,
and to prove that direct images and duality preserve constructibles.
\end{auditbox}

\chapter{Weil Structures and Frobenius on Complete Complexes}
\label{chap:proetale-weil}

Let \(k=\Fq\), let \(\bar k=\Fbar\), and let
\(X_0\) be a separated scheme of finite type over \(k\).  Put
\[
  X=X_0\times_k\bar k.
\]
This chapter adds Frobenius descent data to the complete coefficient
categories.  No estimate on Frobenius eigenvalues is used.

\section{The two Frobenius descriptions}

Write
\[
 \sigma_q:a\longmapsto a^q
 \quad\text{and}\quad
 \varphi_q=\sigma_q^{-1}
\]
for arithmetic and geometric Frobenius in
\(\Gal(\bar k/k)\), respectively.  The element \(\varphi_q\) induces a
semilinear automorphism
\[
  \Fr_X:X\longrightarrow X.
\]
It is the geometric Frobenius automorphism used in every Weil structure
below.

There is also a scheme-theoretic description.  The absolute \(q\)-power
Frobenius
\[
  F^{(q)}_{X_0}:X_0\longrightarrow X_0
\]
is a universal homeomorphism.  Its pullback on the etale and pro-etale topoi
is therefore an equivalence by
\cref{thm:proetale-universal-homeomorphism}.  After geometric base change,
the contravariant pullback action agrees with the geometric-Frobenius
normalization.

\begin{convention}
The map on geometric points induced by raising coordinates to the
\(q\)-th power is called arithmetic Frobenius.  Cohomology is contravariant,
and throughout these notes the resulting operator on stalks and cohomology is
normalized as \emph{geometric} Frobenius \(\varphi_q=\sigma_q^{-1}\).
The invariant test is
\[
  \Fr\mid \Lambda(1)=q^{-1},
  \qquad
  \Fr\mid \Lambda(-1)=q.
\]
Any formula failing this test uses the opposite convention.
\end{convention}

\begin{proposition}[Frobenius pullback is invertible]
\label{prop:proetale-frobenius-equivalence}
Both \(\Fr_X^*\) and pullback by \(F^{(q)}_{X_0}\) are symmetric monoidal
autoequivalences of \(\Dcons\) and commute with derived completion, reduction
modulo \(\mathfrak m^n\), and scalar extension.
\end{proposition}

\begin{proof}
For \(\Fr_X\) this is immediate because it is an automorphism.  For
\(F^{(q)}_{X_0}\), use universal-homeomorphism invariance
\cite[Lemma~5.4.2]{BhattScholzeProEtale}.  The equivalence preserves the
coefficient sheaf and its finite quotients, so it commutes with the inverse
limit defining completion.
\end{proof}

\section{Weil complexes}

\begin{definition}[Integral Weil complex]
\label{def:proetale-weil-complex}
A constructible \(\Lambda\)-Weil complex on \(X_0\) is a pair
\[
  (K,\theta_K),\qquad
  K\in\Dcons(X,\Lambda),\qquad
  \theta_K:\Fr_X^*K\xrightarrow{\sim}K.
\]
A morphism \(u:(K,\theta_K)\to(L,\theta_L)\) satisfies
\[
  u\circ\theta_K
   =\theta_L\circ\Fr_X^*(u).
\]
The resulting category is denoted
\(\Dcons^W(X_0,\Lambda)\).  The same definition is used for
\(E\)- and \(\Qbarl\)-coefficients.
\end{definition}

Iteration of \(\theta_K\) gives compatible isomorphisms
\[
  (\Fr_X^n)^*K\xrightarrow{\sim}K
  \qquad(n\in\mathbf Z).
\]
Thus a Weil structure is a \(\mathbf Z\)-equivariant object.  It is essential
not to replace the isomorphism \(\theta_K\) by the assertion that
\(\Fr_X^*K=K\) definitionally.

\begin{proposition}[Descent and continuity]
\label{prop:weil-descent-continuity}
Pullback from \(X_0\) to \(X\) gives a fully faithful functor from
constructible complexes on \(X_0\) to Weil complexes on \(X_0\).  Its
essential image consists of those Weil complexes for which the
\(\mathbf Z\)-action generated by \(\theta_K\) extends continuously to
\[
  \Gal(\bar k/k)\simeq\widehat{\mathbf Z}.
\]
For an integral constructible complex, continuity can be checked after every
reduction modulo \(\mathfrak m^n\); on a finite stratification it is
equivalent to the finite-level Frobenius actions factoring through finite
quotients.
\end{proposition}

\begin{proof}
This is effective Galois descent on the pro-etale site.  A sheaf on \(X_0\)
pulls back with continuous descent data.  Conversely, continuous descent
data are effective because \(X\to X_0\) is a pro-etale Galois cover with
group \(\widehat{\mathbf Z}\).  By
\cref{thm:proetale-complete-system-equivalence}, an integral object is its
tower of finite reductions.  At each level the constructible Hom groups are
finite, so a continuous action factors through a finite quotient; compatible
finite-level descent reconstructs the complete object.
\end{proof}

\begin{warning}[Weil does not automatically mean continuous over \(E\)]
A single automorphism of an \(E\)-vector space defines a
\(\mathbf Z\)-action, but it need not extend to
\(\widehat{\mathbf Z}\).  For example, multiplication by a nonunit of
\(E\) preserves no \(\mathcal O_E\)-lattice as an automorphism.  Thus the
rational Weil category is larger than the category of objects descended
from \(X_0\).  The distinction disappears only after a Frobenius-stable
lattice or an equivalent continuity condition has been supplied.
\end{warning}

\section{Tate twists}

\begin{definition}[Integral Tate object]
\label{def:proetale-tate-object}
The pro-etale Tate sheaf is
\[
  \Zl(1):=\varprojlim_n \mu_{\ell^n}.
\]
For \(\Lambda=\mathcal O_E\), set
\[
 \Lambda(1)=\Zl(1)
   \widehat\otimes_{\Zl}\Lambda,
 \qquad
 \Lambda(r)=\Lambda(1)^{\widehat\otimes r}
\]
for \(r\ge0\), and define negative twists by the dual object.
For any \(K\), write \(K(r)=K\widehat\otimes\Lambda(r)\).
\end{definition}

\begin{proposition}[Frobenius normalization of twists]
\label{prop:tate-frobenius-normalization}
At a rational geometric point, geometric Frobenius acts on
\(\Lambda(r)\) by \(q^{-r}\).  Consequently twisting by \((r)\) lowers every
Frobenius weight by \(2r\).
\end{proposition}

\begin{proof}
Arithmetic Frobenius sends an \(\ell^n\)-th root of unity \(\zeta\) to
\(\zeta^q\).  Its inverse therefore acts through the inverse of \(q\) in
\((\mathbf Z/\ell^n)^{\times}\).  Passing to the inverse limit gives
\(q^{-1}\) on \(\Zl(1)\); tensor powers and duals give the general formula.
The assertion about weights is the definition of the twist normalization.
\end{proof}

\section{The six operations inherit Weil structures}

Let \(f_0:X_0\to Y_0\) be a separated finite-type morphism and
\(f:X\to Y\) its geometric base change.  The square
\[
\begin{tikzcd}
X \arrow[r,"\Fr_X"] \arrow[d,"f"'] & X \arrow[d,"f"] \\
Y \arrow[r,"\Fr_Y"'] & Y
\end{tikzcd}
\]
commutes.

\begin{theorem}[Weil-equivariant six operations]
\label{thm:proetale-weil-six-compatibility}
Every operation in
\cref{thm:proetale-six-functor-package} lifts canonically to Weil complexes.
For example, the Weil structure on \(\R f_!K\) is the composite
\[
 \Fr_Y^*\R f_!K
   \xrightarrow{\sim}
 \R f_!\Fr_X^*K
   \xrightarrow{\R f_!(\theta_K)}
 \R f_!K,
\]
and the analogous formula holds for \(\R f_*\).  Pullback, \(f^!\), completed
tensor, internal Hom, Verdier duality, Tate twists, base change, projection
formulas, and localization triangles are all Frobenius equivariant.
\end{theorem}

\begin{proof}
Apply the exchange equivalences of the six-functor formalism to the
commutative Frobenius square, then compose with the given Weil structure.
For tensor and internal Hom use the monoidality of \(\Fr^*\).  Coherence of
the six-functor package implies the cocycle and functoriality conditions.
All assertions may also be checked modulo \(\mathfrak m^n\), where they are
the standard torsion Weil formalism, and then recovered by completion.
\end{proof}

\begin{proetalebox}
The same complete object now carries topology, constructibility, all six
operations, and Frobenius.  There is no separate category of adic systems and
no later quotient by negligible towers.  This is the second major
pro-etale simplification.
\end{proetalebox}

\section{Changing coefficients without hiding lattices}

\begin{theorem}[Coefficient tower for Weil complexes]
\label{thm:proetale-weil-coefficient-tower}
Scalar extension gives functors
\[
 \Dcons^W(X_0,\Zl)
   \longrightarrow \Dcons^W(X_0,\mathcal O_E)
   \longrightarrow \Dcons^W(X_0,E)
   \longrightarrow \Dcons^W(X_0,\Qbarl).
\]
The essential image of
\[
 \Dcons^W(X_0,\mathcal O_E)[1/\ell]
   \longrightarrow \Dcons^W(X_0,E)
\]
consists exactly of the rational Weil complexes admitting a
Frobenius-stable \(\mathcal O_E\)-lattice.  Moreover
\[
 \Dcons^W(X_0,\Qbarl)
   =\varinjlim_{E/\Ql\text{ finite}} \Dcons^W(X_0,E).
\]
Every continuous rational Weil complex admits such a lattice after a finite
extension of \(E\), but the lattice is not canonical.
\end{theorem}

\begin{proof}
Scalar extension commutes with Frobenius pullback, so it transports the
structure isomorphism.  An integral Weil model plainly gives a stable
lattice.  Conversely a stable lattice, completed in the
\(\mathfrak m\)-adic topology, is an integral Weil object whose
localization is the given rational object.  A continuous representation has
compact image; a compact subgroup of \(\operatorname{GL}_n(E)\) stabilizes a
lattice after conjugation.  Apply this on finitely many strata and glue by
localization triangles.  The last formula follows from
\cref{eq:qbarl-colimit}, with the Frobenius isomorphism and each morphism
defined over a sufficiently large finite extension.
\end{proof}

\begin{formalization}
The coefficient API should return a proposition saying that a lattice
\emph{exists}; it should not choose one globally.  If a later construction
temporarily chooses a lattice, its output must be shown invariant under
replacing two lattices by a common finite-index sublattice.  The broader
rational Weil category and the stable-lattice subcategory should have
different names.
\end{formalization}

\section{Compactly supported cohomology and Frobenius}

Let \(a:X\to\Spec \bar k\) be the structural morphism.

\begin{theorem}[Adic perfectness]
\label{thm:proetale-weil-perfectness}
\label{thm:adic-perfectness}
For \(K\in\Dcons^W(X_0,\Lambda)\),
\begin{equation}
\label{eq:perfect-compact-cohomology}
 \R\Gamma_c(X,K):=\R a_!K
   \quad\text{belongs to}\quad \operatorname{Perf}(\Lambda).
\end{equation}
After scalar extension to \(E\), every
\(H_c^i(X,K_E)\) is finite-dimensional and vanishes outside a finite range.
The Weil structure induces an automorphism
\[
  F_K:\R\Gamma_c(X,K)\xrightarrow{\sim}\R\Gamma_c(X,K).
\]
If \(X_0\) is proper, \(\R\Gamma_c=\R\Gamma\).
\end{theorem}

\begin{proof}
By \cref{thm:proetale-weil-six-compatibility}, \(\R a_!K\) is a
constructible complete complex on the geometric point.  By
\cref{ex:constructible-point}, that category is \(\operatorname{Perf}(\Lambda)\).  The
Frobenius automorphism is the structure carried by \(\R a_!\); explicitly it
is
\[
 \R a_!K
   \xrightarrow{\sim}
 \R a_!\Fr_X^*K
   \xrightarrow{\R a_!(\theta_K)}
 \R a_!K.
\]
Properness identifies \(a_!\) and \(a_*\).
\end{proof}

\begin{formalization}
In the formalization \(\R a_!K\) is a Weil complex \(M\) on \(\Spec k_0\), and
\(H_c^r(X,K):=H^r(\R\Gamma(\Spec\bar k,M))\).  The Frobenius on it is transported along
the invariance \(\R\Gamma(\Spec\bar k,\sigma^*-)\cong\R\Gamma(\Spec\bar k,-)\) of derived
global sections under the geometric Galois Frobenius \(\sigma\) of \(\Spec\bar k\), which
follows from \(\sigma^*\cong\sigma^{-1}_*\) and the Leray isomorphism
\(\R\Gamma(\Spec\bar k,\R\sigma^{-1}_*-)\cong\R\Gamma(\Spec\bar k,-)\): it is
\[
 H^r(\R\Gamma(\Spec\bar k,M))
   \xrightarrow{\sim}
 H^r(\R\Gamma(\Spec\bar k,\sigma^*M))
   \xrightarrow{H^r(\varphi)}
 H^r(\R\Gamma(\Spec\bar k,M)),
\]
where \(\varphi:\sigma^*M\xrightarrow{\sim}M\) is the Weil structure of \(M\).
\end{formalization}

\begin{definition}[Cohomological trace and determinant]
For a rational Weil complex, define
\begin{align*}
 \Tr(F_K^n\mid\R\Gamma_c(X,K))
   &=\sum_i(-1)^i
     \Tr(F_K^n\mid H_c^i(X,K)),\\
 \det(1-TF_K\mid\R\Gamma_c(X,K))
   &=\prod_i
     \det(1-TF_K\mid H_c^i(X,K))^{(-1)^i}.
\end{align*}
The first expression is a scalar and the second is a normalized rational
function.  Both are well defined because of
\cref{thm:adic-perfectness}.
\end{definition}

\begin{corollary}[Compatibility of traces with coefficients]
\label{cor:weil-trace-coefficient-change}
The Frobenius traces, characteristic polynomials on individual cohomology
groups, and the alternating determinant commute with finite scalar extension
\(E\subset E'\) and hence define the corresponding invariants over
\(\Qbarl\).
\end{corollary}

\begin{proof}
Perfect base change identifies the cohomology after extension of scalars.
Trace and determinant of an endomorphism of a finite-dimensional vector
space are preserved by field extension.
\end{proof}

\begin{auditbox}
Perfectness is a finiteness statement, not a weight statement.  We have not
proved that Frobenius eigenvalues are algebraic, pure, integral, or
independent of \(\ell\), and we have not invoked the trace formula.  Those
assertions begin only after perverse sheaves and weights have been developed.
\end{auditbox}

\section{Why the next layer is perverse}

The six operations act on complexes, but the Weil proof needs an exact
category in which support dimension controls cohomological degree.
The middle-perversity t-structure supplies that category.  Once weights are
attached to its simple objects, the operations \(j_{!*}\), perverse
cohomology, vanishing cycles, and Fourier transform become exact or
one-sided-exact mechanisms that preserve sharp weight bounds.  The next
chapters exploit precisely this extra structure; ordinary cohomology sheaves
cannot provide the same local-to-global control.

\begin{powerbox}
The pro-etale site solves the coefficient problem; perversity solves the
weight-propagation problem.  Neither subsumes the other.  Their combination
is what lets the proof move from a local Frobenius estimate to global purity
without rebuilding adic limits at every step.
\end{powerbox}


\chapter{The Lefschetz--Verdier trace formula}
\label{ch:trace-formula}

The trace formula is first a theorem with finite coefficients and finite Tor
dimension.  Normalization gives the \(\mathbf Z_\ell\)-version and scalar
extension gives the \(\mathbf Q_\ell\)-version used below.  The finite
Tor-dimension hypothesis is essential at finite level: it makes the relevant
objects dualizable and their categorical traces defined.

\section{Traces of perfect complexes}

For a perfect complex \(C\) over a field and an endomorphism \(u\), set
\[
 \operatorname{Tr}(u\mid C)
 =\sum_i(-1)^i\operatorname{Tr}(H^i(u)\mid H^i(C)).
\]

\begin{lemma}[Perfect-complex supertrace]
\label{lem:perfect-supertrace}
Supertrace is invariant under quasi-isomorphism, additive in distinguished
triangles of endomorphisms, multiplicative under derived tensor product, and
cyclic.  Moreover,
\[
 -\log\det(1-Tu\mid C)
 =\sum_{m\geq1}\operatorname{Tr}(u^m\mid C)\frac{T^m}{m},
\]
where
\[
 \det(1-Tu\mid C)
 =\prod_i\det(1-TH^i(u)\mid H^i(C))^{(-1)^i}.
\]
\end{lemma}

\begin{proof}
Represent \(C\) by a bounded finite-dimensional complex and split it into
cohomology and contractible two-term summands.  Contractible summands have
zero supertrace.  This proves invariance and reduces additivity, cyclicity,
and multiplicativity to the corresponding matrix identities.  After a
splitting-field extension, the last identity is the sum of
\(-\log(1-\alpha T)=\sum_{m\geq1}\alpha^mT^m/m\) over the eigenvalues.
\end{proof}

\section{Cohomological correspondences}

Let
\[
 c=(c_1,c_2):C\longrightarrow X\times X
\]
be a correspondence.  For
\(K\in D^b_{ctf}(X,\Lambda_n)\), a cohomological correspondence is a morphism
\[
 u:c_2^*K\longrightarrow Rc_1^!K.
\]
Here \(D^b_{ctf}\) means bounded constructible cohomology and finite Tor
dimension over \(\Lambda_n\).  If \(c_2\) is proper, the unit for
\((c_2^*,Rc_{2*})\), followed by \(u\) and the counit for
\((Rc_{1!},Rc_1^!)\), gives an endomorphism
\[
 R\Gamma_c(u):R\Gamma_c(X,K)\longrightarrow R\Gamma_c(X,K).
\]
The fixed scheme is
\[
 \operatorname{Fix}(c)
 =C\times_{X\times X,c,\Delta_X}X.
\]
Evaluation against the dualizing complex gives a trace class on this fixed
scheme.  Integration over a proper open-and-closed component
\(\beta\) gives its local term \(\operatorname{LT}_\beta(u)\).

\begin{theorem}[Lefschetz--Verdier trace formula]
\label{thm:lefschetz-verdier}
Let \(c_2\) and \(\operatorname{Fix}(c)\) be proper and let
\(K\in D^b_{ctf}(X,\Lambda_n)\).  For
\(u:c_2^*K\to Rc_1^!K\),
\[
 \operatorname{Tr}(R\Gamma_c(u)\mid R\Gamma_c(X,K))
 =
 \sum_{\beta\in\pi_0(\operatorname{Fix}(c))}
 \operatorname{LT}_\beta(u).
\]
The compatible finite-level identities induce the same formula for
finite-Tor-amplitude \(\mathbf Z_\ell\)-complexes and for constructible
\(\mathbf Q_\ell\)-complexes.
\end{theorem}

\begin{proof}[Finite-level proof capsule]
Compactify the correspondence and extend \(K\) by zero.  Localization removes
boundary pieces with no fixed point.  Deformation of the graph to the normal
cone of the diagonal, followed by purity, specializes the diagonal to the
zero section of its normal bundle.  Homotopy invariance and the Thom trace
reduce the vector-bundle case to the identity correspondence on a point,
where the formula is the ordinary categorical trace.

The reduction uses three separately verified identities for trace classes:
proper pushforward, etale pullback, and additivity in localization triangles.
They follow from the units and counits of the six operations, tensor--Hom
adjunction, and proper base change.  They also prove independence of the
compactification.  This is \cite[Expos\'e III]{SGA5}.

No semisimplicity or division by a group order occurs.  In particular, one
does not resolve a \(\Lambda_n[G]\)-module by a finite complex of permutation
modules; this would fail in general when \(\ell\mid|G|\).  The finite-level
theorem is applied directly in \(D^b_{ctf}(X,\Lambda_n)\).  Reduction
compatibility and normalization then give the adic statement.
\end{proof}

\section{Frobenius fixed points and local terms}

Let \(X_0/\mathbf F_q\) be separated of finite type,
\(X=X_0\times\overline{\mathbf F}_q\), and let \((K,\phi)\) be a Weil
complex.  For \(m\geq1\), use the graph correspondence
\[
 C=X,\qquad (c_1,c_2)=(\operatorname{id}_X,F_X^m).
\]
The second projection \(F_X^m\) is finite, hence proper, and the iterated Weil
map is exactly a cohomological correspondence
\[
 u_m:(F_X^m)^*K\longrightarrow K=R(\operatorname{id}_X)^!K.
\]
Its fixed points are \(X_0(\mathbf F_{q^m})\).

At a fixed point \(y\), write \(F_y^m\) for the endomorphism of \(K_{\bar y}\)
induced by \(u_m\).  This is the global \(F^m\) on the fixed-point stalk.  If
the underlying closed point \(x\) has degree \(d\mid m\), then
\[
 F_y^m=F_x^{\,m/d}.
\]

\begin{lemma}[Frobenius local term]
\label{lem:frobenius-local-term}
Every fixed point is isolated and
\[
 \operatorname{LT}_y(F^m,K)
 =\operatorname{Tr}(F_y^m\mid K_{\bar y})
 :=\sum_j(-1)^j
 \operatorname{Tr}(F_y^m\mid\mathcal H^j(K)_{\bar y}).
\]
\end{lemma}

\begin{proof}
First handle singular \(X_0\).  Topological invariance of the etale topos,
together with functoriality of trace classes under the universal homeomorphism
\((X_0)_{\mathrm{red}}\to X_0\), first reduces us to the reduced scheme.
Choose a finite stratification
\[
 X_0=\coprod_a S_a
\]
by smooth locally closed subschemes defined over \(\mathbf F_q\); this exists
because the field is perfect.  Every \(S_a\) is stable under \(F_X^m\).
Repeated localization for a closed stratum and its open complement, together
with additivity of global traces and trace classes, reduces the assertion to
\((S_a,K|_{S_a})\).  Excision identifies the local term at a fixed point in
\(X\) with the one computed in its smooth stratum.

We may therefore assume \(X\) smooth near \(y\).  Choose regular parameters.
The graph-minus-diagonal map has linear part \(1-dF_X^m\).  Since the
\(q\)-power Frobenius has zero differential, this is the identity, so the
fixed point has local intersection multiplicity one.  In the
deformation-to-the-normal-cone definition of the local term only the
categorical trace on the stalk remains.  Devissage through the cohomology
sheaves gives the alternating sum.
\end{proof}

\begin{theorem}[Grothendieck--Lefschetz trace formula]
\label{thm:GL-trace}
For every \(m\geq1\),
\[
 \boxed{\;
 \sum_{y\in X_0(\mathbf F_{q^m})}
 \operatorname{Tr}(F_y^m\mid K_{\bar y})
 =
 \sum_i(-1)^i
 \operatorname{Tr}(F^m\mid H_c^i(X,K)).
 \;}
\]
If \(y\) has underlying closed point \(x\), its summand uses
\(F_x^{m/\deg x}\), not \(F_x^m\).
\end{theorem}

\begin{proof}
Apply Theorem~\ref{thm:lefschetz-verdier} to
\((\operatorname{id},F_X^m,u_m)\).  The fixed locus is finite.
Lemma~\ref{lem:frobenius-local-term} computes every local term, and the global
endomorphism is the \(m\)-fold Weil action on compactly supported cohomology.
\end{proof}

\begin{corollary}[Point-counting trace formula]
\label{cor:point-count-trace}
For the constant sheaf,
\[
 \#X_0(\mathbf F_{q^m})
 =\sum_i(-1)^i\operatorname{Tr}
 (F^m\mid H_c^i(X,\mathbf Q_\ell)).
\]
\end{corollary}

\begin{proof}
Every stalk is one dimensional and its Frobenius trace is \(1\).
\end{proof}

\section{Closed points and rational points}

For \(x\in|X_0|\), let \(d_x=[k(x):\mathbf F_q]\).

\begin{lemma}[Closed-point/rational-point conversion]
\label{lem:closed-rational-traces}
For every \(m\geq1\),
\[
 \sum_{y\in X_0(\mathbf F_{q^m})}
 \operatorname{Tr}(F_y^m\mid K_{\bar y})
 =
 \sum_{\substack{x\in|X_0|\\d_x\mid m}}
 d_x\operatorname{Tr}(F_x^{m/d_x}\mid K_{\bar x}).
\]
\end{lemma}

\begin{proof}
A degree-\(d\) closed point contributes exactly when \(d\mid m\), and then
gives \(d\) rational geometric points.  Global \(F^m\) goes around their
orbit \(m/d\) times, so its stalk action is \(F_x^{m/d}\).  Sum the \(d\)
equal traces.
\end{proof}

\section{Euler products and determinants}

Define
\[
 L_x(K,T)
 =\prod_j\det(1-F_xT^{d_x}\mid
 \mathcal H^j(K)_{\bar x})^{(-1)^{j+1}},
\qquad
 L(X_0,K,T)=\prod_{x\in|X_0|}L_x(K,T).
\]
The Euler product is coefficientwise meaningful because only points of
bounded degree contribute to any fixed coefficient.

\begin{theorem}[Cohomological expression for \(L\)-functions]
\label{thm:cohomological-L}
For every Weil complex \(K\),
\[
 \boxed{\;
 L(X_0,K,T)
 =
 \prod_i\det(1-TF\mid H_c^i(X,K))^{(-1)^{i+1}}.
 \;}
\]
In particular the Euler product is rational.
\end{theorem}

\begin{proof}
The determinant--trace identity gives
\[
\begin{aligned}
\log L(X_0,K,T)
 &=\sum_x\sum_{r\geq1}
   \operatorname{Tr}(F_x^r\mid K_{\bar x})\frac{T^{rd_x}}r\\
 &=\sum_{m\geq1}\frac{T^m}{m}
   \sum_{d_x\mid m}d_x
   \operatorname{Tr}(F_x^{m/d_x}\mid K_{\bar x})\\
 &=\sum_{m\geq1}\frac{T^m}{m}
   \sum_{y\in X_0(\mathbf F_{q^m})}
   \operatorname{Tr}(F_y^m\mid K_{\bar y})\\
 &=\sum_{m\geq1}\frac{T^m}{m}
   \sum_i(-1)^i\operatorname{Tr}
   (F^m\mid H_c^i(X,K)).
\end{aligned}
\]
The second equality uses \(m=rd_x\), the third is
Lemma~\ref{lem:closed-rational-traces}, and the fourth is
Theorem~\ref{thm:GL-trace}.  Apply Lemma~\ref{lem:perfect-supertrace} again
and exponentiate.  Theorem~\ref{thm:adic-perfectness} makes the
cohomological product finite.
\end{proof}

\begin{corollary}[Cohomological zeta formula]
\label{cor:zeta-cohomological}
\[
 Z(X_0,T)
 :=\exp\left(\sum_{m\geq1}
 \#X_0(\mathbf F_{q^m})\frac{T^m}{m}\right)
 =
 \prod_i\det(1-TF\mid H_c^i(X,\mathbf Q_\ell))^{(-1)^{i+1}}.
\]
Thus \(Z(X_0,T)\in\mathbf Q_\ell(T)\).  Rationality over \(\mathbf Q\),
integrality, and independence of \(\ell\) are proved only after weights.
\end{corollary}

\section{Formalization checklist}

\begin{enumerate}
\item Define traces for dualizable objects and use
\(D^b_{ctf}(X,\Lambda_n)\) at finite level.
\item Package a correspondence as \((C,c_1,c_2,u)\) with
\(u:c_2^*K\to Rc_1^!K\) and \(c_2\) proper.
\item For Frobenius use
\((c_1,c_2)=(\operatorname{id},F_X^m)\).
\item Construct local terms, proper pushforward, localization additivity, and
excision as separate theorems.  Singular schemes are reduced by an
\(\mathbf F_q\)-defined smooth stratification before using regular parameters.
\item Do not use a permutation resolution at finite level.
\item Distinguish global \(F_y^m\) from local \(F_x\), and prove
\(F_y^m=F_x^{m/\deg x}\) before manipulating the Euler product.
\item Derive the \(L\)-function identity by equality of formal logarithms.
\end{enumerate}


\chapter{Purity, Verdier duality, and Poincare duality}
\label{ch:duality}

Duality supplies the trace map, reciprocal Frobenius eigenvalues, and the
self-dual parabolic cohomology used in the curve estimate.  We prove the
purity statement needed in this book for regular schemes of finite type over
a perfect field.  Gabber's more general absolute-purity theorem is recorded
separately and is not claimed as a consequence of the proof capsule below.

\section{Cohomology with supports and purity}

Let \(i:Z\hookrightarrow X\) be a closed immersion, with open complement
\(j:U\hookrightarrow X\).  The exceptional inverse image is characterized by
\[
 R\Gamma_Z(X,K)=R\Gamma(Z,Ri^!K),
\]
and localization gives
\[
 i_*Ri^!K\longrightarrow K\longrightarrow Rj_*j^*K\longrightarrow.
\tag{7.1}\label{eq:costalk-localization}
\]

\begin{theorem}[Purity in the geometric case]
\label{thm:absolute-purity}
Let \(k\) be a perfect field of characteristic different from \(\ell\), and
let \(i:Z\hookrightarrow X\) be a regular closed immersion of pure
codimension \(c\), where \(Z\) and \(X\) are regular schemes of finite type
over \(k\).  Equivalently, \(Z\) and \(X\) are smooth over \(k\).  Then the
fundamental class induces a canonical isomorphism
\[
 Ri^!\Lambda_n\xrightarrow{\sim}\Lambda_n(-c)[-2c].
\tag{7.2}\label{eq:absolute-purity}
\]
It is compatible with etale base change, transverse composition of regular
immersions, cup products, and reduction from \(\Lambda_n\) to
\(\Lambda_{n-1}\).

If \(f:X\to Y\) is smooth of pure relative dimension \(d\), then
\[
 Rf^!L\xrightarrow{\sim}f^*L(d)[2d]
\tag{7.3}\label{eq:smooth-purity}
\]
for \(L\in D^b_c(Y,\Lambda_n)\).  The same formulas hold for
\(\mathbf Z_\ell\)- and \(\mathbf Q_\ell\)-coefficients.
\end{theorem}

\begin{proof}[Proof capsule]
For the zero section \(i:\{0\}\hookrightarrow\mathbf A^1_k\), the localization
sequence and Kummer sequence calculate
\[
 H^r_{\{0\}}(\mathbf A^1_k,\Lambda_n)=
 \begin{cases}
 \Lambda_n(-1),&r=2,\\
 0,&r\ne2.
 \end{cases}
\tag{7.4}\label{eq:purity-line}
\]
The boundary of the Kummer class of the coordinate is the generator.
Multiplying the coordinate by a unit does not change its supported boundary.

Kunneth gives the corresponding statement for the zero section of
\(\mathbf A^c\).  Etale descent globalizes the product of the \(c\) Kummer
classes to the Thom class of a rank-\(c\) vector bundle.  For a general
regular immersion form the deformation to the normal cone
\[
 D_ZX=
 \operatorname{Bl}_{Z\times\{0\}}(X\times\mathbf A^1)
 \setminus \operatorname{Bl}_{Z}(X).
\]
Its fiber over \(1\) is \(X\), and its fiber over \(0\) is the normal bundle
\(N_{Z/X}\).  Smoothness of the pair over the perfect field gives the
specialization isomorphism for cohomology with supports.  Under this
isomorphism, the fundamental class of \(Z\subset X\) becomes the Thom class
of the zero section \(Z\subset N_{Z/X}\).  The vector-bundle calculation
therefore proves \eqref{eq:absolute-purity}.  Products of Kummer classes prove
composition compatibility, and base change of the deformation space proves
base-change compatibility.

For a smooth \(f\), the diagonal
\(\Delta_f:X\hookrightarrow X\times_YX\) is a regular immersion of
codimension \(d\).  Apply purity to the diagonal and the two projection
adjunctions.  This identifies \(Rf^!\Lambda_n\) with
\(\Lambda_n(d)[2d]\); the projection formula gives
\eqref{eq:smooth-purity}.  Compatibility in \(n\), normalization, and
inversion of \(\ell\) give the adic assertions.
\end{proof}

\begin{remark}[Gabber's absolute purity]
For a regular closed immersion between excellent regular schemes on which
\(\ell\) is invertible, the same isomorphism
\(Ri^!\Lambda_n\simeq\Lambda_n(-c)[-2c]\) is Gabber's absolute-purity theorem
\cite{SGA4half}.  Its mixed-characteristic proof is genuinely
deeper than deformation to the normal cone and is not supplied here.  The
geometric case above is sufficient for all varieties, smooth morphisms, and
Lefschetz pencils used in this book.
\end{remark}

\begin{corollary}[Gysin morphism]
\label{cor:gysin}
Under the hypotheses of Theorem~\ref{thm:absolute-purity}, adjunction gives
\[
 i_*:H^r(Z,\Lambda_n(s))
 \longrightarrow H^{r+2c}_Z(X,\Lambda_n(s+c)).
\]
Gysin morphisms compose and commute with transverse pullback.
\end{corollary}

\section{Exceptional inverse image and Verdier duality}

\begin{theorem}[Existence of \(Rf^!\)]
\label{thm:exceptional-pullback}
For a separated finite-type morphism \(f:X\to Y\), the functor
\[
 Rf_!:D^b_c(X,\Lambda_n)\longrightarrow D^b_c(Y,\Lambda_n)
\]
has a right adjoint \(Rf^!\), characterized by
\[
 R\operatorname{Hom}_Y(Rf_!K,L)
 \simeq R\operatorname{Hom}_X(K,Rf^!L).
\tag{7.5}\label{eq:fshriek-adjunction}
\]
It is compatible with composition, equals \(f^*\) for etale \(f\), satisfies
\eqref{eq:smooth-purity} for smooth \(f\), and agrees with supported costalks
for a closed immersion.
\end{theorem}

\begin{proof}[Proof capsule]
For a compactification \(f=\bar f j\), set
\(Rf^!=j^*R\bar f^!\).  The right adjoint for a proper map is constructed
first for finite maps by sheaf Hom, for projective space by the trace class,
for projective maps after a closed embedding, and for proper maps by
cohomological descent from projective envelopes.  The adjunction is checked
at each stage, and the uniform cohomological bound preserves the bounded
constructible category.  A common refinement of compactifications and
uniqueness of right adjoints identify the results.  The same uniqueness gives
composition and all coherence isomorphisms; see
\cite[Expos\'e XVIII]{SGA4}.
\end{proof}

Let \(a_X:X\to\operatorname{Spec}k\), with \(k\) separably closed, and define
\[
 \omega_X:=Ra_X^!\mathbf Q_\ell,\qquad
 D_XK:=R\mathcal Hom(K,\omega_X).
\]
If \(X\) is smooth of pure dimension \(d\), then
\[
 \omega_X\simeq\mathbf Q_\ell(d)[2d].
\tag{7.6}\label{eq:smooth-dualizing}
\]

\begin{theorem}[Verdier duality]
\label{thm:verdier-duality}
For \(K\in D^b_c(X,\mathbf Q_\ell)\):
\begin{enumerate}
\item \(K\to D_XD_XK\) is an isomorphism;
\item for separated finite-type \(f:X\to Y\),
\[
 D_YRf_!K\xrightarrow{\sim}Rf_*D_XK;
\tag{7.7}\label{eq:verdier-push}
\]
\item for \(L\in D^b_c(Y,\mathbf Q_\ell)\),
\[
 D_Xf^*L\xrightarrow{\sim}Rf^!D_YL;
\tag{7.8}\label{eq:verdier-pull}
\]
\item these maps commute with composition, base change, and Weil structures.
\end{enumerate}
\end{theorem}

\begin{proof}
Biduality is etale local and additive in localization triangles.  On a smooth
stratum where \(K\) is a bounded complex of finite free lisse sheaves, it is
ordinary biduality for finite-dimensional vector spaces, shifted and twisted
by \eqref{eq:smooth-dualizing}.  Noetherian induction over a finite
stratification proves the global assertion.

To prove \eqref{eq:verdier-push}, test both sides against an arbitrary complex
and use, successively, tensor--Hom adjunction, the projection formula, and
\eqref{eq:fshriek-adjunction}; Yoneda gives the isomorphism.
Equation \eqref{eq:verdier-pull} is its adjoint form.  The maps are uniquely
specified by units, counits, and evaluation, so their coherence and Weil
equivariance follow from the corresponding properties of the six operations.
\end{proof}

\section{Trace and Poincare duality}

For a smooth separated \(d\)-fold \(X/k\), the counit
\(Ra_{X!}Ra_X^!\mathbf Q_\ell\to\mathbf Q_\ell\) is the trace map
\[
 \operatorname{tr}_X:
 H_c^{2d}(X,\mathbf Q_\ell(d))\longrightarrow\mathbf Q_\ell.
\tag{7.9}\label{eq:trace-map}
\]
It is normalized so that the product of the Kummer fundamental classes on
\(\mathbf A^d\) maps to \(1\).

\begin{theorem}[Poincare duality]
\label{thm:poincare-duality}
Let \(X\) be smooth, separated, and pure of dimension \(d\) over a separably
closed field, and let \(L\) be a lisse \(\mathbf Q_\ell\)-sheaf.  Cup product,
evaluation, and trace give a perfect pairing
\[
 H_c^r(X,L)\times H^{2d-r}(X,L^\vee(d))
 \longrightarrow\mathbf Q_\ell.
\tag{7.10}\label{eq:poincare-pairing}
\]
Equivalently,
\[
 R\Gamma_c(X,L)^\vee
 \simeq R\Gamma(X,L^\vee(d))[2d].
\tag{7.11}\label{eq:poincare-derived}
\]
The pairing is functorial and Weil equivariant.
\end{theorem}

\begin{proof}
Apply \eqref{eq:verdier-push} to \(a_X\):
\[
 R\operatorname{Hom}(R\Gamma_c(X,L),\mathbf Q_\ell)
 \simeq R\Gamma(X,D_XL).
\]
Smooth purity identifies \(D_XL\) with \(L^\vee(d)[2d]\).
Finite-dimensionality makes the resulting cohomological pairing perfect.
The construction uses only Weil-equivariant evaluation and adjunction maps.
\end{proof}

\begin{proposition}[Frobenius reciprocity]
\label{prop:frob-duality}
Let \(X_0/\mathbf F_q\) be smooth proper of dimension \(d\).  The untwisted
pairing
\[
 H^r(X,L)\times H^{2d-r}(X,L^\vee)
 \longrightarrow\mathbf Q_\ell(-d)
\tag{7.12}\label{eq:untwisted-pairing}
\]
satisfies
\[
 \langle Fv,Fw\rangle=q^d\langle v,w\rangle.
\tag{7.13}\label{eq:frob-similitude}
\]
Thus an eigenvalue \(\alpha\) on \(H^r(X,L)\) corresponds, with the same
algebraic multiplicity, to \(q^d/\alpha\) on
\(H^{2d-r}(X,L^\vee)\).
\end{proposition}

\begin{proof}
The target is \(\mathbf Q_\ell(-d)\), on which geometric Frobenius acts by
\(q^d\).  Perfectness identifies the operator on the second space with
\(q^d(F^{-1})^{\mathsf t}\), proving both assertions.
\end{proof}

\begin{corollary}[Relative Poincare duality]
\label{cor:relative-poincare}
If \(f:X_0\to Y_0\) is smooth proper of relative dimension \(d\), then
\[
 (R^rf_*L)^\vee\simeq R^{2d-r}f_*L^\vee(d)
\]
as Weil sheaves.
\end{corollary}

\begin{proof}
Apply Poincare duality to every geometric fiber.  Proper and smooth base
change identify the terms with stalks of the displayed sheaves; conservativity
of geometric stalks glues the isomorphism.
\end{proof}

\section{Local duality and middle extension on a curve}

Let \(C\) be a smooth proper geometrically connected curve, let
\(j:U\hookrightarrow C\) be nonempty open, and let \(L\) be lisse on \(U\).
At \(s\in C\setminus U\), write \(I_s\) for inertia and
\(V=L_{\bar\eta}\).  Then
\[
 (j_*L)_{\bar s}=V^{I_s},\qquad
 (R^1j_*L)_{\bar s}=H^1(I_s,V),\qquad
 R^rj_*L=0\quad(r>1).
\tag{7.14}\label{eq:curve-local-stalks}
\]
Indeed, wild inertia is pro-\(p\), while the maximal pro-\(\ell\) quotient of
tame inertia has \(\ell\)-cohomological dimension one.

The intermediate extension in the perverse heart is
\[
 j_{!*}L[1]
 :=\operatorname{im}({}^pj_!L[1]\to{}^pj_*L[1]).
\]
On a smooth curve the stalk calculation above gives
\[
 j_{!*}L[1]\simeq(j_*L)[1].
\tag{7.15}\label{eq:curve-middle-extension}
\]

\begin{theorem}[\(j_*\)-local duality on curves]
\label{thm:curve-middle-duality}
There is a canonical Weil-equivariant isomorphism
\[
 D_C((j_*L)[1])
 \xrightarrow{\sim}(j_*L^\vee(1))[1].
\tag{7.16}\label{eq:jstar-local-duality}
\]
Consequently,
\[
 H^1(C,j_*L)\times H^1(C,j_*L^\vee(1))
 \longrightarrow\mathbf Q_\ell
\tag{7.17}\label{eq:curve-middle-pairing}
\]
is perfect, and
\[
 H^1(C,j_*L)
 =
 \operatorname{im}\bigl(H_c^1(U,L)\to H^1(U,L)\bigr).
\tag{7.18}\label{eq:parabolic-cohomology}
\]
\end{theorem}

\begin{proof}
Verdier duality exchanges \(Rj_!\) and \(Rj_*\) and is \(t\)-exact for the
middle-perverse \(t\)-structure.  Since
\(D_U(L[1])=L^\vee(1)[1]\), it takes the image defining \(j_{!*}\) to the
opposite image and gives
\[
 D_C(j_{!*}L[1])\simeq j_{!*}L^\vee(1)[1].
\]
Use \eqref{eq:curve-middle-extension} on both sides.  The local content is
Tate duality for inertia: evaluation and the residue map pair the invariant
boundary term on one side with the coinvariant boundary term on the other.

Global Verdier duality now gives \eqref{eq:curve-middle-pairing}.  Finally,
the localization long exact sequences show that both
\(H^1(C,j_*L)\) and the image of compactly supported cohomology are the kernel
of
\[
 H^1(U,L)\longrightarrow
 \bigoplus_{s\in C\setminus U}H^1(I_s,V).
\]
This proves \eqref{eq:parabolic-cohomology}.
\end{proof}

\section{Formalization checklist}

\begin{enumerate}
\item State geometric purity with its actual scope: regular finite-type pairs
over a perfect field.  Record Gabber's excellent-regular theorem as a separate
external result, not as a consequence of Popescu approximation.
\item Construct fundamental and Thom classes, deformation to the normal cone,
and their base-change and composition laws.
\item Export \(Rf^!\) with its adjunction, dualizing complexes, biduality, and
the two Verdier exchange isomorphisms.
\item Normalize the trace on affine space and test every shift and twist:
\[
 D_X(\mathbf Q_\ell)=\mathbf Q_\ell(d)[2d],\quad
 H^{2d}(X,\mathbf Q_\ell)=\mathbf Q_\ell(-d).
\]
\item Store both the twisted invariant Poincare pairing and the untwisted
\(q^d\)-similitude pairing; derive
\(\alpha\leftrightarrow q^d/\alpha\).
\item For curves, separate inertia stalk calculations, perverse intermediate
extension, local Tate duality, and the parabolic image theorem.
\end{enumerate}



\part{The perverse engine}
\chapter{The Middle Perverse \texorpdfstring{\(t\)}{t}-Structure}
\label{chap:perverse-t-structure}

Ordinary cohomological degree measures how far a complex is from being a
sheaf.  Perverse degree measures that failure after compensating for the
dimension of the support.  This compensation is exactly what makes a local
system on a smooth \(d\)-fold live in degree \(-d\), makes Verdier duality
\(t\)-exact, and makes intersection cohomology an object of an abelian
category.

Throughout this chapter \(X\) is a finite-type scheme over a field, and
\(\Dbc(X)=\Dbc(X,\Qbarl)\) is the bounded constructible derived category.  The
comparison with the classical constructible adic category is compatible with
stalks, costalks, and recollement, so this is the usual rational middle
perversity, now realized on \(X_{\proet}\).  Integral coefficients require the
distinct \(p\) and \(p+\) structures of \cref{chap:integral-perverse}; we make
no integral self-duality claim in the present chapter.  For a geometric point
\(\bar x\to X\), let
\[
  \delta(x)=\dim\overline{\{x\}}.
\]

\section{Support and cosupport conditions}

\begin{definition}[Middle perversity]
\label{def:middle-perversity}
Define full subcategories of \(\Dbc(X)\) by
\begin{align*}
 {}^p\D^{\leq0}(X)
 &=\{\cK:H^a(i_{\bar x}^*\cK)=0
       \text{ whenever }a>-\delta(x)\},\\
 {}^p\D^{\geq0}(X)
 &=\{\cK:H^a(i_{\bar x}^!\cK)=0
       \text{ whenever }a<-\delta(x)\}.
\end{align*}
Here \(i_{\bar x}:\bar x\to X\) is the geometric-point morphism.  Put
\[
 {}^p\D^{\leq n}={}^p\D^{\leq0}[-n],\qquad
 {}^p\D^{\geq n}={}^p\D^{\geq0}[-n].
\]
\end{definition}

The stalk condition has an equivalent global form:
\[
 \cK\in{}^p\D^{\leq0}
 \quad\Longleftrightarrow\quad
 \dim\Supp\cH^a(\cK)\leq-a\quad\text{for all }a.
\tag{8.1}\label{eq:perverse-support}
\]
The costalk condition is not obtained by replacing stalks with ordinary
stalks.  It records local cohomology transverse to the strata, and is the
Verdier-dual condition.

\begin{theorem}[BBD middle \(t\)-structure]
\label{thm:middle-t-structure}
The pair
\[
  ({}^p\D^{\leq0}(X),{}^p\D^{\geq0}(X))
\]
is a bounded \(t\)-structure on \(\Dbc(X)\).  Its heart
\[
  \Perv(X,\Qbarl)
  ={}^p\D^{\leq0}(X)\cap{}^p\D^{\geq0}(X)
\]
is an abelian category.  The truncation functors are denoted
\(\ptau_{\leq n},\ptau_{\geq n}\), and
\[
  \pH^n(\cK)=\bigl(\ptau_{\geq n}\ptau_{\leq n}\cK\bigr)[n]
\]
is the \(n\)-th perverse cohomology object.
\end{theorem}

\begin{proof}
Fix first a finite algebraic stratification
\(\mathscr S=\{S_\alpha\}\) of \(X\) by smooth, locally closed,
pure-dimensional strata, and let \(\D^b_{\mathscr S}(X)\) be the full
subcategory of complexes whose cohomology sheaves are constructible along
\(\mathscr S\).  We construct the \(t\)-structure on this fixed category by
induction on the number of strata.

If there is one stratum of dimension \(d\), all cohomology sheaves are lisse
and the required \(t\)-structure is the ordinary lisse \(t\)-structure shifted
by \(d\).  In general let \(j:U\hookrightarrow X\) be the union of the open
strata and let \(i:Z\hookrightarrow X\) be its closed complement.  On each
connected component of \(U\) use the ordinary lisse \(t\)-structure shifted by
the dimension of that component; on \(Z\) use the structure already obtained
by induction.  Recollement glues them: the two halves are characterized by
\[
 \begin{split}
 K\in{}^p\D_{\mathscr S}^{\leq0}(X)
 &\Longleftrightarrow
 j^*K\in{}^p\D^{\leq0}(U),\quad i^*K\in{}^p\D^{\leq0}(Z),\\
 K\in{}^p\D_{\mathscr S}^{\geq0}(X)
 &\Longleftrightarrow
 j^*K\in{}^p\D^{\geq0}(U),\quad i^!K\in{}^p\D^{\geq0}(Z).
 \end{split}
\]
The two recollement triangles construct the truncation triangle, while the
adjunctions give
\[
 \Hom({}^p\D^{\leq0}_{\mathscr S},
      {}^p\D^{\geq1}_{\mathscr S})=0.
\]
Induction on the strata identifies these glued conditions with the stated
pointwise stalk and costalk inequalities.

If \(\mathscr S'\) refines \(\mathscr S\), the two constructions agree on
\(\D^b_{\mathscr S}(X)\), because the pointwise inequalities do not mention a
stratification.  Every finite family of constructible complexes is adapted to
a common finite refinement.  We may therefore pass from the fixed-stratum
categories to \(\Dbc(X)\); the orthogonality and truncation axioms are checked
inside such a common category.  Boundedness follows from bounded
constructibility and finite dimension.  Finally, the heart of any bounded
\(t\)-structure is abelian, with kernels and cokernels obtained from the
perverse cohomology of a cone.  This is the gluing construction of
BBD, Sections 1.4 and 2.1--2.2, for the dimension function \(-\delta\);
see \cite{BBD}.
\end{proof}

\begin{perversebox}
The proof uses only constructibility, recollement, and the local dimension
function.  It contains no weights and therefore lies safely upstream of the
Weil conjectures.
\end{perversebox}

\section{The dimension shift}

\begin{proposition}[Local systems are perverse after the dimension shift]
\label{prop:lisse-shift-perverse}
Let \(X\) be smooth and pure of dimension \(d\), and let \(\cL\) be lisse.
Then \(\cL[d]\) is perverse.
\end{proposition}

\begin{proof}
For a point \(x\) of dimension \(e\), the stalk of \(\cL[d]\) is concentrated
in degree \(-d\leq-e\), which gives the support inequality.  Purity of the
smooth local ring and absolute cohomological purity identify the transverse
costalk, up to a lisse factor, with a Tate twist placed in the codimension
predicted by \(d-e\).  Consequently it has no cohomology below \(-e\), giving
the cosupport inequality.  Equivalently, choose a smooth stratification and
verify the two inequalities on normal slices.  Both are equalities on the
generic stratum.
\end{proof}

\begin{example}
On a smooth curve \(C\), a lisse sheaf \(\cL[1]\) is perverse.  At a closed
point its stalk is concentrated in degree \(-1\), and its costalk in degree
\(1\).  A general perverse sheaf on a curve may have stalk cohomology in
degrees \(-1,0\) and costalk cohomology in degrees \(0,1\); the degree-zero
boundary terms detect punctual subobjects and quotients.  This is precisely
the information controlled by intermediate extension.
\end{example}

\begin{example}
For a closed point \(i:x\hookrightarrow X\), the skyscraper \(i_*V\) is
perverse with no shift.  More generally a finite map sends perverse sheaves to
perverse sheaves.  Thus objects supported in different dimensions naturally
occupy different ordinary cohomological degrees but the same perverse degree.
\end{example}

\begin{warning}
If \(X\) is singular, the shifted constant sheaf \(\Qbarl[\dim X]\) need not be
perverse.  The canonical replacement is the intersection complex constructed
in \cref{chap:intermediate-extension}.  The equality
\(\IC_X=\Qbarl[\dim X]\) is a smoothness statement, not a definition.
\end{warning}

\section{Recollement in the heart}

Let \(i:Z\hookrightarrow X\) be closed and
\(j:U=X\setminus Z\hookrightarrow X\) its open complement.  The functors
\(j^*=j^!\) and \(i_*=i_!\) are perverse \(t\)-exact.  The four remaining
functors are one-sided exact:
\[
 \begin{array}{c|c}
 \text{functor}&\text{perverse exactness}\\
 \hline
 i^*,\ j_!&\text{right \(t\)-exact}\\
 i^!,\ Rj_*&\text{left \(t\)-exact}.
 \end{array}
\tag{8.2}\label{eq:recollement-exactness}
\]
Here “right \(t\)-exact” means preserving
\({}^p\D^{\leq0}\); “left \(t\)-exact” means preserving
\({}^p\D^{\geq0}\).  This convention will always be stated when used.

Define heart-level functors
\[
 {}^pj_!=\pH^0j_!,\qquad
 {}^pj_*=\pH^0Rj_*,\qquad
 {}^pi^*=\pH^0i^*,\qquad
 {}^pi^!=\pH^0i^!.
\]
For \(\cP\in\Perv(X)\), the recollement triangles yield morphisms in the
perverse heart
\begin{align}
 {}^pj_!j^*\cP&\longrightarrow\cP\longrightarrow
      i_*{}^pi^*\cP\longrightarrow0,
\tag{8.3}\label{eq:perverse-recollement-one}\\
 0&\longrightarrow i_*{}^pi^!\cP\longrightarrow\cP
      \longrightarrow{}^pj_*j^*\cP.
\tag{8.4}\label{eq:perverse-recollement-two}
\end{align}
Together with their neighboring perverse cohomology terms, these are exact
sequences.  They detect subobjects and quotients supported on the boundary.

\section{Verdier duality}

\begin{theorem}[Self-duality of middle perversity]
\label{thm:perverse-verdier-duality}
Verdier duality exchanges the two halves of the middle \(t\)-structure:
\[
 \D_X({}^p\D^{\leq0})={}^p\D^{\geq0},\qquad
 \D_X({}^p\D^{\geq0})={}^p\D^{\leq0}.
\]
It therefore restricts to an exact contravariant equivalence
\[
 \D_X:\Perv(X)^{\op}\xrightarrow{\sim}\Perv(X)
\]
and satisfies
\[
 \D_X\pH^a(\cK)\simeq\pH^{-a}(\D_X\cK).
\tag{8.5}\label{eq:duality-perverse-cohomology}
\]
\end{theorem}

\begin{proof}
Local Verdier duality exchanges a stalk with the dual of the corresponding
costalk.  The inequality \(a>-\delta(x)\) is therefore changed into
\(-a<-\delta(x)\), precisely the other defining condition.  Biduality gives
equality rather than mere inclusion.  General \(t\)-structure formalism then
gives \eqref{eq:duality-perverse-cohomology}.
\end{proof}

This self-duality is the reason for using the middle perversity.  Ordinary
sheaves are not stable under Verdier duality on a positive-dimensional space.

\section{Functorial amplitudes}

The following list is used repeatedly.  It is best viewed as a table of
interfaces, not as permission to suppress shifts.

\begin{proposition}[Basic perverse exactness]
\label{prop:basic-perverse-exactness}
Let all morphisms be separated and of finite type.
\begin{enumerate}
\item If \(f:X\to Y\) is smooth of pure relative dimension \(r\), then
      \(f^*[r]\) and \(f^![-r]=f^*[r](r)\) are \(t\)-exact.
\item If \(f\) is finite, then \(Rf_*=Rf_!\) is \(t\)-exact.
\item If \(f\) is affine, then \(Rf_*\) is right \(t\)-exact and \(Rf_!\) is
      left \(t\)-exact; equivalently, for \(\cP\in\Perv(X)\),
      \[
        \pH^a(Rf_*\cP)=0\ (a>0),\qquad
        \pH^a(Rf_!\cP)=0\ (a<0).
      \]
      If \(f\) is also quasi-finite, both functors are \(t\)-exact.
\item External product is \(t\)-exact on perverse hearts:
      \(\cP\boxtimes\cQ\in\Perv(X\times Y)\).
\end{enumerate}
\end{proposition}

\begin{proof}
The smooth statement follows by comparing dimensions of strata and using
smooth base change for stalks and costalks.  For a finite map, dimensions of
closures are preserved and proper base change computes both conditions
fiberwise.  The affine statement is the perverse form of Artin vanishing.  One
embeds an affine morphism into affine space over the base, filters by generic
hyperplane sections, and applies ordinary Artin vanishing to the successive
normal slices; Verdier duality gives the compact-support half.  For an affine
quasi-finite morphism the relative support dimension is zero, so the two
one-sided bounds meet.  External products follow from the Kunneth formula for
stalks and costalks and addition of dimensions.  The affine and quasi-finite
assertions are BBD 4.1.1--4.1.3; smooth pullback is BBD 4.2.4, and external
products follow from BBD 4.2.7--4.2.8 (or directly from the stalk/costalk
Kunneth formula).
\end{proof}

\begin{corollary}[Artin vanishing for perverse sheaves]
\label{cor:perverse-artin-vanishing}
If \(X\) is affine over an algebraically closed field and
\(\cP\in\Perv(X)\), then
\[
 H^a(X,\cP)=0\quad(a>0),\qquad
 H_c^a(X,\cP)=0\quad(a<0).
\tag{8.6}\label{eq:perverse-artin}
\]
\end{corollary}

\begin{proof}
Apply \cref{prop:basic-perverse-exactness} to the affine structural morphism
\(X\to\Spec k\).  On a point the perverse and ordinary \(t\)-structures agree.
\end{proof}

\section{Perverse spectral sequences without degeneration}

For any cohomological functor \(T:\Dbc(X)\to\mathcal A\), the perverse
truncation tower gives a spectral sequence
\[
 E_2^{a,b}=T^{a}(\pH^b\cK)
 \Longrightarrow T^{a+b}(\cK).
\tag{8.7}\label{eq:perverse-spectral-sequence}
\]
Because \(\cK\) is bounded, the induced filtration is finite.  Every graded
piece of the abutment is a subquotient of an \(E_2\)-term.  This is all that is
needed for weight inequalities: upper and lower weight bounds pass to
subquotients.  No degeneration assertion is used in the direct-image theorem.

\begin{formalization}
The reusable formalization layer should expose:
\begin{enumerate}
\item a dimension function on points and the two pointwise predicates in
      \cref{def:middle-perversity};
\item the glued bounded \(t\)-structure and its heart;
\item perverse truncation triangles and functorial perverse cohomology;
\item exact recollement functors and the one-sided amplitude table;
\item a natural isomorphism expressing Verdier self-duality;
\item the finite filtration attached to
      \eqref{eq:perverse-spectral-sequence}, separated from any degeneration
      theorem.
\end{enumerate}
The point case and the smooth \(d\)-fold shift should be executable tests:
\(\Perv(\Spec k)\) is the ordinary sheaf category and \(\cL[d]\) is perverse.
\end{formalization}

\chapter{Intermediate Extension and Intersection Complexes}
\label{chap:intermediate-extension}

An extension by zero has too many quotients on the boundary; a direct image
has too many boundary subobjects.  The intermediate extension is the unique
extension with neither defect.  This minimality is what makes it compatible
with duality and what singles out intersection cohomology.

Fix an open immersion \(j:U\hookrightarrow X\), with complementary closed
immersion \(i:Z\hookrightarrow X\).

\section{The image in the perverse heart}

For \(\cP\in\Perv(U)\), adjunction gives a canonical morphism
\[
  {}^pj_!\cP=\pH^0(j_!\cP)
  \longrightarrow
  {}^pj_*\cP=\pH^0(Rj_*\cP).
\tag{9.1}\label{eq:standard-to-costandard}
\]

\begin{definition}[Intermediate extension]
\label{def:intermediate-extension}
The intermediate extension of \(\cP\) is the image in the abelian category
\(\Perv(X)\):
\[
  j_{!*}\cP
  =
  \im_{\Perv(X)}
  \bigl({}^pj_!\cP\longrightarrow{}^pj_*\cP\bigr).
\tag{9.2}\label{eq:intermediate-extension}
\]
\end{definition}

\begin{warning}
The symbol “image” in \eqref{eq:intermediate-extension} is meaningful because
\(\Perv(X)\) is abelian.  A triangulated category has no functorial kernels,
cokernels, or images.  Thus \(j_{!*}\) cannot be defined as an image directly
inside \(\Dbc(X)\).
\end{warning}

Restriction to \(U\) is exact and sends both sides of
\eqref{eq:standard-to-costandard} to \(\cP\).  Consequently
\[
  j^*j_{!*}\cP\simeq\cP.
\tag{9.3}\label{eq:jstar-jmiddle}
\]

\begin{theorem}[Minimal-extension characterization]
\label{thm:minimal-extension-characterization}
Let \(\cQ\in\Perv(X)\) and suppose \(j^*\cQ\simeq\cP\).  Then
\(\cQ\simeq j_{!*}\cP\) if and only if
\begin{enumerate}
\item \(\cQ\) has no nonzero quotient supported on \(Z\);
\item \(\cQ\) has no nonzero subobject supported on \(Z\).
\end{enumerate}
Equivalently,
\[
  {}^pi^*\cQ=0,\qquad {}^pi^!\cQ=0.
\tag{9.4}\label{eq:middle-boundary-vanishing}
\]
The isomorphism extending the identity of \(\cP\) is unique.
\end{theorem}

\begin{proof}
The map \({}^pj_!\cP\to\cQ\) is forced by adjunction.  Since its restriction
to \(U\) is an isomorphism, its cokernel is supported on \(Z\); condition (1)
makes it surjective.  Dually, the adjunction map
\(\cQ\to{}^pj_*\cP\) has kernel supported on \(Z\); condition (2) makes it
injective.  Their composite is the canonical map
\eqref{eq:standard-to-costandard}, so \(\cQ\) identifies with its image.

Conversely, any quotient of that image supported on \(Z\) would be a quotient
of \({}^pj_!\cP\) whose restriction to \(U\) is zero.  Adjunction shows such a
quotient is zero.  The subobject assertion follows dually from
\({}^pj_*\cP\).  Recollement identifies maximal boundary quotients and
subobjects with \(i_*{}^pi^*\cQ\) and \(i_*{}^pi^!\cQ\), proving
\eqref{eq:middle-boundary-vanishing}.  The same argument applied to the graph
of two extensions gives uniqueness.
\end{proof}

\begin{corollary}[Transitivity]
\label{cor:intermediate-extension-transitivity}
If \(V\xrightarrow{k}U\xrightarrow{j}X\) are open immersions, then
\[
  (j\circ k)_{!*}\simeq j_{!*}k_{!*}.
\]
\end{corollary}

\begin{proof}
Both sides restrict to the original perverse sheaf on \(V\), and neither has a
subobject or quotient supported on \(X\setminus V\).  Apply
\cref{thm:minimal-extension-characterization}.
\end{proof}

\begin{warning}
The functor \(j_{!*}\) preserves monomorphisms and epimorphisms: a kernel or
cokernel created at the boundary would contradict the minimal-extension
property.  It is nevertheless not exact in general, because middle exactness
of a short exact sequence can fail.  What is canonical is the image
construction in the perverse heart.
\end{warning}

\section{Compatibility with Verdier duality}

\begin{proposition}[Self-duality of intermediate extension]
\label{prop:duality-intermediate-extension}
There is a canonical isomorphism
\[
  \D_X(j_{!*}\cP)\simeq j_{!*}(\D_U\cP).
\tag{9.5}\label{eq:duality-jmiddle}
\]
\end{proposition}

\begin{proof}
Verdier duality is exact on the perverse heart and exchanges
\({}^pj_!\) with \({}^pj_*\).  It therefore sends the image of
\({}^pj_!\cP\to{}^pj_*\cP\) to the coimage of the dual map.  In an abelian
category image and coimage agree.  Alternatively, duality interchanges
boundary subobjects and boundary quotients, so the result follows immediately
from \cref{thm:minimal-extension-characterization}.
\end{proof}

\section{Finite length and simple perverse sheaves}

\begin{theorem}[BBD finite-length theorem]
\label{thm:perverse-finite-length}
The category \(\Perv(X,\Qbarl)\) is both Noetherian and Artinian.  Every
perverse sheaf has finite length.
\end{theorem}

\begin{proof}
Induct on a stratification of \(X\).  On a smooth connected stratum, a
perverse sheaf is generically a shifted lisse sheaf, and its subobjects give
subrepresentations of a finite-dimensional continuous representation; these
satisfy both chain conditions.  Shrink to make the generic constituent lisse.
The kernel and cokernel of the resulting standard or costandard extension are
supported on the smaller closed complement, where induction applies.  The
recollement exact sequences then prevent an infinite ascending or descending
chain on \(X\).  This is the devissage in BBD 4.3.1.
\end{proof}

\begin{theorem}[Classification of simple objects]
\label{thm:simple-perverse-classification}
Let \(S\subset X\) be a smooth connected locally closed subvariety of
dimension \(d\), let \(\cL\) be an irreducible lisse \(\Qbarl\)-sheaf on
\(S\), and let \(j:S\hookrightarrow\overline S\) and
\(i:\overline S\hookrightarrow X\).  Then
\[
  i_*j_{!*}(\cL[d])
\tag{9.6}\label{eq:simple-IC-object}
\]
is a simple perverse sheaf on \(X\).  Every simple perverse sheaf is obtained
uniquely in this way, up to shrinking \(S\) inside a dense open of its closure
and identifying the resulting local systems.
\end{theorem}

\begin{proof}
Any nonzero subobject or quotient of \(j_{!*}\cL[d]\) restricts to a nonzero
subobject or quotient of \(\cL[d]\).  Irreducibility makes that restriction
all of \(\cL[d]\), and the kernel or cokernel is then supported on the
boundary; minimality makes it zero.  Thus \eqref{eq:simple-IC-object} is
simple.

Conversely, let \(\cP\) be simple and choose a maximal-dimensional stratum
\(S\) on which it is nonzero and lisse after the dimension shift.  Simplicity
forces the generic local system to be irreducible.  The adjunction morphisms
factor \(\cP\) between the standard and costandard extensions, hence produce a
nonzero map \(j_{!*}\cL[d]\to\cP\).  A nonzero map between simple objects is
an isomorphism.  The support and generic local system recover
\(\overline S\) and \(\cL\).
\end{proof}

\section{Intersection complexes}

\begin{definition}[Intersection complex]
\label{def:intersection-complex}
Let \(X\) be irreducible of dimension \(d\), let
\(j:U\hookrightarrow X\) be a smooth dense open, and let \(\cL\) be lisse on
\(U\).  The perverse-normalized intersection complex is
\[
  \IC_X(\cL)=j_{!*}(\cL[d]).
\tag{9.7}\label{eq:IC-definition}
\]
For the constant local system write \(\IC_X=\IC_X(\Qbarl)\).
\end{definition}

Transitivity shows that \(\IC_X(\cL)\) is independent of replacing \(U\) by a
smaller dense smooth open.  When \(\cL\) is irreducible it is a simple
perverse sheaf.  If \(X\) is reducible, define the intersection complex as
the direct sum over irreducible components, with the desired coefficient
local systems on dense smooth opens.

\begin{proposition}[Smooth normalization and duality]
\label{prop:IC-smooth-duality}
\begin{enumerate}
\item If \(X\) is smooth of pure dimension \(d\) and \(\cL\) is lisse on
      all of \(X\), then
      \[
        \IC_X(\cL)=\cL[d].
      \tag{9.8}\label{eq:IC-smooth}
      \]
\item For \(U\) smooth of dimension \(d\),
      \[
        \D_X\IC_X(\cL)
        \simeq \IC_X(\cL^\vee(d)).
      \tag{9.9}\label{eq:IC-duality}
      \]
\end{enumerate}
\end{proposition}

\begin{proof}
For the first statement take \(j=\id_X\).  For the second, smooth Poincare
duality gives
\[
  \D_U(\cL[d])\simeq\cL^\vee(d)[d].
\]
Apply \cref{prop:duality-intermediate-extension}.
\end{proof}

There are two common normalizations.  In these notes \(\IC_X(\cL)\) is
perverse, so on a smooth \(d\)-fold it is \(\cL[d]\).  Some authors shift back
and call \(\IC_X(\cL)[-d]\) the intersection complex.  Every weight formula
below uses the perverse normalization \eqref{eq:IC-definition}.

\section{The curve model}

Let \(C\) be a smooth connected curve and \(j:U\hookrightarrow C\) a dense
open.  For a lisse sheaf \(\cL\) on \(U\),
\[
  \IC_C(\cL)=j_{!*}(\cL[1])\simeq (j_*\cL)[1],
\tag{9.10}\label{eq:curve-IC}
\]
where \(j_*\cL\) on the right is the ordinary sheaf direct image.  At a
boundary point \(x\), its stalk is the inertia-invariant space
\((\cL_{\bar\eta})^{I_x}\).  Formula \eqref{eq:curve-IC} follows by checking
the support and cosupport conditions and the absence of punctual subobjects
and quotients.

Consequently
\[
  \mathbb H^r(C,\IC_C(\cL))
  =H^{r+1}(C,j_*\cL).
\tag{9.11}\label{eq:curve-IH}
\]
The degree-zero intersection cohomology
\(\mathbb H^0(C,\IC_C(\cL))=H^1(C,j_*\cL)\) is the parabolic cohomology that
appears in the curve weight theorem.

\begin{example}[A nodal curve]
Let \(X\) be an irreducible nodal proper curve with normalization
\(\nu:\widetilde X\to X\).  Since \(\nu\) is finite,
\[
  \IC_X\simeq\nu_*\Qbarl[1].
\]
At the node \(x\),
\[
 \cH^{-1}(\nu_*\Qbarl[1])_x\simeq\Qbarl^2,
 \qquad
 \cH^{-1}\bigl((\Qbarl)_X[1]\bigr)_x\simeq\Qbarl.
\]
Thus the shifted constant complex identifies the two branches through the
diagonal gluing condition, whereas the intersection complex retains the two
normalization branches separately.  The latter is the self-dual pure object
appropriate to the singular curve.
\end{example}

\section{Where weights will enter}

Nothing in this chapter proves that \(\IC_X(\cL)\) is pure.  Once weights are
constructed, the correct shift will be
\[
  \cL\text{ pure of weight }w
  \quad\Longrightarrow\quad
  \IC_X(\cL)\text{ pure of weight }w+d.
\tag{9.12}\label{eq:IC-weight-preview}
\]
The dimension term \(d\) is not cosmetic: it is forced by the convention that
ordinary cohomology sheaves of a pure complex of weight \(v\) have weights
\(v+a\) in degree \(a\).  The proof of \eqref{eq:IC-weight-preview} is delayed
until \cref{chap:perverse-weights}; proper cohomological purity is delayed
further, until the Deligne--Laumon direct-image theorem has been proved.

\begin{formalization}
The minimum reusable API consists of:
\begin{enumerate}
\item heart-level standard and costandard extensions;
\item the canonical natural transformation
      \({}^pj_!\Rightarrow{}^pj_*\);
\item \(j_{!*}\) as its abelian image;
\item the equivalent no-boundary-subobject/no-boundary-quotient predicate;
\item transitivity and Verdier compatibility;
\item finite length and the classification of simple objects;
\item \(\IC_X(\cL)\) with its smooth and curve normalization tests.
\end{enumerate}
Choice independence should be an isomorphism characterized by the universal
property, not a definitional equality tied to a selected stratification.
\end{formalization}

\chapter{Integral Perverse Lattices: the \texorpdfstring{\(p/p^+\)}{p/p+} Formalism}
\label{chap:integral-perverse-lattices}
\label{chap:integral-perverse}
\providecommand{\op}{\mathrm{op}}
\providecommand{\im}{\operatorname{im}}

The pro-etale site removes one artificial difficulty from integral
\(\ell\)-adic cohomology: the coefficient ring \(\Zl\), and more generally the
valuation ring of a finite extension of \(\Ql\), is an honest sheaf.  It does
\emph{not} make integral Verdier duality \(t\)-exact for the usual perverse
\(t\)-structure.  Torsion produces a one-degree defect already on a point.
The correct integral formalism therefore contains two adjacent perverse
\(t\)-structures, denoted \(p\) and \(p^+\).  They become identical after
inverting \(\ell\), and Verdier duality exchanges them before inverting
\(\ell\).

This distinction is not optional bookkeeping.  Ignoring it leads to false
claims that integral intermediate extension is self-dual, that reduction
modulo \(\ell\) is perverse \(t\)-exact, or that a rational perverse sheaf has
a canonical integral model.  This chapter gives the exact replacement for
each of those claims.

\section{Coefficient category and notation}

Let \(\mathcal O\) be the valuation ring of a finite extension \(E/\Ql\), let
\(\varpi\) be a uniformizer, and let \(k_{\mathcal O}=\mathcal O/\varpi\).
The special case \((\mathcal O,E)=(\Zl,\Ql)\) is the one needed for the Weil
conjectures.  For a finite-type scheme \(X\) over a field of characteristic
different from \(\ell\), write
\[
  D_{\mathcal O}(X)=D^b_{\cons}(X_{\proet},\mathcal O)
\]
for the bounded constructible, derived-complete coefficient category of
finite Tor-amplitude constructed in the preceding pro-etale chapters.  The
finite Tor-amplitude condition is retained whenever Verdier biduality is
used.  Put
\[
  K_E:=K\otimes_{\mathcal O}^{\mathrm L}E,\qquad
  K_k:=K\otimes_{\mathcal O}^{\mathrm L}k_{\mathcal O}.
\]
Let
\[
  \D_{X,\mathcal O}K
    :=\R\mathcal Hom_{\mathcal O}(K,\omega_{X,\mathcal O})
\]
be Verdier duality.  On the constructible finite-Tor subcategory it is an
anti-equivalence and the biduality morphism is an isomorphism.

The middle perverse \(t\)-structure \(p\) is defined by the same geometric
stalk and costalk inequalities as over \(E\); see
\cref{chap:perverse-t-structure}.  Its heart is denoted
\(\Perv_p(X,\mathcal O)\), and its cohomology functors by
\({}^p\!H^i\).  An object \(T\) of this heart is \(\varpi\)-power torsion if
\(\varpi^nT=0\) for some \(n\).  It is torsion-free if multiplication by
\(\varpi\) is a monomorphism in the perverse heart.

\begin{warning}[Derived completeness is not the perverse issue]
Derived \(\varpi\)-completeness makes inverse limits and the coefficient
sheaf \(\mathcal O\) behave correctly.  The \(p/p^+\) distinction comes from
the homological dimension-one ring \(\mathcal O\), not from a defect of the
pro-etale site.  Passing from the etale site to the pro-etale site therefore
does not erase this distinction.
\end{warning}

\section{The obstruction is visible on a point}

Consider \(X=\Spec\bar k\).  The perverse \(t\)-structure is the ordinary
cohomological \(t\)-structure.  The free module \(\mathcal O\) is fixed by
\(\R\Hom_{\mathcal O}(-,\mathcal O)\), whereas the resolution
\[
  0\longrightarrow\mathcal O\xrightarrow{\ \varpi\ }\mathcal O
    \longrightarrow k_{\mathcal O}\longrightarrow0
\]
gives
\[
  \R\Hom_{\mathcal O}(k_{\mathcal O},\mathcal O)\simeq
  k_{\mathcal O}[-1].
\tag{\thechapter.1}\label{eq:point-torsion-dual-shift}
\]
Thus duality sends a torsion object in degree \(0\) to a torsion object whose
only cohomology lies in degree \(1\).  No self-dual integral \(t\)-structure
can have both free and torsion modules in its heart in degree \(0\).

The remedy is the following tilted \(t\)-structure on a point:
\begin{align*}
 {}^{+}D^{\leq0}
   &=\{K:H^i(K)=0\ (i>1),\ H^1(K)\text{ is torsion}\},\\
 {}^{+}D^{\geq0}
   &=\{K:H^i(K)=0\ (i<0),\ H^0(K)\text{ is torsion-free}\}.
\end{align*}
Its heart consists of two-term objects with torsion-free \(H^0\) and
torsion \(H^1\).  Formula \eqref{eq:point-torsion-dual-shift} says precisely
that integral duality exchanges the ordinary heart with this tilted heart.

\section{The torsion pair and the \texorpdfstring{\(p^+\)}{p+}-structure}

The full subcategories
\[
\begin{aligned}
  \mathsf T_X
    &=\{P\in\Perv_p(X,\mathcal O):P\text{ is }\varpi\text{-power torsion}\},\\
  \mathsf F_X
    &=\{P\in\Perv_p(X,\mathcal O):\varpi:P\to P\text{ is monic}\}.
\end{aligned}
\]
form a torsion pair in \(\Perv_p(X,\mathcal O)\).  Indeed, the ascending
chain of kernels of \(\varpi^n:P\to P\) stabilizes by constructibility and
Noetherianity.  Its stable value \(P_{\mathrm{tor}}\) is the maximal torsion
subobject, and
\[
  0\longrightarrow P_{\mathrm{tor}}\longrightarrow P
    \longrightarrow P_{\mathrm{tf}}\longrightarrow0
\tag{\thechapter.2}\label{eq:perverse-torsion-decomposition}
\]
has \(P_{\mathrm{tf}}\in\mathsf F_X\).  Moreover
\(\Hom(\mathsf T_X,\mathsf F_X)=0\).

\begin{theorem}[The adjacent integral perverse structure]
\label{thm:p-pplus-tilt}
There is a bounded \(t\)-structure \(p^+\) on \(D_{\mathcal O}(X)\), the
right tilt of \(p\) at \((\mathsf T_X,\mathsf F_X)\), characterized by
\begin{equation}
\begin{aligned}
 {}^{p^+}D^{\leq0}
   &=\{K\in{}^pD^{\leq1}:{}^pH^1(K)\in\mathsf T_X\},\\
 {}^{p^+}D^{\geq0}
   &=\{K\in{}^pD^{\geq0}:{}^pH^0(K)\in\mathsf F_X\}.
\end{aligned}
\label{eq:pplus-tilt-formulas}
\end{equation}
Consequently its heart consists exactly of the objects \(K\) for which
\({}^pH^i(K)=0\) outside \(i=0,1\), with
\({}^pH^0(K)\) torsion-free and \({}^pH^1(K)\) torsion.
\end{theorem}

\begin{proof}
For any torsion pair \((\mathsf T,\mathsf F)\) in the heart of a bounded
\(t\)-structure, the two classes in \eqref{eq:pplus-tilt-formulas} satisfy
orthogonality and provide a truncation triangle.  Explicitly, start with the
\(p\)-truncation triangle of \(K\), split \({}^pH^1(K)\) by
\eqref{eq:perverse-torsion-decomposition}, and take the appropriate cone.
The resulting negative piece has only a torsion \({}^pH^1\); the positive
piece has torsion-free \({}^pH^0\).  This constructs the \(p^+\)-truncation
functors and proves boundedness.

On every smooth stratum the construction is the pointwise tilt above,
shifted by the dimension of the stratum.  Recollement therefore glues the
local tilted structures, and refinement of the stratification does not
change them.  This is the Dedekind-coefficient construction of
\cite[\S3.3.4]{BBD}, now realized inside the pro-etale coefficient category.
\end{proof}

\begin{remark}
The inclusions
\[
  {}^pD^{\leq0}\subset{}^{p^+}D^{\leq0}\subset{}^pD^{\leq1},\qquad
  {}^{p^+}D^{\geq0}\subset{}^pD^{\geq0}\subset{}^{p^+}D^{\geq-1}
\]
make the phrase adjacent \(t\)-structures literal.  Their difference is
entirely \(\varpi\)-primary torsion.
\end{remark}

\section{Verdier duality exchanges \texorpdfstring{\(p\) and \(p^+\)}{p and p+}}

\begin{theorem}[Integral perverse duality]
\label{thm:integral-verdier-exchange}
On constructible objects of finite Tor-amplitude, Verdier duality satisfies
\begin{equation}
\begin{aligned}
  \D_{X,\mathcal O}({}^pD^{\leq0})
    &={} ^{p^+}D^{\geq0},&
  \D_{X,\mathcal O}({}^pD^{\geq0})
    &={} ^{p^+}D^{\leq0},\\
  \D_{X,\mathcal O}({}^{p^+}D^{\leq0})
    &={} ^pD^{\geq0},&
  \D_{X,\mathcal O}({}^{p^+}D^{\geq0})
    &={} ^pD^{\leq0}.
\end{aligned}
\label{eq:duality-exchanges-p-pplus}
\end{equation}
In particular,
\[
  \D_{X,\mathcal O}:\Perv_p(X,\mathcal O)^{\op}\xrightarrow{\sim}
  \Perv_{p^+}(X,\mathcal O).
\tag{\thechapter.5}\label{eq:duality-hearts-p-pplus}
\]
\end{theorem}

\begin{proof}
Choose a finite smooth stratification adapted to \(K\).  Under the
six-operation hypotheses fixed in Chapter~4, constructibility and finite
Tor-amplitude make stalks and costalks dualizable.  Smooth absolute
cohomological purity identifies the dualizing complex on a smooth stratum of
dimension \(d\), with the shift and twist dictated by the standing
normalization.  Thus the support and cosupport inequalities are exchanged by
Verdier duality.  On the coefficient complex the only integral defect is the
point calculation \eqref{eq:point-torsion-dual-shift}: finite free modules
remain in degree \(0\), whereas torsion modules move by one degree.
Verdier duality also exchanges \(j_!\) with \(j_*\).  Induction on the number
of strata and the two recollement triangles therefore glue the assertion
over \(X\); biduality gives the second line from the first.  This is the
middle-perversity case of \cite[\S3.3.3--3.3.4]{BBD}.
\end{proof}

\begin{powerbox}
Over \(E\), the middle perverse structure is self-dual.  Integrally, the
correct statement is not that \(\D\) preserves \(p\), but that \(\D\)
exchanges \(p\) and \(p^+\).  It records exactly where torsion can occur in a
dual complex.
\end{powerbox}

\section{Rationalization collapses the two structures}

Because \(E=\mathcal O[1/\varpi]\) is flat, rationalization is a triangulated
functor and kills precisely the \(\varpi\)-power-torsion constructible
objects.  It is compatible with restriction to geometric points and hence
with all support conditions defining perversity.

\begin{theorem}[Rationalization is \(t\)-exact twice]
\label{thm:rationalization-perverse}
The functor
\[
  (-)_E:D_{\mathcal O}(X)\longrightarrow D^b_c(X_{\proet},E)
\]
is \(t\)-exact from the \(p\)-structure to the middle perverse structure
over \(E\), and also from the \(p^+\)-structure to that same structure.  For
every \(K\) and every \(i\) there are functorial isomorphisms
\[
  ({}^pH^iK)_E\simeq{}^pH^i(K_E),\qquad
  ({}^{p^+}H^iK)_E\simeq{}^pH^i(K_E).
\tag{\thechapter.6}\label{eq:rationalization-perverse-cohomology}
\]
The rationalized \(p\)- and \(p^+\)-truncation triangles are canonically
identified.
\end{theorem}

\begin{proof}
For \(p\), flatness commutes with stalks, costalks, and ordinary cohomology,
so the defining vanishing inequalities are preserved.  For \(p^+\), use
\eqref{eq:pplus-tilt-formulas}: the possible \({}^pH^1\) is torsion and
vanishes after tensoring with \(E\), while the torsion-free condition on
\({}^pH^0\) introduces no additional rational cohomology.  Thus the two
integral aisles and coaisles have the same rational images.  A \(t\)-exact
functor commutes with cohomology and truncation, giving
\eqref{eq:rationalization-perverse-cohomology}.
\end{proof}

\begin{proposition}[Duality and the six operations after inverting
\(\varpi\)]
\label{prop:rationalization-six-duality}
For constructible finite-Tor objects, the canonical comparison maps give
\[
  (\D_{X,\mathcal O}K)_E\simeq\D_{X,E}(K_E).
\tag{\thechapter.7}\label{eq:rationalization-duality}
\]
Likewise rationalization commutes with \(f^*,Rf_!,Rf_*,f^!\), tensor
product, and internal Hom whenever those functors are defined on the
constructible categories in question.
\end{proposition}

\begin{proof}
This is the coefficient-change theorem for constructible complete
coefficients, not a formal claim about commuting localization with an
inverse limit.  The equivalence between constructible \(E\)-complexes and
the localization of constructible \(\mathcal O\)-complexes, together with
its tensor and six-operation compatibilities, is
\cite[Proposition~6.8.14 and Remark~6.8.15]{BhattScholzeProEtale}; see also
\cite[Theorem~7.7]{HRSConstructible}.  These comparison theorems give the
statements for \(f^*,Rf_!,Rf_*\), tensor product, and internal Hom.
The \(f^!\)-comparison follows from the adjunction with \(Rf_!\), equivalently
from the duality formula on dualizable constructible objects.  Applying the
internal-Hom comparison to the dualizing object gives
\eqref{eq:rationalization-duality}.
\end{proof}

\begin{warning}
The converse to \cref{thm:rationalization-perverse} is false without an
integral boundedness condition: rationalization cannot see a torsion object
or decide whether it lies in perverse degree \(0\) or \(1\).  Consequently
one must never prove an integral \(t\)-exactness statement merely by
tensoring with \(E\).
\end{warning}

\section{Perverse lattices}

\begin{definition}[Perverse lattice]
\label{def:perverse-lattice}
Let \(P_E\in\Perv(X,E)\).  A \(p\)-perverse lattice in \(P_E\) is an object
\(P_{\mathcal O}\in\Perv_p(X,\mathcal O)\) such that multiplication by
\(\varpi\) is monic and \(P_{\mathcal O,E}\simeq P_E\).  Two lattices are
commensurable if they become comparable after multiplying by powers of
\(\varpi\).
\end{definition}

\begin{lemma}[A torsion-free \(p\)-object lies in both hearts]
\label{lem:lattice-in-both-hearts}
If \(P_{\mathcal O}\) is a perverse lattice, then it belongs to both
\(\Perv_p(X,\mathcal O)\) and \(\Perv_{p^+}(X,\mathcal O)\).
\end{lemma}

\begin{proof}
Its \(p\)-cohomology is concentrated in degree \(0\), and that degree is
torsion-free.  Both conditions in \eqref{eq:pplus-tilt-formulas} follow.
\end{proof}

If \(U\) is smooth of pure dimension \(d\) and \(L_E\) is a lisse
\(E\)-sheaf, the image of the geometric fundamental group is compact.  After
a finite extension of \(E\), it stabilizes a finite free
\(\mathcal O\)-lattice \(L_{\mathcal O}\).  Then
\(L_{\mathcal O}[d]\) is a perverse lattice in \(L_E[d]\), common to the two
hearts.  The choice is not canonical, but any two choices are commensurable.

\section{Two integral intermediate extensions, one rational extension}

Let \(j:U\hookrightarrow X\) be an open immersion and let
\(P_{\mathcal O}\) lie in both integral hearts.  Form the images in the two
different abelian categories:
\begin{equation}
\begin{aligned}
 {}^p j_{!*}P_{\mathcal O}
   &:=\im_{\Perv_p}\bigl({}^pH^0Rj_!P_{\mathcal O}\to
                         {}^pH^0Rj_*P_{\mathcal O}\bigr),\\
 {}^{p^+}j_{!*}P_{\mathcal O}
   &:=\im_{\Perv_{p^+}}\bigl({}^{p^+}H^0Rj_!P_{\mathcal O}\to
                             {}^{p^+}H^0Rj_*P_{\mathcal O}\bigr).
\end{aligned}
\label{eq:two-integral-jmiddle}
\end{equation}

\begin{theorem}[Integral and rational middle extensions]
\label{thm:two-integral-intermediate-extensions}
The two objects in \eqref{eq:two-integral-jmiddle} need not be isomorphic
over \(\mathcal O\).  Their images in the Verdier localization by
\(\varpi\)-primary torsion are canonically isomorphic; equivalently, their
rationalizations are canonically isomorphic:
\[
  ({}^p j_{!*}P_{\mathcal O})_E
  \simeq j_{!*}(P_E)\simeq
  ({}^{p^+}j_{!*}P_{\mathcal O})_E.
\tag{\thechapter.9}\label{eq:two-jmiddle-rationalize}
\]
There is in general no canonical integral comparison morphism between the
two lattices.  If one chooses an integral comparison whose rationalization
is the displayed isomorphism, the cohomology objects of its cone are
\(\varpi\)-primary torsion.  Verdier duality exchanges the two constructions:
\[
  \D_{X,\mathcal O}({}^p j_{!*}P_{\mathcal O})\simeq
  {}^{p^+}j_{!*}(\D_{U,\mathcal O}P_{\mathcal O}).
\tag{\thechapter.10}\label{eq:integral-jmiddle-duality}
\]
\end{theorem}

\begin{proof}
By \cref{thm:rationalization-perverse}, rationalization is exact on either
heart and sends the two maps defining \eqref{eq:two-integral-jmiddle} to the
same morphism
\({}^pH^0Rj_!P_E\to{}^pH^0Rj_*P_E\).  Exact functors preserve images, which
proves \eqref{eq:two-jmiddle-rationalize} and the canonical identification in
the Verdier quotient.  A chosen integral comparison becomes an isomorphism
after tensoring with \(E\); the localization criterion therefore says
precisely that every cohomology object of its cone is torsion.

Duality exchanges \(Rj_!\) and \(Rj_*\), reverses the defining arrow, and by
\cref{thm:integral-verdier-exchange} exchanges the two hearts.  The image of
an arrow in one heart is sent to the image of the reversed dual arrow in the
other heart, giving \eqref{eq:integral-jmiddle-duality}.  Equivalently, use
the minimal-extension characterization: no nonzero boundary subobject or
quotient is carried by duality to the corresponding condition in the other
heart.  Compare \cite[\S1.4.22--1.4.25 and \S3.3.4]{BBD}.
\end{proof}

\begin{warning}[There is no canonical integral intersection complex]
For \(X\) irreducible with smooth dense \(j:U\hookrightarrow X\), a free
lattice \(L_{\mathcal O}\subset L_E\) produces
\[
  {}^p\IC_X(L_{\mathcal O})={}^p j_{!*}L_{\mathcal O}[d],\qquad
  {}^{p^+}\IC_X(L_{\mathcal O})={}
  ^{p^+}j_{!*}L_{\mathcal O}[d].
\]
Both rationalize to \(\IC_X(L_E)\), but they may differ by torsion.  Integral
duality exchanges them.  The rational intersection complex is canonical;
an integral lattice in it is additional data.
\end{warning}

\section{Reduction modulo \texorpdfstring{\(\varpi\)}{the uniformizer}}

Unlike rationalization, reduction is not flat.  The triangle
\[
  K\xrightarrow{\ \varpi\ }K\longrightarrow K_k\longrightarrow K[1]
\tag{\thechapter.11}\label{eq:derived-reduction-triangle}
\]
gives, for \(P\in\Perv_p(X,\mathcal O)\), an exact sequence
\[
  0\longrightarrow{}^pH^{-1}(P_k)\longrightarrow P
  \xrightarrow{\ \varpi\ }P\longrightarrow{}^pH^0(P_k)\longrightarrow0.
\tag{\thechapter.12}\label{eq:perverse-reduction-tor}
\]
Thus
\[
 {}^pH^{-1}(P_k)=P[\varpi],\qquad
 {}^pH^0(P_k)=P/\varpi P.
\]
Reduction of a perverse lattice is perverse because the first term vanishes.
Reduction of an arbitrary integral perverse sheaf may occupy two adjacent
perverse degrees.  This is the modular shadow of the \(p/p^+\) distinction.

\begin{proposition}[Safe lattice route to rational statements]
\label{prop:safe-lattice-route}
For the Weil proof, the following order is valid:
\begin{enumerate}
\item construct the six operations and Verdier duality over \(\mathcal O\);
\item select lattices only where compactness or a geometric construction
      provides them;
\item rationalize to \(E\), where the middle perverse structure is
      self-dual and weights are defined;
\item prove purity and the Weil bounds over \(E\);
\item return to the chosen lattices for integrality arguments.
\end{enumerate}
No step requires choosing one of the two integral intermediate extensions as
the canonical rational one.
\end{proposition}

\begin{proof}
Steps 1--3 are justified by
\cref{thm:integral-verdier-exchange,thm:rationalization-perverse,prop:rationalization-six-duality}.
Weights concern eigenvalues in \(\Qbarl\), so they are rational notions.
Both integral intermediate extensions have the same rationalization by
\cref{thm:two-integral-intermediate-extensions}.  Finally, a
Frobenius-stable lattice makes the characteristic polynomial integral by
Cayley--Hamilton; that later algebraic argument depends on the lattice, not
on a false integral self-duality assertion.
\end{proof}

\begin{auditbox}[title={Dependency firewall for integral coefficients}]
The following implications are forbidden:
\[
\begin{aligned}
 \text{pro-etale}&\Longrightarrow\text{integral self-duality of \(p\)},\\
 K_E\text{ perverse}&\Longrightarrow K\text{ perverse},\\
 P\text{ integral perverse}&\Longrightarrow P_k\text{ perverse},\\
 {}^p j_{!*}P&={} ^{p^+}j_{!*}P.
\end{aligned}
\]
The valid replacements are respectively
\cref{thm:integral-verdier-exchange,thm:rationalization-perverse},
\eqref{eq:perverse-reduction-tor}, and
\eqref{eq:two-jmiddle-rationalize}.
\end{auditbox}

\section{Formalization interface}

\begin{formalization}
A reusable implementation should expose the following objects separately.
\begin{enumerate}
\item A coefficient record \((\mathcal O,\varpi,E,k_{\mathcal O})\), with
      flat localization \(\mathcal O\to E\) and a length-one resolution of
      \(k_{\mathcal O}\).
\item The torsion pair \((\mathsf T_X,\mathsf F_X)\) in the \(p\)-heart.
\item A generic right-tilt construction producing \(p^+\) and proving
      \eqref{eq:pplus-tilt-formulas}.
\item An anti-\(t\)-exactness theorem for Verdier duality between two named
      \(t\)-structures, not an overloaded claim that it preserves one.
\item Two \(t\)-exactness instances for rationalization, both landing in the
      same rational perverse structure.
\item Separate definitions of \({}^p j_{!*}\) and
      \({}^{p^+}j_{!*}\), followed by a theorem identifying only their
      rationalizations.
\end{enumerate}
Regression tests should include the point object
\(\mathcal O/\varpi\): its dual must occur in degree \(1\), and its derived
reduction must exhibit the Tor term in degree \(-1\).
\end{formalization}


\chapter{The Perverse Engine of the Weil Proof}
\label{chap:perverse-engine}
\providecommand{\pr}{\operatorname{pr}}
\providecommand{\FT}{\operatorname{FT}}

Perverse sheaves are not an ornamental reformulation of the final
cohomology groups.  They perform four operations that ordinary sheaves do
not perform cleanly:
\begin{enumerate}
\item they turn extension across a singular boundary into the canonical
      image \(j_{!*}\) inside an abelian category;
\item they put affine cohomology into one-sided degrees by Artin vanishing;
\item they make the Fourier--Deligne transform an exact equivalence of
      abelian categories;
\item after the weight theorem, they convert upper weight bounds and
      duality into proper purity, intersection-cohomology purity, and the
      decomposition theorem.
\end{enumerate}
This chapter gives the entire theorem chain in the order in which it may be
used.  In particular, purity of \(j_{!*}\) and the decomposition theorem are
\emph{outputs}, not inputs, of the proof of the Weil bounds.

\begin{powerbox}[title={The proof in one display}]
\[
\begin{aligned}
 &j_{!*}+\text{Artin vanishing}
      +\text{Fourier support-forgetting}\\
 &\Longrightarrow
      \FT_\psi:\Perv(\A^1)\xrightarrow{\sim}\Perv(\A^1)\\
 &\Longrightarrow \text{the curve weight estimate}\\
 &\Longrightarrow Rf_!\text{ preserves upper weights}\\
 &\xRightarrow[\ f\text{ proper}\ ]{\text{Verdier duality}}
      Rf_*\text{ preserves purity}\\
 &\Longrightarrow H^i(X_{\Fbar},\Ql)
      \text{ is pure of weight }i\\
 &\Longrightarrow \text{the Riemann hypothesis for }Z(X,T)\\
 &\Longrightarrow
   \begin{gathered}
      j_{!*}\text{ preserves purity},\quad \IC\text{ is pure},\\
      \text{decomposition and relative hard Lefschetz}.
   \end{gathered}
\end{aligned}
\]
\end{powerbox}

\section{The noncircular starting package}

We work with \(E=\Qbarl\), or a finite extension large enough to contain the
eigenvalues under consideration.  The integral category of
\cref{chap:integral-perverse-lattices} supplies lattices and the six
operations; \cref{thm:rationalization-perverse} places us in the self-dual
rational perverse category.  At the beginning of the weight proof the
following facts are available.

\begin{enumerate}
\item The middle perverse \(t\)-structure, its recollement triangles, and
      Verdier self-duality; these are geometric and contain no weight
      theorem.
\item For an open immersion \(j:U\hookrightarrow X\),
      \[
        j_{!*}P=\im\bigl({}^pH^0Rj_!P\longrightarrow{}^pH^0Rj_*P\bigr).
      \]
      It is the unique extension having no nonzero subobject or quotient
      supported on \(X-U\).  These are BBD 1.4.22 and 1.4.25.
\item Every simple perverse sheaf is
      \(j_{!*}(L[\dim U])\) for an irreducible lisse \(L\) on a smooth
      irreducible stratum; this is BBD 4.3.1.
\item Perverse Artin vanishing, BBD 4.1.1--4.1.2.
\item The formal six-operation identities, the trace formula, and local
      monodromy theory, but not the general direct-image weight theorem.
\end{enumerate}

\begin{auditbox}[title={What is not yet available}]
At this stage one may not use
\[
  j_{!*}(\text{pure})\text{ is pure},\qquad
  Rf_!(D_{\leq w})\subset D_{\leq w},\qquad
  Rf_*(\text{pure})\text{ is pure},
\]
or the decomposition theorem.  In BBD these occur in
5.3.2, 5.1.13--5.1.14, and 5.4.5 only after the Deligne weight estimate.
The Fourier proof below uses the \emph{minimality} of \(j_{!*}\), never its
later purity-preservation theorem.
\end{auditbox}

\section{\texorpdfstring{\(j_{!*}\)}{Middle extension}: canonical continuation across a boundary}

\begin{proposition}[The minimal extension mechanism]
\label{prop:jmiddle-engine}
Let \(j:U\hookrightarrow X\) have closed complement \(i:Z\hookrightarrow X\)
and let \(P\in\Perv(U,E)\).  The following characterize \(j_{!*}P\):
\begin{enumerate}
\item \(j^*j_{!*}P\simeq P\);
\item \(j_{!*}P\) has no nonzero perverse subobject and no nonzero perverse
      quotient in the essential image of \(i_*\);
\item \(j_{!*}P\) is the image of the standard-to-costandard morphism
      \({}^pj_!P\to{}^pj_*P\).
\end{enumerate}
It preserves monomorphisms and epimorphisms, though it is not in general an
exact functor on short exact sequences.
\end{proposition}

\begin{proof}
The kernel and cokernel of
\({}^pj_!P\to{}^pj_*P\) are boundary objects by recollement.  Its image
therefore has neither a boundary quotient inherited from the source nor a
boundary subobject inherited from the target.  Conversely, if \(Q\) extends
\(P\) and has the stated minimality, the adjunction maps
\({}^pj_!P\to Q\to{}^pj_*P\) force \(Q\) to be their image.  This proves
uniqueness.  Applying the same boundary argument to the kernel of the image
of a monomorphism and to the cokernel of the image of an epimorphism gives
the final assertion.  See BBD 1.4.22--1.4.26 and
\cref{thm:minimal-extension-characterization}.
\end{proof}

For a smooth dense \(U\) of dimension \(d\), the object
\[
  \IC_X(L):=j_{!*}L[d]
\]
therefore remembers the local system on the top-dimensional stratum and
adds no artificial boundary constituent.  This feature, rather than a
weight assertion, is what is used in the Fourier step.

\section{Perverse Artin vanishing and exact cohomology}

\begin{theorem}[Perverse Artin engine]
\label{thm:perverse-artin-engine}
Let \(a:X\to\Spec k\) be affine and \(P\in\Perv(X,E)\).  Then
\[
  H^r(X_{\bar k},P)=0\quad(r>0),\qquad
  H_c^r(X_{\bar k},P)=0\quad(r<0).
\tag{\thechapter.1}\label{eq:artin-engine-vanishing}
\]
More generally, for an affine \(f:X\to Y\), \(Rf_*\) is right perverse
\(t\)-exact and \(Rf_!\) is left perverse \(t\)-exact.
\end{theorem}

\begin{proof}
The relative assertions are BBD 4.1.1 and its Verdier dual, BBD 4.1.2.
Apply them to \(a\).  Since the perverse structure on a point is the
ordinary one, right \(t\)-exactness of \(Ra_*\) is precisely the first
vanishing, and left \(t\)-exactness of \(Ra_!\) is the second.
\end{proof}

This is stronger than a dimension bound.  If \(U\) is an affine curve and
\(L[1]\) is perverse, then the possible cohomology of \(L\) is centered at
the single perverse degree \(0\).  If invariants and coinvariants eliminate
the two endpoint terms, \(R\Gamma_c(U,L)[1]\) is an honest vector space in
degree \(0\).  Consequently a localization triangle or a Fourier stalk
formula becomes a short exact sequence, not a long complex with
uncontrolled connecting maps.

\begin{proposition}[Perverse Leray filtration]
\label{prop:perverse-leray-engine}
For \(f:X\to Y\) and \(K\in D^b_c(X,E)\) there is a functorial spectral
sequence
\[
 E_2^{a,b}
   =H^a\bigl(Y_{\bar k},{}^pH^b(Rf_*K)\bigr)
 \Longrightarrow H^{a+b}(X_{\bar k},K).
\tag{\thechapter.2}\label{eq:perverse-leray-engine}
\]
It is the spectral sequence of the finite perverse truncation tower.  Before
the decomposition theorem, only the resulting finite Frobenius-stable
filtration may be used.
\end{proposition}

\begin{proof}
Apply \(R\Gamma(Y,-)\) to the exact couple of the perverse truncation tower
of \(Rf_*K\).  Bounded constructibility makes the tower finite.  Every
construction commutes with the Weil structure, so the filtration is
Frobenius-stable.  No degeneration assertion has entered.
\end{proof}

\section{Why Fourier--Deligne is perverse \texorpdfstring{\(t\)}{t}-exact}

Let \(A=\A^1_x\), \(A'=\A^1_{x'}\), let
\(\mu(x,x')=xx'\), and let \(\mathcal L_\psi\) be the Artin--Schreier sheaf
for a nontrivial additive character.  Define
\[
 \FT_\psi(K)
  :=R\pr_{2!}\bigl(\pr_1^*K\otimes\mu^*\mathcal L_\psi\bigr)[1].
\tag{\thechapter.3}\label{eq:fourier-engine-definition}
\]

\begin{theorem}[Fourier--Deligne perverse equivalence]
\label{thm:fourier-perverse-equivalence-engine}
The functor \(\FT_\psi\) is \(t\)-exact for the middle perverse
\(t\)-structure and induces an exact equivalence
\[
  \FT_\psi:\Perv(A,E)\xrightarrow{\sim}\Perv(A',E).
\tag{\thechapter.4}\label{eq:fourier-heart-equivalence}
\]
It sends simple perverse sheaves to simple perverse sheaves.
\end{theorem}

\begin{proof}
Let \(P\) be perverse.  The projection \(\pr_1\) is smooth of relative
dimension \(1\), so \(\pr_1^*P[1]\) is perverse; this is the smooth-pullback
case of BBD 4.2.4--4.2.5.  Tensoring with the lisse rank-one sheaf
\(\mu^*\mathcal L_\psi\) preserves perversity.  The projection \(\pr_2\) is
affine.  Hence \(R\pr_{2!}\) gives the lower perverse bound by BBD 4.1.2,
while \(R\pr_{2*}\) gives the upper bound by BBD 4.1.1.

The decisive input is Laumon's support-forgetting theorem
\cite[Thm.~1.3.1.1]{LaumonFourier}: for this oscillatory kernel the natural
map
\[
 R\pr_{2!}\bigl(\pr_1^*P\otimes\mu^*\mathcal L_\psi\bigr)
 \longrightarrow
 R\pr_{2*}\bigl(\pr_1^*P\otimes\mu^*\mathcal L_\psi\bigr)
\]
is an isomorphism.  The same object therefore satisfies both perverse
bounds and is perverse.  Fourier inversion supplies a quasi-inverse, up to
the standard sign pullback and Tate twist.  Thus the induced functor on the
heart is an exact equivalence and preserves simple objects.  This is exactly
\cite[Thm.~1.3.2.3 and Cor.~1.3.2.4]{LaumonFourier}; see also
\cref{thm:fourier-support-forgetting,thm:fourier-perverse-texact-detailed}.
\end{proof}

\begin{powerbox}
The argument uses perversity three times in an essential way:
\(\pr_1^*[1]\) is exactly normalized, the two affine direct images provide
opposite half-exactness bounds, and support-forgetting identifies those two
images.  Ordinary cohomological degrees do not turn these statements into an
exact equivalence of hearts.
\end{powerbox}

\section{The Laumon curve engine}

We now explain the hard one-dimensional estimate in the form in which it
drives all higher-dimensional weights.  Fix
\(\iota:\Qbarl\hookrightarrow\mathbf C\).

\begin{theorem}[Deligne--Laumon curve theorem]
\label{thm:laumon-curve-engine}
Let \(C_0/\Fq\) be smooth projective, let \(j:U_0\hookrightarrow C_0\) be
dense, and let \(L_0\) be lisse and punctually \(\iota\)-pure of weight
\(w\) on \(U_0\).  Then
\[
  H^r(C_{\Fbar},j_{!*}L[1])
  \quad\text{is pure of weight }w+1+r.
\tag{\thechapter.5}\label{eq:curve-ic-purity-engine}
\]
Equivalently, \(H^1(C_{\Fbar},j_*L)\), or parabolic cohomology in degree
\(1\), is pure of weight \(w+1\).
\end{theorem}

\begin{proof}
We separate the reductions from the Fourier step.

\emph{Step 1: reduction to the affine line.}
After a finite extension of the ground field and a finite map
\(C\to\PP^1\), devissage reduces to an irreducible lisse sheaf \(L\) on a
dense open \(U\subset A=\A^1\), non-geometrically-constant and unramified at
\(\infty\).  The reverse reduction is
\cite[Prop.~4.3.2.1]{LaumonFourier}.  Finite pushforward is perverse
\(t\)-exact by BBD 4.1.3, so the reduction preserves the exact perverse
problem.

\emph{Step 2: make the coefficient into one simple object.}
Put
\[
  P=j_{!*}L[1]\in\Perv(A,E).
\]
Minimal extension and irreducibility make \(P\) simple by BBD 4.3.1.  By
\cref{thm:fourier-perverse-equivalence-engine},
\[
  P'=\FT_\psi(P)
\]
is again simple perverse.  On its maximal lisse open it therefore has the
form \(j'_{!*}L'[1]\) with \(L'\) irreducible.  This is the classification
used in \cite[Thm.~1.4.2.1]{LaumonFourier}.

\emph{Step 3: prove that the Fourier local system is pure.}
The sheaf
\[
  L\oplus L^\vee(-w)
\]
is \(\iota\)-real: every local characteristic polynomial is stable under
complex conjugation.  Grothendieck's trace formula and the Fourier trace
identity show that \(L'\) is an irreducible constituent of an
\(\iota\)-real lisse sheaf on its lisse locus.  Deligne's real-sheaf
criterion \cite[Thm.~1.5.1]{DeligneWeilII} says that every irreducible lisse
constituent of an \(\iota\)-real sheaf on a geometrically connected curve is
pure.  Hence \(L'\) is pure of some weight \(w'\).  Notice that no direct
image weight theorem or purity of \(j_{!*}P\) was used.

\emph{Step 4: read global cohomology as local Fourier monodromy.}
Laumon's stationary-phase theorem identifies the local monodromy of \(L'\)
with a finite filtration whose graded pieces are local Fourier transforms
of the local monodromies of \(L\) at the singular points and at infinity.
After taking the relevant inertia invariants, the global group occurs as a
specified Frobenius-stable subquotient:
\[
  H_c^1(U_{\Fbar},L)
  \quad\text{is a subquotient of a local-monodromy filtration on}\quad
  L'_{\bar\eta_{0'}} .
\tag{\thechapter.6}\label{eq:laumon-stationary-filtration}
\]
In the extra normalization used in Laumon \S4.4---irreducible,
nonconstant input with the prescribed behavior at infinity---this
filtration has a two-step portion whose global graded piece is
\(H_c^1(U_{\Fbar},L)\) and whose neighboring graded piece is the
unramified infinity contribution \(L_{\bar\eta_\infty}(-1)\).  We use this
precise graded-piece statement, not a universal short exact sequence for
arbitrary input.  Perverse Artin vanishing is what concentrates the global
contribution in degree \(1\).  The full filtration and its local summands
are constructed in Laumon 2.3.2.1, 2.4.2.1--2.4.3, and 2.5.3.1; the
specialized global calculation is \cite[\S4.4]{LaumonFourier}.  The later
Fourier and curve chapters spell out the corresponding nearby-cycle and
inertia maps.

\emph{Step 5: determine \(w'\) and bound the global term.}
Deligne's local monodromy estimates
\cite[\S1.8, especially Thm.~1.8.4]{DeligneWeilII} compare the weights of
the inertia-invariant and one-step unipotent graded pieces in
\eqref{eq:laumon-stationary-filtration}.  The normalized infinity piece
\(L_{\bar\eta_\infty}(-1)\) has weight \(w+2\).  Applying the opposite local
inequalities to the two adjacent graded pieces forces \(w'=w+1\).
The global graded subquotient \(H_c^1(U_{\Fbar},L)\) consequently has
weights at most \(w+1\).

\emph{Step 6: turn the upper bound into equality.}
The preceding upper bound is not, by itself, a purity statement for all of
\(H_c^1(U_{\Fbar},L)\): boundary classes may have smaller weights.  Pass
instead to the parabolic image
\[
 H_p^1(U,L):=\im\!\left(H_c^1(U_{\Fbar},L)
             \longrightarrow H^1(U_{\Fbar},L)\right)
 \simeq H^1(C_{\Fbar},j_*L).
\]
Poincare--Verdier duality induces a perfect Frobenius-equivariant pairing
between this image and the corresponding parabolic image for
\(L^\vee(1)\).  Since \(L^\vee(1)\) is pure of weight \(-w-2\), applying
the upper bound to it gives the reciprocal lower bound \(w+1\) on
\(H_p^1(U,L)\).  Thus every Frobenius eigenvalue on parabolic cohomology has
absolute value \(q^{(w+1)/2}\).  The endpoint degrees are immediate from
invariants and coinvariants.  This proves \eqref{eq:curve-ic-purity-engine}.

The result is Deligne \cite[Thm.~3.2.3]{DeligneWeilII} and Laumon
\cite[Thm.~4.1.3]{LaumonFourier}.
\end{proof}

\begin{remark}[Where the genuinely deep curve input sits]
The real-sheaf criterion and the local monodromy theorem are proved before
the general direct-image theorem in Weil II.  Their proofs use rationality
of curve \(L\)-functions, positivity, and local monodromy, but not
Deligne 3.2.3 or 3.3.1.  This placement is the essential noncircularity
check for Laumon's proof.
\end{remark}

\section{From curves to arbitrary \texorpdfstring{\(Rf_!\)}{Rf-shriek}}

For a mixed Weil complex \(K_0\), write \(K_0\in D_{\leq w}\) when
\(\mathcal H^a(K_0)\) has weights at most \(w+a\).  BBD 5.4.1 equivalently
tests this condition on perverse cohomology:
\[
 K_0\in D_{\leq w}
 \Longleftrightarrow
 {}^pH^b(K_0)\text{ has weights at most }w+b.
\tag{\thechapter.7}\label{eq:perverse-weight-test-engine}
\]

\begin{theorem}[Perverse direct-image weight bound]
\label{thm:perverse-direct-image-weight-bound}
Let \(f:X_0\to Y_0\) be separated of finite type.  Then
\[
  K_0\in D_{\leq w}(X_0)
  \Longrightarrow
  Rf_!K_0\in D_{\leq w}(Y_0).
\tag{\thechapter.8}\label{eq:rfshriek-upper-engine}
\]
Equivalently, if \(P_0\) is perverse of weights at most \(w\), then
\[
  {}^pH^b(Rf_!P_0)
  \quad\text{has weights at most }w+b.
\tag{\thechapter.9}\label{eq:perverse-rfshriek-bound-engine}
\]
\end{theorem}

\begin{proof}
The detailed induction is carried out in
\cref{thm:deligne-laumon-direct-image}; here we isolate the role of the
perverse machinery.

First devissage by a smooth stratification reduces to shifted irreducible
lisse sheaves on strata.  Open--closed triangles reduce the complement by
Noetherian induction.  After shrinking the base, relative Noether
normalization factors the remaining map through a quasi-finite morphism to
an affine space over the base.  Quasi-finite fibers are zero-dimensional and
preserve the required bound directly.  Project the affine space one
coordinate at a time.  Each relative \(\A^1\)-projection is controlled
fiberwise by \cref{thm:laumon-curve-engine}.

Perverse Artin vanishing confines the cohomology of every affine projection
to the required half-range.  The perverse Leray spectral sequence
\eqref{eq:perverse-leray-engine} then assembles the coordinate projections
and yields a finite Frobenius-stable filtration.  Weight bounds pass to
extensions and subquotients, so no \(E_2\)-degeneration is needed.  The
closed remainder has smaller dimension and is handled by induction.

Translated back through \eqref{eq:perverse-weight-test-engine}, the result is
\eqref{eq:rfshriek-upper-engine}.  In ordinary cohomological notation this
is Deligne \cite[Thm.~3.3.1]{DeligneWeilII}; the derived/perverse formulation
is BBD 5.1.13--5.1.14 and 5.4.1--5.4.2.
\end{proof}

\begin{corollary}[The dual lower bound]
\label{cor:rfstar-lower-engine}
If \(K_0\in D_{\geq w}(X_0)\), then
\[
  Rf_*K_0\in D_{\geq w}(Y_0).
\tag{\thechapter.10}\label{eq:rfstar-lower-engine}
\]
\end{corollary}

\begin{proof}
By definition \(K_0\in D_{\geq w}\) means
\(\D_XK_0\in D_{\leq-w}\).  Apply
\cref{thm:perverse-direct-image-weight-bound} to \(\D_XK_0\) and use
\[
  \D_YRf_*K_0\simeq Rf_!\D_XK_0.
\]
\end{proof}

\section{Proper purity is a two-line theorem once the engine exists}

\begin{theorem}[Proper purity]
\label{thm:proper-purity-engine}
Let \(f:X_0\to Y_0\) be proper and let \(K_0\) be pure of weight \(w\).
Then \(Rf_*K_0\) is pure of weight \(w\), and
\[
  {}^pH^b(Rf_*K_0)
  \quad\text{is pure of weight }w+b.
\tag{\thechapter.11}\label{eq:proper-perverse-purity-engine}
\]
\end{theorem}

\begin{proof}
The upper bound follows from
\cref{thm:perverse-direct-image-weight-bound} because
\(Rf_*=Rf_!\) for proper \(f\).  The lower bound follows from
\cref{cor:rfstar-lower-engine}.  Their intersection is purity.  Apply the
BBD perverse weight criterion, BBD 5.4.4, to obtain
\eqref{eq:proper-perverse-purity-engine}.
\end{proof}

\begin{powerbox}
This is the central payoff.  The difficult part of the Weil proof is the
one-sided theorem for \(Rf_!\).  Properness identifies \(!\) with \(*\), and
Verdier duality reflects the upper bound into a lower bound.  Exact purity
then appears without a second geometric induction.
\end{powerbox}

\section{The smooth proper Weil theorem}

\begin{theorem}[Cohomological Riemann hypothesis]
\label{thm:weil-rh-perverse-engine}
Let \(X_0/\Fq\) be smooth proper of pure dimension \(d\).  Every eigenvalue
\(\alpha\) of geometric Frobenius on
\(H^i(X_{\Fbar},\Ql)\) satisfies, under every complex embedding,
\[
  |\alpha|=q^{i/2}.
\tag{\thechapter.12}\label{eq:weil-rh-engine}
\]
\end{theorem}

\begin{proof}
The shifted constant sheaf
\[
  P_0=\Ql[d]
\]
is perverse and pure of weight \(d\).  Apply
\cref{thm:proper-purity-engine} to the structure morphism
\(a:X_0\to\Spec\Fq\).  Since perverse and ordinary cohomology agree on a
point, the complex \(Ra_*P_0\) is pure of weight \(d\), so
\[
 H^{i-d}(Ra_*P_0)=H^i(X_{\Fbar},\Ql)
\]
is pure of weight
\[
  d+(i-d)=i.
\]
This is exactly \eqref{eq:weil-rh-engine}.
\end{proof}

Together with the trace formula, the determinant--trace identity, duality,
and the integral lattice argument developed elsewhere in these notes,
\cref{thm:weil-rh-perverse-engine} yields rationality, the functional
equation, integral Frobenius polynomials, independence of \(\ell\), and the
full Weil conjectures.  The perverse calculation explains the otherwise
mysterious degree \(i\): it is the sum of the weight \(d\) of
\(\Ql[d]\) and the residual cohomological degree \(i-d\).

\section{Only now: purity of \texorpdfstring{\(j_{!*}\)}{middle extension} and intersection complexes}

\begin{theorem}[Intermediate extension preserves purity]
\label{thm:jmiddle-purity-engine}
Let \(j:U_0\hookrightarrow X_0\) be a locally closed immersion and let
\(P_0\) be a pure perverse sheaf of weight \(w\) on \(U_0\).  Then
\(j_{!*}P_0\) is pure of weight \(w\).
\end{theorem}

\begin{proof}
For an affine open immersion, \(Rj_!\) and \(Rj_*\) have the complementary
perverse exactness properties of BBD 4.1.1--4.1.3.  The upper direct-image
weight theorem puts the relevant standard extension in weights at most
\(w\); BBD 5.3.1 passes this bound to its image \(j_{!*}P_0\).  Duality and
the minimal-extension identity
\(\D j_{!*}\simeq j_{!*}\D\) provide the lower bound.  A general locally
closed immersion is handled by factoring off a closed immersion and by
Noetherian induction on the boundary.  This is BBD 5.3.2 and 5.3.3.
\end{proof}

\begin{theorem}[Purity of intersection cohomology]
\label{thm:ic-purity-engine}
Let \(X_0\) be irreducible of dimension \(d\), let
\(j:U_0\hookrightarrow X_0\) be a smooth dense open, and let \(L_0\) be
lisse and pure of weight \(u\).  Then
\[
  \IC_{X_0}(L_0)=j_{!*}L_0[d]
\]
is a pure perverse sheaf of weight \(u+d\).  If \(X_0\) is proper, then
\[
  IH^i(X_{\Fbar},L)
    :=H^{i-d}(X_{\Fbar},\IC_X(L))
  \quad\text{is pure of weight }u+i.
\tag{\thechapter.13}\label{eq:intersection-purity-engine}
\]
\end{theorem}

\begin{proof}
The lisse perverse sheaf \(L_0[d]\) is pure of weight \(u+d\).  Apply
\cref{thm:jmiddle-purity-engine}.  If \(X_0\) is proper, apply
\cref{thm:proper-purity-engine} to its structure morphism and compute the
shift:
\((u+d)+(i-d)=u+i\).
\end{proof}

Thus perversity extends the Riemann-hypothesis statement from smooth proper
varieties to intersection cohomology of singular proper varieties.  Ordinary
cohomology of a singular proper variety is generally mixed; \(\IC_X\) is the
canonical pure replacement.

\section{Pure complexes split: the decomposition theorem}

\begin{theorem}[Geometric splitting and semisimplicity]
\label{thm:decomposition-engine}
Let \(K_0\) be a pure complex of weight \(w\) on \(Y_0/\Fq\), and write
\(K\) for its pullback to \(Y=Y_0\otimes\Fbar\).  Then
\[
  K\simeq\bigoplus_b{}^pH^b(K)[-b],
\tag{\thechapter.14}\label{eq:pure-splitting-engine}
\]
and every \({}^pH^b(K)\) is a semisimple perverse sheaf.  Consequently, if
\(f:X_0\to Y_0\) is proper and \(P_0\) is pure perverse, then
\[
 Rf_*P\ \simeq\
 \bigoplus_{b,Z,L}
 \IC_{\bar Z}(L)[-b]
\tag{\thechapter.15}\label{eq:decomposition-engine}
\]
after geometric base change, with semisimple lisse \(L\) on smooth dense
opens of the supports \(Z\).
\end{theorem}

\begin{proof}
By BBD 5.4.4, \({}^pH^b(K_0)\) is pure of weight \(w+b\), so the shifted
object \({}^pH^b(K_0)[-b]\) is pure of weight \(w\).  Let \(a\) be the top
nonzero perverse degree.  The upper truncation triangle has the form
\[
 B_0\longrightarrow K_0\longrightarrow
 A_0:={}^pH^a(K_0)[-a]\xrightarrow{\partial_0}B_0[1],
\]
where induction makes \(B_0={}^p\tau_{\leq a-1}K_0\) pure of weight \(w\).
The connecting morphism is an arithmetic class in
\(\Ext^1(A_0,B_0)\).  BBD 5.1.15(iii) states that its image in geometric
\(\Ext^1\) is zero.  Hence the triangle splits after base change to
\(\Fbar\).  Induction on the finite perverse amplitude proves
\eqref{eq:pure-splitting-engine}.  This is BBD 5.4.5.

For semisimplicity, take a pure perverse \(Q_0\).  Over the geometric base,
choose the maximal semisimple subobject \(S\subset Q\); after a finite
extension it descends.  The quotient \(T\) is pure of the same weight by
BBD 5.3.1.  The extension class of
\(0\to S\to Q\to T\to0\) again maps to zero in geometric Ext by
BBD 5.1.15(iii), so it splits.  Maximality forces \(T=0\).  Thus \(Q\) is
semisimple, which is BBD 5.3.8.  Finally, BBD 4.3.1 identifies its simple
summands as intermediate extensions of irreducible lisse sheaves.

For proper \(f\), \cref{thm:proper-purity-engine} makes \(Rf_*P_0\) pure.
Apply the two conclusions above to obtain
\eqref{eq:decomposition-engine}.
\end{proof}

\begin{warning}[Arithmetic Frobenius need not be semisimple]
\Cref{thm:decomposition-engine} is a geometric splitting after passage to
\(\Fbar\).  It does not say that the Frobenius operator on each arithmetic
cohomology group is semisimple.  Pure Weil sheaves can carry nontrivial
unipotent Frobenius blocks.  BBD 5.3.8 asserts geometric semisimplicity of
the perverse sheaf, not semisimplicity of every chosen Weil structure.
\end{warning}

\begin{corollary}[Perverse Leray degeneracy for pure proper maps]
\label{cor:perverse-leray-degeneration-engine}
Under the proper and pure hypotheses of
\cref{thm:decomposition-engine}, the perverse Leray spectral sequence
\eqref{eq:perverse-leray-engine} degenerates at \(E_2\) after geometric base
change, and its filtration admits a noncanonical splitting.
\end{corollary}

\begin{proof}
Apply \(R\Gamma(Y,-)\) to the direct-sum decomposition
\eqref{eq:pure-splitting-engine}.  The resulting cohomology is the direct sum
of the \(E_2\)-terms, which kills all higher differentials.  Noncanonicity is
inherited from the choices in splitting the perverse truncation triangles.
\end{proof}

If \(f\) is projective and \(\eta\) is the Chern class of a relatively ample
line bundle, BBD 5.4.10 further gives relative hard Lefschetz
\[
  \eta^r:{}^pH^{-r}(Rf_*P_0)
  \xrightarrow{\sim}{}^pH^r(Rf_*P_0)(r)
  \qquad(r\geq0).
\tag{\thechapter.16}\label{eq:relative-hard-lefschetz-engine}
\]
This theorem is another downstream benefit of purity.  It must not be used
to establish the weight theorem from which its proof is derived.

\section{What each perverse tool contributes}

\begin{center}
\begin{tabularx}{\textwidth}{>{\raggedright\arraybackslash}p{0.24\textwidth}
                              >{\raggedright\arraybackslash}X}
\toprule
Perverse construction & Irreplaceable role in the proof\\
\midrule
\(j_{!*}\) &
Packages an irreducible local system as a simple global object without
boundary subobjects or quotients; later produces the pure intersection
complex.\\
\addlinespace
Artin vanishing &
Concentrates affine curve cohomology into the degree in which Fourier
stalks and local monodromy form short exact sequences.\\
\addlinespace
Fourier \(t\)-exactness &
Turns an oscillatory derived integral into an exact equivalence of hearts,
so simplicity survives the transform and Deligne's real-sheaf criterion
applies to one irreducible lisse constituent.\\
\addlinespace
Perverse Leray &
Assembles relative one-dimensional estimates through finite
Frobenius-stable filtrations without assuming premature degeneration.\\
\addlinespace
Perverse weight criterion &
Translates the sheafwise \(Rf_!\) estimate into a statement about exact
perverse constituents with the correct degree shift.\\
\addlinespace
Proper purity &
Combines the \(!\)-upper and \(*\)-lower bounds in one object because
\(!=*\) for a proper map.\\
\addlinespace
Pure splitting &
Kills geometric extension classes between equal-weight truncation pieces,
yielding decomposition, semisimplicity, and perverse Leray degeneration.\\
\bottomrule
\end{tabularx}
\end{center}

\section{Primary-source map and dependency audit}

The exact source map used in this chapter is:
\begin{itemize}
\item BBD 1.4.22, 1.4.25, 1.4.26:
      \(j_{!*}\), minimality, and simple objects under recollement;
\item BBD 4.1.1--4.1.3:
      affine Artin vanishing and quasi-finite \(t\)-exactness;
\item BBD 4.2.4--4.2.5 and 4.2.8:
      smooth shifted pullback and external products;
\item BBD 4.3.1:
      finite length and classification of simple perverse sheaves;
\item Deligne 1.5.1 and 1.8.4:
      the real-sheaf purity criterion and local monodromy weights;
\item Laumon 1.3.1.1, 1.3.2.3--1.3.2.4, and 1.4.2.1:
      support-forgetting, Fourier perverse equivalence, and simple types;
\item Laumon 2.3.2.1, 2.4.2.1--2.4.3, 2.5.3.1, and 4.4:
      local Fourier transform, stationary phase, and the curve proof;
\item Deligne 3.2.3 and Laumon 4.1.3:
      the curve purity theorem;
\item Deligne 3.3.1, BBD 5.1.13--5.1.14 and 5.4.1:
      the general upper direct-image bound and its perverse translation;
\item BBD 5.3.1--5.3.2, 5.3.8, 5.4.4--5.4.5:
      subquotient weights, pure \(j_{!*}\), geometric semisimplicity, and
      pure-complex splitting;
\item BBD 5.4.10:
      relative hard Lefschetz after purity.
\end{itemize}

\begin{auditbox}[title={Acyclic theorem order}]
\[
\begin{array}{c}
\text{pro-etale coefficients and six operations}\\
\downarrow\\
\text{perverse recollement, minimal \(j_{!*}\), Artin vanishing}\\
\downarrow\\
\text{Fourier perverse equivalence and local Fourier transform}\\
\downarrow\\
\text{Deligne--Laumon curve theorem}\\
\downarrow\\
\text{general \(Rf_!\) upper weights}\\
\downarrow\\
\text{proper purity and smooth proper Weil RH}\\
\downarrow\\
\text{pure \(j_{!*}\), IC purity, decomposition, relative Lefschetz}.
\end{array}
\]
Every downward arrow is used; no upward arrow is allowed.  In particular,
decomposition is not a premise of the proof of the Weil conjectures in this
volume.
\end{auditbox}

\section{Formalization interface}

\begin{formalization}
The reusable theorem interfaces suggested by the proof are:
\begin{enumerate}
\item a recollement package whose heart-level image defines \(j_{!*}\) and
      proves the no-boundary universal property;
\item affine left/right perverse exactness as two separate theorems, linked
      by Verdier duality;
\item a Fourier kernel record carrying smoothness of \(\pr_1\), affineness
      of \(\pr_2\), support-forgetting, duality, and inversion;
\item a theorem that those fields imply an exact equivalence of perverse
      hearts, rather than postulating Fourier \(t\)-exactness independently;
\item a curve-engine record containing the real-sheaf criterion, local
      monodromy filtration, the stationary-phase filtration with its specified
      global subquotient, and duality;
\item an upper-weight theorem for \(Rf_!\), with the lower theorem derived
      mechanically by duality;
\item proper purity derived from the equality \(Rf_!=Rf_*\), not stored as
      an unrelated axiom;
\item pure splitting derived from the geometric vanishing of the image of
      the relevant arithmetic \(\Ext^1\)-class.
\end{enumerate}
The final test case should reduce
\(\Ql[d]\) along \(X\to\Spec\Fq\) and normalize the weight by the literal
identity \(d+(i-d)=i\).  This catches simultaneously the most common shift,
twist, and Frobenius-convention errors.
\end{formalization}



\part{Weights and the Deligne--Laumon theorem}
\chapter{Weil Sheaves, Weights, and Determinant Weights}
\label{chap:weights}

The purpose of this chapter is to introduce the language in which the
Riemann hypothesis becomes stable under geometric operations.  Two
distinctions must remain visible.  A Weil sheaf carries an action of the
discrete Weil group, whereas a sheaf descended to a finite field carries
continuous descent under a profinite Galois group.  Also, purity relative
to one field isomorphism with \(\mathbf C\) is weaker than the
all-conjugates assertion eventually required for integral Frobenius
polynomials.

Throughout \(k_0=\Fq\), \(k=\Fbar\), and
\[
  X=X_0\times_{k_0}k
\]
for a separated finite-type \(k_0\)-scheme \(X_0\).  We use geometric
Frobenius; it acts on \(\Ql(1)\) by \(q^{-1}\).

\section{Scheme Frobenius, Galois Frobenius, and Weil data}

Let
\[
  \Fr_{X_0}^{(q)}:X_0\longrightarrow X_0
\]
be the \(q\)-power scheme Frobenius: it is the identity on the
underlying topological space and sends a local function \(f\) to
\(f^q\).  Because the ground field is \(\Fq\), this is an
\(\Fq\)-morphism.  Define
\[
  F_X
  =
  \Fr_{X_0}^{(q)}\times_{\Spec\Fq}\id_{\Spec k}
  :X\longrightarrow X.
\tag{8.1}\label{eq:scheme-frobenius}
\]
This is a \(k\)-morphism and a universal homeomorphism; when \(X_0\) is
smooth its differential is zero.  It need not be an isomorphism of
schemes, but pullback by it is an equivalence of the \'etale topos.

This scheme map must not be confused with the semilinear Galois
automorphism.  Let
\[
  \gamma:a\longmapsto a^{1/q}
\]
be geometric Galois Frobenius and let
\(\Gamma_X=\id_{X_0}\times\gamma\).  If
\(\Fr_X^{\mathrm{abs}}\) denotes absolute \(q\)-power Frobenius on
\(X\), then
\[
  F_X=\Fr_X^{\mathrm{abs}}\circ\Gamma_X.
\tag{8.2}\label{eq:two-frobenius-comparison}
\]
Absolute Frobenius induces the identity equivalence on the \'etale
topos.  Consequently a datum
\(F_X^*\cF\simeq\cF\) is canonically the same, on the \'etale topos, as
a datum \(\Gamma_X^*\cF\simeq\cF\).  Under Galois descent the latter is
the action of \emph{geometric} Galois Frobenius.  This comparison is
also why the endomorphism induced by the scheme map \(F_X\) on
cohomology is the geometric-Galois Frobenius action used below.

\begin{definition}
\label{def:weil-sheaf}
A \emph{constructible Weil sheaf} on \(X_0\) is a pair
\((\cF,\varphi)\), where \(\cF\) is a constructible
\(\Qbarl\)-sheaf on \(X\) and
\[
  \varphi:F_X^*\cF\xrightarrow{\sim}\cF.
\]
A Weil complex is defined in the same way in the constructible derived
category.
\end{definition}

Iteration defines, at every degree \(n\) closed point \(x_0\), a
conjugacy class
\[
 F_{x_0}:\cF_{\bar x}\longrightarrow\cF_{\bar x}.
\]
Via \eqref{eq:two-frobenius-comparison}, this is geometric Galois
Frobenius.  On the constant sheaf it is the identity; on
\(\Qbarl(1)\) it is \(q^{-n}\), and on \(\Qbarl(-1)\) it is \(q^n\).
Thus \(\Qbarl(m)\) has weight \(-2m\).

If \(X_0\) is connected, its Weil group fits into
\[
 1\to\pi_1(X,\bar x)\to W(X_0,\bar x)
 \xrightarrow{\deg}\mathbf Z\to0,
\]
with geometric Frobenius of degree \(1\).  Lisse Weil sheaves are
continuous finite-dimensional representations of this group, with the
usual \(\ell\)-adic topology on the geometric subgroup and the discrete
topology on \(\mathbf Z\).

\begin{warning}
\label{warn:weil-versus-etale}
A Weil sheaf need not descend to \(X_0\).  Descent requires the action
to extend continuously from \(\mathbf Z\) to
\(\widehat{\mathbf Z}=\Gal(k/k_0)\).  The Fourier argument uses scalar
Weil twists which need not meet that condition, so the trace formula
and the six operations must be available for Weil objects themselves.
\end{warning}

\section{Punctual purity and mixedness}

Fix \(\iota:\Qbarl\xrightarrow{\sim}\mathbf C\).

\begin{definition}
\label{def:iota-purity}
Let \(\cF_0\) be a constructible Weil sheaf and \(w\in\mathbf R\).
\begin{enumerate}
\item It is \emph{punctually \(\iota\)-pure of weight \(w\)} if every
local eigenvalue \(\alpha\) at a degree \(n\) closed point satisfies
\[
 |\iota(\alpha)|=q^{nw/2}.
\]
\item It is \emph{\(\iota\)-mixed} if it has a finite filtration whose
quotients are punctually pure.
\item It has weights at most \(w\) if all quotient weights are at most
\(w\), and weights at least \(w\) analogously.
\end{enumerate}
\end{definition}

For a complex \(\cK_0\), the derived convention is
\[
 \cK_0\text{ has weights }\le w
 \Longleftrightarrow
 \cH^j(\cK_0)\text{ has weights }\le w+j.
\]

\begin{definition}
\label{def:all-embedding-purity}
An algebraic number \(\alpha\in\Qbar\) is a \(q^n\)-Weil number of
weight \(w\) if
\[
 |\sigma(\alpha)|=q^{nw/2}
\]
for every \(\sigma:\Qbar\hookrightarrow\mathbf C\).  A sheaf is pure
in the all-embeddings sense if all local eigenvalues are algebraic and
satisfy this condition.
\end{definition}

\begin{proposition}[Elementary permanence]
\label{prop:elementary-weight-permanence}
Upper bounds are preserved by subobjects, quotients, extensions, finite
sums, and finite Frobenius-stable filtrations.  Tensor products add pure
weights, duals negate them, a twist by \(m\) subtracts \(2m\), and
pullback preserves punctual weights.
\end{proposition}

\begin{proof}
Stalks are exact and characteristic polynomials multiply in short exact
sequences.  Tensor eigenvalues are products, dual eigenvalues are
inverses, and a twist by \(m\) multiplies a degree \(n\) eigenvalue by
\(q^{-mn}\).  Pullback raises Frobenius to the same power by which it
multiplies the residue degree.
\end{proof}

\section{Rank one: the independent base case}

\begin{foundation}
\label{found:rank-one-purity}
\textbf{Rank-one purity capsule.}
Let \(X_0\) be normal and geometrically connected, and let \(\cL_0\) be
a rank-one lisse Weil sheaf.  For every \(\iota\), it is punctually pure
of some real weight.

The associated character
\(\chi:W(X_0,\bar x)\to\Qbarl^\times\) has finite image on
\(\pi_1(X,\bar x)\) by the rank-one Katz--Lang theorem.  This reduces,
through a sufficiently ample normal curve, to geometric class field
theory and Lang's theorem for a generalized Picard variety; it uses no
higher-rank weight result.  If \(s\) has degree one and \(a=\chi(s)\),
then at a degree \(n\) point
\[
 \chi(F_{x_0})=\zeta_{x_0}a^n
\]
with \(\zeta_{x_0}\) a root of unity.  Hence
\[
 |\iota\chi(F_{x_0})|
 =q^{nw_\iota/2},
 \qquad
 w_\iota=2\log_q|\iota(a)|.
\]
See \cite[Section~1.3]{DeligneWeilII},
\cite[Chapter~I]{KiehlWeissauer}, and
\cite{StanfordL19,StanfordL22}.
\end{foundation}

\section{Determinant weights}

\begin{definition}
\label{def:determinant-weight}
For an irreducible lisse Weil sheaf \(\cF_0\) of rank \(r\), define
\[
 d_\iota(\cF_0)=\frac1r\,w_\iota(\det\cF_0).
\]
For a geometrically semisimple sheaf, take the determinant weights of
its irreducible constituents with multiplicity.
\end{definition}

\begin{lemma}[Determinant-weight arithmetic]
\label{lem:det-weight-arithmetic}
Duality negates determinant weights; a rank-one twist of weight \(c\)
adds \(c\); the largest weight of \(\cF^{\otimes m}\) is \(m\) times
that of \(\cF\); and exterior-power weights are the corresponding sums,
subject to constituent ranks.
\end{lemma}

\begin{proof}
After semisimplification, tensor and exterior constructions are
representations of the product of general linear groups on the
isotypic blocks.  Their central characters are the indicated monomials;
rank-one purity changes products into sums of weights.
\end{proof}

\section{Arithmetic and geometric monodromy}

Let
\[
 G_{\mathrm{geom}}
 =\overline{\rho(\pi_1(X,\bar x))}^{\,\mathrm{Zar}}
 \subseteq\operatorname{GL}(V)
\]
and let \(G_{\mathrm{arith}}\) be generated by this group and a
positive-degree Frobenius element.

\begin{proposition}[Central element of positive degree]
\label{prop:positive-degree-center}
After a connected finite \'etale cover and a finite extension of
\(k_0\), \(H=G_{\mathrm{geom}}^\circ\) is semisimple.  Moreover
\(Z(G_{\mathrm{arith}})(\Qbarl)\) contains an element of positive
degree, and
\[
 Z(G_{\mathrm{arith}})(\Qbarl)\xrightarrow{\deg}\mathbf Z
\]
has finite kernel and cokernel.
\end{proposition}

\begin{proof}
Geometric semisimplicity first makes \(H\) reductive: its unipotent
radical has a fixed vector on every irreducible constituent and
therefore acts trivially; faithfulness kills it.

After taking a power of a positive-degree element \(s\), its action on
\(H\) is inner.  Choose \(g\in H\) inducing that action and put
\(z=g^{-1}s\); then \(z\) centralizes \(H\).  If \(H\) had a nontrivial
algebraic character, the splitting generated by \(H\) and \(z\) would
extend it to a rank-one character of a finite-index Weil group with
infinite geometric image, contradicting
\cref{found:rank-one-purity}.  Thus the connected reductive group \(H\)
has no characters, hence is semisimple and has finite center.

Raise \(z\) so that it acts trivially on the finite component group of
\(G_{\mathrm{geom}}\).  For \(u\in G_{\mathrm{geom}}\),
\[
 c_u=zuz^{-1}u^{-1}\in H.
\]
Since \(z\) centralizes \(H\), direct conjugation shows
\(c_u\in Z(H)\).  This finite-valued cocycle factors through the
component group.  Raising \(z\) to the exponent of \(Z(H)\) kills all
\(c_u\), so a positive power of \(z\) is central in
\(G_{\mathrm{arith}}\).  The kernel of degree is contained in the
finite group \(Z(G_{\mathrm{geom}})\), and the positive degree of \(z\)
makes the cokernel finite.  See
\cite[Chapter~I, Section~3]{KiehlWeissauer} and
\cite{StanfordL22}.
\end{proof}

\begin{corollary}[Dense pullback preserves determinant weights]
\label{cor:dense-pullback-det-weights}
Let \(f:Y_0\to X_0\) have dense image between normal geometrically
irreducible schemes, and let \(\cF_0\) be geometrically semisimple and
lisse.  Then \(\cF_0\) and \(f^*\cF_0\) have the same determinant
weights.
\end{corollary}

\begin{proof}
Choose compatible positive-degree central elements in the two
monodromy groups and replace them by powers until degrees and images
agree.  They act by the same scalar on a constituent and on every
summand after pullback.  The logarithmic absolute value of that scalar,
divided by the degree and normalized by rank, is the determinant
weight.
\end{proof}

\section{The dependency boundary}

\begin{formalization}
\label{form:weights-interface}
\textbf{Reusable interface.}
\begin{enumerate}
\item The \(q\)-power scheme Frobenius \(F_X\), separately from the
semilinear Galois automorphism \(\Gamma_X\), together with
\eqref{eq:two-frobenius-comparison}.
\item Weil objects, effective descent, fixed-\(\iota\) purity, and
all-embeddings purity as distinct notions.
\item Elementary permanence, rank-one purity, determinant weights, and
dense-pullback invariance.
\item The semisimple finite-center proof of a positive-degree central
element.
\end{enumerate}
No result here uses \(Rf_!\)-preservation of weights or smooth-proper
purity.
\end{formalization}

\chapter{Perverse Degree and Weight Bookkeeping}
\label{chap:perverse-weights}

This chapter translates the punctual weights of
\cref{chap:weights} into perverse degree.  It proves only the formal
translation.  Stability of weights under \(Rf_!\), purity of intermediate
extensions, and the decomposition theorem are deliberately not used here;
they are downstream consequences of the Deligne--Laumon estimate.

Fix \(\iota:\Qbarl\simeq\mathbf C\).  All weight assertions in this chapter
are \(\iota\)-weight assertions unless “for every \(\iota\)” is stated.

\section{Weights of complexes}

\begin{definition}[Upper and lower weights]
\label{def:derived-weight-bounds}
A mixed Weil complex \(\cK_0\in\Dbc(X_0,\Qbarl)\) has weights at most \(w\)
if
\[
  \cH^a(\cK_0)\quad\text{has weights at most }w+a
  \quad\text{for every }a.
\tag{10.1}\label{eq:ordinary-derived-upper}
\]
Write \(\cK_0\in D_{\leq w}(X_0)\).  It has weights at least \(w\) if
\[
  \D_{X_0}\cK_0\in D_{\leq-w}(X_0),
\tag{10.2}\label{eq:derived-lower-dual}
\]
and is pure of weight \(w\) if it belongs to both
\(D_{\leq w}\) and \(D_{\geq w}\).
\end{definition}

The definition is engineered so that
\[
 \cK[n]\text{ pure of weight }w+n,\qquad
 \cK(m)\text{ pure of weight }w-2m,\qquad
 \D_X\cK\text{ pure of weight }-w.
\tag{10.3}\label{eq:weight-shift-twist-dual}
\]
Indeed, \(\cH^a(\cK[n])=\cH^{a+n}(\cK)\), geometric Frobenius on
\(\Qbarl(m)\) has weight \(-2m\), and the last assertion is built into
\eqref{eq:derived-lower-dual}.

\begin{warning}
“Punctually pure” is a property of a sheaf and its stalk eigenvalues.
“Pure complex” is the simultaneous upper/lower condition
\cref{def:derived-weight-bounds}.  They agree for a lisse sheaf in ordinary
degree zero on a smooth stratum, but they are not interchangeable words on a
singular space.
\end{warning}

\section{The dimension shift once more}

\begin{proposition}[Perverse normalization of a local system]
\label{prop:lisse-perverse-weight-shift}
Let \(U_0\) be smooth of pure dimension \(d\), and let \(\cL_0\) be lisse and
punctually pure of weight \(u\).  Then
\[
  \cL_0[d]\in\Perv(U_0)
  \quad\text{and is pure of weight }u+d.
\tag{10.4}\label{eq:lisse-perverse-weight}
\]
\end{proposition}

\begin{proof}
Perversity is \cref{prop:lisse-shift-perverse}.  The only ordinary
cohomology sheaf of \(\cL[d]\) is
\(\cH^{-d}(\cL[d])=\cL\).  For a complex of weight \(u+d\), the required
weight in ordinary degree \(-d\) is \((u+d)-d=u\), giving the upper bound.
Smooth duality gives
\[
  \D_U(\cL[d])\simeq\cL^\vee(d)[d].
\]
The dual local system has weight \(-u\), and the twist has weight \(-2d\);
therefore the right side has perverse-normalized weight \(-u-d\).  This is the
lower bound by \eqref{eq:derived-lower-dual}.
\end{proof}

In particular, on a smooth \(d\)-fold
\[
  \Qbarl[d]\quad\text{is perverse and pure of weight }d.
\tag{10.5}\label{eq:constant-perverse-weight}
\]
This single shift is the source of the final identity
\(d+(i-d)=i\).

\section{Ordinary versus perverse cohomology}

For a perverse sheaf \(\cP_0\), say that it has weights at most \(w\) if it
does so as a complex in the sense of \cref{def:derived-weight-bounds};
define lower bounds and purity similarly.

\begin{theorem}[BBD perverse weight criterion]
\label{thm:perverse-weight-criterion}
For a mixed complex \(\cK_0\),
\[
 \cK_0\in D_{\leq w}
 \quad\Longleftrightarrow\quad
 \pH^a(\cK_0)\text{ has weights }\leq w+a
 \text{ for every }a.
\tag{10.6}\label{eq:perverse-upper-equivalence}
\]
The same equivalence holds for lower bounds.  Consequently \(\cK_0\) is pure
of weight \(w\) if and only if
\[
  \pH^a(\cK_0)\quad\text{is pure of weight }w+a
\tag{10.7}\label{eq:pH-pure-shift}
\]
for every \(a\).
\end{theorem}

\begin{proof}
Filter \(\cK_0\) by its perverse truncation tower.  On a smooth stratum of
dimension \(d\), each perverse constituent is a shifted lisse sheaf
\(\cL[d]\), and \cref{prop:lisse-perverse-weight-shift} converts the desired
perverse inequality to the ordinary one.  For the closed complement, use
Noetherian induction and the recollement triangles.  Upper weight bounds pass
through extensions, subobjects, and quotients; duality gives the lower
statement.  The finite tower gives one implication by devissage and the other
by applying perverse cohomology to ordinary truncation triangles.  This is BBD
5.4.1; \eqref{eq:pH-pure-shift} is BBD 5.4.4.
\end{proof}

On the point \(\Spec\Fq\), perverse and ordinary cohomology coincide, so
\eqref{eq:pH-pure-shift} reduces to the familiar rule
\[
  \cK_0\text{ pure of weight }w
  \Longrightarrow H^a(\cK)\text{ pure of weight }w+a.
\tag{10.8}\label{eq:point-weight-shift}
\]

\section{Safe formal consequences}

\begin{lemma}[Weight bounds and finite filtrations]
\label{lem:weight-filtration-subquotients}
Let \(V\) carry geometric Frobenius and a finite Frobenius-stable filtration.
If every graded piece has weights at most \(w\), then \(V\) has weights at
most \(w\).  The analogous statement holds for lower bounds.  Hence every
subquotient of a mixed perverse sheaf of weights at most \(w\) has weights at
most \(w\).
\end{lemma}

\begin{proof}
Characteristic polynomials multiply in a short exact sequence, so the
multiset of eigenvalues is the union of those of the subobject and quotient.
Iterate through the filtration.  Costalks and stalks are exact on the
appropriate perverse ranges; the general assertion follows by the
stratification devissage used in \cref{thm:perverse-weight-criterion}.
\end{proof}

\begin{corollary}[Spectral-sequence rule]
\label{cor:perverse-weight-spectral-rule}
In a bounded Frobenius-equivariant spectral sequence, if every \(E_2^{a,b}\)
has weights at most \(u_{a,b}\), then every graded piece of the abutment has
the maximum of the bounds of the terms that can contribute to it.  No
degeneration is required.
\end{corollary}

\section{The dependency firewall}

The following important statements are \emph{not} available yet:
\begin{align*}
 \cP\text{ pure }w
 &\Longrightarrow j_{!*}\cP\text{ pure }w,\\
 \cK\in D_{\leq w}
 &\Longrightarrow Rf_!\cK\in D_{\leq w},\\
 \cK\text{ pure},\ f\text{ proper}
 &\Longrightarrow Rf_*\cK\text{ pure},\\
 \cK\text{ pure}
 &\Longrightarrow
   \cK_{\Fbar}\simeq\bigoplus_a\pH^a(\cK_{\Fbar})[-a].
\end{align*}
The second line is the main Deligne--Laumon theorem proved by the next four
chapters.  The first and third then follow by combining its upper bound with
duality; the fourth follows from weight separation.  Using any of these four
lines in the Fourier or curve argument would be circular.

\begin{perversebox}
Before the curve theorem, the allowed perverse tools are the \(t\)-structure,
recollement, \(j_{!*}\) as a minimal extension, Artin vanishing, Fourier
\(t\)-exactness, and the bookkeeping formulas
\eqref{eq:lisse-perverse-weight}--\eqref{eq:pH-pure-shift}.  No proper-purity
or decomposition statement is available.
\end{perversebox}

\begin{formalization}
Represent the three notions separately:
\begin{enumerate}
\item punctual weight of a constructible Weil sheaf;
\item upper/lower derived weight predicates;
\item perverse-heart weight predicates transported by the inclusion of the
      heart into the derived category.
\end{enumerate}
Then formalize \cref{thm:perverse-weight-criterion} as two equivalences and
\eqref{eq:pH-pure-shift} as their intersection.  Tests should verify that
shift \([1]\) raises weight by \(1\), twist \((1)\) lowers it by \(2\), and
\(\Qbarl[d]\) has perverse weight \(d\), not \(0\).
\end{formalization}

\chapter{Real Sheaves and Purity}
\label{chap:real-sheaves}

The theorem of this chapter is the engine that prevents the Fourier
argument from being circular.  It extracts punctual purity from a
reality condition by using positivity of even tensor powers.  Its proof
uses the trace formula on a curve, the rank-one theorem of
\cref{found:rank-one-purity}, elementary monodromy, and Poincare duality
in top degree.  It does \emph{not} use the direct-image weight theorem
which it will later help prove.

Fix an isomorphism \(\iota:\Qbarl\to\mathbf C\) throughout this chapter.

\section{Reality and positivity}

\begin{definition}
\label{def:iota-real}
A constructible Weil sheaf \(\cG_0\) on \(X_0\) is
\emph{\(\iota\)-real} if, for every closed point \(x_0\),
\[
  \iota\det(1-TF_{x_0}\mid\cG_{\bar x})
  \in\mathbf R[T].
\]
Equivalently, the multiset of complex numbers obtained from the local
Frobenius eigenvalues is stable under complex conjugation.
\end{definition}

Reality is preserved by direct sums, duals, tensor powers, exterior
powers, subquotients whose complementary factor is also real, and
semisimplification.  It is not in general inherited by an individual
irreducible constituent.  This last observation is why the proof below
needs exterior powers rather than simply applying the irreducible case
constituent by constituent.

\begin{lemma}[Positive local factors]
\label{lem:positive-local-factors}
If \(\cG_0\) is \(\iota\)-real and \(m\geq 1\), then for every closed
point \(x_0\) the series
\[
  \iota\det\left(1-TF_{x_0}\mid
  \cG_{\bar x}^{\otimes 2m}\right)^{-1}
\]
has nonnegative real coefficients.
\end{lemma}

\begin{proof}
For an endomorphism \(A\) of a finite-dimensional vector space,
\[
  \log\det(1-TA)^{-1}
  =\sum_{n\geq1}\Tr(A^n)\frac{T^n}{n}.
\]
Since the characteristic polynomial of \(F_{x_0}\) becomes real under
\(\iota\), all power traces
\(\iota\Tr(F_{x_0}^n\mid\cG_{\bar x})\) are real.  On the tensor power,
\[
  \Tr(F_{x_0}^n\mid\cG_{\bar x}^{\otimes 2m})
  =\Tr(F_{x_0}^n\mid\cG_{\bar x})^{2m},
\]
which becomes a nonnegative real number.  The exponential of a formal
series with zero constant term and nonnegative coefficients again has
nonnegative coefficients.
\end{proof}

We shall use one elementary fact about positive products.

\begin{lemma}[A local factor cannot converge less far]
\label{lem:positive-product-radius}
Let \(f_j(T)\in1+T\mathbf R_{\geq0}[[T]]\) be a locally finite family in
the \(T\)-adic topology, and put \(f=\prod_j f_j\).  The radius of
convergence of each \(f_j\) is at least the radius of convergence of
\(f\).
\end{lemma}

\begin{proof}
Write \(f_j=\sum_na_{j,n}T^n\) and \(f=\sum_na_nT^n\).  On multiplying
series with nonnegative coefficients, the contribution obtained by
taking the degree \(n\) term from \(f_j\) and constant terms from all
other factors shows \(0\leq a_{j,n}\leq a_n\).  The
Cauchy--Hadamard formula gives
\[
  \limsup_n a_{j,n}^{1/n}
  \leq\limsup_n a_n^{1/n},
\]
which is the asserted comparison of radii.
\end{proof}

\section{Top cohomology of an affine curve}

Let \(C_0\) be a smooth geometrically connected affine curve over \(k_0\),
and let \(\cG_0\) be a lisse Weil sheaf corresponding to a representation
\(V\) of \(W(C_0,\bar c)\).

\begin{lemma}
\label{lem:curve-hc0-hc2}
There are Frobenius-equivariant identifications
\[
  H_c^0(C,\cG)=0,\qquad
  H_c^2(C,\cG)\simeq
  V_{\pi_1(C,\bar c)}(-1),
\]
where the subscript denotes geometric coinvariants.  In particular, if
\(\lambda\) is an eigenvalue of geometric Frobenius on \(H_c^2(C,\cG)\),
then
\[
  2\log_q|\iota(\lambda/q)|
\]
is a determinant weight of a geometrically constant quotient of
\(\cG_0\), and hence is among the determinant weights of \(\cG_0\).
\end{lemma}

\begin{proof}
A lisse section with proper support vanishes on a nonempty open set and
therefore vanishes everywhere on the connected curve, proving the first
claim.  Poincare duality gives
\[
 H_c^2(C,\cG)
 \simeq H^0(C,\cG^\vee)^\vee(-1).
\]
The space \(H^0(C,\cG^\vee)\) is
\((V^\vee)^{\pi_1(C,\bar c)}\); its dual is
\(V_{\pi_1(C,\bar c)}\).  The twist \((-1)\) multiplies geometric
Frobenius by \(q\).  The coinvariant representation is geometrically
constant and is a quotient of \(V\).  Its irreducible Weil constituents
are rank one after extension of scalars because the geometric group acts
trivially.  Their weights are determinant weights of constituents of
\(\cG_0\), by \cref{found:rank-one-purity,def:determinant-weight}.
\end{proof}

\begin{corollary}
\label{cor:hc2-upper-det-weight}
If \(\beta\) is the largest determinant weight of a geometrically
semisimple \(\cG_0\), every eigenvalue \(\lambda\) on
\(H_c^2(C,\cG)\) satisfies
\[
  |\iota(\lambda)|\leq q^{(\beta+2)/2}.
\]
For \(\cG_0^{\otimes r}\), the same assertion holds with \(\beta\)
replaced by \(r\beta\).
\end{corollary}

\begin{proof}
The first statement is immediate from
\cref{lem:curve-hc0-hc2}; the second uses
\cref{lem:det-weight-arithmetic}.
\end{proof}

\section{The affine-curve real-sheaf theorem}

\begin{theorem}[Real sheaves on affine curves]
\label{thm:real-sheaf-affine-curve}
Let \(C_0\) be a smooth geometrically connected affine curve, and let
\(\cG_0\) be a lisse \(\iota\)-real Weil sheaf.  Every irreducible
constituent of its semisimplification is punctually
\(\iota\)-pure, of weight equal to its determinant weight.  Consequently
\(\cG_0\) is \(\iota\)-mixed.
\end{theorem}

\begin{proof}
Semisimplification leaves all local characteristic polynomials unchanged,
so we may assume that \(\cG_0\) is semisimple.

\medskip
\noindent\emph{Step 1: a universal upper bound.}
Let \(\beta\) be the largest determinant weight of \(\cG_0\), and fix
\(m\geq1\).  The trace formula for the even tensor power, together with
\cref{lem:curve-hc0-hc2}, gives an identity of rational functions and of
Taylor series near \(0\):
\[
\begin{split}
 L(C_0,\cG_0^{\otimes 2m},T)
 &=
 \prod_{x_0\in|C_0|}
 \det\left(1-T^{\deg(x_0)}F_{x_0}\mid
 \cG_{\bar x}^{\otimes 2m}\right)^{-1}\\
 &=
 \frac{
 \det(1-TF\mid H_c^1(C,\cG^{\otimes 2m}))
 }{
 \det(1-TF\mid H_c^2(C,\cG^{\otimes 2m}))
 }.
\end{split}
\tag{9.1}
\label{eq:real-sheaf-L}
\]
By \cref{cor:hc2-upper-det-weight}, every zero of the denominator on the
right has modulus at least
\[
  q^{-(2m\beta+2)/2}=q^{-m\beta-1}.
\]
The numerator can cancel a pole but cannot create a new one.  Hence the
Taylor series of the rational function in \eqref{eq:real-sheaf-L} has
radius of convergence at least \(q^{-m\beta-1}\).

All local factors on the first line have nonnegative coefficients after
applying \(\iota\), by \cref{lem:positive-local-factors}.  The product is
locally finite because only finitely many closed points have bounded
degree.  By \cref{lem:positive-product-radius}, each local factor has
radius at least \(q^{-m\beta-1}\).

Let \(\alpha\) be an eigenvalue of \(F_{x_0}\) on
\(\cG_{\bar x}\), with \(e=\deg(x_0)\).  Then \(\alpha^{2m}\) is an
eigenvalue on the even tensor power.  The corresponding local reciprocal
factor has a singularity at a complex number \(T\) satisfying
\[
  |T|^e=|\iota(\alpha)|^{-2m}.
\]
The radius estimate therefore gives
\[
  |\iota(\alpha)|^{2m/e}
  \leq q^{m\beta+1},
\]
or
\[
  |\iota(\alpha)|
  \leq q^{e(\beta/2+1/(2m))}.
\]
Letting \(m\) tend to infinity yields the universal bound
\[
  |\iota(\alpha)|\leq q^{e\beta/2}.
\tag{9.2}
\label{eq:universal-real-upper}
\]
Notice that no estimate for \(H_c^1\) was used: it occurs only in the
numerator of \eqref{eq:real-sheaf-L}.  This is the essential
noncircularity of the argument.

\medskip
\noindent\emph{Step 2: the irreducible case.}
Suppose first that \(\cG_0\) is irreducible, so its sole determinant
weight is \(\beta\).  Its dual is again \(\iota\)-real and has determinant
weight \(-\beta\).  Applying \eqref{eq:universal-real-upper} to the
eigenvalue \(\alpha^{-1}\) of the dual gives
\[
  |\iota(\alpha)|^{-1}\leq q^{-e\beta/2}.
\]
Together with \eqref{eq:universal-real-upper}, this proves
\[
  |\iota(\alpha)|=q^{e\beta/2}
\]
at every closed point.  Thus \(\cG_0\) is pure of its determinant weight.

\medskip
\noindent\emph{Step 3: separation by exterior powers.}
An individual constituent of a real sheaf need not be real, so Step 2
cannot simply be applied to each constituent.  Group the irreducible
constituents by strictly decreasing determinant weights
\[
  \gamma_1>\gamma_2>\cdots>\gamma_s,
\]
and let \(r_j\) be the total rank of the constituents of weight
\(\gamma_j\).  Put \(N_j=r_1+\cdots+r_{j-1}\).

The sheaf \(\bigwedge^{N_j+1}\cG_0\) is \(\iota\)-real: its local
characteristic polynomial is a universal polynomial with integer
coefficients in that of \(\cG_0\).  The largest determinant weight of
this exterior power is
\[
  r_1\gamma_1+\cdots+r_{j-1}\gamma_{j-1}+\gamma_j.
\tag{9.3}
\label{eq:exterior-largest}
\]
Indeed, the largest central weight is obtained by taking every available
direction in the higher-weight blocks and one direction in the
\(\gamma_j\)-block.

Fix a closed point of degree \(e\) and an eigenvalue \(\alpha\) belonging
to a constituent of determinant weight \(\gamma_j\).  Multiply it by all
the Frobenius eigenvalues, with multiplicity, in the blocks of weights
\(\gamma_1,\ldots,\gamma_{j-1}\).  The resulting number \(\delta\) is an
eigenvalue on \(\bigwedge^{N_j+1}\cG_0\).  The absolute value of the
product contributed by the higher blocks is known even before their
individual purity: the determinant definition gives
\[
 \left|\iota\prod_{\text{higher blocks}}\alpha_h\right|
 =
 q^{\frac e2(r_1\gamma_1+\cdots+r_{j-1}\gamma_{j-1})}.
\]
Apply the universal upper bound to the exterior power and use
\eqref{eq:exterior-largest}.  Cancellation of the known determinant
factor yields
\[
  |\iota(\alpha)|\leq q^{e\gamma_j/2}.
\tag{9.4}
\]
Now repeat the exterior-power argument for \(\cG_0^\vee\).  Its weight
blocks occur in reverse order and the eigenvalue corresponding to
\(\alpha\) is \(\alpha^{-1}\), of determinant weight \(-\gamma_j\).
The resulting upper bound for \(\alpha^{-1}\) is the reverse inequality
to \eqref{eq:universal-real-upper}.  Hence
\[
  |\iota(\alpha)|=q^{e\gamma_j/2}.
\]
Every constituent is pure of its determinant weight.  A Jordan--Holder
filtration of the original sheaf is therefore a weight filtration, so
the possibly nonsemisimple \(\cG_0\) is mixed.
\end{proof}

\begin{remark}
The tensor-power limit in this proof is a local positivity argument, not
the Kunneth tensor-power argument from Deligne's Weil I proof.  No
Lefschetz pencil, Chebotarev theorem, or coarse middle-cohomology estimate
appears.  The two arguments should not be conflated.
\end{remark}

\section{From curves to arbitrary dimension}

We require a Bertini statement that preserves irreducibility of a given
local system, not merely smoothness of the curve.

\begin{foundation}
\label{found:irreducibility-bertini}
\textbf{Bertini irreducibility capsule.}
Let \(X\) be smooth and geometrically irreducible over an algebraically
closed field, and let \(\cF\) be an irreducible lisse
\(\Qbarl\)-sheaf.  After choosing a sufficiently ample embedding, the
linear sections \(C=X\cap L\) which are smooth geometrically irreducible
curves and for which \(\cF|_C\) remains irreducible form a nonempty open
subset of the relevant Grassmannian.

The geometric input is proved by replacing the representation by a finite
collection of finite quotients that detect invariant subspaces, applying
the fundamental-group Bertini theorem to obtain surjectivity onto those
quotients, and using constructibility of the exceptional locus.  A
compactness argument in the profinite fundamental group passes from the
finite detectors to the \(\ell\)-adic representation.  Over a finite
field, a point of the nonempty Grassmannian open exists after a finite
extension; taking all Galois conjugates descends a dense open on which the
slicing conclusion holds.  This is independent of weights; see
\cite[Appendix~B]{KiehlWeissauer} and \cite{StanfordL22}.
\end{foundation}

\begin{theorem}[Real sheaves are mixed]
\label{thm:real-sheaf-mixed}
Let \(X_0\) be a separated finite-type \(k_0\)-scheme and let
\(\cG_0\) be a constructible \(\iota\)-real Weil sheaf.
\begin{enumerate}
\item The sheaf \(\cG_0\) is \(\iota\)-mixed.
\item If \(X_0\) is normal and geometrically irreducible and \(\cG_0\)
is lisse, every irreducible constituent is punctually
\(\iota\)-pure of its determinant weight.
\end{enumerate}
\end{theorem}

\begin{proof}
We use Noetherian induction on the support.  Nilpotents do not affect the
\'etale topos, so assume \(X_0\) reduced.  Choose a dense open immersion
\(j_0:U_0\hookrightarrow X_0\) such that \(U_0\) is smooth and every
cohomology sheaf under consideration is lisse.  With \(i_0\) the reduced
closed complement, there is an exact sequence
\[
  0\longrightarrow j_{0!}j_0^*\cG_0
  \longrightarrow\cG_0
  \longrightarrow i_{0*}i_0^*\cG_0
  \longrightarrow0.
\tag{9.5}
\label{eq:real-sheaf-stratification}
\]
The closed term is mixed by induction.  It remains to prove purity of the
irreducible constituents on a dense open of each irreducible component.
After a finite extension of \(k_0\) we may suppose the component
geometrically irreducible and the constituents absolutely irreducible;
mixedness descends because the original sheaf is a direct summand of
finite pushforward of its pullback.

Let \(\cF_0\) be one such constituent.  If the dimension is one,
\cref{thm:real-sheaf-affine-curve} applies after deleting finitely many
points.  In higher dimension, use
\cref{found:irreducibility-bertini}.  For every geometric point in a
suitable dense open there is a good curve through that point, defined
after a finite extension, on which \(\cF\) remains irreducible.  The
restriction of the entire \(\cG_0\) is still real, so
\cref{thm:real-sheaf-affine-curve} makes \(\cF|_C\) pure.  By
\cref{cor:dense-pullback-det-weights}, its weight is the determinant
weight of \(\cF_0\), independent of the chosen curve.  Thus all closed
points in the dense open have the required punctual weight.

The Jordan--Holder constituents of \(j_0^*\cG_0\) are therefore pure.
They give a mixed filtration on \(U_0\); extension by zero preserves the
same local Frobenius polynomials on \(U_0\) and has zero stalks outside.
Combining that filtration with
\eqref{eq:real-sheaf-stratification} and the induction filtration on the
complement proves (1).  The same argument without
stratification proves (2).
\end{proof}

\section{Real envelopes}

The Fourier transform of a pure sheaf need not visibly be pure.  We shall
first place it inside a real sheaf and invoke
\cref{thm:real-sheaf-mixed}.

\begin{lemma}[Real envelope of a pure sheaf]
\label{lem:real-envelope}
Let \(\cF_0\) be punctually \(\iota\)-pure of weight \(w\).  There is a
rank-one geometrically constant Weil sheaf \(\cL_b\), with degree-one
Frobenius scalar \(b\), such that
\[
  \cF_0\oplus(\cF_0^\vee\otimes\cL_b)
\]
is \(\iota\)-real and contains \(\cF_0\) as a direct summand.
\end{lemma}

\begin{proof}
Choose \(b\in\Qbarl^\times\) with
\(\iota(b)=q^w\); the existence uses only that \(\iota\) is a field
isomorphism.  Let \(\cL_b\) be the Weil line on which a degree \(n\)
Frobenius acts by \(b^n\).  If \(\alpha\) is a local eigenvalue at a
degree \(n\) point, purity gives
\[
 \overline{\iota(\alpha)}
 =\frac{q^{nw}}{\iota(\alpha)}
 =\iota(b^n/\alpha).
\]
Thus the eigenvalues on \(\cF_0^\vee\otimes\cL_b\) are precisely the
complex conjugates of those on \(\cF_0\), with multiplicity.  The
characteristic polynomial of their direct sum is real.
\end{proof}

\begin{warning}
The line \(\cL_b\) is generally only a Weil sheaf; \(b\) need not define a
continuous representation of \(\widehat{\mathbf Z}\).  This is exactly
why \cref{warn:weil-versus-etale} was necessary.
\end{warning}

\begin{formalization}
\label{form:real-sheaf-interface}
\textbf{Interface exported by the real-sheaf chapter.}

\begin{enumerate}
\item The predicate \(\operatorname{Real}_\iota(\cG)\), characterized
equivalently by real power traces at every closed point.
\item Positivity of the inverse local determinant for even tensor powers.
\item The positive-product radius lemma as a statement about
nonnegative formal power series.
\item The curve calculations \(H_c^0=0\) and
\(H_c^2\simeq V_{\pi_1}(-1)\), Frobenius-equivariantly.
\item The universal determinant-weight upper bound and its limiting
argument.
\item Exterior-power separation of determinant-weight blocks.
\item Bertini slicing as a separately named geometric capability.
\item The terminal theorem
\(\operatorname{Real}_\iota(\cG)\Rightarrow
\operatorname{Mixed}_\iota(\cG)\).
\item A real-envelope constructor for a pure Weil sheaf.
\end{enumerate}

The dependency audit for the terminal theorem contains trace formula,
top-degree duality, rank-one purity, algebraic monodromy, and Bertini.  It
does not contain the curve estimate of \cref{chap:curve-estimate}, Fourier
transform, or the direct-image theorem of
\cref{chap:global-weights}.
\end{formalization}




% Authoritative Fourier chapter.  The implementation is kept in the
% reviewed-v2 file so that earlier review snapshots remain immutable.
\chapter{The Fourier--Deligne Transform}
\label{chap:fourier}

Fourier transform replaces a cohomological trace by a stalk in a family.
We normalize by the shift \([1]\).  Consequently the trace function has
a minus sign, the inverse-character composite has no reflection and
equals a Tate twist, the same-character square has a reflection, and
Plancherel has the factor \(q^n\) over \(\mathbf F_{q^n}\).

\section{The Artin--Schreier kernel}

Fix a nontrivial additive character
\(\psi:k_0\to\Qbarl^\times\).  For the finite \'etale Galois cover
\[
  z\longmapsto z^q-z
\]
of \(\A^1_{k_0}\), let \(\cL_\psi\) be the rank-one isotypic summand
corresponding to \(\psi\).

\begin{proposition}[Trace of the Artin--Schreier sheaf]
\label{prop:as-trace}
The idempotent may be normalized so that, for
\(x\in\mathbf F_{q^n}\),
\[
  t_{\cL_\psi}^{(n)}(x)
  =\psi\!\left(\Tr_{\mathbf F_{q^n}/\Fq}(x)\right).
\]
Moreover
\[
  m^*\cL_\psi\simeq p_1^*\cL_\psi\otimes p_2^*\cL_\psi,
  \qquad
  \cL_\psi^\vee\simeq\cL_{\psi^{-1}}.
\]
\end{proposition}

\begin{proof}
If \(z^q-z=x\), then
\[
 z^{q^n}-z=x+x^q+\cdots+x^{q^{n-1}}
 =\Tr_{\mathbf F_{q^n}/\Fq}(x).
\]
Thus Frobenius acts on the torsor fiber by the displayed translation.
The chosen isotypic projector gives the asserted scalar.  Additivity of
the trace gives the tensor identity, and inversion of the character
gives the dual.
\end{proof}

\begin{convention}
\label{conv:as-sign}
The trace identity in \cref{prop:as-trace} is our sign test.  Replacing
\(\psi\) by \(\psi^{-1}\) changes where a reflection appears in Fourier
inversion but changes no weight estimate.
\end{convention}

For \(a\in\mathbf F_{q^n}\), write
\(\cL_{\psi_a}=m_a^*\cL_\psi\), where \(m_a(x)=ax\).

\section{Definition, stalks, and trace functions}

On \(A_x\times A_y\), with projections \(p_x,p_y\) and
\(\mu(x,y)=xy\), define
\[
 \FT_\psi(\cK_0)
 =
 Rp_{y!}\bigl(p_x^*\cK_0\otimes\mu^*\cL_\psi\bigr)[1].
\tag{10.1}\label{eq:fourier-definition}
\]

\begin{proposition}[Geometric stalk formula]
\label{prop:fourier-stalk}
For a geometric point \(\bar a\to A_y\),
\[
  \FT_\psi(\cK)_{\bar a}
  \simeq
  R\Gamma_c(A_x,\cK\otimes\cL_{\psi_a})[1],
\]
functorially and compatibly with Frobenius.
\end{proposition}

\begin{proof}
This is compact-support base change for \(p_y\); the kernel restricts on
the fiber to \(\cL_{\psi_a}\).
\end{proof}

\begin{corollary}
\label{cor:fourier-cohomology-sheaves}
If \(\cK\) is an ordinary sheaf, then
\[
 \cH^i(\FT_\psi\cK)_{\bar a}
 \simeq H_c^{i+1}(A_x,\cK\otimes\cL_{\psi_a}).
\]
In particular, concentration of the groups on the right in degree \(1\)
makes the transform an ordinary sheaf in degree \(0\).
\end{corollary}

For a Weil complex put
\[
 t_{\cK}^{(n)}(x)
 =\sum_i(-1)^i
 \Tr(F_x\mid\cH^i(\cK)_{\bar x}).
\]

\begin{proposition}[Fourier transform of the trace function]
\label{prop:fourier-trace-function}
For \(y\in\mathbf F_{q^n}\),
\[
 t_{\FT_\psi(\cK)}^{(n)}(y)
 =
 -\sum_{x\in\mathbf F_{q^n}}
 t_{\cK}^{(n)}(x)
 \psi\!\left(\Tr_{\mathbf F_{q^n}/\Fq}(xy)\right).
\tag{10.2}\label{eq:finite-fourier-trace}
\]
\end{proposition}

\begin{proof}
Apply the trace formula to the stalk complex in
\cref{prop:fourier-stalk}.  The shift \([1]\) contributes the minus
sign.
\end{proof}

\section{Fourier inversion}

Let \(a(x)=-x\).

\begin{theorem}[Fourier inversion]
\label{thm:fourier-inversion}
There are functorial Frobenius-equivariant isomorphisms
\[
 \FT_{\psi^{-1}}\FT_\psi(\cK_0)\simeq\cK_0(-1),
\tag{10.3}\label{eq:fourier-inversion}
\]
\[
 \FT_\psi^2(\cK_0)\simeq a^*\cK_0(-1).
\]
\end{theorem}

\begin{proof}
The inverse-character composite has kernel
\[
 Rp_{xz!}\cL_\psi(y(x-z))[2].
\]
For \(x\ne z\), all compactly supported cohomology of the nontrivial
Artin--Schreier sheaf on \(A_y\) vanishes.  For \(x=z\),
\[
 R\Gamma_c(A_y,\Qbarl)=\Qbarl(-1)[-2].
\]
Thus the kernel is the extension by zero of \(\Qbarl(-1)\) from the
diagonal.  With \(\psi\) in both transforms the phase is \(y(x+z)\), so
the kernel is supported on the antidiagonal.  On functions:
\[
 \widehat{\widehat f_\psi}_{\psi^{-1}}(z)=q^nf(z),
 \qquad
 \widehat{\widehat f_\psi}_{\psi}(z)=q^nf(-z).
\]
\end{proof}

\begin{corollary}
\label{cor:fourier-equivalence}
\(\FT_\psi\) is an equivalence with quasi-inverse
\(\FT_{\psi^{-1}}(1)\), hence is conservative.
\end{corollary}

\section{Constructibility and perversity}

\begin{foundation}
\label{found:fourier-perverse}
\textbf{Perverse exactness capsule.}
On the affine line,
\[
 \FT_\psi:\operatorname{Perv}(A_x,\Qbarl)
 \xrightarrow{\sim}\operatorname{Perv}(A_y,\Qbarl)
\]
is exact.  Stalk bounds follow from affine Artin vanishing and
\cref{prop:fourier-stalk}; costalk bounds follow by Verdier duality and
\[
 \mathbb D\FT_\psi(\cK)
 \simeq\FT_{\psi^{-1}}(\mathbb D\cK)(1).
\]
Fourier inversion gives essential surjectivity.  See
\cite{LaumonFourier},
\cite[Chapter~I, Section~5]{KiehlWeissauer}, and
\cite{StanfordL23}.
\end{foundation}

\section{Plancherel}

For \(f:\mathbf F_{q^n}\to\mathbf C\), set
\(\|f\|_{2,n}^2=\sum_x|f(x)|^2\).

\begin{theorem}[Plancherel]
\label{thm:fourier-plancherel}
For every constructible Weil complex,
\[
 \|\iota t_{\FT_\psi(\cK)}^{(n)}\|_{2,n}^2
 =
 q^n\|\iota t_{\cK}^{(n)}\|_{2,n}^2.
\tag{10.4}\label{eq:plancherel}
\]
\end{theorem}

\begin{proof}
Expand the square in \cref{eq:finite-fourier-trace}.  Character
orthogonality gives
\[
 \sum_y\psi_n((x-x')y)
 =\begin{cases}q^n,&x=x',\\0,&x\ne x'.\end{cases}
\]
\end{proof}

\section{\texorpdfstring{\(L^2\)}{L2} growth and the largest weight}

The next result is the determinant-weight radius theorem.  Its delicate
point is not Pringsheim's theorem by itself: positivity does not by
itself prevent cancellation in a rational function.  The required
noncancellation is an isotypic monodromy calculation.

\begin{foundation}
\begin{theorem}[\(L^2\)-growth criterion]
\label{thm:l2-weight-criterion}
Let \(C_0\) be a geometrically connected smooth curve and let
\(\cG_0\ne0\) be a constructible \(\iota\)-mixed Weil sheaf.  Assume
that no maximal-weight subquotient is punctual and that
\(H_c^0(C,\cG)=0\).  If \(w_{\max}\) is the largest constituent weight,
then
\[
 \limsup_{n\to\infty}\frac1n
 \log_q\|\iota t_{\cG}^{(n)}\|_{2,n}^2
 =w_{\max}+1.
\tag{10.5}\label{eq:l2-weight-growth}
\]
The same formula holds after restricting \(n\) to multiples of any
fixed positive integer.
\end{theorem}

\smallskip
\textbf{Determinant-weight proof chain.}
Semisimplification does not change trace functions.  Write its
nonpunctual part as a direct sum of pure irreducible constituents
\(\cM_\lambda\), and form the conjugate sheaf constituentwise:
\[
 \cM_\lambda^c
 =\cM_\lambda^\vee\otimes\cL_{b_\lambda},
 \qquad
 \iota(b_\lambda)=q^{w_\lambda}.
\]
Then the local eigenvalue multiset of \(\cG^c\) is the complex conjugate
of that of \(\cG\).  If
\[
 S_n=\|\iota t_{\cG}^{(n)}\|_{2,n}^2,
\]
the trace formula gives
\[
 \exp\!\left(\sum_{n\ge1}S_n\frac{T^n}{n}\right)
 =L(C_0,\cG_0\otimes\cG_0^c,T).
\tag{10.6}\label{eq:l2-L-function}
\]

Pass to a finite base extension for which
\cref{prop:positive-degree-center} supplies a central degree-one element
\(z\), and decompose the maximal-weight part into its \(z\)-isotypic
blocks.  On a block \(M\), \(z\) acts by a scalar.  In
\[
 M\otimes M^c\simeq\End(M)\otimes L_{b_M},
\]
the geometrically constant quotient is
\[
 \End_{\pi_1(C)}(M)\otimes L_{b_M}.
\]
Its multiplicity
\[
 m_M=\dim\End_{\pi_1(C)}(M)
\]
is strictly positive.  Grouping geometrically isomorphic blocks, the
diagonal multiplicity is the positive Gram number
\[
 \sum_{\Gamma}
 \dim\End_{\pi_1(C)}
 \left(\bigoplus_{\lambda\in\Gamma}M_\lambda\right)>0.
\tag{10.7}\label{eq:positive-diagonal-multiplicity}
\]
Top-degree duality adds \((-1)\).  The diagonal quotient therefore has
cohomological weight \(2w_{\max}+2\), equivalently its Frobenius
eigenvalues have absolute-value exponent \(w_{\max}+1\).

The determinant-weight lemma of Kiehl--Weissauer compares, one
\(z\)-isotypic block at a time, the divisor of the cohomological
determinant in \eqref{eq:l2-L-function}.  Its proof uses the central
decomposition and the exact invariant--coinvariant sequence:
off-diagonal blocks pair with their conjugates, whereas the order in
the diagonal block is the positive integer in
\eqref{eq:positive-diagonal-multiplicity}.  Thus this top-degree factor
survives in the reduced determinant quotient.  The hypotheses exclude
the two extra terms in that exact sequence: a maximal punctual block and
a compactly-supported degree-zero block.  Lower-weight isotypes are
holomorphic on the outermost circle.  Consequently the reduced
logarithmic derivative has radius
\[
 q^{-(w_{\max}+1)}.
\]
This is the noncancellation step; coefficient positivity is not used as
a substitute for it.

Cauchy--Hadamard gives
\[
 \limsup_n S_n^{1/n}=q^{w_{\max}+1}.
\]
Replacing Frobenius by a fixed power proves descent from the finite base
extension and the assertion about multiples.  See
\cite[Chapter~I, Sections~2.12--3.2]{KiehlWeissauer} and
\cite{StanfordL24}.  This result depends on rank-one purity and the
semisimple finite-center monodromy lemma, but not on the curve estimate
or the global direct-image theorem.
\end{foundation}

\begin{warning}
\label{warn:l2-needs-mixed}
\Cref{eq:l2-weight-growth} is false as a general definition of weight.
A punctual constituent lacks the extra \(q^n\) supplied by the number
of points, and without mixedness there is no maximal pure central
isotype.
\end{warning}

\begin{formalization}
\label{form:fourier-interface}
\textbf{Fourier interface.}
\begin{enumerate}
\item Artin--Schreier kernel with its local trace normalization.
\item Fourier transform, stalk base change, and the shifted trace
formula.
\item The inversion identities
\[
 \FT_{\psi^{-1}}\FT_\psi\simeq(-1),
 \qquad
 \FT_\psi^2\simeq a^*(-1).
\]
\item Constructibility, perverse exactness, and Plancherel.
\item The \(L^2\)-growth theorem as a separate determinant-weight
capability, with its explicit nonpunctual and \(H_c^0\) hypotheses.
\end{enumerate}
\end{formalization}

\section{A detailed proof of perverse exactness}
\label{sec:fourier-perverse-exactness-detail}

We now expand the perverse-exactness capsule
\cref{found:fourier-perverse}.  The only input beyond the six operations and
perverse Artin vanishing is the following support theorem.  It is the step
that permits Verdier duality to be used without changing the Fourier
transform from a \(!\)-pushforward to a \(*\)-pushforward.

\begin{theorem}[Laumon's support-forgetting theorem]
\label{thm:fourier-support-forgetting}
For \(\cK\in\Dbc(A_x,\Qbarl)\), the canonical morphism
\[
 Rp_{y!}\bigl(p_x^*\cK\otimes\mu^*\cL_\psi\bigr)
 \longrightarrow
 Rp_{y*}\bigl(p_x^*\cK\otimes\mu^*\cL_\psi\bigr)
\tag{10.15}\label{eq:fourier-support-forgetting}
\]
is an isomorphism.
\end{theorem}

\begin{proof}
Compactify the \(x\)-line by
\(a:\A^1_x\hookrightarrow\mathbf P^1_x\), let
\(\eta_\infty\) be the generic point of the strict henselian trait at
\(\infty_x\), and let \(I_\infty\) be its inertia group.  The localization
triangle for this compactification identifies the cone of
\eqref{eq:fourier-support-forgetting} canonically with the
inertia-cohomology contribution
\[
 R\Gamma\!\left(
 I_\infty,\,
 \cK_{\bar\eta_\infty}\otimes
 R\Psi_{\infty_x}\bigl((a\times1)_!\mu^*\cL_\psi\bigr)
 \right).
\tag{10.17}\label{eq:fourier-boundary-inertia}
\]
Here the nearby-cycle factor is a complex over the \(y\)-line; this is a
global family statement, not a comparison of \(!\)- and \(*\)-cohomology on
each individual fiber.

Laumon's universal local-acyclicity theorem gives
\[
 R\Psi_{\infty_x}\bigl((a\times1)_!\mu^*\cL_\psi\bigr)=0.
\]
Consequently the complex in
\eqref{eq:fourier-boundary-inertia} vanishes, so the support-forgetting cone
vanishes.  This is \cite[Thm.~1.3.1.1]{LaumonFourier}; the displayed
nearby-cycle vanishing is Theorem 1.3.1.2 and Remark 1.3.1.3 there.
\end{proof}

\begin{proposition}[Fourier transform and Verdier duality]
\label{prop:fourier-verdier-detailed}
There is a functorial isomorphism
\[
 \D_{A_y}\FT_\psi(\cK)
 \simeq
 \FT_{\psi^{-1}}\bigl(\D_{A_x}\cK\bigr)(1).
\tag{10.16}\label{eq:fourier-verdier-detailed}
\]
\end{proposition}

\begin{proof}
Put
\(\cM=p_x^*\cK\otimes\mu^*\cL_\psi\).
Relative Verdier duality and the shift in the definition of \(\FT_\psi\)
give
\[
 \D_{A_y}\bigl(Rp_{y!}\cM[1]\bigr)
 \simeq Rp_{y*}\D_{A_x\times A_y}(\cM)[-1].
\]
The projection \(p_x\) is smooth of relative dimension \(1\), and
\(\cL_\psi^\vee=\cL_{\psi^{-1}}\).  Thus
\[
 \D_{A_x\times A_y}(\cM)
 \simeq
 p_x^*(\D_{A_x}\cK)\otimes\mu^*\cL_{\psi^{-1}}[2](1).
\]
Substitution yields
\[
 Rp_{y*}\bigl(
 p_x^*(\D\cK)\otimes\mu^*\cL_{\psi^{-1}}\bigr)[1](1).
\]
Apply \cref{thm:fourier-support-forgetting}, with
\(\psi^{-1}\) in place of \(\psi\), to replace \(Rp_{y*}\) by
\(Rp_{y!}\).  The result is exactly
\eqref{eq:fourier-verdier-detailed}.  This is Laumon
1.3.2.2; the twist is forced by
\(R\Gamma_c(\A^1,\Qbarl)=\Qbarl(-1)[-2]\), the same normalization check used
in Fourier inversion.
\end{proof}

\begin{theorem}[Perverse \(t\)-exactness of Fourier--Deligne transform]
\label{thm:fourier-perverse-texact-detailed}
The functor
\[
 \FT_\psi:\Perv(A_x,\Qbarl)\longrightarrow\Perv(A_y,\Qbarl)
\]
is exact.  Together with Fourier inversion, it is an equivalence of abelian
categories with quasi-inverse \(\FT_{\psi^{-1}}(1)\).
\end{theorem}

\begin{proof}
Let \(\cP\) be perverse on \(A_x\).  Since \(p_x\) is smooth of relative
dimension \(1\), smooth perverse pullback gives
\[
 p_x^*\cP[1]\in\Perv(A_x\times A_y).
\]
Tensoring with the lisse rank-one sheaf
\(\mu^*\cL_\psi\) preserves perversity.  The projection \(p_y\) is affine,
so perverse Artin vanishing says that \(Rp_{y!}\) is left \(t\)-exact.
Using the defining shift inside the displayed perverse object, we obtain
\[
 \FT_\psi(\cP)
 =Rp_{y!}\bigl(p_x^*\cP[1]\otimes\mu^*\cL_\psi\bigr)
 \in{}^p\D^{\geq0}(A_y).
\]

Verdier duality is \(t\)-exact for the middle perverse structure, so
\(\D\cP\) is perverse.  Applying the preceding lower bound with
\(\psi^{-1}\), and then \cref{prop:fourier-verdier-detailed}, gives
\[
 \D\FT_\psi(\cP)
 \simeq\FT_{\psi^{-1}}(\D\cP)(1)
 \in{}^p\D^{\geq0}(A_y).
\]
Because duality exchanges
\({}^p\D^{\leq0}\) and \({}^p\D^{\geq0}\), this is equivalent to
\(\FT_\psi(\cP)\in{}^p\D^{\leq0}(A_y)\).  Hence the transform lies in the
perverse heart.  A triangulated \(t\)-exact functor induces an exact functor
on hearts.  Finally \cref{thm:fourier-inversion} identifies
\(\FT_{\psi^{-1}}(1)\) as its quasi-inverse.  This is
\cite[Thm.~1.3.2.3]{LaumonFourier}.
\end{proof}

\begin{perversebox}
The argument uses both one-sided affine bounds.  The \(!\)-definition gives
the lower perverse bound directly; the upper bound is obtained only after
duality and support-forgetting.  Treating \(Rp_{y!}\) as \(t\)-exact from
affineness alone would be incorrect.
\end{perversebox}


\chapter{The Fundamental Estimate on Curves}
\label{chap:curve-estimate}

We now prove the one-dimensional estimate from which the global
direct-image theorem will follow.  Cohomology in degree \(i\) raises an
upper weight bound by at most \(i\).  The proof uses the real-sheaf
theorem, Fourier inversion, Plancherel, the determinant-weight
\(L^2\)-criterion, local duality, and elementary geometric reductions.

\section{Statement and the intermediate-extension form}

\begin{theorem}[Curve estimate]
\label{thm:curve-estimate}
Let \(C_0\) be a smooth geometrically irreducible curve over \(k_0\),
not necessarily proper, and let \(\cF_0\) be a lisse Weil sheaf
punctually \(\iota\)-pure of weight \(w\).  Then
\[
  H_c^i(C,\cF)
\]
is \(\iota\)-mixed of weights at most \(w+i\) for every \(i\).
\end{theorem}

Let \(\bar C_0\) be the smooth projective completion and
\(j_0:C_0\hookrightarrow\bar C_0\).

\begin{theorem}[Intersection-cohomology form]
\label{thm:curve-intersection-purity}
Under the hypotheses of \cref{thm:curve-estimate},
\[
 H^i(\bar C,j_*\cF)
\]
is punctually \(\iota\)-pure of weight \(w+i\) whenever it is nonzero.
\end{theorem}

\section{Boundary weights and compact support}

The needed boundary assertion is Deligne's semicontinuity theorem.  A
local-monodromy filtration alone does not prove it: that would silently
import a weight--monodromy statement.  The following proof instead uses
an \(L\)-function and a tensor-power limit.

\begin{foundation}
\label{found:boundary-semicontinuity}
\textbf{Semicontinuity at the boundary.}
Let \(j_0:U_0\hookrightarrow C_0\) be a dense open immersion of smooth
geometrically irreducible curves, and let \(\cF_0\) be lisse and pure of
weight \(w\) on \(U_0\).  Then every Frobenius eigenvalue on a boundary
stalk of \(j_{0*}\cF_0\) has weight at most \(w\).

\smallskip
\emph{Proof.}
Fix \(s_0\in C_0-U_0\), of degree \(e\), and an eigenvalue
\(\alpha\) on
\[
  (j_*\cF)_{\bar s}=V^{I_s}.
\]
Delete a closed point of \(U_0\), if necessary, so that \(C_0\) is
affine; this changes neither the boundary stalk nor the assertion.  For
\(m\ge1\), put
\[
  \cG_{m,0}=j_{0*}(\cF_0^{\otimes m}).
\]
The Frobenius-equivariant injection
\[
  (V^{I_s})^{\otimes m}\hookrightarrow(V^{\otimes m})^{I_s}
\]
shows that \(\alpha^m\) occurs on
\((\cG_m)_{\bar s}\).

The adjunction
\(\cG_{m,0}\to j_{0*}j_0^*\cG_{m,0}\) is an isomorphism, so
\(\cG_m\) has no section supported on the boundary.  On the affine
curve this gives \(H_c^0(C,\cG_m)=0\).  Removing finitely many points
does not change \(H_c^2\), and top-degree duality gives
\[
 H_c^2(C,\cG_m)
 \simeq(V^{\otimes m})_{\pi_1(U)}(-1).
\]
The coinvariants are a geometrically constant quotient of the pure
sheaf \(\cF_0^{\otimes m}\), so the eigenvalues on this group have
weights at most \(mw+2\).  Hence the trace formula
\[
 L(C_0,\cG_{m,0},T)
 =
 \frac{\det(1-TF\mid H_c^1(C,\cG_m))}
      {\det(1-TF\mid H_c^2(C,\cG_m))}
\tag{11.1}\label{eq:boundary-trace-ratio}
\]
has no pole for
\[
  |T|<q^{-mw/2-1}.
\tag{11.2}\label{eq:boundary-pole-disk}
\]

Purity on \(U_0\), the logarithmic derivative of its Euler product, and
\(\#U_0(\mathbf F_{q^n})=O(q^n)\) show that
\(L(U_0,\cF_0^{\otimes m},T)\) converges absolutely and is nonzero in
the same disk.  Euler products factor as
\[
 L(C_0,\cG_{m,0},T)
 =
 L(U_0,\cF_0^{\otimes m},T)
 \prod_{t_0\in C_0-U_0}
 \det\!\left(
  1-T^{\deg(t_0)}F_{t_0}\mid(\cG_m)_{\bar t}
 \right)^{-1}.
\tag{11.3}\label{eq:boundary-euler-factorization}
\]
After division by the nonvanishing first factor, the finite boundary
product has no pole in \cref{eq:boundary-pole-disk}.  It has a
denominator but no numerator, so its individual local factors cannot
cancel one another.  Since \(\alpha^m\) occurs in the factor at \(s_0\),
\[
 |\iota(\alpha)|^{m/e}\le q^{mw/2+1}.
\]
Thus
\[
 |\iota(\alpha)|\le q^{e(w/2+1/m)}.
\]
Letting \(m\to\infty\) gives
\(|\iota(\alpha)|\le q^{ew/2}\).

This proof uses only the trace formula, top-degree duality, elementary
Euler-product convergence, and purity on the open; it uses neither
weight--monodromy nor the global direct-image theorem.  See
\cite[Proposition~2.3.0.1]{StanfordL19} and
\cite[Chapter~I, Section~2]{KiehlWeissauer}.
\end{foundation}

Let
\(\cQ_0=j_{0*}\cF_0/j_{0!}\cF_0\).  It is punctual of weights at most
\(w\), and
\[
 0\longrightarrow j_!\cF\longrightarrow j_*\cF
 \longrightarrow\cQ\longrightarrow0.
\tag{11.4}\label{eq:jshriek-jstar}
\]
The long exact sequence contains
\[
 0\to H_c^0(C,\cF)\to H^0(\bar C,j_*\cF)
 \to H^0(\bar C,\cQ)\to H_c^1(C,\cF)
 \to H^1(\bar C,j_*\cF)\to0
\tag{11.5}\label{eq:boundary-long-exact}
\]
and
\[
 H_c^2(C,\cF)\simeq H^2(\bar C,j_*\cF).
\]
Thus \cref{thm:curve-intersection-purity} implies
\cref{thm:curve-estimate}.

\section{Reduction to an open subset of the affine line}

Finite extension of \(k_0\) does not change normalized weights.  Delete
finitely many points so that \(C_0\) is affine.  Noether normalization
gives a finite generically \'etale map to \(\A^1\); after deleting the
branch locus it is finite \'etale.  Finite pushforward preserves purity:
Frobenius permutes the stalks above a point, and each orbit block is
obtained from the Frobenius polynomial upstairs by a power substitution.
Extending the finite map to projective completions gives
\[
 R\Gamma(\bar C,j_*\cF)
 \simeq R\Gamma(\PP^1,j_*g_*\cF).
\]
We may therefore replace the pair by an open
\(U_0\subset\PP^1_{k_0}\) and a pure lisse sheaf on it.

Pass to a Jordan--Holder constituent and, after a finite base extension,
assume it geometrically irreducible.  A geometrically constant
constituent has \(H^1(\PP^1,j_*\cF)=0\), so assume geometric
nonconstancy.  Choose \(s\in U_0(k_0)\), remove it, and send it to
\(\infty\).  Because the original sheaf extended across \(s\), this does
not change \(j_*\cF\), and the resulting sheaf on
\[
 U_0\subset\A^1_{k_0}
\]
is unramified at infinity.  We have reduced to a geometrically
irreducible, geometrically nonconstant sheaf unramified at infinity.

\section{The Fourier transform is a mixed sheaf}

Let \(j:U\hookrightarrow\A^1\), put
\[
 \cK_0=j_!\cF_0,\qquad
 \widehat{\cF}_0=\FT_\psi(\cK_0).
\]

\begin{lemma}
\label{lem:twisted-cohom-concentration}
For every geometric \(a\in\A^1(k)\),
\[
 H_c^i(U,\cF\otimes\cL_{\psi_a})=0\qquad(i\ne1).
\]
Consequently \(\widehat{\cF}_0\) is an ordinary constructible Weil
sheaf in degree \(0\).
\end{lemma}

\begin{proof}
Compactly supported \(H^0\) of a lisse sheaf on the connected affine
curve vanishes.  Top-degree duality identifies \(H_c^2\) with the
geometric coinvariants, twisted by \((-1)\).  Since
\(\cF\otimes\cL_{\psi_a}\) is geometrically irreducible, nonzero
coinvariants would make it geometrically constant.  For \(a=0\) this
contradicts the nonconstancy of \(\cF\).  For \(a\ne0\), it contradicts
the fact that \(\cF\) is unramified at infinity whereas
\(\cL_{\psi_a}\) has nontrivial wild inertia there.  The stalk formula
\cref{cor:fourier-cohomology-sheaves} proves the last assertion.
\end{proof}

\begin{lemma}
\label{lem:fourier-mixed-before-estimate}
The sheaf \(\widehat{\cF}_0\) is \(\iota\)-mixed.
\end{lemma}

\begin{proof}
Let \(\cF_0^c\) be the conjugate summand supplied by
\cref{lem:real-envelope}.  Apply the preceding construction to
\((\cF_0,\psi)\) and \((\cF_0^c,\psi^{-1})\).  Their Fourier trace
functions are complex conjugates, so
\[
 \FT_\psi(j_!\cF_0)\oplus
 \FT_{\psi^{-1}}(j_!\cF_0^c)
\]
is \(\iota\)-real.  It is mixed by
\cref{thm:real-sheaf-mixed}, and so is its direct summand
\(\widehat{\cF}_0\).
\end{proof}

\begin{lemma}
\label{lem:fourier-no-hc0}
\[
 H_c^0(\A^1,\widehat{\cF})=0.
\]
In particular, the semisimplification of \(\widehat{\cF}\) has no
punctual constituent.
\end{lemma}

\begin{proof}
The inverse-transform stalk formula and Fourier inversion give
\[
 H_c^0(\A^1,\widehat{\cF})
 \simeq
 \cH^{-1}\!\left(
  \FT_{\psi^{-1}}(\widehat{\cF})
 \right)_0
 \simeq
 \cH^{-1}\bigl((j_!\cF)(-1)\bigr)_0=0.
\]
A punctual direct summand would have a nonzero compactly supported
section.
\end{proof}

\section{Plancherel gives the upper bound}

The sheaf \(\cK_0=j_!\cF_0\) is pure of weight \(w\), has no punctual
constituent, and has \(H_c^0=0\).  Thus
\cref{thm:l2-weight-criterion} gives
\[
 \limsup_n\frac1n\log_q
 \|\iota t_{\cK}^{(n)}\|_{2,n}^2=w+1.
\tag{11.6}\label{eq:K-l2-growth}
\]
Plancherel gives
\[
 \|\iota t_{\widehat{\cF}}^{(n)}\|_{2,n}^2
 =q^n\|\iota t_{\cK}^{(n)}\|_{2,n}^2.
\]
The left exponent is therefore \(w+2\).  The Fourier sheaf is mixed and
has neither \(H_c^0\) nor punctual constituents, so applying the
\(L^2\)-criterion again gives
\[
  \widehat w_{\max}=w+1.
\tag{11.7}\label{eq:fourier-weight-shift}
\]
At \(a=0\), the stalk formula identifies
\[
 \widehat{\cF}_{\bar0}\simeq H_c^1(U,\cF).
\]

\begin{proposition}[The fundamental upper estimate]
\label{prop:fundamental-curve-upper}
Every eigenvalue on \(H_c^1(U,\cF)\) has weight at most \(w+1\).  The
same holds on its quotient \(H^1(\PP^1,j_*\cF)\).
\end{proposition}

\section{Duality gives equality}

We have
\[
 H^0(\bar C,j_*\cF)=H^0(C,\cF).
\]
Its nonzero constituents are geometrically constant quotients of
\(\cF\), hence have weight \(w\).  Local duality gives
\[
 H^2(\bar C,j_*\cF)
 \simeq H^0(\bar C,j_*\cF^\vee)^\vee(-1),
\]
so its nonzero eigenvalues have weight \(w+2\).

In degree one, local duality gives a perfect pairing
\[
 H^1(\bar C,j_*\cF)\otimes
 H^1(\bar C,j_*\cF^\vee)
 \longrightarrow\Qbarl(-1).
\tag{11.8}\label{eq:curve-local-duality}
\]
If \(\alpha\) occurs in the first factor, \(q/\alpha\) occurs in the
second.  Apply \cref{prop:fundamental-curve-upper} to the dual sheaf,
which has weight \(-w\):
\[
 |\iota(q/\alpha)|\le q^{(-w+1)/2}.
\]
Thus
\[
 |\iota(\alpha)|\ge q^{(w+1)/2}.
\]
The upper estimate for \(\cF\) gives the reverse inequality.  This
proves \cref{thm:curve-intersection-purity}; the boundary exact sequence
then proves \cref{thm:curve-estimate}.

\begin{warning}
\label{warn:curve-proof-noncircular}
The boundary estimate came from an \(L\)-function tensor-power limit,
not from weight--monodromy.  The lower middle-degree estimate came from
the already proved upper estimate applied to the dual, not from
smooth-proper purity.
\end{warning}

\begin{formalization}
\label{form:curve-estimate-interface}
\textbf{Formalization decomposition.}
\begin{enumerate}
\item Boundary stalks as inertia invariants and Deligne's
\(L\)-function/tensor-power semicontinuity proof.
\item The \(j_!\)-to-\(j_*\) exact sequence.
\item Finite-map reduction to \(\PP^1\).
\item Vanishing for every Artin--Schreier twist.
\item Fourier transform as an ordinary mixed sheaf with \(H_c^0=0\).
\item Plancherel plus the determinant-weight \(L^2\)-criterion.
\item Local duality for the reverse bound.
\item The terminal implication
\[
 \operatorname{Pure}(\cF,w)
 \Longrightarrow
 \operatorname{WeightsLE}(H_c^i(C,\cF),w+i).
\]
\end{enumerate}
No step imports the global \(Rf_!\) theorem.
\end{formalization}

\chapter{The Deligne--Laumon Direct-Image Theorem}
\label{chap:global-weights}

We now pass from the estimate on a single curve to the global theorem on
weights.  The reduction has two independent parts.  Geometry and
Noetherian induction reduce an arbitrary morphism to a relative affine
curve.  Proper-support base change then identifies the stalks of its
higher direct images with the compactly supported cohomology treated in
\cref{thm:curve-estimate}.  Mixedness is not inferred merely from
pointwise inequalities: it is obtained from real envelopes and
\cref{thm:real-sheaf-mixed}.  This separation is essential.

\section{The theorem}

\begin{theorem}[Deligne--Laumon direct-image theorem]
\label{thm:deligne-laumon-direct-image}
Let
\[
  f_0:X_0\longrightarrow Y_0
\]
be a separated morphism of finite type between schemes of finite type
over \(k_0\).  If a constructible Weil sheaf \(\cF_0\) on \(X_0\) is
\(\iota\)-mixed of weights at most \(w\), then
\[
  R^i f_{0!}\cF_0
\]
is \(\iota\)-mixed of weights at most \(w+i\) for every \(i\).

More generally, if a constructible Weil complex \(\cK_0\) has weights at
most \(w\), in the sense that
\[
  \cH^j(\cK_0)\text{ has weights at most }w+j,
\]
then \(Rf_{0!}\cK_0\) has weights at most \(w\).
\end{theorem}

The theorem is local on \(Y_0\), but it is not enough to inspect a single
fiber of an arbitrary-dimensional morphism: that would ask for the very
theorem we are proving.  We first reduce the relative dimension to one.

\section{Formal reductions}

We list the reductions as lemmas so that none is hidden in the phrase
``standard devissage''.

\begin{lemma}[Quasi-finite case]
\label{lem:weights-quasifinite}
If \(f_0\) is separated and quasi-finite, then \(R^if_{0!}=0\) for
\(i>0\), and \(f_{0!}\) preserves upper weight bounds.
\end{lemma}

\begin{proof}
By Zariski's main theorem, locally on the target \(f_0\) factors as an
open immersion followed by a finite morphism.  Extension by zero is
exact, and finite pushforward is exact, so the higher direct images
vanish.  At a geometric point of the target, \(f_!\cF\) is the direct
sum of stalks over the finite fiber, with zero contributions from points
outside the chosen open.  Frobenius permutes the summands in finite
orbits.  The characteristic polynomial of one orbit is obtained from
the Frobenius polynomial at the corresponding closed point upstairs by
replacing \(T\) by a power.  As in the finite-map reduction in
\cref{chap:curve-estimate}, the normalized absolute value and the residue
degree are multiplied by the same integer.  Thus weights do not
increase.
\end{proof}

\begin{lemma}[Composition]
\label{lem:weights-composition}
Suppose the conclusion of
\cref{thm:deligne-laumon-direct-image} is known for
\(h_0:X_0\to Z_0\) and \(g_0:Z_0\to Y_0\).  Then it is known for
\(f_0=g_0h_0\).
\end{lemma}

\begin{proof}
The compact-support Leray spectral sequence is Frobenius-equivariant:
\[
  E_2^{a,b}
  =
  R^ag_{0!}R^bh_{0!}\cF_0
  \Longrightarrow
  R^{a+b}f_{0!}\cF_0.
\]
If \(\cF_0\) has weights at most \(w\), the first application gives
weights at most \(w+b\), and the second gives weights at most
\(w+b+a\).  The abutment has a finite Frobenius-stable filtration whose
graded pieces are subquotients of \(E_2^{a,b}\).  Weight bounds pass to
subquotients and extensions by
\cref{prop:elementary-weight-permanence}.  No degeneration of the
spectral sequence is asserted or needed.
\end{proof}

\begin{lemma}[Open--closed devissage]
\label{lem:weights-open-closed}
Let \(j_0:U_0\hookrightarrow X_0\) be open with closed complement
\(i_0:Z_0\hookrightarrow X_0\).  If the theorem is known for the
restrictions of \(f_0\) to \(U_0\) and \(Z_0\), then it is known for
\(f_0\).
\end{lemma}

\begin{proof}
Use the distinguished triangle
\[
  j_{0!}j_0^*\cK_0
  \longrightarrow\cK_0
  \longrightarrow i_{0*}i_0^*\cK_0
  \xrightarrow{+1}.
\]
After applying \(Rf_{0!}\), the long exact cohomology sequence expresses
each cohomology sheaf as an extension of a quotient and a subobject of
the two restricted terms.  Upper weight bounds survive all three
operations.
\end{proof}

\begin{lemma}[Generic relative Noether normalization]
\label{lem:relative-noether-normalization}
After replacing \(X_0\) and \(Y_0\) by dense opens on each irreducible
component, a dominant affine morphism \(f_0:X_0\to Y_0\) of relative
dimension \(r\) admits a factorization
\[
 X_0\xrightarrow{\nu_0}\A^r_{Y_0}
 \xrightarrow{\pi_0}Y_0
\]
with \(\nu_0\) finite.  The projection \(\pi_0\) is a composition of
\(r\) relative affine-line projections.
\end{lemma}

\begin{proof}
Apply Noether normalization to the finitely generated extension of the
function field of \(Y_0\) by that of \(X_0\).  The resulting algebraically
independent functions and the integral dependence relations spread out
after inverting finitely many functions on \(X_0\) and \(Y_0\).  They
define the finite map to affine \(r\)-space.  The last assertion is the
coordinate-by-coordinate factorization of affine space.
\end{proof}

\section{Reduction to a relative affine curve}

We prove the sheaf statement of
\cref{thm:deligne-laumon-direct-image} by simultaneous Noetherian
induction on the supports in source and target.

Use a finite stratification and a weight filtration to reduce to
\(\cF_0\) pure and lisse on a smooth, irreducible, locally closed
stratum of \(X_0\).  Extension by zero and
\cref{lem:weights-open-closed} allow us to replace \(X_0\) by that
stratum.  Shrinking the target to the dense image, and treating the
closed complement by induction, we may assume that \(f_0\) is dominant
and both source and target are affine.  Relative Noether normalization
and \cref{lem:weights-quasifinite,lem:weights-composition} reduce the
problem to a projection
\[
  \pi_0:\A^1_{Y_0}\longrightarrow Y_0.
\]

The finite pushforward from the normalization step is constructible, not
necessarily lisse everywhere.  Stratify \(\A^1_{Y_0}\), use
\cref{lem:weights-open-closed}, and shrink the base once more.  We arrive
at
\[
  h_0:U_0\longrightarrow Y_0,
\]
where \(U_0\) is a dense open of \(\A^1_{Y_0}\), \(h_0\) is smooth and
surjective, all geometric fibers are nonempty geometrically irreducible
affine curves, and \(\cF_0\) is lisse and pure of weight \(w\).

All discarded closed subsets have strictly smaller support and are
covered by the induction.  This proves that it is enough to establish
the theorem in the relative affine-curve case.

\section{Fiberwise upper bounds}

Let \(y_0\) be a closed point of \(Y_0\), with residue field
\(k(y_0)\) of cardinality \(q_y=q^{\deg(y_0)}\), and choose a geometric
point \(\bar y\) above it.  Proper-support base change gives a
Frobenius-equivariant isomorphism
\[
  \left(R^ih_{0!}\cF_0\right)_{\bar y}
  \simeq
  H_c^i(U_{\bar y},\cF|_{U_{\bar y}}).
\tag{12.1}
\label{eq:relative-curve-base-change}
\]
The restriction of \(\cF_0\) to the fiber is pure of weight \(w\)
relative to \(q_y\): a degree \(n\) point over \(k(y_0)\) is a degree
\(n\deg(y_0)\) point over \(k_0\).  Applying
\cref{thm:curve-estimate} over \(k(y_0)\) shows that every eigenvalue in
\eqref{eq:relative-curve-base-change} has weight at most \(w+i\).
Thus \(R^ih_{0!}\cF_0\) is punctually of weights at most \(w+i\).

Punctual bounds do not by themselves construct a mixed filtration on
the base.  We now supply it without invoking the theorem recursively.

\section{Mixedness in the relative-curve case}

\begin{lemma}
\label{lem:relative-curve-mixedness}
In the relative affine-curve situation above,
\[
  R^ih_{0!}\cF_0
\]
is \(\iota\)-mixed of weights at most \(w+i\).
\end{lemma}

\begin{proof}
It is enough first to treat a pure \(\cF_0\).  Form its real envelope
\[
  \cE_0=\cF_0\oplus\cF_0^c
\]
as in \cref{lem:real-envelope}; both summands are pure of weight \(w\),
and \(\cE_0\) is \(\iota\)-real.  Since \(R^ih_!\cF\) is a direct
summand of \(R^ih_!\cE\), it will be enough to prove that the latter is
mixed.

For a geometric affine-curve fiber, \(H_c^0=0\).  Top cohomology is
\[
 H_c^2(U_{\bar y},\cE)
 \simeq
 \left(E_{\bar\eta}\right)_{\pi_1(U_{\bar y})}(-1).
\]
The coinvariant quotient is geometrically constant.  Since \(\cE_0\)
is pure of weight \(w\), every nonzero eigenvalue on this top cohomology
has weight exactly \(w+2\).  By
\cref{thm:curve-estimate}, the eigenvalues on \(H_c^1\) have weights at
most \(w+1\).

The fiber \(L\)-function of the real sheaf \(\cE_0\) has real
coefficients after \(\iota\), because every local factor has real
coefficients.  The trace formula writes it as
\[
  L(U_{y_0},\cE_0,T)
  =
  \frac{
    \det(1-TF_y\mid H_c^1)
  }{
    \det(1-TF_y\mid H_c^2)
  }.
\tag{12.2}
\label{eq:relative-real-fiber-L}
\]
Roots of the numerator have reciprocal absolute value at least
\(q_y^{-(w+1)/2}\), whereas every root of the denominator has reciprocal
absolute value exactly \(q_y^{-(w+2)/2}\).  The two polynomials are
coprime.  Since their quotient is fixed by complex conjugation and both
have constant term \(1\), uniqueness of a reduced fraction shows that
each polynomial has real coefficients.  This remains true after every
finite extension of \(k(y_0)\), so the local characteristic polynomials
of \(R^1h_!\cE\) and \(R^2h_!\cE\) are \(\iota\)-real.

After stratifying \(Y_0\), the higher direct images are lisse on each
stratum.  Apply \cref{thm:real-sheaf-mixed} on each stratum and use
open--closed induction to conclude that they are mixed.  Their punctual
upper bounds were already proved from \cref{eq:relative-curve-base-change},
so their mixed weights are at most \(w+i\).  Direct summands inherit
mixedness and the bound, giving the assertion for \(\cF_0\).

For a general mixed \(\cF_0\), choose a finite weight filtration.  Apply
\(R^ih_!\) to each short exact sequence.  The resulting long exact
sequences, together with induction on the filtration length and
subquotient/extension permanence, prove the claim.
\end{proof}

\begin{warning}
The coprimality in \eqref{eq:relative-real-fiber-L} uses the strict gap
between \(w+1\) and \(w+2\).  Merely knowing that the whole alternating
product is real would not imply that each cohomological factor is real.
This is the only no-cancellation argument needed to pass from punctual
bounds to mixed sheaves in the reduction.
\end{warning}

\section{Completion of the direct-image theorem}

\begin{proof}[Proof of \cref{thm:deligne-laumon-direct-image}]
For a sheaf, the formal reductions of the preceding sections reduce the
claim to the relative affine-curve case, which is
\cref{lem:relative-curve-mixedness}.  Open--closed induction restores the
discarded strata, composition restores the successive affine-line
projections, and the quasi-finite case restores the finite normalization
map.  This proves the sheaf statement.

Let now \(\cK_0\) be a bounded constructible complex of weights at most
\(w\).  The hypercohomology spectral sequence
\[
  E_2^{a,b}
  =
  R^af_{0!}\cH^b(\cK_0)
  \Longrightarrow
  \cH^{a+b}(Rf_{0!}\cK_0)
\tag{12.3}
\label{eq:weight-hypercohomology-ss}
\]
is Frobenius-equivariant and finite.  The sheaf
\(\cH^b(\cK_0)\) has weights at most \(w+b\), so the result just proved
puts weights at most \(w+a+b\) on \(E_2^{a,b}\).  The abutment in total
degree \(n=a+b\) has a finite filtration by subquotients of those terms
and therefore has weights at most \(w+n\).  This is precisely the derived
weight convention from \cref{def:iota-purity}.
\end{proof}

\section{Smooth proper purity}

Upper bounds become exact under smoothness and properness because
relative Poincare duality supplies the opposite bound.

\begin{corollary}[Smooth-proper purity]
\label{cor:smooth-proper-purity}
Let \(f_0:X_0\to Y_0\) be smooth and proper of pure relative dimension
\(d\), and let \(\cF_0\) be a lisse Weil sheaf punctually
\(\iota\)-pure of weight \(w\).  Then \(R^if_{0*}\cF_0\) is lisse and
punctually \(\iota\)-pure of weight \(w+i\) for every \(i\).
\end{corollary}

\begin{proof}
Properness identifies \(Rf_!=Rf_*\), and smooth proper base change makes
each \(R^if_*\cF\) lisse.  The direct-image theorem gives weights at most
\(w+i\).

Relative Poincare duality gives a Frobenius-equivariant isomorphism
\[
  (R^if_*\cF)^\vee
  \simeq
  R^{2d-i}f_*(\cF^\vee)(d).
\tag{12.4}
\label{eq:relative-pd-weight}
\]
The dual sheaf \(\cF^\vee\) is pure of weight \(-w\).  Applying the upper
bound in degree \(2d-i\) and then twisting by \((d)\) gives weights at
most
\[
  -w+(2d-i)-2d=-w-i
\]
on the right side of \eqref{eq:relative-pd-weight}.  Dualizing reverses
weights, so \(R^if_*\cF\) has weights at least \(w+i\).  Combining upper
and lower bounds proves purity.
\end{proof}

\begin{corollary}[Cohomological Riemann hypothesis]
\label{cor:cohomological-rh}
Let \(X_0\) be smooth and proper over \(k_0\).  For every \(i\), every
eigenvalue \(\alpha\) of geometric Frobenius on
\[
  H^i(X,\Qbarl)
\]
satisfies
\[
  |\iota(\alpha)|=q^{i/2}.
\]
The assertion holds for every isomorphism
\(\iota:\Qbarl\to\mathbf C\).
\end{corollary}

\begin{proof}
Apply \cref{cor:smooth-proper-purity} to
\(X_0\to\Spec k_0\) and the constant sheaf, which is pure of weight
\(0\).  The entire argument from \cref{chap:weights} onward was valid
for an arbitrary choice of \(\iota\), so the conclusion may be repeated
for every such choice.
\end{proof}

\begin{remark}
At this stage the Frobenius eigenvalues are elements of \(\Qbarl\) with
the asserted absolute value under every field isomorphism to
\(\mathbf C\).  The next part of the notes combines this with the
rational zeta function to prove that they are algebraic over
\(\mathbf Q\), that every complex conjugate has the same modulus, and
that the individual polynomials \(P_i(T)\) lie in \(\mathbf Z[T]\) and
are independent of \(\ell\).  Those arithmetic conclusions have not
been used here.
\end{remark}

\section{Exact dependency graph}

\[
\begin{tikzcd}[column sep=small,row sep=small]
\substack{\text{six operations, base change,}\\\text{trace, and duality}}
  \arrow[r] & \substack{\text{rank-one and}\\\text{determinant weights}}
  \arrow[r] & \substack{\text{real-sheaf}\\\text{purity}}
  \arrow[d] \\
\substack{\text{Artin--Schreier}\\\text{orthogonality}}
  \arrow[r] & \substack{\text{Fourier inversion}\\\text{and Plancherel}}
  \arrow[r] & \substack{\text{\(L^2\)-weight}\\\text{criterion}}
  \arrow[d] \\
& \text{curve estimate}
  \arrow[r] & \substack{\text{relative affine-curve}\\\text{case}}
  \arrow[d] \\
& \substack{\text{Noether normalization,}\\\text{excision, and Leray}}
  \arrow[r] & \substack{Rf_!\text{ upper}\\\text{weights}}
  \arrow[d] \\
& \substack{\text{relative Poincare}\\\text{duality}}
  \arrow[r] & \substack{\text{smooth-proper}\\\text{purity.}}
\end{tikzcd}
\]

There is no arrow from smooth-proper purity, integral Frobenius
polynomials, hard Lefschetz, a curve Riemann hypothesis, Chebotarev, or
Picard--Lefschetz back into the graph.

\begin{formalization}
\label{form:global-weight-interface}
\textbf{Terminal interfaces.}

\begin{enumerate}
\item A sheaf-level theorem:
\[
\operatorname{WeightsLE}_\iota(\cF,w)
\Longrightarrow
\operatorname{WeightsLE}_\iota(R^if_!\cF,w+i).
\]
\item A derived theorem using the convention
\(\operatorname{WeightsLE}(\cH^j\cK,w+j)\).
\item Separate reduction lemmas for quasi-finite maps, composition,
open--closed triangles, relative Noether normalization, and affine-line
projections.
\item A Frobenius-equivariant proper-support base-change isomorphism at
every closed point.
\item A mixedness lemma based on real envelopes and reduced real fiber
\(L\)-functions.
\item A spectral-sequence interface requiring only a finite filtration
and subquotient maps, never degeneration.
\item The terminal theorem
\cref{thm:deligne-laumon-direct-image}.
\item Relative duality in the exact form
\eqref{eq:relative-pd-weight}, and the terminal purity corollary
\cref{cor:smooth-proper-purity}.
\end{enumerate}

The architecture follows the reductions in
\cite{DeligneWeilII,LaumonFourier,KiehlWeissauer,StanfordL24}.
\end{formalization}




\part{Proper purity and the Weil conjectures}
\chapter{Proper Purity and the Decomposition Theorem}
\label{chap:decomposition}

The arithmetic estimate is now complete:
\cref{thm:deligne-laumon-direct-image} proves that \(Rf_!\) preserves upper
weights.  We may therefore turn the formal perverse machinery into purity,
intersection-cohomology statements, and decomposition.  This order is the
logical heart of the notes.

\section{The direct-image theorem in perverse form}

\begin{corollary}[Perverse upper direct image]
\label{cor:perverse-upper-direct-image}
Let \(f_0:X_0\to Y_0\) be separated of finite type.  If
\(\cP_0\in\Perv(X_0)\) has weights at most \(w\), then
\[
  \pH^a(Rf_{0!}\cP_0)
  \quad\text{has weights at most }w+a
\tag{16.1}\label{eq:perverse-rfshriek-upper}
\]
for every \(a\).
\end{corollary}

\begin{proof}
The derived statement of
\cref{thm:deligne-laumon-direct-image} says
\(Rf_{0!}\cP_0\in D_{\leq w}(Y_0)\).  Apply the perverse weight criterion
\cref{thm:perverse-weight-criterion}.
\end{proof}

\begin{corollary}[Dual lower direct image]
\label{cor:perverse-lower-direct-image}
If \(\cK_0\in D_{\geq w}(X_0)\), then
\[
  Rf_{0*}\cK_0\in D_{\geq w}(Y_0).
\tag{16.2}\label{eq:rfstar-lower}
\]
Equivalently, if a perverse sheaf \(\cP_0\) has weights at least \(w\), then
\[
  \pH^a(Rf_{0*}\cP_0)
  \quad\text{has weights at least }w+a.
\tag{16.3}\label{eq:perverse-rfstar-lower}
\]
\end{corollary}

\begin{proof}
By definition \(\D_{X_0}\cK_0\in D_{\leq-w}\).  Apply the upper theorem and
the Verdier-duality isomorphism
\[
  \D_{Y_0}Rf_{0*}\cK_0
  \simeq Rf_{0!}\D_{X_0}\cK_0.
\]
The right side has weights at most \(-w\); dualizing is exactly
\eqref{eq:rfstar-lower}.  Formula
\eqref{eq:perverse-rfstar-lower} follows from
\cref{thm:perverse-weight-criterion}.
\end{proof}

\begin{theorem}[Proper direct-image purity]
\label{thm:proper-direct-image-purity}
Let \(f_0:X_0\to Y_0\) be proper and let \(\cK_0\) be pure of weight \(w\).
Then \(Rf_{0*}\cK_0\) is pure of weight \(w\).  In perverse degrees,
\[
  \pH^a(Rf_{0*}\cK_0)
  \quad\text{is pure of weight }w+a.
\tag{16.4}\label{eq:proper-perverse-purity}
\]
\end{theorem}

\begin{proof}
Properness gives \(Rf_{0!}=Rf_{0*}\).  The upper direct-image theorem gives
weights at most \(w\); \cref{cor:perverse-lower-direct-image} gives weights at
least \(w\).  Their intersection is purity.  Apply
\eqref{eq:pH-pure-shift} for the last assertion.
\end{proof}

\begin{perversebox}
Upper weights come from the Fourier/curve argument.  Lower weights come from
Verdier duality.  Properness is the point at which the two functors become the
same and the two inequalities meet.  Merely knowing the upper theorem never
implies purity for a nonproper map.
\end{perversebox}

\section{Purity of intermediate extensions}

\begin{theorem}[Middle extension preserves purity]
\label{thm:jmiddle-purity}
Let \(j_0:U_0\hookrightarrow X_0\) be an open immersion and let
\(\cP_0\in\Perv(U_0)\) be pure of weight \(w\).  Then
\[
  j_{0!*}\cP_0
  \quad\text{is pure of weight }w.
\tag{16.5}\label{eq:jmiddle-pure}
\]
\end{theorem}

\begin{proof}
The upper direct-image theorem applied to \(j_0\) shows that
\(j_{0!}\cP_0\) has weights at most \(w\).  Hence
\({}^pj_{0!}\cP_0=\pH^0(j_{0!}\cP_0)\) has weights at most \(w\).
Duality and the same theorem show that
\({}^pj_{0*}\cP_0\) has weights at least \(w\).

By definition \(j_{0!*}\cP_0\) is a quotient of
\({}^pj_{0!}\cP_0\) and a subobject of
\({}^pj_{0*}\cP_0\).  Upper bounds pass to quotients and lower bounds pass to
subobjects.  Thus it has both bounds and is pure.
\end{proof}

\begin{corollary}[Purity of intersection complexes]
\label{cor:IC-purity}
Let \(X_0\) be irreducible, let \(j_0:U_0\hookrightarrow X_0\) be a smooth
dense open of dimension \(d\), and let \(\cL_0\) be lisse and punctually pure
of weight \(u\).  Then
\[
  \IC_{X_0}(\cL_0)=j_{0!*}(\cL_0[d])
  \quad\text{is pure of weight }u+d.
\tag{16.6}\label{eq:IC-pure-weight}
\]
\end{corollary}

\begin{proof}
Use \cref{prop:lisse-perverse-weight-shift,thm:jmiddle-purity}.
\end{proof}

\section{Pure intersection cohomology}

\begin{corollary}[Proper intersection-cohomology purity]
\label{cor:proper-IH-purity}
Let \(X_0\) be proper of pure dimension \(d\), and let \(\cL_0\) be as in
\cref{cor:IC-purity}.  Then
\[
  IH^r(X,\cL)
  :=\mathbb H^r(X,\IC_X(\cL))
\]
is pure of weight
\[
  u+d+r.
\tag{16.7}\label{eq:IH-pure}
\]
\end{corollary}

\begin{proof}
Apply \cref{thm:proper-direct-image-purity} to the structural morphism
\(a_0:X_0\to\Spec\Fq\) and the pure complex
\(\IC_{X_0}(\cL_0)\) of weight \(u+d\).  On a point, perverse and ordinary
cohomology agree; \eqref{eq:point-weight-shift} gives
\eqref{eq:IH-pure}.
\end{proof}

\begin{corollary}[Smooth-proper cohomological Riemann hypothesis]
\label{cor:perverse-smooth-proper-purity}
If \(X_0/\Fq\) is smooth and proper of pure dimension \(d\), then
\[
  H^i(X,\Qbarl)\quad\text{is pure of weight }i
\tag{16.8}\label{eq:perverse-smooth-purity}
\]
for every \(i\).
\end{corollary}

\begin{proof}
Smoothness gives \(\IC_{X_0}=\Qbarl[d]\), pure of weight \(d\).  Apply
\cref{thm:proper-direct-image-purity} to \(a_0:X_0\to\Spec\Fq\):
\[
  Ra_{0*}\Qbarl[d]\quad\text{is pure of weight }d.
\]
Its degree-\(r\) cohomology is
\[
  H^r(Ra_{0*}\Qbarl[d])=H^{r+d}(X,\Qbarl)
\]
and has weight \(d+r\).  Put \(i=d+r\).
\end{proof}

\section{Geometric semisimplicity}

\begin{theorem}[BBD geometric semisimplicity]
\label{thm:pure-perverse-geometric-semisimple}
Let \(\cP_0\) be a pure perverse Weil sheaf on \(X_0\).  After base change to
\(\Fbar\) and forgetting the Frobenius structure, \(\cP\) is a finite direct
sum of simple perverse sheaves:
\[
  \cP\simeq
  \bigoplus_\lambda
  \IC_{\overline{S_\lambda}}(\cL_\lambda)^{\oplus m_\lambda}.
\tag{16.9}\label{eq:pure-perverse-semisimple}
\]
\end{theorem}

\begin{proof}
Finite length is \cref{thm:perverse-finite-length}.  Let \(\cP'\subset\cP\)
be a maximal semisimple geometric subobject and put
\(\cP''=\cP/\cP'\).  After replacing the ground field by a finite extension,
the object, subobject, quotient, and exact sequence all descend:
\[
 0\longrightarrow\cP'_0\longrightarrow\cP_0
 \longrightarrow\cP''_0\longrightarrow0.
\]
Because \(\cP_0\) is pure of weight \(w\), the quotient \(\cP''_0\) has
weights \(\leq w\), while the subobject \(\cP'_0\) has weights \(\geq w\).
The sequence defines
\[
 e_0\in\Ext^1_{X_0}(\cP''_0,\cP'_0).
\]
BBD 5.1.15(iii) says that, for a first argument of weights \(\leq w\) and a
second argument of weights \(\geq w\), every positive-degree arithmetic
extension maps to zero after geometric base change.  Hence the image of
\(e_0\) in \(\Ext^1_X(\cP'',\cP')\) vanishes and the geometric sequence
splits.  If \(\cP''\neq0\), a simple subobject of the finite-length object
\(\cP''\) would enlarge \(\cP'\), contradicting maximality.  Thus
\(\cP''=0\).  The simple-object classification
\cref{thm:simple-perverse-classification} now gives
\eqref{eq:pure-perverse-semisimple}.  This is BBD 5.3.8.
\end{proof}

\begin{warning}[Geometric is not Frobenius semisimplicity]
\label{warn:frobenius-not-semisimple}
The theorem forgets the Weil structure.  It does not say that Frobenius acts
semisimply.  On a point, take a two-dimensional representation on which
Frobenius is \(\alpha\begin{psmallmatrix}1&1\\0&1\end{psmallmatrix}\), with
\(\alpha\) pure.  All eigenvalues have the same weight, but the Frobenius
operator has a nontrivial Jordan block.
\end{warning}

\section{Splitting a pure complex}

\begin{theorem}[BBD splitting of pure complexes]
\label{thm:pure-complex-splitting}
If \(\cK_0\in\Dbc(X_0,\Qbarl)\) is pure, then after base change to \(\Fbar\)
there exists an isomorphism
\[
  \cK\simeq\bigoplus_a\pH^a(\cK)[-a].
\tag{16.10}\label{eq:pure-complex-splitting}
\]
The objects \(\pH^a(\cK)\) are canonical; the splitting generally is not.
\end{theorem}

\begin{proof}
Assume \(\cK_0\) is pure of weight \(w\), and let \(a\) be its largest
nonzero perverse degree.  Put
\[
 A_0=\pH^a(\cK_0)[-a],
 \qquad
 B_0=\ptau_{\leq a-1}\cK_0.
\]
The perverse weight criterion makes \(A_0\) pure of weight \(w\).  It makes
\(B_0\) pure of weight \(w\) as well, since its perverse cohomology objects
are precisely \(\pH^r(\cK_0)\), pure of weight \(w+r\), for \(r<a\).
The truncation triangle is
\[
 B_0\longrightarrow\cK_0\longrightarrow A_0
 \xrightarrow{\partial_0}B_0[1].
\tag{16.11}\label{eq:perverse-truncation-obstruction}
\]
Its connecting morphism is an arithmetic extension class
\[
 [\partial_0]\in\Ext^1_{X_0}(A_0,B_0).
\]
Apply BBD 5.1.15(iii) with \(A_0\) of weights \(\leq w\) and \(B_0\) of
weights \(\geq w\): the image of this class in
\(\Ext^1_X(A,B)\) after geometric base change is zero.  Therefore the
geometric truncation triangle splits.  Induction on the finite perverse
amplitude of \(\cK_0\) gives \eqref{eq:pure-complex-splitting}.  This is
BBD 5.4.5.
\end{proof}

\begin{theorem}[Decomposition theorem over a finite field]
\label{thm:decomposition-finite-field}
Let \(f_0:X_0\to Y_0\) be proper and let \(\cK_0\) be pure.  After base change
to \(\Fbar\),
\[
  Rf_*\cK\simeq
  \bigoplus_a \pH^a(Rf_*\cK)[-a],
\tag{16.12}\label{eq:decomposition-theorem}
\]
and every perverse cohomology object on the right is semisimple.
In particular, for \(X_0\) pure dimensional and \(\cL_0\) pure on a smooth
dense open, this applies to \(\cK_0=\IC_{X_0}(\cL_0)\).
\end{theorem}

\begin{proof}
Proper direct-image purity is
\cref{thm:proper-direct-image-purity}.  Apply
\cref{thm:pure-complex-splitting} to \(Rf_{0*}\cK_0\), and
\cref{thm:pure-perverse-geometric-semisimple} to each perverse cohomology
object.
\end{proof}

\begin{remark}[Relative hard Lefschetz]
For a projective morphism and an ample class \(\eta\), the stronger theorem
gives
\[
  \eta^r:\pH^{-r}(Rf_*\IC_X)
  \xrightarrow{\sim}\pH^r(Rf_*\IC_X)(r).
\]
Together with primitive decomposition it makes the splitting more structured.
It is not needed for the classical Weil conjectures and is not used anywhere
upstream.  Keeping it outside the dependency chain avoids concealing an
additional deep theorem behind the word “decomposition.”
\end{remark}

\begin{formalization}
The theorem interfaces should be layered:
\[
\begin{array}{c}
 Rf_!\text{ upper weights}\\
 \Downarrow\ \text{duality}\\
 Rf_*\text{ lower weights}\\
 \Downarrow\ \text{properness}\\
 Rf_*\text{ purity}\\
 \Downarrow\\
 \text{IC/IH purity, geometric semisimplicity, decomposition.}
\end{array}
\]
The decomposition is an existence of an isomorphism after base change, not a
chosen definitional equality and not a Frobenius-equivariant semisimplicity
record.  The smooth test must reduce
\(Ra_*\Qbarl[d]\) to \(H^{r+d}\) with weight \(d+r\).
\end{formalization}

\chapter{The Weil Conjectures from the Perverse Package}
\label{chap:weil-assembly}

Let \(X_0\) be smooth, projective, and geometrically connected of dimension
\(d\) over \(k_0=\Fq\), and put
\[
  X=X_0\times_{\Fq}\Fbar,\qquad
  a_0:X_0\longrightarrow\Spec\Fq.
\]
Fix an auxiliary prime \(\ell\ne\operatorname{char}\Fq\), and let \(F\) be
geometric Frobenius.

\section{The one-line perverse proof of purity}

\begin{theorem}[Terminal perverse calculation]
\label{thm:terminal-perverse-calculation}
Every eigenvalue \(\alpha\) of \(F\) on
\[
  V_{i,\ell}=H^i_{\et}(X,\Ql)
\]
has \(\iota\)-absolute value \(q^{i/2}\), for every field isomorphism
\(\iota:\Qbarl\simeq\mathbf C\).
\end{theorem}

\begin{proof}
The constant sheaf is punctually pure of weight \(0\).  Since \(X_0\) is
smooth of dimension \(d\),
\[
  \IC_{X_0}=\Qbarl[d]
\]
is perverse and pure of weight \(d\), by
\cref{eq:IC-smooth,eq:constant-perverse-weight}.  The structural morphism is
proper, so \cref{thm:proper-direct-image-purity} makes
\[
  Ra_{0*}\Qbarl[d]
\]
pure of weight \(d\).  On the point, perverse and ordinary cohomology agree:
\[
  \pH^r(Ra_{0*}\Qbarl[d])
  =H^r(Ra_{0*}\Qbarl[d])
  =H^{r+d}(X,\Qbarl).
\]
The left side is pure of weight \(d+r\).  Taking \(r=i-d\) gives weight
\(i\), which means \(|\iota(\alpha)|=q^{i/2}\).

The argument in the weight chapters was carried out for an arbitrary
\(\iota\); repeating it for every \(\iota\) gives the universal-realization
assertion.  No rationality or algebraicity of the individual factor has yet
been assumed.
\end{proof}

\begin{perversebox}
The decomposition theorem is compatible with this calculation but not needed
for it.  Proper purity is the essential input.  Logically, decomposition was
proved only after proper purity in \cref{chap:decomposition}.
\end{perversebox}

\section{Trace formula and rationality}

Finiteness gives \(V_{i,\ell}=0\) outside \(0\leq i\leq2d\), and all these
spaces are finite dimensional.  Properness identifies ordinary and
compactly supported cohomology.  The trace formula applied to \(F^n\) gives
\[
  \#X_0(\mathbf F_{q^n})
  =
  \sum_{i=0}^{2d}(-1)^i
  \Tr(F^n\mid V_{i,\ell}).
\tag{17.1}\label{eq:trace-X}
\]
Put
\[
  P_{i,\ell}(T)=\det(1-TF\mid V_{i,\ell})
  \in1+T\Ql[T].
\tag{17.2}\label{eq:Pi-definition}
\]

\begin{theorem}[Cohomological rationality]
\label{thm:cohomological-rationality}
For every \(\ell\ne\operatorname{char}\Fq\),
\[
  Z(X_0,T)
  =
  \prod_{i=0}^{2d}P_{i,\ell}(T)^{(-1)^{i+1}}
  \quad\text{in }\Ql[[T]].
\tag{17.3}\label{eq:cohomological-zeta}
\]
In particular \(Z(X_0,T)\in\mathbf Q(T)\).
\end{theorem}

\begin{proof}
Insert \eqref{eq:trace-X} into
\[
 Z(X_0,T)=
 \exp\!\left(\sum_{n\geq1}
 \#X_0(\mathbf F_{q^n})\frac{T^n}{n}\right).
\]
The alternating determinant--trace identity
\cref{cor:alternating-det-trace} changes the result into
\eqref{eq:cohomological-zeta}.  Thus \(Z\) is rational over \(\Ql\).

The point-counting definition lies in \(\mathbf Q[[T]]\), and the closed-point
Euler product \cref{prop:euler-product} lies in \(\mathbf Z[[T]]\).
Hankel descent \cref{thm:hankel-descent}, applied to
\(\mathbf Q\subset\Ql\), yields \(Z\in\mathbf Q(T)\).
\end{proof}

\begin{warning}
At this point only the alternating product is rational.  It would be false in
general to infer that each \(P_{i,\ell}\) is rational: numerator and
denominator factors can cancel.  Universal purity will exclude such
cancellation by separating their absolute values.
\end{warning}

\section{Endpoint factors}

\begin{proposition}[Degree zero and top degree]
\label{prop:endpoint-factors}
One has
\[
  P_{0,\ell}(T)=1-T,\qquad
  P_{2d,\ell}(T)=1-q^dT.
\tag{17.4}\label{eq:endpoint-factors}
\]
\end{proposition}

\begin{proof}
Geometric connectedness gives
\(H^0(X,\Ql)=\Ql\), with Frobenius acting trivially.  Poincare duality
identifies \(H^{2d}(X,\Ql)\) with \(\Ql(-d)\), on which geometric Frobenius
acts by \(q^d\).  These are exactly the two determinants.
\end{proof}

\section{Duality and reciprocal eigenvalues}

Poincare duality gives perfect \(F\)-equivariant pairings
\[
  V_{i,\ell}\times V_{2d-i,\ell}
  \longrightarrow\Ql(-d)
\tag{17.5}\label{eq:poincare-X}
\]
and therefore
\[
  \langle Fx,Fy\rangle=q^d\langle x,y\rangle.
\tag{17.6}\label{eq:poincare-similitude}
\]
If \(\alpha_1,\ldots,\alpha_{b_i}\) are the eigenvalues on \(V_{i,\ell}\),
then those on \(V_{2d-i,\ell}\) are
\[
  \frac{q^d}{\alpha_1},\ldots,\frac{q^d}{\alpha_{b_i}}.
\tag{17.7}\label{eq:reciprocal-eigenvalues}
\]
In particular \(b_i=b_{2d-i}\).

\begin{proposition}[Purity in every complex realization]
\label{prop:all-realization-purity}
Every eigenvalue \(\alpha\in\Qbarl^\times\) of geometric Frobenius on
\(V_{i,\ell}\) satisfies
\[
  |\iota(\alpha)|=q^{i/2}
\tag{17.8}\label{eq:smooth-proper-purity}
\]
for every isomorphism \(\iota:\Qbarl\simeq\mathbf C\).
\end{proposition}

\begin{proof}
This is \cref{thm:terminal-perverse-calculation}.  The word “eigenvalue” here
means a root of the characteristic polynomial of \(F\); the roots of
\(P_{i,\ell}(T)\) as a polynomial in \(T\) are their reciprocals.
\end{proof}

\section{The functional equation}

Let
\[
  \chi_\ell=\sum_{i=0}^{2d}(-1)^i\dim_{\Ql}V_{i,\ell}.
\]

\begin{theorem}[Functional equation]
\label{thm:zeta-functional-equation}
There exists \(\varepsilon_\ell\in\{1,-1\}\) such that
\[
  Z(X_0,q^{-d}T^{-1})
  =
  \varepsilon_\ell q^{d\chi_\ell/2}
  T^{\chi_\ell}Z(X_0,T).
\tag{17.9}\label{eq:zeta-functional-equation}
\]
The integer \(\chi_\ell\) and the sign \(\varepsilon_\ell\) are independent
of \(\ell\).
\end{theorem}

\begin{proof}
Apply the abstract determinant computation
\cref{thm:abstract-functional-equation} to the graded spaces
\(V_{i,\ell}\), the operator \(F\), and the perfect pairings
\eqref{eq:poincare-X}, whose multiplier is \(q^d\).  Its alternating
determinant is \(Z(X_0,T)\) by
\cref{thm:cohomological-rationality}.  This proves
\eqref{eq:zeta-functional-equation}; the middle-degree sign is controlled by
the \((-1)^d\)-symmetry of cup product.

If two equations of the displayed shape hold for the same nonzero rational
function \(Z\), their quotient makes
\(T^{\chi_\ell-\chi_{\ell'}}\) constant.  Hence the exponents agree, and then
the constants and signs agree.  Thus both are determined by the
point-counting zeta function and are independent of \(\ell\).
\end{proof}

\section{What remains arithmetic}

\begin{corollary}[Weil theorem before separation of factors]
\label{cor:weil-before-integrality}
The zeta function is rational over \(\mathbf Q\), has the cohomological
factorization \eqref{eq:cohomological-zeta}, has endpoint factors
\eqref{eq:endpoint-factors}, satisfies the functional equation
\eqref{eq:zeta-functional-equation}, and its degree-\(i\) Frobenius
eigenvalues have universal weight \(i\).
\end{corollary}

Three linked assertions remain: every eigenvalue must be shown algebraic over
\(\mathbf Q\), each \(P_{i,\ell}\) must lie in \(\mathbf Z[T]\), and the
polynomial must be independent of \(\ell\).  The next chapter derives all
three from the already proved statements in the safe order
\[
\begin{gathered}
  Z\in\mathbf Z[[T]]\cap\mathbf Q(T)
  \Longrightarrow\text{integral reduced factors}\\
  \Longrightarrow\text{weight separation}\\
  \Longrightarrow P_i\in\mathbf Z[T]\\
  \text{and \(P_i\) is \(\ell\)-independent}.
\end{gathered}
\]

\begin{formalization}
The final assembly has four independent interfaces:
\begin{enumerate}
\item \texttt{perverseProperPurityToCohomology}, implementing
      \[
      \Qbarl[d]\mapsto Ra_*\Qbarl[d]\mapsto
      H^{i-d}(Ra_*\Qbarl[d])=H^i(X);
      \]
\item \texttt{traceFormulaToZeta}, using the determinant--trace identity;
\item \texttt{dualityToFunctionalEquation}, using the similitude pairing;
\item \texttt{universalPurityToSeparatedFactors}, developed in the next
      chapter.
\end{enumerate}
The release theorem must import the proved Deligne--Laumon theorem and proper
purity, not an axiom named “Weil purity.”
\end{formalization}

\chapter{Weight separation, integrality, and independence of \texorpdfstring{\(\ell\)}{ell}}
\label{chap:integrality}

The zeta function has integral Taylor coefficients for the elementary
reason that it is an Euler product over closed points.  Its
cohomological factors were constructed over \(\Ql\), and their roots
initially live only in \(\Qbarl\).  This chapter proves, in the required
noncircular order, that those roots are algebraic over \(\mathbf Q\), that
the individual factors lie in \(\mathbf Z[T]\), and that they are
independent of \(\ell\):
\[
\begin{split}
 &\text{purity for every }\iota:\Qbarl\simeq\mathbf C\\
 &\quad\Longrightarrow\text{no cancellation between distinct weights}\\
 &\quad\Longrightarrow\text{roots occur in the reduced integral
 numerator or denominator of }Z\\
 &\quad\Longrightarrow\text{algebraicity over }\mathbf Q
 \text{ and Galois-stable weight pieces}\\
 &\quad\Longrightarrow\mathbf Q\text{-rational factors}
 \Longrightarrow\mathbf Z\text{-integral factors and }\ell\text{-independence}.
\end{split}
\tag{14.1}\label{eq:arithmetic-order}
\]
In particular, algebraicity over \(\mathbf Q\) is a conclusion in the
middle of the argument, not part of the initial definition of weight.

\section{Universal realization-purity}

\begin{definition}
Let \(\alpha\in\Qbarl^\times\), let \(q>1\), and let \(w\in\mathbf Z\).
We say that \(\alpha\) is \emph{universally \(\iota\)-pure of weight
\(w\)} if
\[
  |\iota(\alpha)|=q^{w/2}
\]
for every abstract field isomorphism
\(\iota\colon\Qbarl\xrightarrow{\sim}\mathbf C\).
\end{definition}

The two algebraically closed fields in this definition have
characteristic \(0\), cardinality \(2^{\aleph_0}\), and transcendence
degree \(2^{\aleph_0}\) over \(\mathbf Q\), so such isomorphisms exist.
No topology is imposed on them; a continuous map from the \(\ell\)-adic
topology to the complex topology would be the wrong notion.

\begin{lemma}[Invariance and disjointness]
\label{lem:universal-weight-invariance}
Let \(\alpha\in\Qbarl^\times\).
\begin{enumerate}
\item If \(\alpha\) is universally \(\iota\)-pure of weight \(w\) and
      \(\sigma\in\operatorname{Aut}(\Qbarl/\mathbf Q)\), then
      \(\sigma(\alpha)\) is universally \(\iota\)-pure of weight \(w\).
\item The same nonzero element cannot be universally
      \(\iota\)-pure of two distinct weights.
\end{enumerate}
\end{lemma}

\begin{proof}
For every \(\iota\colon\Qbarl\simeq\mathbf C\), the composite
\(\iota\circ\sigma\) is another such isomorphism.  Hence
\[
 |\iota(\sigma\alpha)|
 =|(\iota\circ\sigma)(\alpha)|=q^{w/2}.
\]
This proves the first assertion.  If \(\alpha\) had weights \(w\) and
\(w'\), evaluation under any one isomorphism would give
\(q^{w/2}=q^{w'/2}\).  Since \(q>1\), this implies \(w=w'\).
\end{proof}

\begin{lemma}[Extension and fixed fields]
\label{lem:qbarl-automorphism-fixed-field}
The following field-theoretic facts hold.
\begin{enumerate}
\item Every \(\mathbf Q\)-embedding of a finite algebraic subextension
      \(E\subset\Qbarl\) into \(\mathbf C\) extends to an isomorphism
      \(\Qbarl\xrightarrow{\sim}\mathbf C\).
\item Every \(\mathbf Q\)-embedding between finite algebraic
      subextensions of \(\Qbarl\) extends to an element of
      \(\operatorname{Aut}(\Qbarl/\mathbf Q)\).
\item The fixed field of
      \(\operatorname{Aut}(\Qbarl/\mathbf Q)\) is \(\mathbf Q\).
\end{enumerate}
\end{lemma}

\begin{proof}
For (1), first extend the given embedding over the relative algebraic
closure of \(\mathbf Q\) in \(\Qbarl\).  Choose transcendence bases on
the source and target complements.  Both bases have cardinality
\(2^{\aleph_0}\), so a bijection between them, together with the
algebraic embedding already chosen, gives an isomorphism of the purely
transcendental subfields.  Since source and target are algebraically
closed, the isomorphism extends across their algebraic closures.  The
same argument with target \(\Qbarl\) proves (2).

For (3), let \(x\in\Qbarl\setminus\mathbf Q\).  If \(x\) is algebraic
over \(\mathbf Q\), then either it has a distinct conjugate, in which case
(2) moves it, or its minimal polynomial is linear, which would put it in
\(\mathbf Q\).  If \(x\) is transcendental, include it in a
transcendence basis and permute that basis; the permutation extends to an
automorphism of the algebraic closure and moves \(x\).  Thus no element
outside \(\mathbf Q\) is fixed by every automorphism.
\end{proof}

\section{No cancellation before algebraicity}

Fix an auxiliary prime \(\ell\), and set
\[
  O_\ell(T)=\prod_{\substack{0\leq i\leq2d\\ i\ {\rm odd}}}
       P_{i,\ell}(T),\qquad
  E_\ell(T)=\prod_{\substack{0\leq i\leq2d\\ i\ {\rm even}}}
       P_{i,\ell}(T).
\]
These polynomials lie in \(1+T\Ql[T]\), and
\[
  Z(X_0,T)=\frac{O_\ell(T)}{E_\ell(T)}.
  \tag{14.2}\label{eq:odd-even-zeta-final}
\]

\begin{lemma}[Cohomological numerator and denominator are coprime]
\label{lem:odd-even-coprime-before-Q}
The polynomials \(O_\ell\) and \(E_\ell\) are coprime in
\(\Qbarl[T]\).  More generally, \(P_{i,\ell}\) and \(P_{j,\ell}\) have
no common root when \(i\neq j\).
\end{lemma}

\begin{proof}
Suppose \(P_{i,\ell}\) and \(P_{j,\ell}\) shared the ordinary root
\(\rho\).  Frobenius is invertible, so \(\rho\neq0\), and
\(\alpha=\rho^{-1}\) would be a reciprocal root of both polynomials.
By \cref{prop:all-realization-purity}, it would be universally
\(\iota\)-pure of weights \(i\) and \(j\).  The disjointness part of
\cref{lem:universal-weight-invariance} forces \(i=j\).  In particular an
odd factor and an even factor cannot meet, proving the first assertion.
\end{proof}

Now use integrality of the \emph{whole} zeta series.  By
\cref{prop:euler-product,thm:cohomological-rationality},
\[
 Z(X_0,T)\in1+T\mathbf Z[[T]]\cap\mathbf Q(T).
\]
The reduced integral Fatou presentation
\cref{cor:reduced-integral-presentation} gives coprime
\(A,B\in\mathbf Z[T]\) with
\[
  Z(X_0,T)=A(T)/B(T),\qquad A(0)=B(0)=1.
  \tag{14.3}\label{eq:reduced-integral-zeta}
\]

\begin{proposition}[Parity products are the integral reduced factors]
\label{prop:parity-products-integral}
For every \(\ell\),
\[
  O_\ell(T)=A(T),\qquad E_\ell(T)=B(T)
\]
inside \(\Qbarl[T]\).
\end{proposition}

\begin{proof}
The fraction \(O_\ell/E_\ell\) is reduced over the algebraically closed
field \(\Qbarl\) by
\cref{lem:odd-even-coprime-before-Q}.  The fraction \(A/B\) remains
reduced after scalar extension from \(\mathbf Q\), since Bezout
coprimality is preserved.  Reduced presentations over a field are unique
up to a common nonzero scalar.  All four polynomials have constant term
\(1\), so that scalar is \(1\).
\end{proof}

\begin{corollary}[Arithmetic algebraicity of eigenvalues]
\label{cor:eigenvalues-algebraic-over-Q}
Every Frobenius eigenvalue on every \(V_{i,\ell}\) is algebraic over
\(\mathbf Q\).
\end{corollary}

\begin{proof}
Every reciprocal root of every \(P_{i,\ell}\) occurs in the product
\(O_\ell\) or \(E_\ell\).  By
\cref{prop:parity-products-integral}, it is therefore a reciprocal root
of the polynomial \(A\) or \(B\) with rational coefficients.  Its
inverse is algebraic over \(\mathbf Q\).
\end{proof}

This is the precise point where it becomes legitimate to regard the
eigenvalues as algebraic numbers.  Nothing in the proof of no
cancellation used that conclusion.

\section{Recovery of the rational weight factors}

The roots of \(A\) and \(B\), viewed inside \(\Qbarl\), are algebraic over
\(\mathbf Q\).  Partition their reciprocal roots by universal
\(\iota\)-weight.  By
\cref{lem:universal-weight-invariance}, every part is stable under
\(\operatorname{Aut}(\Qbarl/\mathbf Q)\), and distinct parts are
disjoint.  The part of weight \(i\), with multiplicity, is exactly the
reciprocal-root multiset of \(P_{i,\ell}\).

\begin{theorem}[Rationality of individual factors]
\label{thm:rational-individual-factors}
For every \(i\) and every \(\ell\neq\operatorname{char}(\Fq)\),
\[
  P_{i,\ell}(T)\in\mathbf Q[T].
\]
Moreover, if \(\alpha\) is one of its reciprocal roots, then for every
embedding
\(\tau\colon\mathbf Q(\alpha)\hookrightarrow\mathbf C\),
\[
  |\tau(\alpha)|=q^{i/2}.
  \tag{14.4}\label{eq:classical-all-embeddings}
\]
\end{theorem}

\begin{proof}
An automorphism of \(\Qbarl\) over \(\mathbf Q\) permutes the roots of
\(A\) and \(B\), because those polynomials have rational coefficients.
It preserves each universal-weight part by
\cref{lem:universal-weight-invariance}.  Hence it preserves the
reciprocal-root multiset of \(P_{i,\ell}\).  Since
\(P_{i,\ell}(0)=1\), the polynomial itself is fixed coefficientwise by
every such automorphism.  The fixed-field assertion of
\cref{lem:qbarl-automorphism-fixed-field} gives
\(P_{i,\ell}\in\mathbf Q[T]\).

Let \(\tau\) be any complex embedding of the now-known number field
\(\mathbf Q(\alpha)\).  By
\cref{lem:qbarl-automorphism-fixed-field}(1), it extends to an isomorphism
\(\iota\colon\Qbarl\simeq\mathbf C\).  Universal realization-purity then
gives
\[
 |\tau(\alpha)|=|\iota(\alpha)|=q^{i/2}.
\]
Thus the universal \(\ell\)-adic formulation has yielded the classical
all-number-field-embeddings formulation.
\end{proof}

\section{Gauss's lemma and integral factors}

\begin{theorem}[Integrality of the individual factors]
\label{thm:integral-frobenius-factors}
For every \(i\) and every auxiliary prime \(\ell\),
\[
  P_{i,\ell}(T)\in\mathbf Z[T].
\]
Consequently every Frobenius eigenvalue is an algebraic integer.
\end{theorem}

\begin{proof}
According to the parity of \(i\), the rational polynomial
\(P_{i,\ell}\) is a factor of \(A=O_\ell\) or \(B=E_\ell\).  Its
complement is the product of the remaining factors of the same parity.
Both factor and complement have rational coefficients by
\cref{thm:rational-individual-factors} and constant term \(1\).  The
normalized Gauss lemma \cref{lem:normalized-gauss-factor} therefore
places both in \(\mathbf Z[T]\).

Write
\[
 P_{i,\ell}(T)=1+c_1T+\cdots+c_bT^b
              =\prod_{j=1}^b(1-\alpha_jT).
\]
Each \(\alpha_j\) is a root of the monic integral polynomial
\[
 U^b+c_1U^{b-1}+\cdots+c_b.
\]
It is therefore an algebraic integer.
\end{proof}

\begin{remark}
The Fatou lemma alone provides integral \(A\) and \(B\), but does not
identify their cohomological subfactors.  Universal purity first proves
that the cohomological parity fraction is reduced; only then does
\cref{prop:parity-products-integral} identify it with \(A/B\).
Automorphism-stable weight pieces then give rationality, after which
Gauss's lemma applies.  This order contains no hidden use of individual
integrality.
\end{remark}

\section{Independence of the auxiliary prime}

\begin{theorem}[\(\ell\)-independence]
\label{thm:ell-independence}
For each \(i\), the polynomial \(P_{i,\ell}(T)\) is independent of
\(\ell\neq\operatorname{char}(\Fq)\).  There is therefore a unique
\[
  P_i(T)\in\mathbf Z[T]
\]
such that
\[
  P_i(T)=\det(1-TF\mid H^i_{\et}(X,\Ql))
\]
for every auxiliary prime.  In particular
\[
  b_i=\deg P_i=\dim_{\Ql}H^i_{\et}(X,\Ql)
\]
is independent of \(\ell\).
\end{theorem}

\begin{proof}
The reduced integral polynomials \(A\) and \(B\) in
\eqref{eq:reduced-integral-zeta} are determined by the point counts and
are independent of \(\ell\).  By
\cref{prop:parity-products-integral}, their reciprocal-root multisets are
the union of the odd, respectively even, Frobenius eigenvalue multisets
for every auxiliary prime.

Now work with the roots of the fixed polynomials \(A\) and \(B\) in
\(\Qbar\).  By \cref{thm:rational-individual-factors}, each root has a
classical, all-complex-embeddings weight, and
\cref{lem:universal-weight-invariance} shows distinct weights cannot
meet.  Thus the normalized factor whose reciprocal roots have weight
\(i\) is uniquely determined by \(A/B\), \(q\), and \(i\), with no
reference to \(\ell\).  For every \(\ell\), that unique factor is
\(P_{i,\ell}\).  Hence the polynomials and their degrees agree.
\end{proof}

\begin{corollary}[Betti comparison when a lift exists]
\label{cor:betti-comparison-lift}
Suppose \(X_0\) is the special fiber of a smooth proper scheme over a
mixed-characteristic discrete valuation ring, and let
\(X_{\mathbf C}\) be a complex fiber obtained from an embedding of the
fraction field into \(\mathbf C\).  Then
\[
 b_i=\dim_{\mathbf Q}H^i_{\mathrm{sing}}
       (X_{\mathbf C}(\mathbf C),\mathbf Q).
\]
\end{corollary}

\begin{proof}
Smooth proper base change identifies the etale cohomology of geometric
special and generic fibers.  The characteristic-zero comparison theorem
identifies the latter, after tensoring with \(\Ql\), with the singular
cohomology of the complex analytic fiber.  Taking dimensions proves the
claim.
\end{proof}

\section{Completion of the proof}

\begin{theorem}[The classical Weil conjectures]
\label{thm:weil-conjectures-proved}
All assertions of \cref{thm:weil-conjectures-statement} hold.
\end{theorem}

\begin{proof}
Cohomological rationality is
\cref{thm:cohomological-rationality}; the endpoint factors are
\cref{prop:endpoint-factors}; and the functional equation is
\cref{thm:zeta-functional-equation}.  Universal smooth-proper purity,
followed by \cref{cor:eigenvalues-algebraic-over-Q} and
\cref{thm:rational-individual-factors}, gives the all-complex-embeddings
Riemann-hypothesis statement for algebraic numbers.
\Cref{thm:integral-frobenius-factors} adds algebraic integrality, and
\cref{thm:ell-independence} gives common polynomials and Betti numbers.
These are precisely the conclusions in
\cref{thm:weil-conjectures-statement}.
\end{proof}

\begin{formalization}
The formal development should keep three predicates distinct:
\begin{enumerate}
\item `UniversallyRealizationPure alpha w`, quantified over all field
      isomorphisms \(\Qbarl\simeq\mathbf C\);
\item `IsAlgebraic Q alpha`, proved only after
      \cref{prop:parity-products-integral};
\item `IsIntegral Q alpha`, proved only after the normalized Gauss
      factor theorem.
\end{enumerate}
The reusable arithmetic pipeline is:
\[
\begin{array}{c}
\text{universal weights}+\text{rational integral global series}\\
\Downarrow\\
\text{coprime parity products}=\text{reduced }A/B\\
\Downarrow\\
\text{algebraic roots}+\operatorname{Aut}(\Qbarl/\mathbf Q)
  \text{-stable weight strata}\\
\Downarrow\\
\mathbf Q\text{-factors}\to\mathbf Z\text{-factors}
\to\text{realization independence}.
\end{array}
\]
Keeping these as separate theorem boundaries prevents a later
formalization from smuggling arithmetic algebraicity into the definition
of a Weil sheaf.
\end{formalization}


\chapter{Examples and convention tests}
\label{chap:examples}

The examples in this chapter are computational proofs as well as
regression tests.  Each isolates a convention that can otherwise remain
incorrect for many pages: geometric versus arithmetic Frobenius, the
alternating determinant sign, compact supports, the direction of a Tate
twist, geometric connectedness, tensor-product weights, or base-field
Frobenius.  Every answer is computed independently from point counts and
from cohomology.

\section{The rational point}

Let \(X_0=\Spec\Fq\).  There is one point over every finite extension, so
\[
 N_n=1,\qquad
 Z(X_0,T)=\exp\left(\sum_{n\geq1}\frac{T^n}{n}\right)
          =\frac1{1-T}.
 \tag{15.1}\label{eq:zeta-point}
\]
The closed-point calculation has one degree-one point and gives the same
Euler factor.

Cohomologically,
\[
 H^0(X,\Ql)=\Ql,\qquad H^i(X,\Ql)=0\quad(i>0),
\]
and geometric Frobenius acts trivially on constants.  Hence
\[
 P_0(T)=1-T.
\]
Because \(d=0\), degree zero is also top degree; the endpoint factor is
one polynomial, not two copies of it.

The Euler characteristic is \(1\), and the functional equation becomes
\[
 Z(T^{-1})=-T Z(T).
\]
Indeed,
\[
 \frac1{1-T^{-1}}=-\frac{T}{1-T}.
\]
Thus the sign \(\varepsilon\) in the general theorem can be negative even
for the simplest smooth proper variety.

\section{The affine line and compact supports}

For \(X_0=\A^1_{\Fq}\),
\[
 N_n=q^n,\qquad
 Z(\A^1,T)=\frac1{1-qT}.
 \tag{15.2}\label{eq:zeta-A1}
\]
The scheme is smooth but not proper.  Ordinary cohomology sees only
\(H^0=\Ql\) and would produce \(1/(1-T)\), which is wrong.  Compactly
supported cohomology is
\[
 H_c^i(\A^1_{\Fbar},\Ql)=
 \begin{cases}
   \Ql(-1),&i=2,\\
   0,&i\neq2.
 \end{cases}
\]
Geometric Frobenius acts on \(\Ql(-1)\) by \(q\), and the even
cohomological degree appears in the denominator:
\[
 \det(1-TF\mid H_c^2)^{-1}=(1-qT)^{-1}.
\]

\begin{convention}
The calculation \eqref{eq:zeta-A1} tests three choices simultaneously:
\begin{itemize}
\item the trace formula uses geometric Frobenius;
\item \(F\) acts by \(q\) on \(\Ql(-1)\);
\item the exponent in the zeta product is \((-1)^{i+1}\).
\end{itemize}
Reversing any one choice gives either \(1/(1-q^{-1}T)\) or
\(1-qT\), both incompatible with the elementary point count.
\end{convention}

\section{Projective spaces}

The decomposition of \(\PP^r(\mathbf F_{q^n})\) into affine cells gives
\[
 \#\PP^r(\mathbf F_{q^n})=1+q^n+\cdots+q^{rn}.
\]
Taking the exponential term by term,
\[
 Z(\PP^r,T)=\prod_{j=0}^r(1-q^jT)^{-1}.
 \tag{15.3}\label{eq:zeta-projective-space}
\]
Etale cohomology gives the identical answer:
\[
 H^{2j}(\PP^r_{\Fbar},\Ql)=\Ql(-j)
 \quad(0\leq j\leq r),\qquad
 H^{2j+1}=0.
\]
Thus geometric Frobenius has eigenvalue \(q^j\) in degree \(2j\), of
weight \(2j\), and every nonzero group contributes a denominator factor.

For \(r=1\),
\[
 Z(\PP^1,T)=\frac1{(1-T)(1-qT)}.
\]
Here \(d=1\), \(b_0=b_2=1\), and \(\chi=2\).  A direct substitution gives
\[
 Z(\PP^1,q^{-1}T^{-1})=qT^2Z(\PP^1,T).
\]
This tests both endpoint factors and the positive sign for the
projective line.

For general \(r\), the transformation \(j\mapsto r-j\) pairs the factors
\(1-q^jT\) and \(1-q^{r-j}T\), visibly reproducing Poincare duality and
the multiplier \(q^r\).

\section{Finite etale schemes and component permutations}

Every finite etale \(\Fq\)-scheme has a unique presentation
\[
 Y_0=\coprod_{a=1}^s\Spec\mathbf F_{q^{m_a}}.
\]
There is one closed point of degree \(m_a\) for each component, so
\[
 Z(Y_0,T)=\prod_{a=1}^s(1-T^{m_a})^{-1}.
 \tag{15.4}\label{eq:zeta-finite-etale}
\]
Equivalently,
\[
 \#Y_0(\mathbf F_{q^n})
 =\sum_{a:m_a\mid n}m_a,
\]
because a field embedding
\(\mathbf F_{q^{m_a}}\hookrightarrow\mathbf F_{q^n}\) exists precisely
when \(m_a\mid n\), and then there are \(m_a\) of them.

The geometric fiber is a finite set \(S\).  Its cohomology is concentrated
in degree zero:
\[
 H^0(Y,\Ql)=\Ql^S.
\]
On the points arising from \(\mathbf F_{q^m}\), Frobenius is a cycle of
length \(m\).  The eigenvalues of its permutation matrix are all
\(m\)-th roots of unity, once each.  Therefore
\[
 \det(1-TF\mid\Ql^{m})
 =\prod_{\zeta^m=1}(1-\zeta T)=1-T^m.
\]
Multiplication over the cycles recovers
\eqref{eq:zeta-finite-etale}.

This is the cleanest test of geometric connectedness.
\(\Spec\mathbf F_{q^m}\) is connected as an \(\Fq\)-scheme, but its
geometric fiber consists of \(m\) points.  For \(m>1\),
\[
 P_0(T)=1-T^m\neq1-T.
\]
A finite etale scheme over \(\Fq\) is geometrically connected only in
the one-point case \(\Spec\Fq\).

\section{The multiplicative group}

For \(\Gm=\A^1-\{0\}\),
\[
 \#\Gm(\mathbf F_{q^n})=q^n-1,
\qquad
 Z(\Gm,T)=\frac{1-T}{1-qT}.
 \tag{15.5}\label{eq:zeta-Gm}
\]
We now derive both compact-support groups rather than merely quote them.
For the open immersion \(j:\Gm\hookrightarrow\A^1\) and closed
complement \(i:\{0\}\hookrightarrow\A^1\), the localization triangle
\[
 j_!\Ql\longrightarrow\Ql\longrightarrow i_*\Ql\longrightarrow
\]
gives a long exact compact-support sequence.  Since
\[
 H_c^\bullet(\A^1)=\Ql(-1)[-2],
\qquad
 H_c^\bullet(\{0\})=\Ql[0],
\]
the connecting map identifies
\[
 H_c^1(\Gm,\Ql)=\Ql,
\qquad
 H_c^2(\Gm,\Ql)=\Ql(-1),
\]
and all other groups vanish.  The construction is Frobenius equivariant.
Frobenius therefore acts by \(1\) in degree \(1\) and by \(q\) in degree
\(2\).  The alternating determinant is
\[
 \frac{\det(1-TF\mid H_c^1)}
      {\det(1-TF\mid H_c^2)}
 =\frac{1-T}{1-qT},
\]
which is \eqref{eq:zeta-Gm}.  The example proves that an odd
compact-support group can create a numerator factor even for a smooth
one-dimensional variety.

\section{Smooth proper curves}

Let \(C_0/\Fq\) be smooth, proper, geometrically connected, of genus
\(g\).  Its etale cohomology has dimensions
\[
 b_0=1,\qquad b_1=2g,\qquad b_2=1.
\]
Consequently
\[
 Z(C_0,T)=\frac{P_1(T)}{(1-T)(1-qT)},
\quad
 P_1(T)=\det(1-TF\mid H^1(C,\Ql)).
 \tag{15.6}\label{eq:zeta-curve}
\]
The theorem gives \(P_1\in\mathbf Z[T]\), independent of \(\ell\), and
every reciprocal root has absolute value \(q^{1/2}\) under every complex
embedding.

The cup-product pairing on \(H^1\) is alternating and satisfies
\[
 \langle Fx,Fy\rangle=q\langle x,y\rangle.
\]
Thus eigenvalues occur in pairs
\[
 \alpha,\quad q/\alpha.
\]
The determinant of a \(2g\)-dimensional symplectic similitude of
multiplier \(q\) is \(q^g\).  Applying
\cref{prop:reciprocal-characteristic} to the middle pairing gives
\[
 P_1(q^{-1}T^{-1})=q^{-g}T^{-2g}P_1(T).
 \tag{15.7}\label{eq:curve-P1-reciprocal}
\]
Substitution into \eqref{eq:zeta-curve}, with
\(\chi=2-2g\), yields
\[
 Z(C_0,q^{-1}T^{-1})
 =q^{1-g}T^{2-2g}Z(C_0,T).
\]
Thus \(\varepsilon=1\) for every smooth proper geometrically connected
curve.

The trace formula in degree one is
\[
 \#C_0(\Fq)=1+q-\Tr(F\mid H^1).
\]
If the eigenvalues are \(\alpha_1,\ldots,\alpha_{2g}\), then every
complex embedding has \(|\alpha_j|=q^{1/2}\).  The triangle inequality
therefore gives
\[
 \left|\#C_0(\Fq)-(q+1)\right|\leq2gq^{1/2}.
 \tag{15.8}\label{eq:hasse-weil-bound}
\]
For \(g=0\) this recovers the projective-line count.  For \(g=1\), write
\[
 P_1(T)=1-aT+qT^2,
\]
where \(a=q+1-\#C_0(\Fq)\); then
\eqref{eq:hasse-weil-bound} is \(|a|\leq2\sqrt q\).

In this book the curve bound is a consequence of the terminal
smooth-proper purity theorem.  It is not used as the first step of a
Lefschetz-pencil proof.

\section{Products and Kunneth weights}

Let \(X_0\) and \(Y_0\) be smooth and proper.  Pointwise,
\[
 \#(X_0\times Y_0)(\mathbf F_{q^n})
 =\#X_0(\mathbf F_{q^n})\#Y_0(\mathbf F_{q^n}).
 \tag{15.9}\label{eq:product-counts}
\]
This does not make zeta functions multiply: the logarithmic coefficient
of a zeta function is the point count, and
\eqref{eq:product-counts} is a product of those coefficients.

The Kunneth isomorphism is Frobenius equivariant:
\[
 H^r(X\times Y,\Ql)
 \simeq\bigoplus_{i+j=r}
 H^i(X,\Ql)\otimes H^j(Y,\Ql).
 \tag{15.10}\label{eq:kunneth-example}
\]
If \(\alpha\) and \(\beta\) are eigenvalues on the two factors, the
tensor eigenvalue is \(\alpha\beta\).  If their weights are \(i\) and
\(j\), then for every complex realization
\[
 |\iota(\alpha\beta)|
 =q^{i/2}q^{j/2}=q^{(i+j)/2}.
\]
Thus weights add exactly as cohomological degrees add.  Taking
dimensions with alternating signs in \eqref{eq:kunneth-example} gives
\[
 \chi(X_0\times Y_0)=\chi(X_0)\chi(Y_0).
\]
Poincare multipliers multiply as well:
\(q^{\dim X}q^{\dim Y}=q^{\dim X+\dim Y}\).

For example,
\[
 Z(\PP^1\times\PP^1,T)
 =\frac1{(1-T)(1-qT)^2(1-q^2T)}.
\]
The square on \(1-qT\) records the two generators of \(H^2\), one from
each ruling.  This is a useful multiplicity and tensor-product test.

\section{Extension of the finite base field}

Fix \(m\geq1\), and put
\[
 X_m=X_0\times_{\Fq}\mathbf F_{q^m}.
\]
Geometric Frobenius for the new ground field is \(F^m\).  If the
eigenvalues of \(F\) on \(H^i\) are \(\alpha_{ij}\), then
\[
 P_i(X_m,T)
 =\det(1-TF^m\mid H^i(X,\Ql))
 =\prod_j(1-\alpha_{ij}^mT).
 \tag{15.11}\label{eq:base-change-polynomial}
\]
Purity is preserved with the same numerical weight:
\[
 |\iota(\alpha_{ij}^m)|=q^{mi/2}=(q^m)^{i/2}.
\]
The top eigenvalue becomes \(q^{md}\), and the functional equation uses
\((q^m)^{-d}\).  Betti numbers and Euler characteristic do not change.

There is also a point-count identity.  Let \(\mu_m\) be the set of
complex \(m\)-th roots of unity.  Then
\[
 Z(X_m,T^m)=\prod_{\zeta\in\mu_m}Z(X_0,\zeta T).
 \tag{15.12}\label{eq:base-change-zeta}
\]
Taking logarithms, the right side becomes
\[
\begin{aligned}
 \sum_{\zeta^m=1}\sum_{n\geq1}
 \#X_0(\mathbf F_{q^n})\frac{\zeta^nT^n}{n}
 &=\sum_{r\geq1}
 \#X_0(\mathbf F_{q^{mr}})\frac{T^{mr}}r,
\end{aligned}
\]
because the root-of-unity sum is \(m\) when \(m\mid n\) and zero
otherwise.  The last expression is \(\log Z(X_m,T^m)\), proving
\eqref{eq:base-change-zeta}.

\begin{formalization}
These calculations should be executable tests with the following
expected outputs:
\begin{enumerate}
\item `Spec Fq`: \(P_0=1-T\), and the dimension-zero functional-equation
      sign is \(-1\).
\item `A1`: \(H_c^2=\Ql(-1)\), geometric Frobenius has eigenvalue \(q\),
      and \(Z=(1-qT)^{-1}\).
\item `ProjectiveSpace r`: the even-degree eigenvalues are
      \(1,q,\ldots,q^r\), giving \eqref{eq:zeta-projective-space}.
\item finite etale schemes: a cycle of length \(m\) has determinant
      \(1-T^m\); ordinary connectedness does not force \(P_0=1-T\).
\item `Gm`: the localization sequence produces an odd numerator
      \(1-T\) and even denominator \(1-qT\).
\item curves: a \(2g\)-dimensional alternating similitude has determinant
      \(q^g\) and satisfies \eqref{eq:curve-P1-reciprocal}.
\item products: tensor eigenvalues multiply, weights add, and Euler
      characteristics multiply.
\item field extension: Frobenius becomes \(F^m\), and both
      \eqref{eq:base-change-polynomial} and
      \eqref{eq:base-change-zeta} hold.
\end{enumerate}
These tests should depend only on proved cohomology computations and the
algebraic core, and should import no project-specific axioms.
\end{formalization}


\appendix
\part{Foundations, audits, and formalization}
\chapter{Foundational Theorem Dossier}
\label{app:proetale-foundations}

This appendix is a ledger of the large theorems that support the proof.  Its
purpose is to prevent two opposite errors: pretending that the pro-etale site
makes finiteness, duality, or the trace formula formal, and treating the
Deligne--Laumon weight theorem as an unexplained black box.  Every entry has a
status, an exact interface, and a declared use.

\section{Status notation and global hypotheses}

The four statuses are:
\[
\begin{array}{c|l}
\mathbf F&
 \text{foundational theorem imported with an exact primary source},\\
\mathbf C&
 \text{construction or comparison carried out in the pro-etale chapters},\\
\mathbf P&
 \text{theorem proved in the main mathematical proof spine},\\
\mathbf A&
 \text{formal assembly from results already proved}.
\end{array}
\tag{A.1}\label{eq:dossier-statuses}
\]

Unless a statement says otherwise, schemes in the proof spine are
quasi-excellent, noetherian, and of finite type over a finite field;
morphisms are separated and of finite type.  In particular they are qcqs,
their structural maps are finitely presented, and \(\ell\) is invertible.
These hypotheses are stronger than necessary but simultaneously satisfy the
finiteness assumptions in the torsion and pro-etale six-operation theorems.

For the integral theory, \(E/\Ql\) is finite,
\(\mathcal O=\mathcal O_E\), \(\varpi\) is a uniformizer, and
\[
 \widehat{\mathcal O}_X
   =\varprojlim_n(\mathcal O/\varpi^n)_X
\]
is the sheaf associated with the topological ring \(\mathcal O\) on
\(X_{\proet}\).  The rational proof is performed over \(E\), followed by
finite coefficient extension when needed.

\section{Pro-etale topology and complete coefficients}

\begin{longtable}{@{}
 >{\raggedright\arraybackslash}p{0.12\textwidth}
 >{\raggedright\arraybackslash}p{0.24\textwidth}
 >{\raggedright\arraybackslash}p{0.39\textwidth}
 >{\raggedright\arraybackslash}p{0.15\textwidth}@{}}
\toprule
Status & Input & Exact interface & Use\\
\midrule
\endhead
\(\mathbf F\) &
Weakly etale topology &
A map is weakly etale when it and its diagonal are flat.  Weakly etale
objects with fpqc covering families define \(X_{\proet}\); etale maps are
exactly finitely presented weakly etale maps. &
Definition of the site and functoriality.\\

\(\mathbf F\) &
Weakly contractible basis &
Every scheme has a pro-etale cover by affine \(w\)-contractible objects.
Evaluation on such an object is exact and commutes with limits. &
Derived completion and local perfectness.\\

\(\mathbf F\) &
Repleteness &
For a countable sequential tower
\(F_0\leftarrow F_1\leftarrow\cdots\) of sheaves on the pro-etale topos in
which every transition is an epimorphism,
\(\varprojlim F_n\simeq\Rlim F_n\).  General countable sequential inverse
limits have cohomological dimension at most one, and the derived category is
left complete. &
Convergence of coefficient towers.\\

\(\mathbf F\) &
Etale comparison &
For torsion coefficients, pullback along
\(\nu:X_{\proet}\to X_{\et}\) is fully faithful on the bounded-below derived
category and computes the same cohomology. &
Transport of finite-level theorems.\\

\(\mathbf C\) &
Complete coefficient sheaf &
The sheaf associated with the topological ring \(\mathcal O\) is
\(\widehat{\mathcal O}_X=\varprojlim_n\mathcal O/\varpi^n\), and it is
derived complete. &
An honest integral coefficient object.\\

\(\mathbf C\) &
Constructible category &
\(K\in\Dcons(X_{\proet},\widehat{\mathcal O})\) exactly when \(K\) is
derived complete and \(K\otimes^L\mathcal O/\varpi\) is constructible of
finite Tor amplitude.  On topologically noetherian \(X\), this is equivalent
to local perfectness along a finite stratification. &
Primary derived category of the notes.\\

\(\mathbf C\) &
Coefficient change &
\[
 \begin{gathered}
  \Dcons(X,\widehat{\mathcal O})[\varpi^{-1}]\\[-0.2em]
  {}\simeq \Dcons(X,E),\\
  K\simeq\Rlim_nK/\varpi^n.
 \end{gathered}
\]
Every \(\Qbarl\)-constructible object and every finite diagram involving it
descends to some finite \(E/\Ql\). &
Lattices and rational weights.\\
\bottomrule
\end{longtable}

The primary source for the first six rows is
\href{https://people.mpim-bonn.mpg.de/scholze/proetale.pdf}
{Bhatt--Scholze, \emph{The pro-etale topology for schemes}},
especially Definition 1.2, Theorem 1.5, Proposition 3.1.10,
Corollary 5.1.6, Definition 6.5.1, Proposition 6.6.11, and
Proposition 6.8.14.  The characterization
\emph{lisse means dualizable, constructible means lisse on a
stratification} is developed systematically in
\href{https://arxiv.org/abs/2305.18131}
{Hemo--Richarz--Scholbach, \emph{Constructible sheaves on schemes}}.

\begin{warning}
Geometric points are conservative for the constructible category by
finite-level comparison and stratification.  The notes do not require, and
do not assert, that the usual set of geometric points detects every object
of the entire unbounded pro-etale derived category.
\end{warning}

\section{The six operations}

There are two logically separate claims: the functors exist on an ambient or
constructible derived category, and they preserve bounded constructibility.
The second contains the deep geometric input.

\begin{foundation}[title={Active architecture: Route B}]
The ambient category in this edition is
\(\Dsolid(X,\Lambda)\), the reflective localization of solid pro-etale
complexes.  It has the six operations before constructibility is imposed,
by solid pro-etale motives and rigidity
\cite[Theorem~3.3.7]{RuimyTubachWolfSolid}.  Constructible complexes then
embed as a full subcategory.  The table below records the proof that the
ambient operations preserve that subcategory and their comparison with the
classical adic theory; it is not a levelwise definition of the ambient
functors.
\end{foundation}

\begin{longtable}{@{}
 >{\raggedright\arraybackslash}p{0.11\textwidth}
 >{\raggedright\arraybackslash}p{0.20\textwidth}
 >{\raggedright\arraybackslash}p{0.45\textwidth}
 >{\raggedright\arraybackslash}p{0.14\textwidth}@{}}
\toprule
Status & Theorem & Form used & Dependency\\
\midrule
\endhead
\(\mathbf F\) &
Torsion finiteness &
For \(\Lambda_n=\mathcal O/\varpi^n\), finite-type direct images of bounded
constructible complexes have bounded constructible cohomology; the
cohomological amplitude is uniformly bounded in \(n\). &
SGA 4 and SGA \(4\frac12\).\\

\(\mathbf C\) &
Completed pullback &
The left adjoint on derived-complete objects is
\[
 \begin{aligned}
  Lf_{\comp}^*K&=\bigl(Lf^*K\bigr)^\wedge,\\
  (Lf_{\comp}^*K)/\varpi^n
    &\simeq Lf^*(K/\varpi^n).
 \end{aligned}
\]
It preserves constructibility. &
Completion plus torsion comparison.\\

\(\mathbf C\) &
Direct image &
\(Rf_*\) preserves derived completeness.  Under the stated finiteness
hypotheses it preserves constructibility and reduces to finite-level
\(Rf_*\) modulo every \(\varpi^n\). &
Uniform torsion finiteness.\\

\(\mathbf C\) &
Exceptional direct image &
For \(f=\bar f j\), with \(j\) open and \(\bar f\) proper, set
\[
 Rf_!K=R\bar f_*j_!K.
\]
Common refinements give canonical comparison isomorphisms satisfying the
cocycle condition. &
Nagata compactification and proper base change.\\

\(\mathbf C\) &
Exceptional pullback &
\(Rf^!\) is the right adjoint of \(Rf_!\) on the constructible category and
\[
 (Rf^!K)/\varpi^n\simeq Rf_n^!(K/\varpi^n).
\]
For smooth \(f\) of pure relative dimension \(d\),
\(Rf^!K\simeq Lf^*K(d)[2d]\). &
Finite-level duality and purity.\\

\(\mathbf C\) &
Tensor and internal Hom &
Constructible objects are closed under derived tensor and internal Hom.
The integral tensor is already derived complete on constructible objects;
both operations commute with finite reduction in their finite-Tor range. &
Local perfectness.\\

\(\mathbf F/\mathbf C\) &
Base change and projection formula &
Proper base change for \(Rf_*\), arbitrary base change for \(Rf_!\), smooth
base change, projection formula, localization, and Kunneth hold at finite
level and lift conservatively through all reductions. &
SGA finite-level theorems, then completion.\\
\bottomrule
\end{longtable}

The active ambient construction is
\cite[Theorems~2.5.4 and~3.3.7]{RuimyTubachWolfSolid}.  Its proof of proper
base change uses torsion constructible generators, so the torsion row of the
table remains a genuine foundational dependency.  Constructible preservation
and comparison with classical adic complexes are in Bhatt--Scholze
Section~6.7 and Hemo--Richarz--Scholbach Theorem~7.7.  The particularly
important nonexample remains Bhatt--Scholze Remark~6.7.8: \(Rf_!\) is not, in
general, the derived functor of a degree-zero abelian \(f_!\).  Route B uses
the coherent ambient compactification construction and never that invalid
shortcut.

\section{Duality, purity, and trace}

\begin{longtable}{@{}
 >{\raggedright\arraybackslash}p{0.11\textwidth}
 >{\raggedright\arraybackslash}p{0.21\textwidth}
 >{\raggedright\arraybackslash}p{0.45\textwidth}
 >{\raggedright\arraybackslash}p{0.13\textwidth}@{}}
\toprule
Status & Theorem & Exact form consumed & Source\\
\midrule
\endhead
\(\mathbf F\) &
Dualizing complex &
For excellent \(\ell\)-coprime \(X\) with a dimension function, there is
\(\omega_X\in\Dcons(X,\widehat{\mathcal O})\) such that
\[
 \D_XK=R\mathcal Hom(K,\omega_X),\qquad
 K\xrightarrow{\sim}\D_X\D_XK
\]
on constructible objects. &
ILO; Bhatt--Scholze 6.7.20.\\

\(\mathbf F\) &
Absolute cohomological purity &
For a regular closed immersion \(i:Z\hookrightarrow X\) of codimension
\(c\), with regular source and target and \(\ell\) invertible,
\[
 Ri^!\widehat{\mathcal O}_X
   \simeq\widehat{\mathcal O}_Z(-c)[-2c].
\]
&
Gabber purity.\\

\(\mathbf F/\mathbf C\) &
Verdier duality &
\[
 \D_YRf_!\simeq Rf_*\D_X,\qquad Rf_!\dashv Rf^!.
\]
compatible with composition, localization, and coefficient reduction. &
Finite level, then completion.\\

\(\mathbf F/\mathbf C\) &
Perfectness &
For \(X\) separated of finite type over a separably closed or finite field
and \(K\) constructible of finite Tor amplitude,
\(R\Gamma_c(X,K)\) is a perfect coefficient complex. &
Finiteness and topological Nakayama.\\

\(\mathbf F\) &
Lefschetz--Verdier trace &
For a constructible Weil complex and a proper fixed-point correspondence,
the categorical trace equals the sum of local terms.  For geometric
Frobenius on a finite-type scheme,
\[
 \begin{aligned}
 &\sum_{x\in X_0(\mathbf F_{q^n})}
   \Tr(F_x^n\mid K_{\bar x})\\[-0.2em]
 &\qquad =\sum_i(-1)^i
   \Tr(F^n\mid H_c^i(X,K)).
 \end{aligned}
\]
&
SGA 5 and SGA \(4\frac12\).\\
\bottomrule
\end{longtable}

In the rational-point formula, \(F_x^n\) means the action at the chosen
\(\mathbf F_{q^n}\)-point.  If it is rewritten as a sum over closed points
of degree \(e\mid n\), the operator is \(F_x^{\,n/e}\).  This conversion is
recorded separately so that no exponent convention is hidden.

\section{Perverse foundations and the integral warning}

The middle perverse \(t\)-structure over \(E\) or \(\Qbarl\) consumes the
following non-arithmetic statements from BBD \cite{BBD}:

\begin{longtable}{@{}
 >{\raggedright\arraybackslash}p{0.12\textwidth}
 >{\raggedright\arraybackslash}p{0.31\textwidth}
 >{\raggedright\arraybackslash}p{0.31\textwidth}
 >{\raggedright\arraybackslash}p{0.16\textwidth}@{}}
\toprule
Status & Input & Form used & Primary location\\
\midrule
\endhead
\(\mathbf F/\mathbf P\) &
Gluing of the middle \(t\)-structure &
Support and cosupport conditions glue over a finite smooth
stratification, are independent of refinement, and give a bounded
\(t\)-structure. &
BBD 1.4, 2.1--2.2.\\

\(\mathbf F/\mathbf P\) &
Affine Artin vanishing &
For affine \(f\), \(Rf_*\) is right perverse \(t\)-exact and \(Rf_!\) is
left perverse \(t\)-exact; quasi-finite affine maps are \(t\)-exact. &
BBD 4.1.1--4.1.3.\\

\(\mathbf F/\mathbf P\) &
Functorial amplitudes &
Smooth pullback shifted by relative dimension, finite direct image, and
external product are perverse \(t\)-exact. &
BBD 4.2.4, 4.2.7--4.2.8.\\

\(\mathbf P\) &
Intermediate extension &
\(j_{!*}P=\im({}^pj_!P\to{}^pj_*P)\), characterized by the absence of
boundary subobjects and quotients. &
BBD 1.4.23 and 4.3.1.\\

\(\mathbf F/\mathbf P\) &
Finite length and simple objects &
The rational perverse heart has finite length; its simple objects are
intersection complexes of irreducible lisse sheaves on smooth strata. &
BBD 4.3.1.\\
\bottomrule
\end{longtable}

Over \(\mathcal O\), Verdier duality does not preserve the ordinary
\(p\)-perverse heart.  It exchanges \(p\) and \(p^+\), the tilt by the
torsion/torsion-free pair.  This is BBD Section 3.3.  The notes prove the
comparison with rationalization and use the self-dual middle perversity only
after passage to \(E\).  No integral decomposition theorem is assumed.

\section{Arithmetic inputs upstream of the main weight theorem}

The following inputs are permitted before the general direct-image theorem.
None is smooth-proper purity in arbitrary dimension.

\begin{longtable}{@{}
 >{\raggedright\arraybackslash}p{0.11\textwidth}
 >{\raggedright\arraybackslash}p{0.24\textwidth}
 >{\raggedright\arraybackslash}p{0.43\textwidth}
 >{\raggedright\arraybackslash}p{0.12\textwidth}@{}}
\toprule
Status & Input & Exact role & Source\\
\midrule
\endhead
\(\mathbf F\) &
Rank-one theory &
Class field theory and determinant arguments control weights for rank-one
lisse Weil sheaves. &
Weil II 1.2--1.4.\\

\(\mathbf F\) &
Local monodromy weights &
For a pure lisse sheaf on a punctured trait, the graded pieces of the
monodromy filtration have the prescribed shifted weights. &
Weil II 1.8.4.\\

\(\mathbf F/\mathbf P\) &
Real-sheaf criterion &
For an \(\iota\)-real lisse Weil sheaf, every irreducible constituent is
\(\iota\)-pure, with the compatible pairing of weights in Deligne's
criterion.  A finite-order determinant after a Tate or constant twist is
used to normalize the determinant and the weight; that condition alone does
not imply purity. &
Weil II 1.5.1.\\

\(\mathbf F\) &
Fourier support-forgetting &
For the Artin--Schreier Fourier kernel, the canonical map from
\(!\)-pushforward to \(*\)-pushforward is an isomorphism. &
Laumon 1.3.1.1.\\

\(\mathbf P\) &
Fourier formalism &
Inversion, duality with the inverse character, Plancherel, and perverse
\(t\)-exactness are derived from the six operations, character
orthogonality, support-forgetting, and affine exactness. &
Laumon 1.2--1.3.\\

\(\mathbf P\) &
Curve estimate &
For lisse input of weights at most \(w\), the relevant
\(H_c^1\) eigenvalues have weights at most \(w+1\), with local boundary
terms controlled separately. &
Laumon 4; Weil II 3.2.3.\\
\bottomrule
\end{longtable}

The Artin--Schreier kernel is defined over a finite \(E/\Ql\) containing the
values of the chosen additive character and has an
\(\mathcal O_E\)-lattice.  It is not silently assumed to be defined over
\(\Zl\).

\section{Outputs proved in the notes}

The following statements are not foundational imports:

\begin{enumerate}
\item the perverse Fourier mechanism and the curve weight estimate;
\item the perverse-first theorem
      \[
       \operatorname{WeightsLE}(\cP,w)
       \Longrightarrow
       \operatorname{WeightsLE}(\pH^aRf_!\cP,w+a);
      \]
\item the derived \(Rf_!\) upper-weight theorem obtained from the perverse
      truncation tower;
\item the dual \(Rf_*\) lower bound and proper direct-image purity;
\item purity of \(j_{!*}\), intersection complexes, and proper intersection
      cohomology;
\item geometric semisimplicity of pure perverse sheaves and the
      decomposition theorem after geometric base change;
\item purity of \(H^i\) for smooth proper varieties;
\item rationality, the functional equation, integrality of the individual
      factors, and independence of \(\ell\).
\end{enumerate}

\section{Terminal dependency audit}

\[
\begin{tikzcd}[column sep=small,row sep=small]
\substack{\mathbf C:\ \text{solidification and}\\
                 \text{pro-etale motives}}
 \arrow[r] &
\substack{\mathbf C:\ \text{solid rigidity}}
 \arrow[r] &
\substack{\mathbf C:\ \text{ambient coherent}\\
                 \text{six operations}} \\
\substack{\mathbf F:\ \text{torsion base change}\\
                 \text{and rigidity}}
 \arrow[ur] \arrow[r] &
\substack{\mathbf F:\ \text{finiteness, duality,}\\
                 \text{purity, and trace}}
 \arrow[r] &
\substack{\mathbf C:\ \text{restriction to}\\
                 \Dcons}
 \arrow[u] \arrow[d] \\
\substack{\mathbf A:\ \text{classical}\\
                 \text{Weil conjectures}}
 &
\substack{\mathbf P:\ \text{proper purity,}\\
                 \text{IC and decomposition}}
 \arrow[l] &
\substack{\mathbf P:\ \text{perverse Fourier}\\
                 \text{weight engine}}
 \arrow[l]
\end{tikzcd}
\tag{A.2}\label{eq:dossier-terminal-dependencies}
\]

There is no reverse arrow from smooth-proper purity, decomposition, hard
Lefschetz, integral Frobenius factors, or \(\ell\)-independence to the
Fourier or direct-image nodes.

\chapter{The Pro-etale--Classical Adic Dictionary}
\label{app:proetale-classical-dictionary}

The pro-etale category does not define a new Weil cohomology theory.  It gives
a direct model for the classical \(\ell\)-adic category.  This appendix states
the comparison in the precise form used by the notes and records which
expressions correspond only by canonical isomorphism rather than by literal
equality.

\section{The two integral models}

Let \(X\) be topologically noetherian, let
\(\mathcal O\) be a complete discrete valuation ring with uniformizer
\(\varpi\), and put \(\Lambda_n=\mathcal O/\varpi^n\).

The classical model begins with projective systems
\[
  K_\bullet=(K_n)_{n\geq1},
  \qquad
  K_n\in\Dbc(X_{\et},\Lambda_n),
\]
equipped with derived normalization isomorphisms
\[
  K_{n+1}\otimes^L_{\Lambda_{n+1}}\Lambda_n
     \xrightarrow{\sim}K_n.
\tag{B.1}\label{eq:dictionary-normalized-system}
\]
One then localizes an appropriate derived category of systems at AR-null
complexes.  We denote the result schematically by
\(\mathrm D^b_{c,\mathrm{adic}}(X,\mathcal O)\).

The pro-etale model uses the actual sheaf
\[
  \widehat{\mathcal O}_X
     =\varprojlim_n(\Lambda_n)_X
\]
and the full category
\[
  \Dcons(X_{\proet},\widehat{\mathcal O}_X)
  \subset
  \D(X_{\proet},\widehat{\mathcal O}_X)
\]
of derived-complete constructible complexes.

\begin{theorem}[Integral comparison dictionary]
\label{thm:dictionary-integral-comparison}
There is a canonical equivalence
\[
 \Phi_X:
 \mathrm D^b_{c,\mathrm{adic}}(X,\mathcal O)
 \xrightarrow{\sim}
 \Dcons(X_{\proet},\widehat{\mathcal O}_X)
\tag{B.2}\label{eq:dictionary-integral-equivalence}
\]
characterized on a normalized system by
\[
 \Phi_X(K_\bullet)
   \simeq
 \Rlim_n\nu^*K_n,
\tag{B.3}\label{eq:dictionary-realization}
\]
where \(\nu:X_{\proet}\to X_{\et}\).  Its inverse sends \(K\) to the
compatible reductions
\[
 K_n=K\otimes^L_{\widehat{\mathcal O}_X}\Lambda_n.
\tag{B.4}\label{eq:dictionary-reductions}
\]
After inverting \(\varpi\), it gives
\[
 \mathrm D^b_{c,\mathrm{adic}}(X,E)
 \xrightarrow{\sim}
 \Dcons(X_{\proet},E_X).
\tag{B.5}\label{eq:dictionary-rational-equivalence}
\]
\end{theorem}

\begin{proof}
Derived completeness gives
\[
 K\simeq\Rlim_n K_n.
\]
Every \(K_n\) is classical by the torsion comparison
\(\D^+(X_{\et},\Lambda_n)\to\D^+(X_{\proet},\Lambda_n)\).
Conversely a normalized constructible system has a derived inverse limit;
repleteness and the uniform Tor-amplitude bound show that reduction of this
limit is \(K_n\).  A morphism is an isomorphism exactly when all finite
reductions are isomorphisms.  Hence the two constructions are inverse after
the classical localization by pro-zero errors.  Inverting \(\varpi\) gives
\eqref{eq:dictionary-rational-equivalence}.  This is the content of the
comparison developed in Bhatt--Scholze Sections 6.5--6.8.
\end{proof}

\begin{proetalebox}
The equivalence explains the simplification.  AR-null complexes have not
been declared nonexistent; their localization has been absorbed into
derived completion in a replete topos.  The public object is \(K\), and its
coefficient tower is recovered functorially as \((K/\varpi^n)_n\).
\end{proetalebox}

\section{Object-by-object dictionary}

\begin{longtable}{
 >{\raggedright\arraybackslash}p{0.23\textwidth}
 >{\raggedright\arraybackslash}p{0.31\textwidth}
 >{\raggedright\arraybackslash}p{0.36\textwidth}}
\toprule
Concept & Classical adic model & Pro-etale model\\
\midrule
\endhead
Coefficient ring &
The tower \((\Lambda_n)_n\). &
The sheaf
\(\widehat{\mathcal O}_X=\varprojlim_n(\Lambda_n)_X\).\\

Integral complex &
Normalized constructible system, modulo AR-null systems. &
Derived-complete \(K\) with constructible perfect reduction.\\

Reduction modulo \(\varpi^n\) &
The \(n\)-th normalized component \(K_n\). &
\(K\otimes^L_{\widehat{\mathcal O}}\Lambda_n\).\\

Completion &
Renormalization of the coefficient tower. &
The idempotent functor
\[
 K^\wedge=\Rlim_n(K\otimes^L\Lambda_n).
\]\\

Lisse integral object &
A compatible tower of finite locally constant
\(\Lambda_n\)-modules, locally finite free. &
A locally finite-free
\(\widehat{\mathcal O}\)-sheaf; derived lisse objects are locally perfect.\\

Constructible object &
Lisse normalized systems along a finite stratification. &
Locally perfect derived objects along a finite stratification.\\

Rational object &
An integral lattice with morphisms tensored by \(E\), up to commensurability. &
An actual constructible \(E_X\)-complex; equivalently an integral
constructible object after inverting \(\varpi\).\\

\(\Qbarl\)-object &
Extension from a sufficiently large finite coefficient field. &
The same finite-field descent, expressed using the sheaf associated with the
topological algebraic extension.\\

Stalk &
The compatible system \((K_{n,\bar x})_n\), then inverse limit. &
The pullback of the actual complete complex to the geometric point.\\

Cohomology &
\(\Rlim_nR\Gamma(X_{\et},K_n)\). &
\(R\Gamma(X_{\proet},K)\) literally.\\

Continuous representation &
A compatible system of finite continuous representations of the classical
profinite etale fundamental group. &
For connected locally noetherian \(X\) with a geometric base point \(\bar x\),
a locally constant pro-etale sheaf, equivalently a continuous representation
of the Noohi group \(\pi_1^{\proet}(X,\bar x)\).\\
\bottomrule
\end{longtable}

The base-pointed representation statement in the last row is made only under
the displayed connectedness and local-noetherian hypotheses.  Finite
discrete representations, and those \(E\)-representations that stabilize an
\(\mathcal O_E\)-lattice, factor through the classical profinite
\(\pi_1^{\et}(X,\bar x)\).  On geometrically unibranch schemes this captures
the lisse \(E\)-objects relevant to the proof.  For a general nonnormal
scheme, however, the Noohi group \(\pi_1^{\proet}\) can be strictly larger,
and a continuous finite-dimensional \(E\)-representation need not preserve a
lattice.  The Noohi and classical statements must therefore not be identified
globally.

\section{The six-functor dictionary}

Write \(K_n=K\otimes^L\Lambda_n\).  Under
\eqref{eq:dictionary-integral-equivalence}, every operation is characterized
by finite reduction and its adjunction.

\begin{longtable}{
 >{\raggedright\arraybackslash}p{0.19\textwidth}
 >{\raggedright\arraybackslash}p{0.34\textwidth}
 >{\raggedright\arraybackslash}p{0.37\textwidth}}
\toprule
Operation & Pro-etale construction & Comparison formula\\
\midrule
\endhead
Pullback &
Completed derived pullback
\(Lf_{\comp}^*K=(Lf^*K)^\wedge\). &
\[
 (Lf_{\comp}^*K)/\varpi^n
   \simeq Lf^*(K_n).
\]\\

Direct image &
Derived pushforward \(Rf_*\), which preserves completeness. &
\[
 (Rf_*K)/\varpi^n
   \simeq Rf_*(K_n)
\]
in the bounded constructible range.\\

Exceptional direct image &
Compactify \(f=\bar f j\) and set
\(Rf_!=R\bar f_*j_!\). &
\[
 (Rf_!K)/\varpi^n
   \simeq Rf_!(K_n).
\]\\

Exceptional pullback &
Right adjoint \(Rf^!\) of \(Rf_!\). &
\[
 (Rf^!L)/\varpi^n
   \simeq Rf_n^!(L_n).
\]\\

Tensor &
Derived tensor over \(\widehat{\mathcal O}_X\); on constructible objects it
is already complete. &
\[
 (K\otimes^LL)/\varpi^n
 \simeq K_n\otimes^L_{\Lambda_n}L_n.
\]\\

Internal Hom &
Derived sheaf Hom in the finite-Tor constructible category. &
{\small
\[
 \mathrm{RHom}(K,L)/\varpi^n
 \simeq \mathrm{RHom}(K_n,L_n).
\]
}\\

Duality &
\(\D_XK=R\mathcal Hom(K,\omega_X)\). &
\[
 (\D_XK)/\varpi^n
   \simeq\D_{X,n}(K_n).
\]\\
\bottomrule
\end{longtable}

These formulas are canonical isomorphisms of functors.  They are not
definitional equalities between selected resolutions.  Their coherence under
composition is part of the comparison theorem.

\begin{warning}[The exceptional direct image]
\label{warn:dictionary-fshriek}
Even on the pro-etale site, the constructible \(Rf_!\) is not generally the
right or left derived functor of the degree-zero operation on all abelian
sheaves.  Compactification is part of its definition.  Any formalization
should therefore construct it at the derived level and prove independence of
compactification before exposing the notation \(Rf_!\).
\end{warning}

\section{Perverse structures under the comparison}

For field coefficients \(E\), the equivalence
\eqref{eq:dictionary-rational-equivalence} preserves geometric stalks,
costalks, dimensions of supports, and Verdier duality.  It is therefore
\(t\)-exact for the middle perverse structures:
\[
 \Phi_X({}^p\D^{\leq0}_{\mathrm{adic}})
   ={}^p\D^{\leq0}_{\proet},
\qquad
 \Phi_X({}^p\D^{\geq0}_{\mathrm{adic}})
   ={}^p\D^{\geq0}_{\proet}.
\tag{B.6}\label{eq:dictionary-perverse-texact}
\]
Consequently it commutes canonically with perverse truncation, perverse
cohomology, \(j_{!*}\), and intersection complexes.

Integral coefficients require a two-column dictionary:

\begin{center}
\begin{tabularx}{0.94\textwidth}{>{\raggedright\arraybackslash}p{0.26\textwidth}XX}
\toprule
Operation & \(p\)-heart & \(p^+\)-heart\\
\midrule
Verdier duality &
\(\D:{}^p\Perv(X,\mathcal O)^{\op}\to{}^{p^+}\Perv(X,\mathcal O)\) &
\(\D:{}^{p^+}\Perv(X,\mathcal O)^{\op}\to{}^p\Perv(X,\mathcal O)\)\\
Intermediate extension &
\({}^pj_{!*}\) &
\({}^{p^+}j_{!*}\)\\
Rationalization &
\multicolumn{2}{c}{Both map to the same
\(j_{!*}\) in \(\Perv(X,E)\).}\\
\bottomrule
\end{tabularx}
\end{center}

The two integral middle extensions may differ by boundary
\(\varpi\)-torsion.  Rationalization kills that difference.  The rational
perverse heart is Verdier self-dual and is the category in which weights,
simple objects, pure semisimplicity, and decomposition are used.

\section{Weil structures and Frobenius}

Let \(X_0/\Fq\), let \(X=X_0\times_{\Fq}\Fbar\), and let
\(F_X:X\to X\) be geometric Frobenius in the convention of these notes.  A
Weil complex is
\[
  (K,\phi),\qquad
  \phi:F_X^*K\xrightarrow{\sim}K.
\tag{B.7}\label{eq:dictionary-weil-object}
\]
Under the integral comparison, this is the same as a compatible system of
finite-level Weil structures on \(K_n\).  The comparison isomorphisms for all
six operations commute with \(\phi\), so traces and local terms agree in the
two models.

The Tate object has the three matching presentations
\[
 \Zl(1)=\varprojlim_n\mu_{\ell^n},
 \qquad
 \widehat{\mathcal O}(1)
   =\Zl(1)\widehat{\otimes}_{\Zl}\mathcal O,
\qquad
 E(1)=\widehat{\mathcal O}(1)[\varpi^{-1}],
\qquad
 F\mid E(1)=q^{-1}.
\tag{B.8}\label{eq:dictionary-tate}
\]
The last equality uses geometric Frobenius.  It becomes \(q\) for arithmetic
Frobenius.

\section{What is and is not gained}

\begin{longtable}{
 >{\raggedright\arraybackslash}p{0.44\textwidth}
 >{\raggedright\arraybackslash}p{0.46\textwidth}}
\toprule
Eliminated from the proof spine & Still mathematically required\\
\midrule
\endhead
AR-null objects in every public definition. &
Derived completeness and conservativity of finite reduction.\\

Repeated choices of normalized representatives. &
Canonical comparison isomorphisms and their coherence.\\

Pretending that \(\Ql\) is a sheaf on the small etale site. &
Finite-level etale finiteness and base-change theorems.\\

Defining \(R\Gamma(X,\Ql)\) as an external inverse limit. &
The theorem identifying pro-etale cohomology with continuous etale
cohomology.\\

Coefficient towers in the statements of perverse and Fourier functors. &
Integral \(p/p^+\) bookkeeping and finite coefficient-field descent.\\
\bottomrule
\end{longtable}

Thus the pro-etale model shortens the foundational interface without
weakening it.  The classical model remains available as a comparison target
and as the source of finite-level geometric theorems.

\chapter{Noncircularity and Dependency Audit}
\label{app:perverse-dependency-audit}

This appendix makes the proof graph explicit.  An arrow means that the target
uses the source.  There is no implicit appeal to a theorem later in the graph.

\section{The main graph}

\[
\begin{tikzcd}[column sep=small,row sep=large]
\substack{\text{six operations}\\\text{base change, trace, duality}}
 \arrow[r] \arrow[dr] &
\substack{\text{perverse \(t\)-structure}\\j_{!*},\ \IC,\ \text{Artin vanishing}}
 \arrow[d] \\
&
\substack{\text{Fourier inversion}\\\text{\(t\)-exactness, Plancherel}}
 \arrow[d]\\
\substack{\text{rank-one purity}\\\text{real-sheaf theorem}}
 \arrow[r]&
\substack{\text{curve weight estimate}\\H_c^a\text{ upper weights}}
 \arrow[d]\\
\substack{\text{Noether normalization}\\\text{excision, Leray filtration}}
 \arrow[r]&
\substack{Rf_!\text{ preserves}\\\text{upper weights}}
 \arrow[d]\\
\substack{\text{Verdier duality}\\Rf_!=Rf_*\text{ for proper }f}
 \arrow[r]&
\substack{\text{proper purity}\\\text{IC/IH purity}}
 \arrow[d]\\
\substack{\text{trace and Poincare}\\\text{determinant algebra}}
 \arrow[r]&
\substack{\text{smooth-proper purity}\\\text{zeta and functional equation}}
 \arrow[d]\\
\substack{\text{universal-\(\iota\) weights}\\\text{Fatou and Gauss}}
 \arrow[r]&
\substack{P_i\in\mathbf Z[T]\\\ell\text{-independence}.}
\end{tikzcd}
\tag{B.1}\label{eq:dependency-graph}
\]

The decomposition theorem branches from proper purity and points forward:
\[
  \text{proper purity}
  \Longrightarrow
  \text{pure perverse semisimplicity}
  \Longrightarrow
  \text{pure-complex splitting}
  \Longrightarrow
  \text{decomposition}.
\tag{B.2}\label{eq:decomposition-branch}
\]
There is no arrow from \eqref{eq:decomposition-branch} back to the curve or
direct-image theorem.

\section{Audit of tempting circular shortcuts}

\begin{longtable}{
 >{\raggedright\arraybackslash}p{0.31\textwidth}
 >{\raggedright\arraybackslash}p{0.30\textwidth}
 >{\raggedright\arraybackslash}p{0.29\textwidth}}
\toprule
Tempting shortcut & Why it is circular or false & Replacement used here\\
\midrule
\endhead
“\(Rf_*\) of a pure complex is pure by BBD.” &
BBD proper purity uses Deligne's direct-image weight theorem, whose
smooth-proper specialization is the RH being proved. &
Prove the curve estimate, globalize \(Rf_!\) upper weights, then dualize.\\

“The decomposition theorem proves smooth-proper purity.” &
The finite-field decomposition theorem is proved from purity and weight
separation. &
Purity is proved before decomposition; the latter is recorded as a
corollary.\\

“Fourier transform raises weights by one.” &
The Fourier transform is an \(Rp_!\) construction, so this assertion is a
case of the direct-image theorem. &
Before the curve theorem use only stalk formulas, inversion,
\(t\)-exactness, real mixedness, and Plancherel.\\

“The perverse spectral sequence degenerates.” &
No such general degeneration is available at that stage. &
Use only its finite Frobenius-stable filtration and subquotients.\\

“Rationality of \(Z\) makes each \(P_i\) rational.” &
Alternating numerator and denominator factors can cancel. &
First prove universal weights; distinct moduli rule out cancellation.\\

“Pure perverse means Frobenius semisimple.” &
Purity controls eigenvalue moduli, not Jordan blocks. &
State only geometric semisimplicity after forgetting Frobenius.\\

“\(\Qbarl[d]\) is perverse on every pure \(d\)-fold.” &
Singularities may violate cosupport. &
Use \(\IC_X\); identify it with \(\Qbarl[d]\) only when \(X\) is smooth.\\
\bottomrule
\end{longtable}

\section{The weight-shift audit}

All shifts can be reconstructed from four tests:
\[
\begin{array}{c|c}
\text{object}&\text{weight}\\
\hline
\Qbarl&0\\
\Qbarl(1)&-2\\
\cL[d]\text{ for \(\cL\) pure \(u\)}&u+d\\
\pH^a(\cK)\text{ for \(\cK\) pure \(w\)}&w+a.
\end{array}
\tag{B.3}\label{eq:weight-shift-audit}
\]
For \(X\) smooth proper of dimension \(d\),
\[
  \Qbarl[d]\xrightarrow{Ra_*}
  Ra_*\Qbarl[d]
\xrightarrow{H^{\,i-d}}
  H^i(X,\Qbarl)
\]
has successive weight labels
\[
  d,\qquad d,\qquad d+(i-d)=i.
\tag{B.4}\label{eq:terminal-shift-audit}
\]
Any proof producing \(i-d\), \(i+d\), or \(2d-i\) at the last stage has lost a
shift or dualized the wrong object.

\section{Frobenius audit}

The notes use geometric Galois Frobenius on stalks and cohomology.  Therefore
\[
 F|\Qbarl(1)=q^{-1},\qquad F|\Qbarl(-1)=q,
\]
\[
 P_i(T)=\det(1-TF\mid H^i),\qquad
 |\iota(\alpha)|=q^{i/2}
\]
for an eigenvalue \(\alpha\) of \(F\).  The polynomial roots in the variable
\(T\) are \(\alpha^{-1}\).  Poincare duality pairs eigenvalues
\(\alpha\) and \(q^d/\alpha\).  These conventions are tested on
\(\PP^n\), whose zeta function is
\[
  Z(\PP^n,T)=\prod_{r=0}^{n}(1-q^rT)^{-1}.
\]

\section{Use of all complex realizations}

For each chosen
\(\iota:\Qbarl\simeq\mathbf C\), the entire weight proof yields
\(\iota\)-purity.  Only after repeating it for all \(\iota\) do we use the
arithmetic argument
\[
\begin{gathered}
Z\in\mathbf Z[[T]]\cap\mathbf Q(T)\\
\Downarrow\\
\text{integral reduced numerator and denominator}\\
\Downarrow\quad\text{universal modulus separates weights}\\
P_i(T)\in\mathbf Z[T]\text{ and is independent of }\ell.
\end{gathered}
\]
A fixed embedding gives an analytic modulus but not Galois stability of the
weight strata.

\section{Terminal theorem audit}

The final theorem depends on the following proved objects, not theorem
constants introduced as assumptions:

\begin{enumerate}
\item finite-dimensional \(H^i\) with geometric Frobenius;
\item trace formula for all powers of Frobenius;
\item Poincare duality with multiplier \(q^d\);
\item the derived \(Rf_!\) upper-weight theorem;
\item its Verdier-dual lower theorem and proper-purity corollary;
\item the perverse identity \(\IC_X=\Qbarl[d]\) for smooth \(X\);
\item the algebraic integral-factor and separation lemmas.
\end{enumerate}

Deleting the chapters on semisimplicity and decomposition would not break the
classical Weil theorem.  Deleting proper purity would.  This is the simplest
mechanical check that decomposition has not been used circularly.

\chapter{Formalization Map: Reusable Pro-etale and Perverse APIs}
\label{app:proetale-formalization}

This appendix is a theorem-boundary specification, not Lean code.  It records
an implementation order in which each layer is reusable outside the Weil
conjectures and in which the final theorem cannot acquire a hidden
project-specific axiom.

\section{The starting point and the first audit}

The publicly documented mathlib endpoint already contains the small
pro-etale site and an \(\ell\)-adic coefficient sheaf with its sheaf
cohomology; see
\href{https://leanprover-community.github.io/mathlib4_docs/Mathlib/AlgebraicGeometry/Sites/ElladicCohomology.html}
{mathlib's \texttt{ElladicCohomology} module}.
This is a genuine head start, but it is not yet the constructible
six-functor category used by the proof.

At the pinned project commit, the first mechanical audit should produce a
minimal-import file checking at least:

\begin{enumerate}
\item weakly etale morphisms, the pro-etale site, and its functoriality;
\item the exact type and universe of the existing coefficient sheaf;
\item the ambient sheaf category and its Grothendieck/abelian instances;
\item the available derived category, \(t\)-structure, and localization
      infrastructure;
\item the definitions of sheaf cohomology and the existing
      \(\ell\)-adic cohomology endpoint.
\end{enumerate}

No declaration should be inferred from a similarly named documentation page.
The inventory consists of successful \(\#check\) commands under the minimal
imports used by the eventual file.

\section{Layer 0: coefficient-independent categorical infrastructure}

The first layer should not mention schemes, finite fields, or weights.  It
exports:

The completion identities in this layer are parameterized by hypotheses;
they are not unconditional facts about an arbitrary ringed topos.  The
reconstruction theorem used in the proof spine assumes a replete topos, a
noetherian coefficient ring, and a finitely generated ideal.  In the more
general algebraic API, the ideal must be weakly proregular for the standard
derived-completion formulas.  Comparisons with compatible systems in an
unbounded derived category require the stronger regular-ideal hypothesis
appearing in Bhatt--Scholze, Lemma 3.5.7(5).  For the complete DVR used here,
the principal ideal \((\varpi)\) satisfies these conditions.  No global
instance should silently infer them for arbitrary rings or topoi.

\begin{enumerate}
\item derived completion for a commutative ring object \(R\) and an ideal
      generated by \(\varpi\):
      \[
       \operatorname{complete}_\varpi(K)
        =\Rlim_n(K\otimes_R^LR/\varpi^n);
      \]
\item the predicate \(\operatorname{DerivedComplete}_\varpi(K)\);
\item idempotence of completion and the universal property as a left
      adjoint to inclusion of complete objects;
\item conservativity of all finite reductions on complete objects satisfying
      a separatedness or constructibility hypothesis;
\item completed pullback and completed tensor as generic constructions;
\item dualizable objects in a closed symmetric monoidal derived category;
\item transport of dualizability and completion across equivalences.
\end{enumerate}

\begin{formalization}
Suggested public interfaces include
{\small
\[
\begin{aligned}
\texttt{derivedComplete}(K)&:\operatorname{Prop},\\
\texttt{completion}(K)&:\D(R),\\
\texttt{completionIsoRlim}(K;
  \texttt{Replete},\texttt{Noetherian},\texttt{Ideal.FG})&:
 K^\wedge\simeq\Rlim_n(K\otimes^LR/\varpi^n),\\
\texttt{unboundedSystemComparison}
  (\texttt{Ideal.IsRegular})&:\operatorname{Equivalence},\\
\texttt{isIsoOfReductionIsIso}(f)&:
 \bigl(\forall n,\operatorname{IsIso}(f_n)\bigr)
 \to\operatorname{IsIso}(f).
\end{aligned}
\]
}
A generic commutative-algebra version of the first equivalence should instead
take \texttt{Ideal.WeaklyProregular}.  The last theorem must state the
completeness and separatedness or constructibility hypotheses under which
finite reductions are conservative; it is not true for arbitrary
noncomplete objects.
\end{formalization}

\section{Layer 1: the pro-etale coefficient sheaf}

For a scheme \(X\) and complete noetherian local ring \(\mathcal O\), define
the sheaf associated with the topological ring:
\[
 \widehat{\mathcal O}_X(U)
   =C_{\mathrm{cont}}(\pi_0(U),\mathcal O)
\]
on the weakly contractible basis, followed by sheafification or the
equivalent topological constant-sheaf construction.  Export canonical
isomorphisms
\[
 \widehat{\mathcal O}_X
   \simeq\varprojlim_n(\mathcal O/\varpi^n)_X,
\qquad
 f^{-1}\widehat{\mathcal O}_Y
   \longrightarrow\widehat{\mathcal O}_X.
\tag{C.1}\label{eq:formalization-coefficient-sheaf}
\]

The existing \(\Zl\)-valued sheaf should be identified with the specialization
of this generic construction, not replaced by a second unrelated constant
sheaf.  The coefficient interface should accept \(\mathcal O_E\) as well as
\(\Zl\); otherwise the Artin--Schreier kernel will force a redesign later.

\section{Layer 2: lisse and constructible objects}

The reusable derived definition is
\[
 \operatorname{Lisse}(K)
   :\Longleftrightarrow
   \operatorname{IsDualizable}(K).
\tag{C.2}\label{eq:formalization-lisse-dualizable}
\]
On a weakly contractible affine this must be compared with a perfect complex
of coefficient modules.  Define constructibility by the existence of a
finite locally closed stratification on which all restrictions are lisse.
The restriction of the ambient standard \(t\)-structure is installed only
under the hypotheses of Hemo--Richarz--Scholbach, Theorem 6.2: the condensed
coefficient ring is \(t\)-admissible; for lisse objects, \(X\) has Zariski
locally finitely many irreducible components; and for constructible objects,
every constructible subset of \(X\) has Zariski locally finitely many
irreducible components.  The coefficient-limit and comparison interfaces
separately retain their qcqs and transition-map hypotheses.  Locally
noetherian schemes and the coefficients \(\mathcal O_E,E\) used in these
notes meet the selected conditions, but there is no such instance for an
arbitrary scheme and arbitrary condensed ring.

For the integral category, package the equivalent interface
\[
\begin{aligned}
\operatorname{IntegralConstructible}(K)
 :\Longleftrightarrow{}&
 \operatorname{DerivedComplete}_\varpi(K)\\
 &{}\land
 \operatorname{Constructible}
   (K\otimes^L\mathcal O/\varpi)\\
 &{}\land
 \operatorname{FiniteTorAmplitude}(K).
\end{aligned}
\tag{C.3}\label{eq:formalization-integral-constructible}
\]
The equivalence with stratified local perfectness is a theorem, not a second
definition exposed to downstream users.

\begin{formalization}
The layer should export:
\begin{enumerate}
\item stability under shifts, cones, retracts, tensor, and internal Hom;
\item restriction to locally closed subschemes;
\item common refinement for finitely many constructible objects;
\item geometric-stalk conservativity on constructible objects;
\item finite-level reduction and reconstruction by \(\Rlim\);
\item coefficient extension
      \(\mathcal O\to E\to E'\) and descent of finite diagrams from
      \(\Qbarl\) to finite \(E\).
\end{enumerate}
Do not store a chosen stratification in the public object.  Store its
existence and expose choice-independent restriction theorems.
\end{formalization}

\section{Layer 3: comparison with classical adic systems}

The primary API remains the pro-etale category.  The classical tower model is
a comparison target:
\[
 \texttt{adicSystemRealization}:
 \mathrm D^b_{c,\mathrm{adic}}(X,\mathcal O)
 \xrightarrow{\sim}
 \Dcons(X_{\proet},\widehat{\mathcal O}).
\tag{C.4}\label{eq:formalization-adic-comparison}
\]

Its unit and counit are the reconstruction maps
\[
 \Rlim_n\nu^*K_n\longrightarrow K,
\qquad
 \bigl(\Rlim_n\nu^*K_n\bigr)/\varpi^m\longrightarrow\nu^*K_m.
\]
All later compatibility statements should be phrased as natural
isomorphisms involving this equivalence.  The project should not implement
the AR-null quotient as the primary public category merely because it is the
historical model.

\section{Layer 4: ambient solid operations and constructible restriction}

Route B separates ambient existence from geometric preservation:

\begin{enumerate}
\item the unrestricted pro-etale category and the solidification reflector;
\item a presentable closed monoidal category \(\Dsolid(X,\Lambda)\);
\item the correspondence-level ambient operations
      \(f^*,f_*,f_!,f^!,\otimes,R\mathcal Hom\);
\item the full subcategory \(\Dindcons\) and the theorem
      \(\Dcons=\Dindcons^\omega\);
\item an intrinsic constructibility predicate, defined by finite
      stratifications and dualizability;
\item separate geometric theorems that every ambient operation preserves
      constructibility under explicit hypotheses.
\end{enumerate}

The ambient functor record must not contain constructibility hypotheses.
Those belong to restriction lemmas whose conclusions are functors between
the full constructible subcategories.

The exceptional direct image requires a choice-management structure:
\[
\begin{array}{l}
\texttt{Compactification}(f),\\
\texttt{rfShriekUsing}(c,K),\\
\texttt{rfShriekComparison}(c_1,c_2,K),\\
\texttt{rfShriekComparison\_cocycle},\\
\texttt{rfShriek}(f,K)
\quad\text{with a canonical-isomorphism specification}.
\end{array}
\tag{C.5}\label{eq:formalization-rfshriek}
\]
If a canonical choice of compactification is made for implementation, no
theorem statement may depend on that choice.

Integral constructible functors should still export finite-reduction
compatibilities.  For example,
\[
\begin{aligned}
\texttt{reduction\_rfStar}:&
 (Rf_*K)_n\simeq Rf_*(K_n),\\
\texttt{reduction\_rfShriek}:&
 (Rf_!K)_n\simeq Rf_!(K_n),\\
\texttt{reduction\_rfUpperShriek}:&
 (Rf^!L)_n\simeq Rf_n^!(L_n).
\end{aligned}
\tag{C.6}\label{eq:formalization-six-reduction}
\]
These are comparison theorems, not definitions.  Base change, projection
formula, Kunneth, localization, and their higher coherence are first fields
or theorems of the ambient correspondence formalism and then restrict to
constructibles.  The finite-reduction API remains useful as a conservative
test and as the bridge to classical torsion geometry; it is not the sole
coherence engine.

\begin{auditbox}
There must be no declaration identifying \(Rf_!\) with a derived functor of
an abelian pro-etale \(f_!\).  There must also be no theorem asserting proper
base change on the unrestricted naive pro-etale category, identifying
solidity with derived completion, or asserting
\(\Dcons=(\Dsolid)^\omega\).  Ambient compactification coherence and every
constructible-preservation theorem are separate release gates.  The
finite-coefficient theorems remain dependencies of the known proof even
though the ambient operations are not defined levelwise.
\end{auditbox}

\section{Layer 5: Weil structures, duality, and traces}

A Weil object should be a structure on a geometric-base-change object:
\[
 \texttt{WeilObject}(K)
   =\{\phi:F_X^*K\xrightarrow{\sim}K\}.
\tag{C.7}\label{eq:formalization-weil-object}
\]
Iteration must be an ordered composite, with lemmas for zero, successor, and
addition.  The six operations lift to Weil objects using explicit exchange
isomorphisms.

Keep distinct types or named constructions for:

\begin{enumerate}
\item arithmetic Galois Frobenius;
\item geometric Galois Frobenius;
\item the scheme Frobenius correspondence;
\item local closed-point Frobenius;
\item the induced endomorphism of a stalk or cohomology object.
\end{enumerate}

The trace layer should first be generic for dualizable objects in a symmetric
monoidal category, then specialized to perfect coefficient complexes.
The geometric trace formula returns equality between the categorical global
trace and local terms.  The point-counting corollary is a separate theorem
that performs the fixed-point and closed-point exponent conversion.

\section{Layer 6: perverse structures, including the integral pair}

Build the abstract machinery before mentioning sheaves:

\begin{enumerate}
\item bounded \(t\)-structures and their hearts;
\item gluing along recollement;
\item one-sided and two-sided \(t\)-exact functors;
\item truncation towers and finite filtrations;
\item images in a heart and the generic intermediate-extension construction;
\item tilting a \(t\)-structure by a torsion pair.
\end{enumerate}

For a noetherian space with dimension function \(\delta\), define the rational
support and cosupport predicates and prove that they produce the middle
perverse structure.  The terminal identities are
\[
\begin{aligned}
\texttt{perverseLE}(K)&\leftrightarrow
 \forall x\,a,\ a>-\delta(x)\to H^a(i_x^*K)=0,\\
\texttt{perverseGE}(K)&\leftrightarrow
 \forall x\,a,\ a<-\delta(x)\to H^a(i_x^!K)=0.
\end{aligned}
\tag{C.8}\label{eq:formalization-perverse-predicates}
\]

Over \(\mathcal O\), define both \(p\) and \(p^+\).  The required release
theorems include:
\[
\begin{aligned}
\texttt{verdierPToPPlus}&:
 ({}^p\Perv(X,\mathcal O))^{\op}
   \xrightarrow{\sim}{}^{p^+}\Perv(X,\mathcal O),\\
\texttt{verdierPPlusToP}&:
 ({}^{p^+}\Perv(X,\mathcal O))^{\op}
   \xrightarrow{\sim}{}^p\Perv(X,\mathcal O),\\
\texttt{rationalization\_perverse}&:
 {}^p\Perv(\mathcal O)\to\Perv(E),\\
\texttt{middleExtension\_p\_rationalize}&:
 ({}^pj_{!*}P)\otimes E\simeq j_{!*}(P\otimes E),\\
\texttt{middleExtension\_pplus\_rationalize}&:
 ({}^{p^+}j_{!*}P)\otimes E\simeq j_{!*}(P\otimes E).
\end{aligned}
\tag{C.9}\label{eq:formalization-integral-perverse}
\]
Do not install a Verdier self-duality instance for the integral \(p\)-heart.

\section{Layer 7: weights and the perverse engine}

Keep the following predicates distinct:
\[
\begin{array}{ll}
\texttt{PunctuallyPure}(\iota,w,\cF),&
\texttt{WeightsLE}(\iota,w,K),\\
\texttt{WeightsGE}(\iota,w,K),&
\texttt{Pure}(\iota,w,K),\\
\texttt{UniversallyPure}(w,K).&
\end{array}
\tag{C.10}\label{eq:formalization-weight-predicates}
\]

The ordinary/perverse bridge is a first-class theorem:
\[
\operatorname{WeightsLE}(K,w)
\Longleftrightarrow
\forall a,\ 
\operatorname{WeightsLE}(\pH^aK,w+a),
\tag{C.11}\label{eq:formalization-perverse-weight-bridge}
\]
and similarly for lower bounds and purity.  Tests must verify:
\[
 \wt(K[1])=\wt(K)+1,\qquad
 \wt(K(1))=\wt(K)-2,\qquad
 \wt(E_X[d])=d
\]
when \(X\) is smooth of pure dimension \(d\).

The main theorem should be exported first in perverse form:
\[
\texttt{perverseRfShriekWeightsLE}:
\operatorname{WeightsLE}(\cP,w)
\to
\operatorname{WeightsLE}(\pH^aRf_!\cP,w+a).
\tag{C.12}\label{eq:formalization-perverse-direct-image}
\]
Its dependencies are separate declarations for quasi-finite exactness,
open--closed devissage, relative Noether normalization, affine-line
projection, Fourier exactness, and the curve estimate.  The derived theorem
is then obtained from the finite perverse truncation tower.

\begin{powerbox}
This API order exposes the mathematical gain: the hard theorem returns
canonical perverse cohomology objects with support and boundary behavior.
Ordinary cohomology is recovered only after applying the perverse weight
criterion.
\end{powerbox}

\section{Layer 8: proper purity, IC, and decomposition}

Implement the following chain as separate theorems:
\[
\begin{aligned}
\texttt{rfStarWeightsGE}&:
 \operatorname{WeightsGE}(K,w)\to
 \operatorname{WeightsGE}(Rf_*K,w),\\
\texttt{properPushforwardPure}&:
 \operatorname{Proper}(f)\to
 \operatorname{Pure}(K,w)\to
 \operatorname{Pure}(Rf_*K,w),\\
\texttt{middleExtensionPure}&:
 \operatorname{Pure}(P,w)\to
 \operatorname{Pure}(j_{!*}P,w),\\
\texttt{intersectionComplexPure}&:
 \operatorname{PunctuallyPure}(L,u)\to
 \operatorname{Pure}(\IC_X(L),u+d),\\
\texttt{properIntersectionCohomologyPure}&:
 \operatorname{Proper}(X)\to
 \operatorname{Pure}(IH^r(X,L),u+d+r).
\end{aligned}
\tag{C.13}\label{eq:formalization-proper-purity-chain}
\]

Only after these declarations may the library add geometric semisimplicity,
splitting of a pure complex, and the decomposition theorem.  The splitting
returns existence of an isomorphism after geometric base change; it is
neither canonical nor a claim that Frobenius acts semisimply.

\section{Layer 9: final assembly}

The terminal purity term for a smooth proper \(d\)-fold should visibly
compose
\[
\begin{tikzcd}
 E_X[d]
 \arrow[r,"\text{proper purity}"]&
 Ra_*E_X[d]
 \arrow[r,"H^{i-d}"]&
 H^i(X,E).
\end{tikzcd}
\tag{C.14}\label{eq:formalization-terminal-term}
\]
The full Weil theorem combines this with three independent, reusable
geometry-free modules:

\begin{enumerate}
\item trace-to-zeta via the determinant--trace identity;
\item duality-to-functional-equation via perfect similitude pairings;
\item rational-series integrality and universal-weight separation.
\end{enumerate}

\section{Milestones and release gates}

\begin{longtable}{
 >{\raggedright\arraybackslash}p{0.11\textwidth}
 >{\raggedright\arraybackslash}p{0.49\textwidth}
 >{\raggedright\arraybackslash}p{0.30\textwidth}}
\toprule
Gate & Required artifact & Mandatory tests\\
\midrule
\endhead
M0 &
Pinned inventory and minimal imports. &
All \(\#check\) declarations compile.\\

M1 &
Generic completion and coefficient sheaf. &
\(\Zl/\ell^n\simeq\mathbf Z/\ell^n\);
completion idempotence.\\

M2 &
Constructible integral and rational categories. &
Point, finite etale scheme, and lisse representation examples.\\

M3 &
Completed four operations and localization. &
Finite reduction, composition, open--closed triangles.\\

M4 &
\(Rf_!,Rf^!\), base change, duality. &
Compactification independence and cocycle; affine-line compact support.\\

M5 &
Weil objects, twists, and trace. &
Eigenvalues \(1,q^{-1},q\) on a point, \(E(1)\), and
\(H^2(\PP^1,E)\).\\

M6 &
Rational and integral perverse structures. &
\(p/p^+\) dual exchange; \(E[d]\) perverse on a smooth \(d\)-fold.\\

M7 &
Fourier and curve engine. &
Delta/constant Fourier transforms, inversion shifts and twists.\\

M8 &
Direct-image weights and proper purity. &
No import from decomposition or final Weil files.\\

M9 &
IC purity, decomposition, and final assembly. &
All final declarations pass \(\#print\ \mathrm{axioms}\).\\
\bottomrule
\end{longtable}

\section{Import and axiom discipline}

Use four one-way import layers:
\[
\begin{tikzcd}[column sep=small]
\texttt{Generic}
 \arrow[r] &
\texttt{ProEtale}
 \arrow[r] &
\texttt{PerverseWeights}
 \arrow[r] &
\texttt{Weil}.
\end{tikzcd}
\tag{C.15}\label{eq:formalization-import-layers}
\]
No file in the first three layers may import the terminal Weil theorem.
Interface records used for parallel development must live under a visibly
temporary namespace and must not occur in the release import cone.

For every terminal theorem:

\begin{enumerate}
\item search the transitive imports for \(\texttt{sorry}\),
      \(\texttt{admit}\), and project-specific axioms;
\item run \(\#print\ \mathrm{axioms}\);
\item verify that all compactification, stratification, lattice, and finite
      coefficient-field choices have independence theorems;
\item test Frobenius and Tate conventions on the three basic examples above;
\item ensure that no definitional equality is being used where the
      mathematics supplies only a canonical isomorphism.
\end{enumerate}

This architecture makes the pro-etale and perverse infrastructure useful for
many other formalization projects: constructible sheaves with topological
coefficients, intersection cohomology, Fourier transforms, and decomposition
for proper maps all sit below the final Weil-specific layer.

\chapter{Primary-Source Guide and Theorem Crosswalk}
\label{app:proetale-references}

This appendix identifies the primary source of each deep theorem and records
the convention under which it is used.  Expository notes are valuable route
maps, but when a shift, twist, hypothesis, or logical dependency is delicate,
the primary theorem listed here governs this text.

\section{Bhatt--Scholze: the pro-etale model}

The primary coefficient source is Bhatt--Scholze,
\href{https://people.mpim-bonn.mpg.de/scholze/proetale.pdf}
{\emph{The pro-etale topology for schemes}}, Astérisque 369 (2015).

\begin{longtable}{
 >{\raggedright\arraybackslash}p{0.28\textwidth}
 >{\raggedright\arraybackslash}p{0.20\textwidth}
 >{\raggedright\arraybackslash}p{0.42\textwidth}}
\toprule
Topic & Primary location & Use in these notes\\
\midrule
\endhead
Weakly etale and pro-etale &
Definition 1.2; Theorem 1.3 &
Definition of \(X_{\proet}\), and comparison of weakly etale with
ind-etale maps.\\

Weakly contractible basis &
Definition 1.4; Theorem 1.5; Theorem 1.8 &
\cref{thm:proetale-wcontractible-cover}; exact evaluation and control of
inverse limits.\\

Replete topoi and inverse limits &
Propositions 3.1.10--3.1.11; Section 3.3 &
Surjective systems satisfy
\(\varprojlim\simeq\Rlim\); countable inverse limit has cohomological
dimension one; left completeness.\\

Derived completion in a replete topos &
Proposition 3.5.1; Lemma 3.5.7 &
For a replete topos, a noetherian ring, and an ideal \(\mathfrak m\), derived
completeness is reconstruction from the reductions.  Complete complexes are
compared with compatible systems in the bounded-above range; the unbounded
comparison additionally assumes that \(\mathfrak m\) is regular.\\

Comparison with the etale site &
Corollaries 5.1.5--5.1.6 &
Discrete local systems agree, and bounded-below torsion complexes have the
same cohomology after pullback by \(\nu\).\\

Complete constructible objects &
Definition 6.5.1; Lemmas 6.5.3--6.5.9 &
\cref{thm:proetale-complete-system-equivalence}; derived completeness,
finite reduction, tensor closure, and locally closed functors.\\

Stratified characterization &
Proposition 6.6.11 &
\cref{thm:proetale-constructible-characterization}: constructible exactly
means locally perfect along a finite stratification when \(X\) is
topologically noetherian.\\

Pushforward and base change &
Theorem 6.7.1; Lemmas 6.7.2, 6.7.4--6.7.5 &
Constructibility of \(Rf_*\), smooth base change, and proper base change
under the stated finiteness hypotheses.\\

Exceptional direct image &
Definition 6.7.6; Lemma 6.7.7; Lemma 6.7.10 &
Compactification definition of \(Rf_!\), independence, constructibility,
and base change.\\

Failure of the naive derived \(f_!\) &
Remarks 6.7.8--6.7.9 &
\Cref{warn:dictionary-fshriek}; prevents an incorrect abelian-derived
construction.\\

Tensor, internal Hom, projection &
Lemmas 6.7.12--6.7.14 &
Closed monoidal constructible category and projection formula.\\

Exceptional inverse image &
Lemmas 6.7.18--6.7.19 &
Construction of \(Rf^!\) from compatible finite-level adjoints.\\

Dualizing complexes &
Lemma 6.7.20 &
Integral Verdier duality, built from the finite-level dualizing complexes.\\

Integral structures for lisse \(E\)-sheaves &
Lemma 6.8.13 &
On a topologically noetherian scheme, every lisse \(E\)-sheaf admits a
constructible \(\mathcal O_E\)-model.  The model need not be lisse globally;
the proof first obtains lattices after a finite stratification.\\

\(\mathcal O_E,E,\Qbarl\)-coefficients &
Definitions 6.8.1--6.8.3; Proposition 6.8.14; Remark 6.8.15 &
\cref{thm:proetale-coefficient-tower,%
thm:proetale-six-functor-package}; integral lattices, inversion of the
uniformizer, finite coefficient-field descent, and six operations.\\

Pro-etale fundamental group &
Theorem 1.10; Section 7 &
Continuous representations on nonnormal schemes; warning against replacing
\(\pi_1^{\proet}\) everywhere by its profinite completion.\\
\bottomrule
\end{longtable}

The hypotheses of Section 6.7 are part of every use: complete noetherian
local coefficients with finite residue field of characteristic \(\ell\);
qcqs and finitely presented maps; and either properness or the relevant
quasi-excellent, \(\ell\)-coprime hypothesis.  Finite-type schemes over
\(\Fq\), with \(\ell\nmid q\), lie in this range.

\section{Modern constructibility by dualizability}

Hemo--Richarz--Scholbach,
\href{https://arxiv.org/abs/2305.18131}
{\emph{Constructible sheaves on schemes}}, Advances in Mathematics 429
(2023), gives the coefficient-flexible categorical formulation.

\begin{longtable}{
 >{\raggedright\arraybackslash}p{0.29\textwidth}
 >{\raggedright\arraybackslash}p{0.20\textwidth}
 >{\raggedright\arraybackslash}p{0.41\textwidth}}
\toprule
Topic & Location & Formalization role\\
\midrule
\endhead
Lisse as dualizable &
Definition 1.1 &
The reusable predicate
\(\operatorname{Lisse}(K)\leftrightarrow
 \operatorname{IsDualizable}(K)\).\\

Weakly contractible calculation &
Lemma 1.2, referring to the detailed theorem in Section 4 &
On a weakly contractible affine, lisse objects are perfect coefficient
complexes.\\

Pro-etale descent &
Lemma 1.3 and the detailed descent results &
Lisse and constructible categories form hypersheaves; choices of
trivializing covers do not enter public objects.\\

Sequential coefficient limits &
Proposition 5.1 &
If \(\Lambda=\varprojlim_n\Lambda_n\) is sequential, with surjective
transitions having locally nilpotent kernels, then lisse and constructible
categories are the corresponding inverse limits.\\

Filtered coefficient colimits &
Proposition 5.2 &
For qcqs \(X\) and \(\Lambda=\varinjlim_i\Lambda_i\), lisse and constructible
categories are the corresponding filtered colimits.\\

Localization of coefficients &
Proposition 5.5; Remark 5.7 &
For qcqs \(X\), localization along a multiplicative subset gives fully
faithful functors on lisse and constructible categories.  Essential
surjectivity is not automatic for lisse objects, reflecting the distinction
between the etale and pro-etale fundamental groups.\\

Natural \(t\)-structure &
Theorem 6.2 &
The coefficient ring must be \(t\)-admissible.  Lisse objects additionally
require that \(X\) have Zariski locally finitely many irreducible components;
constructible objects require this for every constructible subset of \(X\).
These hypotheses are part of the interface.\\

Classical adic comparison &
Theorem 7.7 &
For qcqs \(X\), the pro-etale constructible categories with
\(\mathcal O_E/\mathfrak m^n\), \(\mathcal O_E\), and \(E\) coefficients
agree with the corresponding classical finite-Tor and adic categories.\\

Algebraic coefficient closure &
Lemma 7.9 &
For qcqs \(X\), filtered colimits over finite extensions
\(E/\Ql\) recover the constructible categories with \(\overline{\mathbf
Z}_\ell\) and \(\Qbarl\) coefficients.\\
\bottomrule
\end{longtable}

The notes use Bhatt--Scholze's concrete complete-DVR model for the proof
spine and the dualizability formulation as the reusable public abstraction.

\section{Finite-level six operations, purity, and trace}

The pro-etale formalism lifts finite-level geometric theorems; it does not
replace their proofs.

\begin{longtable}{
 >{\raggedright\arraybackslash}p{0.25\textwidth}
 >{\raggedright\arraybackslash}p{0.27\textwidth}
 >{\raggedright\arraybackslash}p{0.38\textwidth}}
\toprule
Result & Primary source & Use\\
\midrule
\endhead
Etale topoi and torsion cohomology &
SGA 4 \cite{SGA4}, Exposes V--XII &
Finite coefficient derived categories, stalks, cohomological dimension,
and descent.\\

Finiteness and base change &
SGA 4 \cite{SGA4}, Exposes XIV--XVII; SGA \(4\frac12\)
\cite{SGA4half} &
Constructibility and finite amplitude, proper/smooth base change, and
compact support.\\

Exceptional operations &
SGA 4 \cite{SGA4}, Expose XVII &
Compactification, \(Rf_!\), composition, and base change at finite level.\\

Duality and quasi-excellent refinements &
\href{https://www.numdam.org/item/AST_2014__363-364__R1_0/}
{Illusie--Laszlo--Orgogozo, \emph{Travaux de Gabber}},
especially Expose XVII &
Finite-level \(Rf^!\), dualizing complexes, and compatibility in the range
used by Bhatt--Scholze 6.7.18--6.7.20.\\

Absolute cohomological purity &
Gabber's absolute purity theorem, as presented in the same volumes &
Costalk normalization, smooth purity, and Poincare duality.\\

Lefschetz--Verdier trace &
SGA 5 \cite{SGA5}; SGA \(4\frac12\) \cite{SGA4half} &
\Cref{ch:trace-formula}; local terms, Frobenius fixed points, and the
cohomological zeta formula.\\
\bottomrule
\end{longtable}

\section{BBD: perverse sheaves and weights}

The categorical and weight-theoretic source is
Beilinson--Bernstein--Deligne, \emph{Faisceaux pervers} \cite{BBD}.

\begin{longtable}{
 >{\raggedright\arraybackslash}p{0.29\textwidth}
 >{\raggedright\arraybackslash}p{0.19\textwidth}
 >{\raggedright\arraybackslash}p{0.42\textwidth}}
\toprule
Result & BBD location & Use here\\
\midrule
\endhead
Gluing and middle perversity &
Sections 1.4, 2.1--2.2 &
\cref{def:middle-perversity,thm:middle-t-structure}.\\

Integral \(p\) and \(p^+\) &
Section 3.3 &
\cref{thm:p-pplus-tilt,thm:integral-verdier-exchange};
duality exchanges the two hearts.\\

Affine and quasi-finite exactness &
Theorem 4.1.1 and Corollaries 4.1.2--4.1.3 &
\cref{prop:basic-perverse-exactness,thm:perverse-artin-engine}.\\

Smooth and product exactness &
4.2.4 and 4.2.7--4.2.8 &
Dimension-shifted smooth pullback and external products.\\

Finite length and simple objects &
Theorem 4.3.1 &
\cref{thm:perverse-finite-length,thm:simple-perverse-classification}.\\

Geometric positive-degree Ext vanishing &
5.1.15(iii) &
Splitting arithmetic extensions after geometric base change; no claim of
Frobenius semisimplicity.\\

Intermediate extension and weights &
Corollary 5.3.2 &
Purity of \(j_{!*}\), used only after the upper direct-image theorem.\\

Pure perverse semisimplicity &
Theorem 5.3.8 &
Geometric semisimplicity of the rational perverse object.\\

Ordinary/perverse weight criterion &
Theorem 5.4.1 and Corollary 5.4.4 &
Conversion between derived weights and weights on
\(\pH^a\).\\

Splitting of pure complexes &
Theorem 5.4.5 &
Noncanonical splitting after geometric base change.\\

Decomposition theorem &
Theorem 6.2.5 &
\cref{thm:decomposition-engine,thm:decomposition-finite-field};
a consequence of purity, not an input to it.\\
\bottomrule
\end{longtable}

For an accessible discussion of the integral failure of self-duality and
the \(p/p^+\) pair, see Juteau--Mautner--Williamson,
\href{https://arxiv.org/abs/0901.3322}
{\emph{Perverse sheaves and modular representation theory}}, Section 2.3,
which points back to BBD 3.3.

\section{Deligne: the arithmetic weight theorem}

Deligne's \emph{La conjecture de Weil II} \cite{DeligneWeilII} supplies the
arithmetic estimates and their exact noncircular order.

\begin{longtable}{
 >{\raggedright\arraybackslash}p{0.28\textwidth}
 >{\raggedright\arraybackslash}p{0.18\textwidth}
 >{\raggedright\arraybackslash}p{0.44\textwidth}}
\toprule
Result & Location & Role\\
\midrule
\endhead
Weights and determinant weights &
Sections 1.2--1.3 &
Definitions, tensor/dual/twist bookkeeping, and rank-one inputs.\\

Real-sheaf criterion &
Theorem 1.5.1 &
Mixedness mechanism used before the general direct-image theorem.\\

Local monodromy weights &
Theorem 1.8.4 &
Boundary terms in the local Fourier and curve arguments.\\

Curve theorem &
Theorem 3.2.3 &
\cref{thm:laumon-curve-engine}; reconstructed through perverse Fourier
transform.\\

General direct-image upper bound &
Theorem 3.3.1 &
\cref{thm:perverse-direct-image-weight-bound,%
thm:deligne-laumon-direct-image}.\\

Cohomological upper and lower bounds &
Corollaries 3.3.4--3.3.6 &
Corollary 3.3.4 says that \(H_c^i\) sends weights \(\leq n\) to weights
\(\leq n+i\).  Corollary 3.3.5 gives the dual lower bound for \(H^i\) on a
smooth scheme with lisse coefficients.  Corollary 3.3.6 says that, for a
pure lisse sheaf of weight \(n\) on a smooth scheme, the image of
\(H_c^i\to H^i\) is pure of weight \(n+i\).  Proper purity follows in the
proper case because this map is an isomorphism; it is not the literal
statement of all three corollaries.\\

Smooth-proper consequence &
Corollary 3.3.9 &
\cref{thm:weil-rh-perverse-engine,%
cor:perverse-smooth-proper-purity}.\\
\bottomrule
\end{longtable}

\section{Laumon: Fourier transform as a perverse equivalence}

The Fourier route is from Laumon,
\emph{Transformation de Fourier, constantes d'equations fonctionnelles et
conjecture de Weil} \cite{LaumonFourier}.

\begin{longtable}{
 >{\raggedright\arraybackslash}p{0.28\textwidth}
 >{\raggedright\arraybackslash}p{0.20\textwidth}
 >{\raggedright\arraybackslash}p{0.42\textwidth}}
\toprule
Result & Location & Normalization/use\\
\midrule
\endhead
Definition &
1.2.1.1 &
\[
 \FT_\psi(K)
  =Rp_{2!}
   \bigl(p_1^*K\otimes\cL_\psi(\langle x,\xi\rangle)\bigr)[r].
\]\\

Inversion &
Theorem 1.2.2.1 &
Inverse character gives the inverse transform with the prescribed Tate
twist; using the same character inserts pullback by the sign map.\\

Support-forgetting &
Theorem 1.3.1.1 &
The Fourier-kernel map \(Rp_!\to Rp_*\) is an isomorphism.\\

Duality &
Corollary 1.3.2.2 &
Verdier duality replaces \(\psi\) by \(\psi^{-1}\) and contributes the
displayed twist.\\

Perverse \(t\)-exactness &
Theorem 1.3.2.3 and Corollary 1.3.2.4 &
\cref{thm:fourier-perverse-equivalence-engine,%
thm:fourier-perverse-texact-detailed}; exact equivalence of perverse hearts.\\

Local Fourier transforms &
Sections 2.3--2.5 &
Boundary monodromy and conductor contributions.\\

Curve-weight proof &
Sections 4.1--4.4 &
The perverse curve engine and the reduction to the global direct-image
theorem.\\
\bottomrule
\end{longtable}

The additive character is allowed to enlarge the finite coefficient field
\(E\).  The Fourier kernel has an \(\mathcal O_E\)-lattice; it is not
asserted to be defined over bare \(\Zl\).

\section{Classical Weil assembly}

The trace-to-zeta and duality-to-functional-equation steps follow the
classical cohomological architecture in SGA 5 \cite{SGA5}, Deligne
\cite{DeligneWeilI,DeligneWeilII}, and standard references such as
Freitag--Kiehl \cite{FreitagKiehl}.  In these notes, however, their algebraic
content is proved independently in the algebraic-core chapter:

\begin{center}
\begin{tabularx}{0.94\textwidth}{>{\raggedright\arraybackslash}p{0.32\textwidth}X}
\toprule
Geometric input & Geometry-free output\\
\midrule
Trace formula &
Alternating determinant expression for \(Z(X,T)\).\\
Perfect Poincare pairing &
Reciprocal characteristic polynomials and functional equation.\\
Universal-realization purity &
Separation of factors of distinct weights.\\
Integral rational series &
Integral individual factors by the Fatou lattice and Gauss lemma.\\
\bottomrule
\end{tabularx}
\end{center}

\section{Convention crosswalk}

\begin{longtable}{
 >{\raggedright\arraybackslash}p{0.30\textwidth}
 >{\raggedright\arraybackslash}p{0.29\textwidth}
 >{\raggedright\arraybackslash}p{0.31\textwidth}}
\toprule
Convention in these notes & Consequence & Conversion warning\\
\midrule
\endhead
Geometric Frobenius &
\(F\mid E(1)=q^{-1}\). &
Arithmetic Frobenius inverts every eigenvalue.\\

Derived weight convention &
\(\cH^aK\) has weight \(w+a\) when \(K\) is pure of weight \(w\). &
Thus \(K[n]\) is pure of weight \(w+n\).\\

Perverse-normalized IC &
\(\IC_X(\cL)=j_{!*}(\cL[d])\). &
An unshifted intersection complex differs by \([-d]\).\\

Integral perversity &
Verdier duality sends \(p\) to \(p^+\). &
Self-duality is asserted only after rationalization.\\

Fourier normalization &
The rank-\(r\) transform includes \([r]\). &
Without this shift it is not perverse \(t\)-exact.\\

Fourier duality &
\(\psi\) changes to \(\psi^{-1}\). &
Using the same character also changes the sign-map convention.\\

Frobenius polynomial &
\(P_i(T)=\det(1-TF\mid H^i)\). &
Polynomial roots are reciprocals of Frobenius eigenvalues.\\

Comparison categories &
Classical and pro-etale objects correspond by an equivalence. &
The equivalence and its compatibilities are not definitional equalities.\\
\bottomrule
\end{longtable}

These choices determine every shift, twist, exponent, and absolute value in
the terminal theorem.

\section{Expository route maps}

Kiehl--Weissauer \cite{KiehlWeissauer} remains the closest book-length guide
to the perverse Fourier proof.  The Stanford Weil II seminar
\cite{StanfordWeilIISeminar}, especially Lectures 18--24
\cite{StanfordL18,StanfordL19,StanfordL22,StanfordL23,StanfordL24}, gives a
useful reading sequence.  These sources predate or do not center the direct
pro-etale coefficient model; Bhatt--Scholze governs that replacement.


\backmatter
\bibliographystyle{alpha}
\bibliography{bibliography}

\end{document}
